%\newpage
\section{Introduction} 
The analysis presented in this chapter is based on Ref.~\cite{Bhattacharyya:2023kbh}.
Detection of gravitational waves (GW) \cite{LIGOScientific:2016aoc, LIGOScientific:2016sjg, LIGOScientific:2016vlm, LIGOScientific:2017bnn, Kokeyama:2020dkg, LIGOScientific:2019hgc} opens a new pathway for testing different theories of gravity. Gravitational waves are generated primarily by binary systems consisting of black holes and neutron stars. These achievements motivate increasingly detailed probes of the structure of the source. The evolution of such binaries has three phases: \textit{inspiral}, \textit{merger}, and \textit{ringdown}. Probing the merger phase, where gravity is strong, requires non-perturbative methods based on numerical simulations \cite{Lehner:2014asa, LIGOScientific:2014oec}. However, the inspiral phase can be studied using perturbative analytic techniques because the orbital velocity of the binaries is small compared to the speed of light, i.e. $\frac{v}{c}\ll 1$. \textcolor{black}{Signals detected by ground-based detectors derive much of their constraining power from the long duration of the inspiral phase}, so this binary-inspiral problem can be studied using the Post-Newtonian approximation \cite{Blanchet:2013haa}. %One can solve the problem by investigating the PN approximation involving the energy-momentum tensor. 
\par
The gravitational-wave radiation from non-spinning binaries has been computed up to 3PN order in \cite{Tagoshi:2000zg, Faye:2006gx, Blanchet:2006gy, Blanchet:2004ek, Damour:2000ni, Itoh:2003fy,Boetzel:2019nfw,Faye:2012we,Mishra:2013rna} using traditional methods based on Green's functions. Recently, this computation has been extended up to 4.5PN order in \cite{Fujita:2010xj,Faye:2014fra, Blanchet:2023sbv,Blanchet:2023bwj} \footnote{Interestingly, both UV and IR divergences occur during the computation. The UV divergences appear because the binaries are modelled as point-like particles without internal structure. On the other hand, IR divergences occur due to the insertion of a typically divergent (at long distance) PN expansion in the near zone while defining the source moments. Interested readers are referred to \cite{Larrouturou:2021gqo, Blanchet:2004bb,Blanchet:2023soy} for a detailed discussion of the origin of these divergences and how to regularize them.}. Using these waveform results, parameter estimation based on the Fisher information matrix has been carried out in \cite{Arun:2004hn}. Later, these studies were extended by taking into account the effect of the binary's spin at leading PN order in \cite{PhysRevD.12.329, PhysRevD.2.1428, Kidder:1992fr}. Several subsequent works pushed this analysis to higher PN orders \cite{Cho:2022syn}. The Hamiltonians generating the binary dynamics, including spin effects up to 3PN order, were also derived using the ADM formalism in \cite{Steinhoff:2007mb, Steinhoff:2008ji, Hergt:2008jn, Hergt:2010pa,Porto:2012as}. Gravitational-wave radiation for eccentric orbits has likewise been studied in \cite{Enoki:2006kj, Favata:2011qi, Munna:2019fjz}. Beyond general relativity (GR), these computations are also being extended to alternative theories of gravity \cite{Zhang:2017srh,AbhishekChowdhuri:2022ora, Zhang:2018prg,Saffer:2018jmx, Lin:2018ken,Li:2022grj,Shiralilou:2021mfl,Julie:2019sab}.
\par
%Around the time when only upto 2.5 PN spin corrections were known \cite{Tagoshi:2000zg, Faye:2006gx, Blanchet:2006gy}, 
In parallel with the traditional method, a more efficient and systematic methodology based on \textit{\textcolor{black}{Worldline} Effective Field Theory} (WEFT) emerged for computing PN corrections. This framework is sometimes called ``Non-Relativistic General Relativity'' (NRGR) \cite{Goldberger:2004jt}. It has been used to compute gravitational-wave radiation at leading order in \cite{Goldberger:2004jt}. \textcolor{black}{In this approach, one can decouple and treat the conservative and radiative degrees of freedom separately. Integrating out the potential modes (heavy modes) yields an effective action for the radiative modes (light modes). The real part of this effective action determines the effective potential and hence the binary dynamics, while the imaginary part gives the total radiated power through an optical theorem.}\footnote{Like the traditional method, various divergences occur while computing various multipole moments of the source using the EFT method. Interested readers are referred to \cite{Porto:2017dgs,Porto:2007pw} for more details about how to handle these divergences (in particular, the IR ones). Also, the renormalizability of EFTs of binary systems and the necessity of a time-dependent mass counterterm have been discussed in \cite{Goldberger:2012kf,Galley:2015kus}.} Subsequently, the method was developed further to include effects of the binary's spin (up to quadratic order) in the radiative multipole moments up to \cite{Porto:2010zg} 
\textcolor{black}{3PN order}, and then it has been extended to higher PN  orders in \cite{Levi:2011eq, Levi:2015uxa, Levi:2015ixa, Levi:2014sba}. Apart from these, this EFT based approach has been used to study various phenomena like self-force effect on the curved spacetime, gravitational tail effects and also inspecting conservative and radiative dynamics of the binary inspiral problem (via computing effective Hamiltonian/Lagrangian) \cite{Maia:2017gxn,Maia:2017yok,Foffa:2019yfl,Kalin:2022hph,Porto:2008jj,Porto:2008tb,Porto:2007px,Mandal:2022nty,Mandal:2022ufb,Goldberger:2009qd,Ross:2012fc,Goldberger:2020fot,Goldberger:2017ogt,Goldberger:2006bd,Goldberger:2005cd}
\footnote{The list is by no means exhaustive. Interested readers are referred to the following reviews \cite{Goldberger:2022rqf,Goldberger:2022ebt} and references therein for more details.}
 Finally, the EFT approach has been extended to compute the conservative dynamics as well as power radiation due to a massive (as well as massless) scalar field minimally coupled to GR up to 1PN order \cite{Kuntz:2019zef,Huang:2018pbu}, and to study the conservative dynamics of the electromagnetic field minimally coupled to GR in \cite{Patil:2020dme,Gupta:2022spq}.
\par
Ultra-light bosonic fields are promising candidates for dark-matter \cite{GrillidiCortona:2015jxo, Sanchis-Gual:2022ooi, Goldstein:2022pxu, Schutz:2020jox} and  dynamical dark-energy \cite{Kamionkowski:2014zda}
\footnote{The list is by no means exhaustive.  Interested readers are referred to this Snowmass review \cite{Adams:2022pbo} and references therein for more details.}. One famous example of such ultra-light bosonic fields is the QCD axion. It arises from the Peccei--Quinn mechanism \cite{Fukuda:2021drn,CAPP:2020utb,Sakhelashvili:2021eid, Peccei:2006as, RevModPhys.82.557,PhysRevLett.40.223,PhysRevLett.40.279,Preskill:1982cy,Gorghetto:2018ocs,DiLuzio:2021pxd,Berezhiani:2000gh,Fukuda:2015ana,Dimopoulos:2016lvn,Gherghetta:2016fhp,Kim:1998va,DiLuzio:2020wdo,Conlon:2006tq,Chakraborty:2021fkp,Harigaya:2019qnl}, which was proposed as a solution to the strong CP problem \cite{PhysRevLett.38.1440,PhysRevD.16.1791,PhysRevLett.40.279,Clowe:2006eq,Hsu:2004mf,Agrawal:2017ksf,Gupta:2020vxb,PhysRevLett.43.103}. Experimental bounds on the neutron electric dipole moment imply that the strong CP angle must be much smaller than $10^{-10}$, while the CP angle in the CKM matrix is known to be $\mathcal{O}(1)$ \cite{Baker:2006ts}. This tension can be resolved by introducing the axionic coupling \cite{Huang:2018pbu, Zhang:2021mks}.$$\frac{a}{f_a}\frac{g_{a\gamma\gamma}^2}{32\pi^2}F^{\mu\nu}\tilde{F}^{\mu\nu},$$
where $g_{a\gamma\gamma}$ is the strong coupling constant, $\tilde{F}^{\mu\nu}$ is the dual field strength tensor, and $f_a$ is the axionic decay constant. QCD Axions have several constraints.  For example, the ADMX experiment showed that the first constraint on the QCD axion parameter space in the $\mu$eV mass range \cite{PhysRevLett.120.151301}.  Also, it was suggested that one could possibly obtain constraints on the parameter space for QCD axions by studying its effects on the stellar configuration of a white dwarf \cite{Balkin:2022qer}.  There are other ultra-light bosons with similar properties to the QCD axion, with only the difference that their mass is not related to the decay constant.  They are often termed as \textit{Axion-like particles (ALP)}. They can arise from different considerations, e.g., from string compactifications  \cite{Cicoli:2013ana,Svrcek:2006hf,Hiramatsu:2010yu}.  There exist phenomenological constraints on the mass of these ultra-light bosons stemming from dark-matter phenomenology, black hole superradiance  \cite{Arias:2012az, DiLuzio:2021pxd, QuilezLasanta:2021yzt, Ishii:2022lwc, 
 Cardoso:2018tly, Brito:2015oca,Ghosh:2023tyz} as well as on its anomalous coupling with electromagnetic field ($g_{a\gamma\gamma}$) by using the induced oscillating electric dipole moment of the electron as advocated in \cite{Hill:2015vma}.
Besides axions (ultra-light boson fields), \textit{ultra-light vector fields} (ULV) are also possible candidates for dark matter \cite{Cardoso:2018tly}.  One example of ultra-light vector fields is the so-called \textit{dark photon} \cite{Flambaum:2019cih}.  This can often be thought of as a \textit{hidden U(1)} massive vector field which can interact with the photon \cite{Cardoso:2018tly}.
\par
Both the ultra-light boson and vector fields can couple with gravity, e.g., axions can couple with gravity minimally and non-minimally.  The effect of axions on
the detected waveform is negligible until the binaries are separated by roughly a Compton wavelength of the axion \cite{Zhang:2019eid}. As the orbital length scale decays to the Compton wavelength, scalar radiation may become an important source
of orbital energy loss \cite{Huang:2018pbu}, especially for large Compton wavelengths.  Scalar radiation is detectable for both NS-NS and NS-BH binaries \cite{Dar:2018dra}.  Furthermore, one can write down Chern-Simons type terms both in the electromagnetic and gravitational sectors by which the massive scalar (axion) can couple with both electromagnetic and gravitational fields.  This is natural in the context when we consider the axion to play the role of dark matter \cite{Yoshida:2017cjl}.  If the axions interact with gravity via the Chern-Simons term, GW waves induce axion decays into gravitons \cite{Yoshida:2017cjl}.  Apart from this, in \cite{Machado:2018nqk, Machado:2019xuc}, it was shown that if axions couple with the dark photon with an unbroken $U(1)$ symmetry, there is a possibility of the generation of a stochastic gravitational wave when it experiences a tachyonic instability.  The resulting signal can be detected by ground-based GW detectors  \cite{Machado:2018nqk, Machado:2019xuc,10.21468/SciPostPhys.12.5.171}.  Last but not least, as advocated in \cite{ Ejlli:2022zah, Nagano:2021kwx}, polarimetry experiments on GW signals may provide a novel way to probe the axion-like particles.
\par 
Previously, several works have taken steps toward constraining \textcolor{black}{ALP couplings} using GW observations \cite{Huang:2018pbu, Zhang:2021mks}. This is done by considering a massive scalar field minimally coupled to GR and then computing the additional contribution to the total radiated power at leading PN order, together with the correction to the phase of the emitted GW due to the presence of the scalar field. In this paper, we consider a more generic scenario. Motivated by the dark-matter model discussed in \cite{Cardoso:2018tly}, we study a setup containing both the ALP, in the form of a massive scalar field with an anomalous coupling to the electromagnetic field via a Chern--Simons-type term, and the ULV, in the form of a massive vector field (Proca field) that also couples to the electromagnetic field through a kinetic-mixing term. For simplicity, we neglect the coupling of the ALP through the gravitational Chern--Simons term. Our broader goal is to investigate whether GW observations can constrain, in addition to the masses of the scalar and Proca fields, the axion-photon coupling parameter $g_{a\gamma\gamma}$ as well as the Proca-electromagnetic coupling constant $\gamma$ arising from kinetic mixing. To that end, we take a first step by computing the power radiation from a binary system in this dark-matter environment, which is essential for determining the phase of the gravitational waveform within the WEFT approach, and we eventually comment on the possibility of detecting these couplings.
\par
%\vspace{0.3 cm}
This chapter is organised as follows: In Sec.~(\ref{Sec2}), we briefly review the general machinery of the EFT approach. We also discuss various length scales associated with the EFT of inspiralling binaries, which is crucial for separating various degrees of freedom. In Sec.~(\ref{Sec3}) we discuss how the binaries can be modelled by a point particle action. We review the procedure of integrating out the potential modes to get effective action for Einstein's gravity. We also give the EFT power counting for the bound sector. In Sec.~(\ref{Sec4}), we introduce our model, which consists of a scalar field (axion), electromagnetic (photon), and Proca field (dark photon) minimally coupled with Einstein gravity and an interaction term between photon and dark photon as well between axion and photon. Then we focus on the bound sector, which gives rise to the corrections to the effective potential. \textit{Interestingly, we observe that the correction due to the axion-photon coupling ($g_{a\gamma\gamma}$) enters at the 2.5PN order.} In Sec.~(\ref{Sec5}), we discuss the dissipative sector. We used the familiar optical theorem to calculate the power radiation. We again discuss the EFT power counting for the radiative sector. We describe the corrections to the gravitational radiation coming from different field vertices. We reproduce the well-known expression for the power radiation coming from the pure gravitational and scalar sectors at the leading order. \par 
We extend the computation to the $N^{(4)}LO$ and also showed that the $g_{a\gamma\gamma}$ appears at $N^{(5)}LO$ for the  scalar sector. Due to the presence of other fields, we obtain new contributions to the power radiation from the gravitational sector at different orders of perturbations. We observe that the contribution due to the axion-photon coupling ($g_{a\gamma\gamma}$) enters at the $N^{(7)}LO$ in the expression for the power radiation from the gravitational sector. Furthermore, we also compute the contribution of electromagnetic and Proca fields to the total power radiation at the leading order.  We also comment on the possibility of detecting axion coupling ($g_{a\gamma\gamma}$) through the gravitational flux. Lastly, the kinetic mixing constant $\gamma$ starts contributing from $1PN$ in conservative sector and appears in radiative sector at $N^{(2)}LO$ and $N^{(4)}LO$ through gravitational radiation.\par
Finally, in Sec.~(\ref{disc}), we summarise our main findings and conclude with some future directions. Some details regarding the computation of a few integrals are given in Appendix ~(\ref{ch1:app:A}), (\ref{ch1:app:B}), (\ref{ch1:app:C}) and (\ref{ch1:app:E}), while comments on the middle vertex contribution to radiation are collected in Appendix~(\ref{ch1:app:D}).

For the reader, the chapter is easiest to follow in three passes. First, the introductory sections identify the model, the relevant scales, and the EFT power counting. Second, the conservative and radiative sections explain how each physical sector modifies the orbit or the emitted power. Third, the appendices collect the more repetitive integral reductions and tensor algebra. In the main text, we therefore emphasise the physical role of each class of diagrams, work out one representative computation whenever a new contraction pattern appears, and quote the remaining contributions in a compact form.
\subsection*{\textbf{ \textit{Notations and conventions:}}}

\begin{multicols}{2}
\begin{itemize}
    \item Metric signature: (-,+,+,+).
    \item $\rmint \mathcal{D}\hat\xi \,e^{iS_{\text{quad}}}\rightarrow \rmint \Bar{\mathcal{D}}\hat\xi$.
    \item $\rmint \frac{d^3\boldsymbol{k}}{(2\pi)^3}\rightarrow \rmint_{\boldsymbol{k}}.$
    \item ${\boldsymbol{x}}_0=\frac{\boldsymbol{x}_1-\boldsymbol{x}_2}{2}.$
    \item Reduced mass: $\mu=\frac{m_1m_2}{M}$\,\,\text{with},\,\,$M=m_1+m_2$.
    \item Mass ratio: $\nu=\frac{\mu}{m_1+m_2}$.
    \item Planck mass: $m_p:=\frac{1}{\sqrt{8\,\pi\, G}},\textstyle{with\, (\hbar,\,c)=1.}$
\end{itemize}
\end{multicols}
\begin{itemize}
\item Levi-Civita symbol convention: $\epsilon^{\mu\nu\rho\sigma}\equiv\frac{\Hat{\epsilon}^{\,\mu\nu\rho\sigma}}{\sqrt{-g}}$, with $\Hat{\epsilon}^{\,0123}=1$.
\item $nPN$ in conservative sector symbolizes $\sim \mathcal{O}(Lv^{2n})$, where $L$ denotes angular momentum.
 $LO$ (Leading Order) symbolizes $\sim\mathcal{O}$($L^{1/2}v^{1/2}$) in radiation sector. Then the subsequent orders are denoted by $N^{(n)}LO\sim \mathcal{O}(L^{1/2}v^{n+1/2}).$

\end{itemize}

\section{Brief review of machinery of EFTs}\label{Sec2}
Consider a QFT with a general tensor field $\xi_{\alpha\beta..}$ in 3+1 dimensions described by the action $S[\xi]$. In order to compute the Effective action, one could first write down the generating functional and then take the logarithm. Usually, in QFT, one writes down the generating functions corresponding to the connected diagrams and, by Legendre transformation, can systematically evaluate the effective action as,
\begin{eqnarray}
        \mathcal{Z}[J]&=&\rmint \mathcal{D}\xi_{\alpha\beta..} \, e^{iS[\xi]+\rmint J\cdot \xi},\nonumber\\
      e^{i W[J]}&=& \mathcal{Z}[J],\\
      S_{\text{eff}}[\xi_{\text{cl}}]&=& W[J]-\rmint d^4x J\,\xi_{\text{cl}}.\nonumber \label{2.3m}
\end{eqnarray}
Here $W[J]$ have contributions only from connected diagrams and $\xi_{\text{cl}}$ is defined as $\xi_{\text{cl}}=\frac{\delta W[J]}{\delta J}$. 
 Effective action in (\ref{2.3m}) is complex in general. \textcolor{black}{The real part gives the conservative dynamics of the system, and the imaginary part gives the radiative dynamics of the system. Hence, for a binary inspiral problem, the real part gives the equation of motion, i.e., the orbit equation of the binary system, and the imaginary part gives the power radiation from the system.} The perturbative computation in $\hbar$ also leads to effective action at the classical and quantum levels.
 \begin{align}
     \begin{split}
S_{\text{eff}}^{\text{tot}}=\sum \underbrace{\text{Tree-level diagrams}}_{\mathcal{O}(\hbar^0)(\text{classical contribution})}+ \sum \underbrace{\text{loop level diagram}}_{\mathcal{O}(\hbar^n)(\text{quantum corrections})}. \label{Stotal}
     \end{split}
 \end{align}
 \par
 At this point, we assume that the theory has two different energy scales, a high energy scale ($\Lambda$) and a low energy scale ($\epsilon$). As we are interested in the low energy dynamics of the quantum field theory, it is convenient to decompose the field in those two modes: $\xi_{\alpha\beta...}=\hat{\xi}_{\alpha\beta..}+\Bar{\xi}_{\alpha\beta..}$ such that,
\begin{itemize}
    \item $\Bar{\xi}_{\alpha\beta..}$ are the light modes with energy scales $\epsilon$.
    \item  $\hat{\xi}_{\alpha\beta..}$ are the heavy modes with energy scale $\Lambda>>\epsilon$.
\end{itemize}
%The heavy modes $\hat{\xi}_{\alpha\beta..}$ can then be integrated out and  
The effective theory of the light modes $\Bar{\xi}_{\alpha\beta..}$ can then be determined by integrating out the heavy modes $\hat{\xi}_{\alpha\beta..}$. This can be done using a path integral in the following way,
\begin{align}
    \begin{split}
    \mathcal{Z}[\Bar{\xi}_{\alpha\beta..}]:=    e^{i\mathcal{S}_{\text{eff}}[\Bar{\xi}_{\alpha\beta..}]}=\rmint \mathcal{D}\hat{\xi}_{\alpha\beta..}\,e^{i\,\mathcal{S}^{\mathcal{O}(\hbar^0)}_{\text{eff}}[\Bar{\xi}_{\alpha\beta..},\hat{\xi}_{\alpha\beta..}]}\,,
        \end{split} \label{stotal1}
\end{align}
where $S^{\mathcal{O}(\hbar^0)}_{eff}$ is the classical part of the action mentioned in (\ref{Stotal}).
The path integral eventually gives the effective action for the light modes as follows,
\begin{align}
    \begin{split}
       \mathcal{S}_{\text{eff}}[\Bar{\xi}_{\alpha\beta..}]= \rmint d^4x\Big[\frac{1}{2}\partial_{\mu}\Bar{\xi}_{\alpha\beta..}\partial^{\mu}\Bar{\xi}^{\alpha\beta..}+
        \sum_{n}C_{n}\mathcal{O}_{n}(\Bar{\xi}_{\alpha\beta..})
\Big]\,.
    \end{split}\label{2.4m}
\end{align}
The coefficients $C_{n}$ are sometimes called Wilson coefficients, and they contain information about the UV sector of the theory. This approach is called the \textit{top-down} approach. One can also proceed in the opposite direction. In that case, one first writes down an effective action with the $C_{n}'s$ left unfixed, and then determines them by comparing with experimental data or with the computation from the \textit{top-down} approach. This is called the \textit{bottom-up} approach. We must therefore compute the effective action perturbatively in the small parameter $\frac{\epsilon}{\Lambda}$ and determine how many $\mathcal{O}_n$ contribute at a given order. This is fixed by the EFT power-counting rules, which tell us how the relevant quantities scale with $\frac{\epsilon}{\Lambda}$. This is discussed systematically in Sec.~(\ref{Sec4}) for our case. \par 
\par
 \begin{figure}
     \centering
    \thesisdiagrampanel{\scalebox{0.36}{\begin{feynman}
    \fermion[lineWidth=4, showArrow=false]{5.40, 4.00}{5.80, 5.20}
    \fermion[showArrow=false, lineWidth=4]{5.80, 5.20}{7.00, 4.00}
    \fermion[lineWidth=4, showArrow=false]{11.40, 5.80}{11.40, 6.40}
    \electroweak[flip=true, lineWidth=4]{10.80, 5.20}{11.40, 5.80}
    \fermion[lineWidth=4]{9.80, 6.40}{13.20, 6.40}
    \fermion[showArrow=false, lineWidth=4, label=$\mathcal{O}(\hbar^0)$]{5.80, 5.20}{5.80, 6.40}
    \fermion[lineWidth=4]{4.00, 4.00}{7.60, 4.00}
    \fermion[showArrow=false, lineWidth=4]{4.60, 4.00}{5.80, 5.20}
    \fermion[lineWidth=4, showArrow=false]{6.20, 4.00}{5.80, 5.20}
    \fermion[lineWidth=4]{9.80, 4.00}{13.20, 4.00}
    \electroweak[lineWidth=4]{11.40, 4.80}{10.80, 5.20}
    \fermion[lineWidth=4]{4.00, 6.40}{7.60, 6.40}
    \electroweak[label=$\mathcal{O}(\hbar)$, lineWidth=4]{11.40, 5.80}{12.00, 5.20}
    \electroweak[flip=true, lineWidth=4]{11.40, 4.80}{12.00, 5.20}
    \fermion[showArrow=false, lineWidth=4]{11.40, 4.00}{11.40, 4.80}
\end{feynman}
}}
     \caption{Tree level diagram (left) vs. loop level (one loop) (right) diagram in WEFT}
     \label{mfig}
 \end{figure}
At this point, we are not interested in the quantum effective action, since the quantum corrections are highly subleading. Hence, we neglect the $\mathcal{O}(\hbar^n)$ terms on the right-hand side of (\ref{stotal1}) and focus only on the tree-level diagrams, i.e. the $\mathcal{O}(\hbar^0)$ diagrams shown in Fig.~(\ref{mfig}). Furthermore, we expand the classical effective action in powers of $(v/c)$. This expansion is known as the Post-Newtonian (PN) expansion, and the $n$-PN order scales as $\mathcal{O}(\frac{v}{c})^{2n}$. It is worth noting that, in order to compute the tree-level, i.e. $\mathcal{O}(\hbar^0)$, PN diagrams, one still needs to evaluate several Feynman loop integrals. These tree-level PN diagrams are sometimes called ``Classical Loops.'' For more details, interested readers are referred to \cite{Burgess:2020tbq,Porto:2016pyg}. As explained before, EFT approaches are especially useful when there is a clear separation of scales, which is precisely the case for the binary inspiral problem. The relevant length scales are:
\begin{itemize}
    \item The size of the astrophysical object $R\sim G\,M$.
    \item The orbital radius r.
    \item The wavelength of radiation: $\lambda$.
\end{itemize}
These three scales are not independent; rather, they are related through the relative velocity of the binaries, with $\frac{R}{r}\sim v^2$. Keeping in mind the sensitivity of gravitational-wave detectors, one may roughly assume $\frac{1}{\lambda}\sim \omega_{GW}\sim \frac{v}{r}=\Omega_{orbit}$. These relations show that the different orders in the PN expansion probe different physical length scales, with the hierarchy $R\ll r \ll\lambda$. Hence, in EFT language, the heavy and light field modes are decoupled, and their dynamics can be treated separately.  %In Sec.~(\ref{Sec3}) we discuss how one can construct an Effective Field theory of a binary system. 
\par
\textcolor{black}{In the Worldline Effective Field Theory (WEFT) approach used to study the binary inspiral problem in this paper, apart from field degrees of freedom we do have worldline degrees of freedom $\{x^{\mu}_{a}\}$. We treat the worldline degrees of freedom as non-dynamical DOFs. Then integrating out the heavy modes, we get the total effective action: $\mathcal{S}_{\text{eff}}^{\text{tot}}[\boldsymbol{x}_a,\Bar{\xi}]$ which has the following form,
\begin{align}
    \begin{split}
        \mathcal{S}_{\text{eff}}^{\text{tot}}[\boldsymbol{x}_a,\Bar{\xi}]=-i\,\log\,\mathcal{Z}[\boldsymbol{x}_a,\Bar{\xi}]=\underbrace{\mathcal{S}_{\text{eff}}^{\text{cons}}[\boldsymbol{x}_a]}_{\text{real}}+\underbrace{\mathcal{S}_{\text{eff}}^{\text{rad}}[\boldsymbol{x}_a,\bar\xi]}_{\text{complex}}.
    \end{split}
\end{align}
Extremizing $\mathcal{S}_{\text{eff}}^{\text{cons}}[\boldsymbol{x}_a]$ we get the conservative dynamics of the binary system:
\begin{align}
    \begin{split}
     &    \frac{\delta}{\delta \boldsymbol{x}_a(t)}\,\mathcal{S}_{\text{eff}}^{\text{cons}}[\boldsymbol{x}_a]\Big |_{\mathcal{O}(\hbar^0)}=0 \implies \text{Orbit equation of the binary system.}
    \end{split}
\end{align}
and, $ \mathcal{S}_{\text{eff}}^{\text{rad}}[\boldsymbol{x}_a,\bar\xi]$ is in general complex and related to the radiated power from the system via optical theorem as,
\begin{align}
    \begin{split}
        \,\text{Im}[\gamma_{\text{eff}}[\boldsymbol{x}_a]]\equiv \text{Im}[-i\log\rmint \mathcal{D}\bar\xi \,e^{i\mathcal{S}_{\text{eff}}^{\text{rad}}[\boldsymbol{x}_a,\bar\xi]}]=\frac{\mathcal{T}}{2}\rmint dE\,d\Omega \frac{d^2\Gamma}{dEd\Omega},\,\text{with radiated power},\, dP=E\,d\Gamma.
    \end{split}
\end{align}}

The hierarchy of EFT for the binary inspiral problem is shown in Fig.~(\ref{newfig}).  
\begin{figure}
    \centering
    \includegraphics[scale=0.20]{eft_main.jpg}
    \caption{Different scales in the binary inspiral problem. We reproduce the figure from \cite{Porto:2016pyg}.}
    \label{newfig}
\end{figure}
\section{Review of worldline effective field theory for a point particle in GR}\label{Sec3}
In this section, we briefly discuss the action of an inspiralling non-spinning black hole binary in Einstein theory. The background metric is chosen to be in the Kaluza-Klein form, consisting of a non-relativistic gravitational field. %\textcolor{red}{ Point particle from the start ? We first write down the point particle action along with the finite size effect terms in this background and expand it perturbatively in field strength.} From this, we extract the worldline coupling of the electromagnetic field and the axion field.

%The action for a relativistic point particle in Einstein's gravity is given by, 
We start with the following action, 
\begin{equation}
    S=S_{\text{EH}}+S_{\text{pp}} \label{2.1}
\end{equation}
where the Einstein-Hilbert (EH) action is given by,
\begin{eqnarray}
    S_{EH}=\frac{1}{16\pi G}\rmint d^4x \sqrt{-g}\,R
\end{eqnarray}
and describes the dynamics of the graviton. $S_{\text{pp}}$ models the binary (determines the dynamics of the two-body system (black holes or any other compact objects)) and takes the following form, 
\begin{eqnarray} \label{eqnew}
    S_{pp}=\sum_{a=1}^2 m_a \rmint d\tau_{a}+\sum_{a=1}^2 C^{\text{R}}_a \rmint d\tau_{a} R({x_a})+\sum_{a=1}^2 C_{a}^{\text{V}}\rmint d\tau_a R_{\mu\nu}(x_a)\,\dot{x}_a^\mu \,\dot{x}_\nu^a+.....\,.
\end{eqnarray}

Here, $\tau_a$ denotes the proper time along the particle's world line and  $C^R$, $C^V$ are the Wilson coefficients. In this paper, we will only consider non-spinning binaries, hence we can ignore the finite size effects. Henceforth, we will not consider the proportional to $C^R$, $C^V$ from (\ref{eqnew}).  \par
A gravitational wave detector such as LIGO measures the gravitational waveform. The first step towards computing such waveforms is to calculate the radiated energy flux from the binary system. Given the action in (\ref{2.1}) we can calculate this quantity by considering the perturbation around a flat metric, i.e., $g_{\mu\nu}=\eta_{\mu\nu}+h_{\mu\nu}$. Integrating out the graviton field, one can obtain the effective action for the point particles,
\begin{eqnarray}
    e^{i\,\mathcal{S}_{\text{eff}}}=\rmint \mathcal{D}h_{\mu\nu} \exp{[iS_{EH}+iS_{pp}]}\,.
\end{eqnarray}\label{2.4}
The real part of the effective action describes the dynamics of the two-body system, and the imaginary part measures the graviton emission from the binary system over a large time. In principle, one can directly evaluate the full relativistic path integral using the Feynman rules for the Einstein-Hilbert action, which involves tree-level as well as loop-level diagrams. However, after calculating the amplitudes, it is not easy to take the non-relativistic limit. In order to overcome the problem, following \cite{Kol:2007bc}, one can decompose the metric in terms of the non-relativistic gravitational fields and then perform the calculations. We will only focus on the tree-level diagrams relevant to classical results. The methodology is discussed in detail in the subsequent sections. \par

\subsubsection*{Non relativistic general relativity}
In this section, we will briefly discuss how one can arrive at the effective action in a non-relativistic regime \cite{Goldberger:2004jt,Levi:2018nxp,Porto:2016pyg}. We expand the point particle action in powers of the particle's three-velocities.
\begin{eqnarray}
    S_{pp}=\sum_{a=1}^2\Big[\frac{1}{2}v_a^2-\frac{1}{2}h_{00}-h_{0i}\,v^a_{i}-\frac{1}{2}h_{00}v_{a}^2-\frac{1}{2}h_{ij}v^{i}_a\,v^{j}_a+\frac{1}{8}v_a^4+.......\Big]\,.
\end{eqnarray}
Here we should remember that all components of the metric perturbations $h_{\mu\nu}$ are evaluated at the particular particle world line.\par 
However, at this point the propagator for $h_{\mu\nu}$ is fully relativistic. Hence we can not distinguish between the potential gravitons (heavy modes) and the long-wavelength radiation gravitons (light modes). We are not interested in potential gravitons as the gravitational wave detector can not detect them. So due to our interest in finding out the effective action in the non-relativistic regime, we integrate over the potential gravitons. Now, to connect with the discussion in Sec.~(\ref{Sec2}) about integrating out the heavy modes we decompose the metric perturbation as follows,
\begin{eqnarray}
    h_{\mu\nu}=\Bar{h}_{\mu\nu}+
\mathcal{H}_{\mu\nu}\,,\label{3.6n}
\end{eqnarray}
where $\mathcal{H}_{\mu\nu}$ is the potential gravitons with the properties $\partial_{i}\mathcal{H}_{\mu\nu}\sim\frac{\mathcal{H}_{\mu\nu}}{r},\partial_{0}\mathcal{H}_{\mu\nu}\sim{\mathcal{H}_{\mu\nu}}\frac{v}{r}$. And $\Bar{h}_{\mu\nu}$ is the radiation gravitons with $\partial_{\beta}\bar{h}_{\mu\nu}\sim{\bar{h}_{\mu\nu}}\frac{v}{r}$. It is more convenient to express $\mathcal{H}_{\mu\nu}$ as,
\begin{eqnarray}
    \mathcal{H}_{\mu\nu}(x^\mu)=\rmint\frac{d^3\boldsymbol{k}}{(2\pi)^3}\,e^{i\boldsymbol{k}.\boldsymbol{x}}\,\mathcal{H}^{k}_{\mu\nu}(x^0)\,.
\end{eqnarray}
Now, the conservative (or radiative) effective action for the binaries $S_{\text{NRGR}}$ can be computed by integrating over the non-relativistic potential gravitons,
\begin{eqnarray}
    e^{i\mathcal{S}_{\text{eff}}[x_a,\Bar{h}_{\mu\nu}]}=\rmint \mathcal{D}\mathcal{H}_{\mu\nu}\,\exp{[iS[\Bar{h}+\mathcal{H},x_a]+iS_{GF}]}\label{2.8}
\end{eqnarray}
where, $S_{GF}$ is the gauge fixing term. While evaluating the path integral, ghost terms arise in the effective action, but at the classical level, they do not contribute. Now, to preserve the gauge invariance, we choose the gauge fixing term as \cite{Goldberger:2004jt},
\begin{eqnarray}
    S_{GF}=\rmint d^4x\sqrt{\Bar{g}}\,\Gamma_\mu \Gamma^\mu
\end{eqnarray}
which is invariant under the general coordinate transformation of the background metric $\Bar{g}_{\mu\nu}=\eta_{\mu\nu}+\Bar{h}_{\mu\nu}$ and $\Gamma_\mu=\nabla_{\alpha}\mathcal{H}^{\alpha}_{\mu}-\frac{1}{2}\nabla_{\mu}\mathcal{H}_{\alpha}^{\alpha}.$ From the EH Lagrangian (up to $\mathcal{O}(\mathcal{H}^2)$) in momentum space, one can write down the \textcolor{black}{orbital} graviton propagator as \cite{Goldberger:2004jt},
\begin{eqnarray}
    \Big\langle\mathcal{H}^{k}_{\mu\nu}(x^0)\mathcal{H}^{q}_{\alpha\beta}(0)\Big\rangle=-(2\pi)^3 \delta^{(3)}(\boldsymbol{k}+\boldsymbol{q})\delta(x^0)\frac{i}{\boldsymbol{k}^2}P_{\mu\nu,\alpha\beta}\label{2.10}
\end{eqnarray}
with, $P_{\mu\nu,\alpha\beta}=\frac{1}{2}[\eta_{\mu\alpha}\eta_{\nu\beta}+\eta_{\mu\beta}\eta_{\nu\alpha}-\frac{2}{d-2}\eta_{\mu\nu}\eta_{\alpha\beta}]$. So in order to find $S_{NRGR}$ one have to compute the functional integral in (\ref{2.8}) using the (\ref{2.10}). To do this, one has to sum over the relevant Feynman diagram, which has the following topological properties \cite{Goldberger:2022ebt, Goldberger:2004jt},
\begin{itemize}
    \item All diagrams must be connected with stripped-off particle worldlines.
    \item Diagram can contain internal graviton potential modes ($\mathcal{H}_{\mu\nu}$) but can not have external potential mode.
    \item Radiative graviton modes ($\Bar{h}_{\mu\nu}$) are only appear in the external lines.
\end{itemize}
As there is a gauge redundancy, it is better to choose a particular gauge from the beginning and use the proper decomposition of the non-relativistic graviton field \cite{Kol:2007bc}. This helps us to identify the correct diagrams at each PN order and avoid any confusion arising due to the mixing of various components of $h_{\mu\nu}$. Hence, following \cite{Kol:2007bc}, we choose the standard Kaluza-Klein (KK) decomposition,
\begin{align}
    \begin{split}
        ds^2=g_{\mu\nu}dx^{\mu}dx^{\nu}=-e^{2\psi}(dt-\tilde{\mathcal{A}_i}dx^{i})^2+e^{-2{\psi}}\gamma_{ij}dx^{i}dx^{j}\,.
    \end{split} \label{KK}
\end{align}
In the component form:
$$g_{\mu\nu}=
\begin{pmatrix}
-e^{2\psi} & e^{2\psi} \tilde{\mathcal{A}}_{j}\\
e^{2\psi} \tilde{\mathcal{A}}_{i} & \,\,e^{-2\psi}\gamma_{ij}-e^{2\psi}\tilde{\mathcal{A}}_{i}\tilde{\mathcal{A}}_{j}\\
\end{pmatrix}$$
where $\psi,\tilde{\mathcal{A}_i},\gamma_{ij}=\delta_{ij}+\sigma_{ij}$ are the set of \textit{non-relativistic graviton field} (NRG), sometimes called as Kol-Smolkin variables, are Newtonian potential, gravitomagnetic potential, and metric 3-tensor, respectively \cite{Kol:2007bc}.
For our future computations to make all NRG fields have mass-dimension one, one can rescale the fields as follows,
\begin{align}
    \begin{split}
        \psi\rightarrow\frac{\psi}{m_p},\,\,\tilde{\mathcal{A}}_{i}\rightarrow\frac{\tilde{\mathcal{A}}_{i}}{m_p},\,\,\sigma_{ij}\rightarrow \frac{\sigma_{ij}}{m_p}.
    \end{split}
\end{align}
\section{Conservative dynamics of binary black holes in a theory of ALP and ULV }\label{Sec4}
In this section, we describe the action for inspiralling binary black holes in a theory of \textit{Axion-like particles and Ultra-light vectors}. Apart from the usual EH term, it contains a Maxwell field, a Proca field, and a massive scalar field.  We have the following action for this theory \cite{Cardoso:2018tly}:
\begin{align}
\begin{split}
S_{\text{field}}=\rmint d^4x\sqrt{-g}&\Big(\frac{m_{p}^2}{2}R-\frac{1}{4}F_{\mu\nu}F^{\mu\nu}-\frac{1}{2}\partial_\mu\phi\partial^\mu\phi-\frac{1}{2}m^2\phi^2+\frac{g_{a\gamma\gamma}}{4\,m_{p}}\phi F_{\mu\nu}^*F^{\mu\nu}-\frac{1}{4}B_{\mu\nu}B^{\mu\nu}\\ &
+\frac{\gamma}{2}F_{\mu\nu}B^{\mu\nu}
-\frac{1}{2}\mu_{\gamma}^2B_{\mu}B^{\mu}\Big)\,.
\label{2.11}
\end{split}
\end{align}
The action in (\ref{2.11}) describes the Axion-like-particles (ALPs) (scalar), massless photon and massive (often known as \textit{dark photon}) gauge field (Proca) minimally coupled with gravity. In (\ref{2.11}) $R,\, \phi $ are the Ricci scalar and ALP. Also, 
$$F_{\mu\nu}=\partial_{\mu} {A}_{\nu}-\partial_{\nu} {A}_{\mu},\, B_{\mu\nu}=\partial_{\mu} {B}_{\nu}-\partial_{\nu} {B}_{\mu},$$ where ${A}_{\mu}$ and ${B}_{\mu}$ are Electromagnetic and Proca field respectively. $\gamma$ is the coupling constant between photon and dark photon. Furthermore, $g_{a\gamma\gamma}$ is the coupling constant for axion-photon, and it has a negative mass dimension, implying the theory is not renormalizable, and the term $\phi F^{*} F$ is sometimes called the \textit{Theta term}. However, that is not a problem for our case as we finally intend to write down a classical effective action for the light modes in powers of $(v/c)$ as mentioned in Sec.~(\ref{Sec2}). A good EFT does not need to be renormalized. We are only concerned about the tree-level diagrams as we want classical results. Eventually, we will encounter loop integrals while computing the effective action, and we call them classical loops as they are proportional to $\hbar ^0$. \par

Next we compute the effective action at $\mathcal{O}{(\hbar^0)}$ and organize it as an expansion in $\mathcal{O}(\frac{v}{c})^n$. We therefore work in the non-relativistic regime, choose the background metric in KK form as mentioned in (\ref{KK}), and keep the non-trivial worldline couplings of the non-relativistic gravitational (NRG) fields. Our conservative analysis extends up to 1PN order in the gravitational, electromagnetic, and Proca sectors, and up to 2PN order in the scalar sector. We will also show that the leading correction from the CP-violating $g_{a\gamma\gamma}$ term in (\ref{2.11}) enters the conservative dynamics only at \textit{2.5PN}. \par

The calculation follows a uniform pattern throughout the chapter. We first identify the quadratic action and the relevant worldline operators, which determine the propagators and the vertices. We then use EFT power counting to decide which connected diagrams contribute at a given PN order. Whenever a new contraction pattern appears, we display one explicit derivation; for the remaining diagrams of the same type we quote the reduced amplitudes directly so that the physical structure of the result stays visible in the main text. The EFT of radiation fields, in KK parametrization, is obtained by decomposing fields into potential and radiation sectors as $\psi=\Tilde{\psi}+\Bar{\psi},\Tilde{\mathcal{A}}_i=\Hat{\mathcal{A}}_i+\Bar{\mathcal{A}}_i,\sigma_{ij}=\zeta_{ij}+\Bar{\sigma}_{ij}\,,\phi=\varphi+\Bar{\phi}, A_{\mu}=\mathcal{A}_{\mu}+\Bar{a}_{\mu}$ and $B_{\mu}=\mathcal{B}_{\mu}+\Bar{b}_{\mu}$ and then integrating over the potential modes. The partition function takes the form
\begin{align}
\begin{split}
  \mathcal{Z}[\boldsymbol{x}_a,\Bar{\psi},\Bar{{\mathcal{A}}}_i,\Bar{\sigma}_{ij},\Bar{a}_{\mu},\Bar{b}_{\mu}]  &:=e^{i\mathcal{S}_{\text{eff}}[\boldsymbol{x}_a,\Bar{\psi},\Bar{{\mathcal{A}}}_i,\Bar{\sigma}_{ij},\Bar{a}_{\mu},\Bar{b}_{\mu}]}\,,\\ &
  =\rmint \underbrace{\mathcal{D}\tilde{\psi} \,\mathcal{D}\hat{\mathcal{A}}_{i}\,\mathcal{D}\varphi \,\mathcal{D}\mathcal{A}_{\mu}\,\mathcal{D}\mathcal{B}_{\nu}\,
  \mathcal{D}\boldsymbol{\zeta}_{ij}\,}_{\text{integrating over potential modes}\,(\mathcal{D}\hat{\xi})}\,e^{iS_{\text{field}}+iS_{\text{pp}}}\,.\label{ch1:2.13}
\end{split}
\end{align}
We can compute the effective action in (\ref{ch1:2.13}) by evaluating the  amplitudes of the suitable Feynman diagrams by taking the logarithm as,
\begin{align}
    \begin{split}
        i\mathcal{S}_{\text{eff}}[\boldsymbol{x}_a,\Bar{\xi}]&=\log\Big[{\rmint\mathcal{D}\hat\xi\,e^{iS_{\text{field}}+iS_{\text{pp}}}}\Big]\,.\label{4.3mm}
    \end{split}
\end{align}
One can compute (\ref{4.3mm}) perturbatively by evaluating the connected Feynman diagrams. We need two kinds of vertices, worldline vertices and field vertices. To get the worldline vertices, we need to consider the point particle action, including all the fields.
Now we can write down the field action by, $S_{\text{field}}=S_{\text{quad}}+S_{\text{int}}$ and use,
\begin{align}
    \begin{split}
        \Big\langle\mathcal{O}_{1}(\boldsymbol{x}_1(t_1))...\mathcal{O}_{n}(\boldsymbol{x}_n(t_n))\Big\rangle=\rmint\mathcal{D}\hat\xi \,e^{iS_{\text{quad}}}\,\underbrace{\mathcal{O}_{1}(\boldsymbol{x}_1(t_1),\hat\xi)...\mathcal{O}_{n}(\boldsymbol{x}_n(t_n),\hat\xi)}_{\text{got expanding \,$e^{iS_{\text{pp}}+iS_{\text{int}}}$}}
    \end{split}
\end{align}
and the wick contraction to reduce the higher point correlation function into the two-point functions to compute the effective action.
\par
In order to evaluate the connected diagrams, one next needs the Feynman propagators (two-point functions). We therefore record the quadratic actions for the relevant field configurations and the resulting non-relativistic propagators. The purpose of the next formulas is mainly to collect the ingredients that repeatedly enter the later diagram evaluation, rather than to dwell on algebra that is identical from sector to sector.


%\vspace{0.3 cm}
\textbf{Propagators of different fields:} We list the two-point functions, i.e. the propagators of different fields in our theory. We also do the conventional non-relativistic expansion of the propagator, assuming $|k_0|\ll |\boldsymbol{k}|$. In this limit, we expand the relativistic Feynman propagator, which contains an infinite sum of the delta function and its time derivatives (even). We only consider the non-relativistic contribution and the first-order relativistic corrections.

\newpage
\underline{\textit{Scalar propagator:}}

 In KK parametrization the scalar action has the following form:
 \begin{align}
    \begin{split}
        S_{\phi}&=\rmint d^4x \sqrt{-g}\Big[-\frac{1}{2}\partial_{\mu}\phi\partial^{\mu}\phi-\frac{1}{2}m^2 \phi^2\Big]\,,\\ &
        =\rmint d^3x \,dt \sqrt{\gamma}\Big[\frac{e^{-4\psi}}{2}\partial_{0}\phi\partial_{0}\phi-2\Hat{\mathcal{A}}^{i}\partial_{i}\phi\partial_{0}\phi-\frac{1}{2}\gamma^{ij}\partial_{i}\phi\partial_{j}\phi-\frac{1}{2}e^{-2\psi}m^2\phi^2 \nonumber
        \Big] \,,\end{split}
\end{align}
\begin{align}
    \begin{split}
\,\,\,\,\,\,\,\,\,\,\,\,\,\,\,\,\,\,\,\,\,\,\,\,\,\,\,=\rmint d^3x dt \,\Big(1+\frac{\text{Tr}(\sigma)}{2}\Big)&\Big[\frac{1}{2}\partial_{0}\phi\partial_{0}\phi-2\psi \partial_{0}\phi\partial_{0}\phi-2\Hat{\mathcal{A}}^{i}\partial_{i}\phi\partial_{0}\phi-\frac{1}{2}\delta^{ij}\partial_{i}\phi\partial_{j}\phi \\&+ \frac{1}{2}\sigma^{ij}\partial_{i}\phi\partial_{j}\phi-\frac{1}{2}(1-2\psi)m^2\phi^2\Big]\,.\label{24m}
    \end{split}
\end{align}
From the action (\ref{24m}) one can identify the scalar propagator as follows,
 \begin{align}
    \begin{split}
\Big\langle\varphi(x_1)\varphi(x_2)\Big\rangle &=\rmint \frac{d^4 k}{(2\pi)^4} \frac{e^{ik\cdot(x_1-x_2)}}{-k^2-m^2}\\&=\rmint \frac{dk_0}{2\pi}\,e^{-ik_0 (t_1-t_2)}\rmint \frac{d^3k}{(2\pi)^3}\frac{e^{i\boldsymbol{k}\cdot(\boldsymbol{x}_1-\boldsymbol{x}_2)}}{\boldsymbol{k}^2+m^2}\Big[1+\frac{k_0^2}{\boldsymbol{k}^2+m^2}+...\Big]
        \\ &
        =\delta(t_1-t_2)\rmint \frac{d^3k}{(2\pi)^3}\frac{e^{i\boldsymbol{k}\cdot\boldsymbol{r}}}{\boldsymbol{k}^2+m^2}+\partial_{t_1}\partial_{t_2}\delta(t_1-t_2)\rmint \frac{d^3k}{(2\pi)^3}\frac{e^{i\boldsymbol{k}\cdot\boldsymbol{r}}}{(\boldsymbol{k}^2+m^2)^2}+...\,. 
    \end{split}
\end{align}
\underline{\textit{ Electromagnetic propagator:}}

The action for free Maxwell theory is given by,
\begin{align}
    \begin{split}
        S_{\text{EM}}=-\frac{1}{4}\rmint d^3x dt \sqrt{\gamma}&\Big[-F_{0i}F_{0i}+e^{2\psi}\gamma^{ij}\gamma^{kl}F_{ik}F_{jl}+e^{2\psi}\Hat{\mathcal{A}}^{i}\Hat{\mathcal{A}}^{j}F_{0j}F_{i0}+e^{2\psi}\Hat{\mathcal{A}}^{i}\gamma^{jk}F_{0j}F_{ik}\Big]\\ &
        +(\text{higher order in fields})\,.\label{25m}
    \end{split}
\end{align}
From the quadratic part of the action (\ref{25m}), one can identify the propagator of the electromagnetic field as follows,
\begin{align}
    \begin{split}
&\Big\langle\mathcal{A}_{0}(x_1)\mathcal{A}_0(x_2)\Big\rangle=\textcolor{black}{-}\Big[\delta(t_1-t_2)\rmint \frac{d^3k}{(2\pi)^3}\frac{e^{i\boldsymbol{k}\cdot\boldsymbol{r}}}{\boldsymbol{k}^2}+\partial_{t_1}\partial_{t_2}\delta(t_1-t_2)\rmint \frac{d^3k}{(2\pi)^3}\frac{e^{i\boldsymbol{k}\cdot\boldsymbol{r}}}{\boldsymbol{k}^4}...\Big]\,,\\ &
        \Big\langle\mathcal{A}_{i}(x_1)\mathcal{A}_k(x_2)\Big\rangle=\Big[\delta(t_1-t_2)\rmint \frac{d^3k}{(2\pi)^3}\frac{e^{i\boldsymbol{k}\cdot\boldsymbol{r}}}{\boldsymbol{k}^2}+\partial_{t_1}\partial_{t_2}\delta(t_1-t_2)\rmint \frac{d^3k}{(2\pi)^3}\frac{e^{i\boldsymbol{k}\cdot\boldsymbol{r}}}{\boldsymbol{k}^4}...\Big]\delta_{ik}\,.\\ &
    \end{split}
\end{align}
\underline{\textit{ Proca propagator:}}

The action for the Proca field is the same as (\ref{25m}) with a mass term, hence the propagator has the following form:
\begin{align}
    \begin{split}
    &\Big\langle\mathcal{B}_{0}(x_1)\mathcal{B}_0(x_2)\Big\rangle=\textcolor{black}{-}\Big[\delta(t_1-t_2)\rmint \frac{d^3k}{(2\pi)^3}\frac{e^{i\boldsymbol{k}\cdot\boldsymbol{r}}}{\boldsymbol{k}^2+\mu_{\gamma}^2}+\partial_{t_1}\partial_{t_2}\delta(t_1-t_2)\rmint \frac{d^3k}{(2\pi)^3}\frac{e^{i\boldsymbol{k}\cdot\boldsymbol{r}}}{(\boldsymbol{k}^2+\mu_{\gamma}^2)^2}+.....\Big]\,.\\ &
\Big\langle\mathcal{B}_{i}(x_1)\mathcal{B}_k(x_2)\Big\rangle=\Big[\delta(t_1-t_2)\rmint \frac{d^3k}{(2\pi)^3}\frac{e^{i\boldsymbol{k}\cdot\boldsymbol{r}}}{\boldsymbol{k}^2+\mu_{\gamma}^2}+\partial_{t_1}\partial_{t_2}\delta(t_1-t_2)\rmint \frac{d^3k}{(2\pi)^3}\frac{e^{i\boldsymbol{k}\cdot\boldsymbol{r}}}{(\boldsymbol{k}^2+\mu_{\gamma}^2)^2}...\Big]\delta_{ik}\,.\label{2.28jjj}
    \end{split}
\end{align}
The propagator given in \eqref{2.28jjj} is not the full expression; in general, it also contains a correction term of the form \(\frac{k^\mu k^\nu}{\mu_{\gamma}^2}\). However, this additional piece contributes only at subleading order to both the potential and the waveform. Since, for the massive modes, our interest is restricted to the leading contribution, we neglect these extra terms.\\\\
\underline{\textit{Graviton propagators:}}

Lastly, the gravitational action takes the following form in KK parametrization:
\begin{align}
    \begin{split}
        S_{\text{G}}=\frac{m_{p}^2}{2}\rmint d^3x dt \,\sqrt{\gamma}&\Big[R[\gamma]-2\gamma^{ij}\partial_{i}\psi\partial_{j}\psi+\frac{e^{4\psi}}{4}\Hat{F}_{ij}\Hat{F}_{kl}\gamma^{ik}\gamma^{jl}+\frac{e^{-4\psi}}{4}[\Dot{\gamma}_{ij}\Dot{\gamma}_{kl}\gamma^{ik}\gamma^{jl}-(\gamma^{ij}\Dot{\gamma}_{ij})^2]\\ &
-4\gamma^{ij}\Dot{\Hat{\mathcal{A}}}\,\partial_{j}\psi
        +e^{-4\psi}(2\Dot{\psi}\gamma^{ij}\Dot{\gamma}_{ij}-6\Dot{\psi}^2)\Big]\,.\label{23m}
    \end{split}
\end{align}
For the action in (\ref{23m}) one can identify the static graviton propagator and its relativistic time corrections as:
\begin{align}
  \begin{split}
&\Big\langle\tilde{\psi}(x_1)\tilde{\psi}(x_2)\Big\rangle=\frac{1}{2}\Big[\delta(t_1-t_2)\rmint \frac{d^3k}{(2\pi)^3}\frac{e^{i\boldsymbol{k}\cdot\boldsymbol{r}}}{\boldsymbol{k}^2}+\partial_{t_1}\partial_{t_2}\delta(t_1-t_2)\rmint \frac{d^3k}{(2\pi)^3}\frac{e^{i\boldsymbol{k}\cdot\boldsymbol{r}}}{\boldsymbol{k}^4}...\Big]\,,\\ &
        \Big\langle\Hat{\mathcal{A}}_{i}(x_1)\Hat{\mathcal{A}}_k(x_2)\Big\rangle=-2\Big[\delta(t_1-t_2)\rmint \frac{d^3k}{(2\pi)^3}\frac{e^{i\boldsymbol{k}\cdot\boldsymbol{r}}}{\boldsymbol{k}^2}+\partial_{t_1}\partial_{t_2}\delta(t_1-t_2)\rmint \frac{d^3k}{(2\pi)^3}\frac{e^{i\boldsymbol{k}\cdot\boldsymbol{r}}}{\boldsymbol{k}^4}...\Big]\delta_{ik}\,,\\ &
        \Big\langle\boldsymbol{\zeta}_{ij}(x_1)\boldsymbol{\zeta}_{kl}(x_2)\Big\rangle=4\Big[\delta(t_1-t_2)\rmint \frac{d^3k}{(2\pi)^3}\frac{e^{i\boldsymbol{k}\cdot\boldsymbol{r}}}{\boldsymbol{k}^2}+\partial_{t_1}\partial_{t_2}\delta(t_1-t_2)\rmint \frac{d^3k}{(2\pi)^3}\frac{e^{i\boldsymbol{k}\cdot\boldsymbol{r}}}{\boldsymbol{k}^4}...\Big]P_{ij,kl} \,.
    \end{split}
\end{align}
We denote the relativistic time correction to the non-relativistic propagator as $\Big\langle....\Big\rangle_{\varoslash}$.

\textbf{Point Particle action:} In our analysis, we model the binaries as charged point particles, ignoring any internal structure of the astrophysical objects under consideration. Our first task is to identify the worldline operators. For that, we first write down the point particle worldline action. As we have a massive scalar field in the theory, we expect the objects' masses to depend on the scalar field \cite{Dyadina:2018ryl}. We also assume that the binary possesses electromagnetic and Proca charges (dark charges). So we include the coupling between the charges (electromagnetic (EM) and Proca) with the relevant fields. Here, we ignore the finite-size effects terms. %which appear at the \textcolor{red}{$\mathcal{O}(\partial^4g)$.} %These finite-size effects include Gravitoelectric and Gravitomagnetic tensors. \textcolor{red}{ However, we do not consider finite-size effects because of the point particle assumption, considering the effective length scale of the problem.}
The action is therefore given (at leading order in gauge fields) by,
\begin{align}
    \begin{split}
        S_{pp}=&-\sum_{a}\rmint m_a(\phi)d\tau_{a}+\sum_a\rmint\,Q_a \,d\tau_a \textcolor{black}{A^\mu}\,v_\mu+\sum_a\rmint\,Q_a'\, d\tau_a B^\mu\,v_\mu \,,\\ &
        =\sum_{a}\, m_a\rmint dt\Big[-1+\frac{1}{2}v_a^2-\frac{\psi}{m_{p}}-\frac{\hat{\mathcal{A}}_{i}}{m_{p}}\,v^i_{a}-\frac{\psi^2}{2m_{p}^2}-\frac{3}{2}\frac{\psi}{m_{p}} v_{a}^2\textcolor{black}{+}\frac{1}{2}\frac{\sigma_{ij}}{m_{p}}v^{i}_a\,v^{j}_a+\frac{1}{8}v_a^4+\mathcal{O}(v^6)\Big]\Tilde{\mathcal{L}} \\ &
        +\sum_a \rmint dt \Big[Q_a\,A_{0}(1-\frac{\psi}{m_{p}}-\frac{\psi^2}{2m_{p}^2}) +Q_a\,v^{i}_a A_{i}(1-\frac{\psi}{m_{p}}-\frac{\psi^2}{2m_p^2})+\cdots\Big]\\ &
        \sum_a \rmint dt \Big[Q_a'\,B_{0}(1-\frac{\psi}{m_{p}}-\frac{\psi^2}{2m_{p}^2}) +Q_a'\,v^{i}_a B_{i}(1-\frac{\psi}{m_{p}}-\frac{\psi^2}{2m_p^2})+\cdots\Big]\,, \label{4.10}
    \end{split}
 \end{align}
 where, $\Tilde{\mathcal{L}}=\Big[1+s_a\,\frac{\phi}{m_{p}}+g_a\,\frac{\phi^2}{m_{p}^2}\Big].$ Also, $a=1,2$ denotes the two objects in the binary system. We have expanded $m_a(\phi)$ in a Taylor series up to second order, with $s_a=\frac{1}{m_a}\frac{\partial m_a(\phi)}{\partial \phi}|_{\phi=0}$ and $g_a=\frac{1}{m_a}\frac{\partial^2 m_a(\phi)}{\partial \phi^2}|_{\phi=0} $. \textit{ $Q_1'$ and $Q_2'$ denote two dark charges. They are not associated with an electric or magnetic field. They are sometimes also referred to as `hidden' charges because they do not interact with ordinary matter in the same way that electric and magnetic charges do.}
\footnote{Terms with time derivatives in fields will contribute as corrections to the non-relativistic propagator.  One can show that taking the 2-point self-interacting temporal vertices produces the same corrections to the non-relativistic propagator. 
}



\begin{table}[htb!]
\centering
\scalebox{0.70}{
\setlength{\arrayrulewidth}{0.3mm}
\setlength{\tabcolsep}{20pt}
\renewcommand{\arraystretch}{2.2}
\begin{tabular}{|p{3cm}|p{3cm}|p{3cm}|p{3cm}|p{3cm}|}
\hline
\multicolumn{4}{|c|}{EFT PN COUNTING} \\
\hline
Fields & Worldline-Coupling & Diagram &order \\
\hline
$\hat{\mathcal{A}}_i$& -$\frac{m_a}{m_{p}} \rmint dt\,\Hat{\mathcal{A}}^i\,v_i^a$ & \scalebox{0.3}{\begin{feynman}
    \gluon[lineWidth=4,color=fcb900,label=$\Hat{\mathcal{A}}^i$]{4.00, 5.20}{5.80, 5.20}
    \fermion[lineWidth=4, showArrow=true, flip=false]{4.00, 4.00}{4.00, 6.40}
   
\end{feynman}}& $\sim v$\\
\hline
$ \mathcal{A}_i$& $Q\rmint dt \mathcal{A}_i v^i$& \scalebox{0.3}{\begin{feynman}
    \fermion[lineWidth=4, showArrow=false]{4.00, 4.00}{4.00, 6.20}
    \electroweak[color=9900ef, lineWidth=4, label=$\mathcal{A}_i$]{4.00, 5.00}{5.40, 5.00}
\end{feynman}}& $\sim v$\\
\hline
$ \mathcal{B}_i$& $Q'\rmint dt\, \mathcal{B}_i \,v^i$& \scalebox{0.2}{\begin{feynman}
    \gluon[lineWidth=4, flip=true, endcaps=false, color=eb144c, label=$\mathcal{B}_i$]{4.00, 5.20}{5.60, 5.20}
    \fermion[lineWidth=4, showArrow=true, flip=false]{4.00, 4.00}{4.00, 6.40}
    
\end{feynman}}& $\sim v$\\
\hline
$ \mathcal{A}_0$& $Q\rmint dt\,  \mathcal{A}_0$& \scalebox{0.3}{\begin{feynman}
    \dashed[lineWidth=4, showArrow=false, label=$\mathcal{A}_0$,flip=true]{4.00, 5.00}{6.00, 5.00}
    \fermion[lineWidth=4]{4.00, 4.00}{4.00, 6.00}
    
\end{feynman}}& $\sim v^0$\\
\hline
$\tilde{\psi}$& $\frac{-m_a}{m_{p}}\rmint dt\, \,\tilde{\psi}$& \scalebox{0.3}{\begin{feynman}
    \electroweak[lineWidth=4, label=$\boldsymbol{\tilde{\psi}}$]{4.00, 5.20}{6.00, 5.20}
    \fermion[lineWidth=4]{4.00, 4.00}{4.00, 6.20}
\end{feynman}
}& $\sim v^0$\\
\hline
$ \mathcal{B}_0$& $Q'\,\rmint dt\,  \mathcal{B}_0$ & \scalebox{0.3}{\begin{feynman}
    \dashed[lineWidth=4, showArrow=false, color=9900ef, label=$\mathcal{B}_0$]{4.00, 5.20}{5.40, 5.20}
    \fermion[lineWidth=4, showArrow=true, flip=false]{4.00, 4.00}{4.00, 6.40}
   
\end{feynman}}& $\sim v^0$\\
\hline
$\varphi$& -$s_a\, \frac{m_a}{m_{p}} \rmint dt\, \varphi $ & \scalebox{0.3}{\begin{feynman}
    \electroweak[lineWidth=4, color=0693e3, label=$\boldsymbol{\varphi}$, flip=true]{4.00, 5.00}{5.40, 5.00}
    \fermion[lineWidth=4]{4.00, 4.00}{4.00, 6.20}
\end{feynman}}& $\sim v^0$\\
\hline
\end{tabular}
}
\caption{Table showing the different worldline couplings and the order at which they contribute.}
\label{table1m}
\end{table}



\textbf{EFT counting for different fields:}
As discussed in Sec.~(\ref{Sec2}), we write the effective action for the light modes up to a given PN order. This leads directly to EFT power counting. Hence, we need to know how the different field couplings to the worldline scale with $v$. We use Table~(\ref{table1m}) to determine how the different terms in the conservative effective action scale with the velocities of the black holes. %The power counting or the scaling of fields with different parameters have implications in calculating radiation of different sectors at a given post-newtonian order. we have non-relativistic gravitational fields $\tilde{\psi}, \hat{\mathcal{A}_i}, \sigma $. These fields when couple to wordlines gives the following sort of couplings ,\par
The contribution of effective action corresponding to every diagram has the following form $Lv^{2n}$, which we call the $n PN$ diagram.
\subsection{The Gravitational bound sector}
With the propagators and power counting rules in hand, the remaining task is to identify the connected diagrams that contribute to the real part of the effective action. That real part determines the corrected orbit. In what follows, we therefore organize the calculation sector by sector. For readability, we work out one representative Wick contraction whenever a new structure appears and then collect the remaining diagrams in their reduced form. We begin with the gravitational bound sector. Its 1PN result is standard \cite{Kuntz:2019zef}, but it also provides the template for the electromagnetic and Proca sectors discussed next. Figure~(\ref{fig:my_labela}) shows the four graphs that contribute at this order.\footnote{To do the integrals, we take the help of the master integrals given in Appendix~(\ref{ch1:app:E}).}
\begin{figure}
    \centering
    \begin{subfigure}[t]{0.22\textwidth}
        \centering
\thesisdiagrampanel{\scalebox{0.28}{\begin{feynman}
    \electroweak[lineWidth=6, label=$\tilde{\psi}$]{5.60, 4.00}{5.60, 5.00}
    \fermion[lineWidth=6]{4.00, 6.40}{7.20, 6.40}
    \electroweak[label=$\tilde{\psi}$, lineWidth=6]{5.60, 5.40}{5.60, 6.40}
    \fermion[lineWidth=6]{4.00, 4.00}{7.20, 4.00}
    \parton{5.60,5.20}{0.20}
\end{feynman}
}}
        \caption{}
        \label{fig1a}
    \end{subfigure}
    \begin{subfigure}[t]{0.22\textwidth}
        \centering
\thesisdiagrampanel{\scalebox{0.36}{\begin{feynman}
    \electroweak[label=$\Tilde{\psi}$, lineWidth=4]{4.40, 4.00}{5.20, 5.60}
    \fermion[lineWidth=4]{4.00, 5.60}{6.40, 5.60}
    \electroweak[flip=true, label=$\Tilde{\psi}$, lineWidth=4]{6.00, 4.00}{5.20, 5.60}
    \fermion[lineWidth=4]{4.00, 4.00}{6.40, 4.00}
\end{feynman}}}
        \caption{}
        \label{fig1b}
    \end{subfigure}
    \begin{subfigure}[t]{0.22\textwidth}
        \centering
\thesisdiagrampanel{\scalebox{0.36}{\begin{feynman}
    \gluon[label=$\mathcal{A}_i$, color=fcb900, lineWidth=4]{5.20, 4.00}{5.20, 5.60}
    \fermion[lineWidth=4]{4.00, 5.60}{6.40, 5.60}
    \fermion[lineWidth=4]{4.00, 4.00}{6.40, 4.00}
\end{feynman}}}
        \caption{}
        \label{fig1c}
    \end{subfigure}
    \begin{subfigure}[t]{0.22\textwidth}
        \centering
\thesisdiagrampanel{\scalebox{0.36}{\begin{feynman}
    \fermion[lineWidth=4]{4.00, 5.60}{6.60, 5.60}
    \electroweak[lineWidth=4, label=$\Tilde{\psi}$]{5.20, 4.00}{5.20, 5.40}
    \fermion[lineWidth=4]{4.00, 4.00}{6.60, 4.00}
    \parton{5.20,5.60}{0.20}
\end{feynman}
}}
        \caption{}
        \label{fig1d}
    \end{subfigure}
    \caption{Diagrams contributing to the gravitational bound sector at 1PN. }
    \label{fig:my_labela}
\end{figure}
\noindent
The four graphs in Fig.~(\ref{fig:my_labela}) capture four distinct 1PN effects: the relativistic correction to the $\tilde{\psi}$ propagator, the quadratic worldline coupling $\tilde{\psi}^2$, exchange of the gravito-magnetic field $\hat{\mathcal{A}}_i$, and the velocity correction to the $\tilde{\psi}$ worldline vertex. The first graph is representative of the basic EFT step:
\begin{align}
    \begin{split}
   \mathcal{S}_{\text{eff}}\Big|_{\text{fig}(\ref{fig1a})}& =\frac{m_1 m_2}{m_p^2}\rmint \Bar{\mathcal{D}}\hat\xi\rmint dt_1\,\tilde{\psi}(\boldsymbol{x}_1
    (t_1))\rmint dt_2\,\tilde{\psi}(\boldsymbol{x}_2(t_2))
  \\ & =\frac{m_1m_2}{m_{p}^2}\rmint dt_1dt_2\Big{\langle}\tilde{\psi}(\boldsymbol{x}_1
    (t_1))\tilde{\psi}(\boldsymbol{x}_2(t_2))\Big{\rangle}_{\varoslash}\,.\\ &
    =\frac{ m_1m_2}{16 \pi m_p^2}\rmint dt\, \frac{(\boldsymbol{v}_1\cdot \boldsymbol{v}_2)-(\boldsymbol{v}_1\cdot \boldsymbol{r})(\boldsymbol{v}_2\cdot \boldsymbol{r})}{r}\sim \mathcal{O}(Lv^2)\,. \label{4.13a}
    \end{split} 
\end{align}
Here the passage from the first line to the second is simply the Wick contraction of the two potential gravitons. The remaining three diagrams use the same logic and reduce to
\begin{align}
    \mathcal{S}_{\text{eff}}\Big|_{\text{fig}(\ref{fig1b})}
    &=-\frac{m_1m_2^2}{128 \pi^4 m_p^4}\rmint \frac{ dt}{r^2}+(1\leftrightarrow 2)\,\sim \mathcal{O}(Lv^2)\,,\\
    \mathcal{S}_{\text{eff}}\Big|_{\text{fig}(\ref{fig1c})}
    &=-\frac{m_1m_2}{2\pi m_p^2}\rmint dt \,\frac{(\boldsymbol{v}_1\cdot \boldsymbol{v}_2)}{r}+(1\leftrightarrow 2)\,\sim \mathcal{O}(Lv^2)\,,\\
    \mathcal{S}_{\text{eff}}\Big|_{\text{fig}(\ref{fig1d})}
    &=-\frac{3  m_1m_2}{16\pi m_p^2}\rmint dt\, \frac{(v_1^2+v_2^2)}{r}\,\sim \mathcal{O}(Lv^2)\,.
\end{align}
Together, these four pieces reproduce the familiar 1PN gravitational correction to the conservative action.
\subsection{The Electromagnetic (EM) bound sector} \label{EM section}
Now, the EM and Proca sectors are left to be investigated. The EM bound sector consists of the following terms :
\begin{align}
    \begin{split}
  F_{\mu\nu}F^{\mu\nu}&=(\nabla_{\mu} A_{\nu}-\nabla_{\nu} A_{\mu})(\nabla^{\mu} A^{\nu}-\nabla^{\nu} A^{\mu})\\ &
%=g^{\mu a }g^{\nu b}(\nabla_\mu A_\nu-\nabla_\nu A_\mu)(\nabla_{a} A_{b}-\nabla_{b} A_{a})\\ &
%=g^{\mu a }g^{\nu b}[\nabla_{\mu} A_{\nu}\nabla_{a} A_{b}-\nabla_{\mu} A_{\nu}\nabla_{b} A_{a}]\\ &
=(\partial_0 A_i\partial_0 A_i-\partial_0 A_i\partial_iA_0+\partial_lA_i\partial_l A_i-\partial_l A_i\partial_iA_l-\partial_k A_0\partial_k A_0+\partial_kA_0\partial_0 A_k)\,.
      \end{split}
\end{align}


\begin{figure}
    \centering
    \begin{subfigure}[t]{0.22\textwidth}
        \centering
\thesisdiagrampanel{\scalebox{0.36}{\begin{feynman}
    \dashed[lineWidth=4, showArrow=false, label=$\mathcal{A}_{0}$]{5.20, 4.00}{5.20, 5.60}
    \fermion[lineWidth=4]{4.00, 4.00}{6.60, 4.00}
    \fermion[lineWidth=4]{4.00, 5.60}{6.60, 5.60}
\end{feynman}}}
        \caption{}
        \label{mfig4a}
    \end{subfigure}
      \begin{subfigure}[t]{0.22\textwidth}
        \centering
\thesisdiagrampanel{\scalebox{0.36}{\begin{feynman}
    \electroweak[lineWidth=4, label=$\mathcal{A}_{i}$, color=9900ef]{5.20, 4.00}{5.20, 5.60}
    \fermion[lineWidth=4]{4.00, 4.00}{6.60, 4.00}
    \fermion[lineWidth=4]{4.00, 5.60}{6.60, 5.60}
\end{feynman}}}
        \caption{}
        \label{mfig4b}
    \end{subfigure}
      \begin{subfigure}[t]{0.22\textwidth}
        \centering
\thesisdiagrampanel{\scalebox{0.34}{\begin{feynman}
    \dashed[showArrow=false, lineWidth=4, label=$\mathcal{A}_{0}$]{4.40, 4.00}{5.20, 5.60}
    \electroweak[lineWidth=4, label=$\Tilde{\psi}$]{6.20, 4.00}{5.20, 5.40}
    \fermion[lineWidth=4]{4.00, 4.00}{6.60, 4.00}
    \fermion[lineWidth=4]{4.00, 5.60}{6.60, 5.60}
\end{feynman}}}
        \caption{}
        \label{mfig4c}
    \end{subfigure}
    \begin{subfigure}[t]{0.22\textwidth}
        \centering
\thesisdiagrampanel{\scalebox{0.34}{\begin{feynman}
    \dashed[showArrow=false, lineWidth=4, label=$\mathcal{A}_{0}$]{5.20, 4.00}{5.20, 5.60}
    \fermion[lineWidth=4]{4.00, 4.00}{6.60, 4.00}
    \fermion[lineWidth=4]{4.00, 5.60}{6.60, 5.60}
    \parton{5.20,4.80}{0.20}
\end{feynman}}}
        \caption{}
        \label{mfig4d}
    \end{subfigure}
    \caption{Diagrams contributing to the electromagnetic bound sector up to 1PN.}
    \label{ch1:fig:em-bound}
\end{figure}
Remember that terms quadratic in a single field species belong to the propagator and do not appear again as independent interaction vertices. The gauge-fixing term $(\partial_\mu A^\mu)^2$ likewise does not contribute at 1PN order. The electromagnetic sector, therefore, contains one Coulomb exchange and three 1PN corrections. Since the Wick contractions are the same as in the gravitational example, it is enough to display the leading exchange explicitly,
\begin{align}
    \begin{split}
     \mathcal{S}_{\text{eff}}\Big|_{\text{fig}(\ref{mfig4a})}
     &=Q_1Q_2\rmint dt_1 dt_2\,\Big\langle\mathcal{A}_{0}(\boldsymbol{x}_1(t_1))\mathcal{A}_{0}(\boldsymbol{x}_2(t_2))\Big\rangle\\ &
     =-\frac{Q_1Q_2}{4\pi}\rmint \frac{dt}{r}+(1\leftrightarrow 2)\sim \mathcal{O}(Lv^0)\,.
    \end{split}
\end{align}
and summarize the remaining 1PN contributions as
\begin{align}
    \mathcal{S}_{\text{eff}}\Big|_{\text{fig}(\ref{mfig4b})}
    &=\frac{Q_1 Q_2}{4\pi}\rmint dt\, \frac{\boldsymbol{v}_1\cdot \boldsymbol{v}_2}{r}+(1\leftrightarrow 2)\sim \mathcal{O}(Lv^2)\,,\\
    \mathcal{S}_{\text{eff}}\Big|_{\text{fig}(\ref{mfig4c})}
    &=-\frac{m_2Q_1Q_2}{32\pi^2 m_p^2}\rmint \frac{dt}{r^2}+(1\leftrightarrow 2)\sim \mathcal{O}(Lv^2)\,,\\
    \mathcal{S}_{\text{eff}}\Big|_{\text{fig}(\ref{mfig4d})}
    &=\frac{Q_1Q_2}{8\pi}\rmint dt\,\Big[\frac{\boldsymbol{v}_1\cdot \boldsymbol{v}_2}{r}-\frac{(\boldsymbol{v}_1\cdot \Hat{n})(\boldsymbol{v}_2\cdot \Hat{n})}{r}\Big]\sim \mathcal{O}(Lv^2)\,.
\end{align}
The three 1PN pieces can be read respectively as the spatial-photon exchange, the graviton dressing of the Coulomb interaction, and the relativistic correction to the $\mathcal{A}_0$ propagator.
\subsection{The Proca  bound sector}
The Proca bound sector consists of the following term:
\begin{align}
    \begin{split}
  B_{\mu\nu}B^{\mu\nu}&=\frac{1}{2}(\nabla_\mu B_\nu-\nabla_\nu B_\mu)(\nabla^\mu B^\nu-\nabla^\nu B^\mu)\\ &
%=\frac{\gamma}{2}g^{\mu a }g^{\nu b}(\nabla_\mu B_\nu-\nabla_\nu B_\mu)(\nabla_a B_b-\nabla_b B_a)\\ &
%=\gamma\,g^{\mu a }g^{\nu b}[\nabla_\mu B_\nu\nabla_a B_b-\nabla_\mu B_\nu\nabla_bB_a]\\ &
=(\partial_0B_i\partial_0B_i-\partial_0B_i\partial_iB_0+\partial_lB_i\partial_lB_i-\partial_lB_i\partial_iB_l-\partial_kB_0\partial_kB_0+\partial_kB_0\partial_0B_k)\,.
      \end{split}
\end{align}

We show the effective action for the Proca field (coupled to gravity) up to 1PN order below. There are a total of four contributing diagrams. The computation is similar to that of Sec.~(\ref{EM section}).
\begin{figure}
    \centering
    \begin{subfigure}[t]{0.22\textwidth}
        \centering
\thesisdiagrampanel{\scalebox{0.35}{\begin{feynman}
    \dashed[showArrow=false, lineWidth=4, label=$\mathcal{B}_{0}$, color=9900ef]{5.20, 4.00}{5.20, 5.60}
    \fermion[lineWidth=4]{4.00, 4.00}{6.60, 4.00}
    \fermion[lineWidth=4]{4.00, 5.60}{6.60, 5.60}
\end{feynman}}}
        \caption{}
        \label{mfig5a}
    \end{subfigure}
         \begin{subfigure}[t]{0.22\textwidth}
        \centering
\thesisdiagrampanel{\scalebox{0.35}{\begin{feynman}
    \fermion[lineWidth=4]{4.00, 4.00}{6.60, 4.00}
    \gluon[lineWidth=4, color=eb144c, label=$\mathcal{B}_i$]{5.20, 4.00}{5.20, 5.60}
    \fermion[lineWidth=4]{4.00, 5.60}{6.60, 5.60}
\end{feynman}}}
        \caption{}
        \label{mfig5b}
    \end{subfigure}
      \begin{subfigure}[t]{0.20\textwidth}
        \centering
\thesisdiagrampanel{\scalebox{0.35}{\begin{feynman}
    \electroweak[lineWidth=4, label=$\Tilde{\psi}$]{6.20, 4.00}{5.20, 5.60}
    \fermion[lineWidth=4]{4.00, 4.00}{6.60, 4.00}
    \dashed[lineWidth=4, showArrow=false, color=9900ef, label=$\mathcal{B}_{0}$]{4.40, 4.00}{5.20, 5.60}
    \fermion[lineWidth=4]{4.00, 5.60}{6.60, 5.60}
\end{feynman}}}
        \caption{}
        \label{mfig5c}
    \end{subfigure}
      \begin{subfigure}[t]{0.20\textwidth}
        \centering
\thesisdiagrampanel{\scalebox{0.35}{\begin{feynman}
    \fermion[lineWidth=4]{4.00, 4.00}{6.60, 4.00}
    \dashed[lineWidth=4, showArrow=false, color=9900ef, label=$\mathcal{B}_{0}$]{5.20, 4.00}{5.20, 5.60}
    \fermion[lineWidth=4]{4.00, 5.60}{6.60, 5.60}
    \parton{5.20,4.80}{0.20}
\end{feynman}}}
        \caption{}
        \label{mfig5d}
    \end{subfigure}
    \caption{Diagrams contributing to the Proca bound sector up to 1PN.}
    \label{ch1:fig:proca-bound}
\end{figure}
The Proca sector is structurally identical to the electromagnetic one, with the massless propagator simply replaced by the Yukawa propagator. It is therefore convenient to keep only the leading Yukawa exchange explicitly,
\begin{align}
    \begin{split}
     \mathcal{S}_{\text{eff}}\Big|_{\text{fig}(\ref{mfig5a})}
     &=Q_1'Q_2'\rmint dt_1 dt_2\,\Big\langle\mathcal{B}_{0}(\boldsymbol{x}_1(t_1))\mathcal{B}_{0}(\boldsymbol{x}_2(t_2))\Big\rangle\\ &
     =-\frac{Q_1'Q_2'}{4\pi}\rmint dt\,\frac{e^{-\mu_{\gamma}r}}{r}+(1\leftrightarrow 2)\sim \mathcal{O}(Lv^0)\,,
    \end{split}
\end{align}
while the remaining 1PN terms are
\begin{align}
    \mathcal{S}_{\text{eff}}\Big|_{\text{fig}(\ref{mfig5b})}
    &=\frac{Q_1' Q_2'}{4\pi}\rmint dt\, \frac{(\boldsymbol{v}_1\cdot \boldsymbol{v}_2)e^{-\mu_{\gamma}r}}{ r}+(1\leftrightarrow 2)\sim \mathcal{O}(Lv^2)\,,\\
    \mathcal{S}_{\text{eff}}\Big|_{\text{fig}(\ref{mfig5c})}
    &=-\frac{m_2Q_1'Q_2'}{32\pi^2 m_p^2}\rmint dt \,\frac{e^{-\mu_{\gamma}r}}{r^2}+(1\leftrightarrow 2)\sim \mathcal{O}(Lv^2)\,,\\
    \mathcal{S}_{\text{eff}}\Big|_{\text{fig}(\ref{mfig5d})}
    &=-\frac{Q_1'Q_2'}{8\pi}\rmint dt \,e^{-\mu_{\gamma}r}\Big[\frac{\boldsymbol{v}_1\cdot\boldsymbol{v}_2}{r}-\frac{(\boldsymbol{v}_1\cdot\Hat{n})(\boldsymbol{v}_2\cdot\Hat{n})}{r}(1+\mu_{\gamma}r)\Big]+(1\leftrightarrow 2)\sim\mathcal{O}(Lv^2)\,.
\end{align}
This makes the close analogy between the electromagnetic and Proca sectors manifest: the new ingredient is the Yukawa suppression and the extra $(1+\mu_{\gamma}r)$ structure that accompanies the longitudinal part of the massive propagator.

\begin{figure}[htb!]
    \centering 
    \begin{subfigure}[t]{0.2\textwidth}
   \thesisdiagrampanel{\scalebox{0.2}{ \begin{feynman}
    \electroweak[flip=true, lineWidth=8, label=$\mathcal{A}_{i}$, color=9900ef]{5.20, 4.00}{6.40, 5.80}
    \electroweak[lineWidth=8, label=$\mathcal{A}_{i}$, color=9900ef]{6.40, 5.80}{7.60, 4.00}
    \fermion[lineWidth=8]{4.00, 7.60}{9.40, 7.60}
    \electroweak[lineWidth=8, label=$\varphi$, color=0693e3]{6.40, 5.80}{6.40, 7.60}
    \fermion[lineWidth=8]{4.00, 4.00}{9.40, 4.00}
\end{feynman}}}
\caption{}
\label{fig4an}
\end{subfigure}
\begin{subfigure}[t]{0.2\textwidth}
    \centering
     \thesisdiagrampanel{\scalebox{0.2}{\begin{feynman}
    \electroweak[flip=true, lineWidth=8, label=$\mathcal{A}_{i}$, color=9900ef]{5.20, 4.00}{6.40, 5.80}
    \fermion[lineWidth=8]{4.00, 7.60}{9.40, 7.60}
    \electroweak[lineWidth=8, label=$\varphi$, color=0693e3]{6.40, 5.80}{6.40, 7.60}
    \fermion[lineWidth=8]{4.00, 4.00}{9.40, 4.00}
    \dashed[showArrow=false, lineWidth=8, label=$\mathcal{A}_0$]{6.40, 5.80}{7.60, 4.00}
\end{feynman}}}
    \caption{}
    \label{fig4cn}
\end{subfigure}
\begin{subfigure}[t]{0.2\textwidth}
    \centering 
\thesisdiagrampanel{\scalebox{0.2}{\begin{feynman}
    \electroweak[flip=true, lineWidth=8, label=$\mathcal{A}_{i}$, color=9900ef]{6.40, 4.00}{6.40, 5.80}
    \electroweak[lineWidth=8, label=$\mathcal{A}_i$, color=9900ef]{6.40, 5.80}{7.60, 7.60}
    \fermion[lineWidth=8]{4.00, 7.60}{9.40, 7.60}
    \electroweak[lineWidth=8, label=$\varphi$, color=0693e3]{6.40, 5.80}{5.20, 7.60}
    \fermion[lineWidth=8]{4.00, 4.00}{9.40, 4.00}
\end{feynman}
}}
\caption{}
\label{fig4bn}
\end{subfigure}
    \caption{Scalar-Electromagnetic interaction diagrams which contribute to the bound sector.}
    \label{figtheta1}
\end{figure}

\subsection{The EM-Axion bound sector}\label{Sec4.2}
In this section, we mainly focus on the \textit{Theta term} as we want to investigate the effect of $g_{a\gamma\gamma}$ on the bound potential as well as the power radiation as discussed in Sec.~(\ref{sec5.5}). \textcolor{black}{The bulk interaction vertex which is relevant for this is the following,
\begin{equation}
    S_{int}^{\text{Theta}}=
\frac{g_{a\gamma\gamma}}{4m_p}\rmint d^4x \sqrt{-g}\,\phi\,F_{\mu\nu}\,\Tilde{F}^{\mu\nu}. 
\end{equation}
}
The contribution comes from this term at 2.5 PN order and has the following form\footnote{From the subsequent sections we replace $\hat\epsilon\rightarrow \epsilon$.},
\begin{equation}
    S_{int}^{\text{Theta}}\Big |_{2.5PN}= \frac{g_{a\gamma\gamma}}{m_p}\rmint d^4x [\hat\epsilon^{\,0ikm}\phi\,\partial_{0}\mathcal{A}_{i}\partial_{k}\mathcal{A}_{m}+ \hat\epsilon^{\,i0km}\phi\,\partial_{i}\mathcal{A}_{0}\partial_{k}\mathcal{A}_{m}]\,.
\end{equation}
 Three possible diagrams will contribute to this specific interaction. The contributing diagrams have the two-photon line and one dynamical scalar line as shown in Fig.~(\ref{figtheta1}).
\begin{itemize}
\item The first diagram that contributes has two electromagnetic propagators from the same worldline as shown in Fig.~(\ref{fig4an}) and Fig.~(\ref{fig4cn}) \textcolor{black}{with vertex factors $\phi\,\partial_{0}\mathcal{A}_{i}\partial_{k}\mathcal{A}_{m}$ and $\phi\,\partial_{i}\mathcal{A}_{0}\partial_{k}\mathcal{A}_{m}$ respectively}. The amplitude is given by, $  \mathcal{S}_{\text{eff}}=\frac{g_{a\gamma\gamma}}{m_p}Q_2^2\,m_1\,s_1\epsilon^{i0km}(  \mathcal{S}^{\text{eff}}_{ikm}+\Bar{  \mathcal{S}}^{\text{eff}}_{ikm}).$ The first and second terms are defined in (\ref{2.15}) and (\ref{4.18n})\,.
\begin{align}
   \begin{split}
      \mathcal{S}_{\text{eff}}^{ikm}\Big|_{\text{fig}(\ref{fig4an})}&=\rmint \Bar{\mathcal{D}}\hat{\xi}\rmint dt_2  \mathcal{A}_{j}(\boldsymbol{x}_2(t_2))v_{2}^{j}(t_2)\rmint dt_{3}  \mathcal{A}_{l}(\boldsymbol{x}_2(t_3))v_2^l(t_3)\rmint dt_1  \frac{\varphi(\boldsymbol{x}_1(t_1))}{m_p}\\ &
      \hspace{1.2 cm}\rmint d^4x\,  \varphi\,\partial_0  \mathcal{A}_{i}\partial_{k}\mathcal{A}_{m}(x)\,, \\ &
    =\frac{1}{m_p}\rmint \prod_{i=1}^{3}dt_i\,  dt\, v_{2}^{j}(t_2) v_2^l(t_3) \rmint d^3x\, \partial_{0}\Big\langle \mathcal{A}_{j}(\boldsymbol{x}_2(t_2))\mathcal{A}_{i}(x)\Big\rangle\,\partial_{k}\boldsymbol{\Big\langle}\mathcal{A}_{l}(\boldsymbol{x}_2(t_3))\mathcal{A}_{m}(x)\Big\rangle\,\\&\hspace{6.6cm}\Big\langle\varphi(\boldsymbol{x}_1(t_1))\varphi({x})\Big\rangle,\\ &
     =-\frac{i}{m_p}\rmint dt \,v_2^{m}(t)\Big[a_{2}^{i}(t)\rmint_{k_1,k_2}\frac{k_2^{k}}{\boldsymbol{k}_1^2\boldsymbol{k}_2^2[(\boldsymbol{k}_1+\boldsymbol{k}_2)^2+m^2]}e^{i(\boldsymbol{k}_1+\boldsymbol{k}_2)\cdot \boldsymbol{r}}\\ &
   \,\,\,\,\,\,\,\,\,\,\,\,\,\,  +v_2^{i}(t)v_2^{a}(t)\rmint_{k_1,k_2}\frac{i\,k_2^k\,k_1^a\,e^{i(\boldsymbol{k}_1+\boldsymbol{k}_2)\cdot \boldsymbol{r}}}{\boldsymbol{k}_1^2\boldsymbol{k}_2^2[(\boldsymbol{k}_1+\boldsymbol{k}_2)^2+m^2]}\Big]
   +(1\leftrightarrow 2)
   \label{2.15}
    \end{split}
\end{align}
and
\begin{align}
    \begin{split}
        \mathcal{ \Bar{S}_{\text{eff}}}^{ikm}\Big |_{\text{fig}(\ref{fig4cn})}&=-\rmint\Bar{\mathcal{D}}\hat{\xi}\rmint dt_2 \mathcal{A}_{0}(\boldsymbol{x}_2(t_2))\rmint dt_3 \mathcal{A}_{j}(\boldsymbol{x}_2(t_3))v^{j}(t_3)\rmint dt_1\,\varphi(\boldsymbol{x}_1(t_1))\\ &
        \hspace{4.3 cm}\rmint d^4x \,\varphi(x)\partial_{i}\mathcal{A}_{0}\partial_{k}\mathcal{A}_{m}\,,\\ &
        =-\rmint dt\, dt_1 dt_2 dt_3 \,v^{j}(t_3)\rmint d^3 x\,\partial_{i}\Big\langle \mathcal{A}_{0}(\boldsymbol{x}_2(t_2))\mathcal{A}_{0}(x)\Big\rangle \partial_{k}\Big\langle \mathcal{A}_{j}(\boldsymbol{x}_2(t_3))\mathcal{A}_{m}(x) \Big\rangle\\ &
       \hspace{5.8 cm} \Big\langle 
\varphi(\boldsymbol{x}_1(t_1))\varphi(x)\Big\rangle\,,\\&
        =\rmint dt \,v^{m}(t)\rmint_{\boldsymbol{k}_1,\boldsymbol{k}_2}\frac{k_1^i\,k_2^k\,e^{i(\boldsymbol{k}_1+\boldsymbol{k}_2)\cdot\boldsymbol{r}(t)}}{\boldsymbol{k}_1^2\boldsymbol{k}_2^2[(\boldsymbol{k}_1+\boldsymbol{k}_2)^2+m^2]}\,,\\ &
        =-\frac{1}{64}\rmint dt\, v^m(t)\Big[f(r)n^in^k+g(r)\delta^{ik}\Big]\,.\label{4.18n}
    \end{split}
\end{align}
Now due to the symmetric nature of $n^in^k$ and $\delta^{ik},$  $\mathcal{ \Bar{S}_{\text{eff}}}^{ikm}\Big |_{\text{eff}(\ref{fig4cn})}$ does not contribute to the amplitude after getting contracted with the $\epsilon_{i0km}.$
%\vspace{-0.45cm}
%\begin{align}
  %  \begin{split}
%      \mathcal{\Bar{S}}_{\text{eff}}&=g_{a\gamma\gamma}Q_1Q_2\,\epsilon_{i0km}\Bar{  \mathcal{S}}_{\text{eff}}^{ikm}=0\,.
%    \end{split}
%\end{align}
Therefore only $ \mathcal{S}_{\text{eff}}^{ikm}\Big|_{\text{fig}(\ref{fig4an})}$ contributes to the amplitude. The details of the computations are given in Appendix~(\ref{ch1:app:A}).
\begin{align}
    \begin{split}
          \mathcal{S}_{\text{eff}}\Big|_{\text{fig}(\ref{fig4an})}&=\frac{g_{a\gamma\gamma}}{m_{p}^2}\,Q_2^2 \,m_1\,s_1\,\epsilon^{i0km}\,  \mathcal{S}^{\text{eff}}_{ikm}\,,\\ &
= \frac{\pi^{3/2}\, g_{a\gamma\gamma}Q_2^2m_1s_1 }{8m_p^2}\Bigg[\rmint dt\, \Vec{a_2}\cdot(\hat{n}\times\boldsymbol{v}_2)\Bigg\{\frac{1}{2} \, m \,G_{1,3}^{2,1}\left(\frac{m^2 r^2}{4}\Big|
\begin{array}{c}
 -\frac{1}{2} \\
 -\frac{1}{2},-\frac{1}{2},0 \\
\end{array}
\right)\\&\hspace{5.3cm}-\frac{\, G_{1,3}^{2,1}\left(\frac{m^2 r^2}{4}\Big|
\begin{array}{c}
 \frac{1}{2} \\
 \frac{1}{2},\frac{1}{2},0 \\
\end{array}
\right)}{m\, r^2}\Bigg\}\Bigg] +1\leftrightarrow 2
\sim 
 \mathcal{O}(Lv^5).
%\,\,\,\,\frac{1}{8}\rmint dt\,[\boldsymbol{v}_1\cdot (\hat{n}\times\boldsymbol{v}_2)](\hat{n}.\boldsymbol{v}_2)\Bigg\{-\frac{m G_{1,3}^{2,1}\left(\frac{m^2 r^2}{4}|
%\begin{array}{c}
% -\frac{1}{2} \\
 %-\frac{1}{2},-\frac{1}{2},0 \\
%\end{array}
%\right)}{r}+\frac{G_{1,3}^{2,1}\left(\frac{m^2 r^2}{4}|
%\begin{array}{c}
 %\frac{1}{2} \\
% \frac{1}{2},\frac{1}{2},0 \\
%\end{array}
%\right)}{m r^3}+\\ &
%\,\,\,\,\,\,\frac{1}{2} m^3 r G_{1,3}^{2,1}\left(\frac{m^2 r^2}{4}|
%\begin{array}{c}
% -\frac{3}{2} \\
% -\frac{3}{2},-\frac{3}{2},0 \\
%\end{array}
%\right)-\frac{m G_{1,3}^{2,1}\left(\frac{m^2 r^2}{4}|
%\begin{array}{c}
 %-\frac{1}{2} \\
% -\frac{1}{2},-\frac{1}{2},0 \\
%\end{array}
%\right)}{r}+\frac{2 G_{1,3}^{2,1}\left(\frac{m^2 r^2}{4}|
%\begin{array}{c}
% \frac{1}{2} \\
 %\frac{1}{2},\frac{1}{2},0 \\
%\end{array}
%\right)}{m r^3}
%  \Bigg]+1\longleftrightarrow2
\
\label{4.31}
\end{split}
\end{align}
%The second part will not contribute because of the anti-symmetry of the Levi-Civita symbol. 
Also we can easily check that, when one takes $m\rightarrow 0$ limit, (\ref{4.31}) smoothly goes to the following, 
%\begin{align}
%  \begin{split}
%      \mathcal{S}^{\text{eff}}_{ikm}(m=0)&=-i\rmint dt\,v_{2m}(t)\Big[a_{2i}(t)\rmint_{\boldsymbol{k}_1\Vec{{k_2}}} \frac{k_2^k}{\boldsymbol{k}_1^2\boldsymbol{k}_2^2(\boldsymbol{k}_1^2+\boldsymbol{k}_2^2)}e^{i(\boldsymbol{k}_1+\boldsymbol{k}_2).\boldsymbol{r}(t)}\Big]\,,\\ &
 %   =\frac{1}{32\pi^2}\rmint dt\, v_{2m}(t)\,a_{2i}(t)\frac{n^k}{r(t)}\,.\\ &
%    \end{split}
%\end{align}

%Therefore from (\ref{2.15}) after contracting the indices appropriately one can write down the amplitude as,
\begin{align}
    \begin{split}
          \mathcal{S}_{\text{eff}}\Big|_{\text{fig}(\ref{fig4an})}&
          %=g_{a\gamma\gamma}Q_2^2\,m_1\,s_1\,\epsilon^{i0km}  \mathcal{S}_{\text{eff}1}^{ikm}\,,\\ &
        = \frac{ g_{a\gamma\gamma}\,Q_2^2\,m_1\,s_1\,}{32\pi^2\,m_p^2}\rmint dt\,\frac{\Vec{a_2}.(\hat{n}\times\boldsymbol{v}_2)}{r}+(1\leftrightarrow2)\sim \mathcal{O}(Lv^5)\,.
    \end{split}
\end{align}
\par 


\item  There is another diagram that contributes  consists of two electromagnetic fields coupled with two different worldline vertices  $\rmint dt Q_a v^{i}_{a}\mathcal{A}_{i}.$ The corresponding \textcolor{black}{amplitude} is given by,
\begin{align}
    \begin{split}
\mathcal{S}_{\text{eff}}^{ikm}\Big |_{\text{fig}(\ref{fig4bn})}&=\frac{g_{a\gamma\gamma Q_1Q_2m_1s_1}}{m_p^2}{\rmint\Bar{\mathcal{D}}\Hat{\xi}\rmint dt_1\,v_1^{j}(t_1) \mathcal{A}_{j}(\boldsymbol{x}_1(t_1))\rmint dt_2\,\varphi(\boldsymbol{x}_1(t_2))\rmint dt_3 v_2^{l}(t_1)\mathcal{A}_{l}(\boldsymbol{x}_2(t_3))}\\ & \hspace{7 cm}
\rmint d^4x \,\varphi \partial_{0}  \mathcal{A}_{i}\partial_{k}\mathcal{A}_{m}(x)\,,\\ &
        =\frac{g_{a\gamma\gamma}}{m_p^2}\rmint dt\prod_{i=1}^{3}dt_{i}\,v_1^{j}(t_1)v_2^{l}(t_3)\rmint d^3x \partial_{t}\Big\langle \mathcal{A}_{j}(\boldsymbol{x}_1(t_1))\mathcal{A}_{i}(x)\Big\rangle\partial_{k}\Big\langle \mathcal{A}_{l}(\boldsymbol{x}_2(t_3))\mathcal{A}_{m}(x)\Big\rangle \\& \hspace{5cm} \Big\langle\varphi(\boldsymbol{x}_1(t_2))\varphi(x,t)\Big\rangle\,,\\ &
        =\frac{2ig_{a\gamma\gamma}Q_1Q_2m_1s_1}{m_p^2}\rmint dt v_2^{m}(t)a_1^{i}(t)\rmint_{k}\frac{k^k\,e^{i\boldsymbol{k}\cdot\boldsymbol{r}}}{\boldsymbol{k}^2}\textcolor{black}{\rmint_{k_3}\frac{1}{(\boldsymbol{k}-\boldsymbol{k}_3)^2(\boldsymbol{k}_3^2+m^2)}}\,,\\ &
  =\frac{ g_{a\gamma\gamma}Q_1Q_2m_1s_1}{2\pi m_p^2}\rmint dt \, v_2^m (t)a_1^i(t)\,\partial_r\alpha(m;r)n^k\,,\label{4.33mm}
    \end{split}
\end{align}
Therefore, the contracted amplitude contributing to the effective action is given by,
\begin{align}
    \begin{split}
        \mathcal{S}_{\text{eff}}\Big|_{\text{fig}(\ref{fig4bn})}=\frac{ g_{a\gamma\gamma}Q_1Q_2m_1s_1}{2\pi m_p^2}\rmint dt \,(\Vec{a}_1 \times \Hat{n})\cdot \boldsymbol{v}_2\, \partial_{r}\alpha(m;r)\sim \mathcal{O}(Lv^5).
    \end{split}
\end{align}
\end{itemize}
\textcolor{black}{where the function $\alpha(m;r)$ is defined in (\ref{C3m}) of Appendix(\ref{ch1:app:C}).}
%Now come to the massive contribution to the amplitude. From (\ref{2.14}),
%\begin{align}
   % \begin{split}
     %     \mathcal{S}_{\text{eff}}_{2}^{ikm}(\mu)&=\mu^2\,\iiiint dt_1 dt_2 dt_3 dt\, v_{1}^{j}(t_2) v_2^l(t_3)\rmint d^3x \, \partial^{t}\delt)a(t-t_2)\delta_{ij}\rmint d^3k_1\frac{e^{i\boldsymbol{k}_1.[\boldsymbol{x}_2(t_2)-\boldsymbol{x}(t)]}}{\boldsymbol{k}_1^2}\delta(t-t_3)\delta_{ml}\\ &
    %\,\,\,\,\,\rmint d^3k_2\,\partial^{k}\,\frac{e^{i\boldsymbol{k}_2.[\boldsymbol{x}_2(t_3)-\boldsymbol{x}(t)]}}{\boldsymbol{k}_2^2}\delta(t-t_1)\rmint d^3k_{3}\delta(t-t_3)\rmint d^3k \frac{e^{i\boldsymbol{k}_3.[\boldsymbol{x}_1(t_1)-\boldsymbol{x}(t)]}}{\boldsymbol{k}_3^4}\\ &
   % =-i\,\mu^2\rmint dt\,v_2^{m}(t)\Big[a_{1}^{i}(t)\rmint_{\boldsymbol{k}_1\Vec{{k_2}}} \frac{k_2^k}{\boldsymbol{k}_1^2\boldsymbol{k}_2^2(\boldsymbol{k}_1+\boldsymbol{k}_2)^4}e^{i(\boldsymbol{k}_1+\boldsymbol{k}_2).\boldsymbol{r}(t)}+v_1^{i}(t) v_2^{a}(t)\rmint_{\boldsymbol{k}_1\boldsymbol{k}_2}\frac{i\,k_2^k\, k_1^a\,e^{i(\boldsymbol{k}_1+\boldsymbol{k}_2).\boldsymbol{r}(t)}}{\boldsymbol{k}_1^2\boldsymbol{k}_2^2(\boldsymbol{k}_1+\boldsymbol{k}_2)^4}\Big]
  % \end{split}
%\end{align}
%Adopting the same treatment discussed earlier in this section we can calculate this two momentum integrals. The only difference is that in the denominator there is a fourth power of the momentum. Then using (\ref{}) and (\ref{}) we will get,
%\textcolor{black}{\begin{align}
 %   \begin{split}
  %      \rmint_{\boldsymbol{k}_1\Vec{{k_2}}} \frac{k_2^k}{\boldsymbol{k}_1^2\boldsymbol{k}_2^2(\boldsymbol{k}_1+\boldsymbol{k}_2)^4}e^{i(\boldsymbol{k}_1+\boldsymbol{k}_2).\boldsymbol{r}(t)}\propto\Gamma(0
   %     )\sim \infty
    %\end{split}
%%}
%$\implies$ $  \mathcal{S}_{\text{eff}}_2(\mu)\rightarrow \infty$\par
%\textcolor{red}{QUESTION: ANY MASSIVE TERM IN THE SCALAR PROPAGATOR GIVES DIVERGENCE IN THE AMPLITUDE.WHY SO? IS THERE ANY PHYSICS BEHIND THAT? CAN'T A MASSIVE SCALAR FIELD COUPLE WITH MASSLESS ELECTROMAGNETIC FIELD IN NON RELATIVISTIC REGIME?}

%If we choose the scalar as massive and don't want to expand it for small mass then we need to compute the following intgerals,
%\begin{align}
    %\begin{split}
   %     \mathcal{I}(r)=\rmint_{\boldsymbol{k}_1\boldsymbol{k}_2} \frac{k_2^k}{\boldsymbol{k}_1^2\boldsymbol{k}_2^2[(\boldsymbol{k}_1+\boldsymbol{k}_2)^2+m^2]}e^{i(\boldsymbol{k}_1+\boldsymbol{k}_2).\boldsymbol{r}(t)},\,\tilde{\mathcal{I}}(r)=\rmint_{\boldsymbol{k}_1\boldsymbol{k}_2}\frac{\,k_2^k\, k_1^a\,e^{i(\boldsymbol{k}_1+\boldsymbol{k}_2).\boldsymbol{r}(t)}}{\boldsymbol{k}_1^2\boldsymbol{k}_2^2[(\boldsymbol{k}_1+\boldsymbol{k}_2)^2+m^2]}
  %  \end{split}
%\end{align}
%Let us loop upon the first term,
%\begin{align}
    %\begin{split}
     %   \mathcal{I}(r)&=\iint_{\boldsymbol{k},\boldsymbol{k}_1}e^{i\boldsymbol{k}.\boldsymbol{r}}\,\Big(\frac{k^k-k_1^k}{\boldsymbol{k}_1^2\,(\boldsymbol{k}^2+m^2)\,(\boldsymbol{k}_1-\boldsymbol{k})^2}\Big)\\ &
   % =\iint_{\boldsymbol{k},\boldsymbol{k}_1} e^{i\boldsymbol{k}.\boldsymbol{r}}\frac{k^k}{\boldsymbol{k}_1^2\,(\boldsymbol{k}^2+m^2)\,(\boldsymbol{k}_1-\boldsymbol{k})^2}-\iint_{\boldsymbol{k},\boldsymbol{k}_1} e^{i\boldsymbol{k}.\boldsymbol{r}}\frac{k_1^k}{\boldsymbol{k}_1^2\,(\boldsymbol{k}^2+m^2)\,(\boldsymbol{k}_1-\boldsymbol{k})^2}\\ &
   % =\iint_{\boldsymbol{k},\boldsymbol{k}_1}e^{i\boldsymbol{k}_1.\boldsymbol{r}}\frac{k_1^k}{(\boldsymbol{k}_1^2+m^2)\,\boldsymbol{k}^2\,(\boldsymbol{k}_1-\boldsymbol{k})^2}-\iint_{\boldsymbol{k},\boldsymbol{k}_1}e^{i\boldsymbol{k}_1.\boldsymbol{r}}\frac{k^k}{(\boldsymbol{k}_1^2+m^2)\,\boldsymbol{k}^2\,(\boldsymbol{k}_1-\boldsymbol{k})^2},\,\,(\boldsymbol{k}\longleftrightarrow \boldsymbol{k}_1)\\ &
   % =\rmint_{\Vec{{k_1}}}\frac{k_1^k\,e^{i\boldsymbol{k}_1.\boldsymbol{r}}}{\boldsymbol{k}_1^2+m^2}\rmint_{\boldsymbol{k}}\frac{1}{\boldsymbol{k}^2(\boldsymbol{k}-\boldsymbol{k}_1)^2}-\rmint_{\Vec{{k_1}}}\frac{e^{i\boldsymbol{k}_1.\boldsymbol{r}}}{\boldsymbol{k}_1^2+m^2}\rmint_{\boldsymbol{k}}\frac{k^k}{\boldsymbol{k}^2(\boldsymbol{k}-\boldsymbol{k}_1)^2}\\ &
   % =\rmint_{\Vec{{k_1}}}\frac{k_1^k\,e^{i\boldsymbol{k}_1.\boldsymbol{r}}}{\boldsymbol{k}_1^2+m^2}\Big[\frac{1}{(4\pi)^{3/2}}\Gamma(1/2)^3\frac{1}{|\boldsymbol{k}_1|}\Big]-\rmint_{\Vec{{k_1}}}\frac{e^{i\boldsymbol{k}_1.\boldsymbol{r}}}{\boldsymbol{k}_1^2+m^2}\Big[\frac{1}{(4\pi)^{3/2
%}}\frac{\Gamma(1/2)\Gamma(3/2)\Gamma(1/2)}{\Gamma(2)}\frac{1}%{|\boldsymbol{k}_1|}k_1^k\Big]\\ &
%=\frac{1}{8}\rmint_{\boldsymbol{k}_1}\frac{k_1^k \,e^{i\boldsymbol{k}_1.\boldsymbol{r}}}{\Vec{|k_1|}(\boldsymbol{k}_1^2+m^2)}-\frac{1}{16}\rmint_{\boldsymbol{k}_1}\frac{k_1^k \,e^{i\boldsymbol{k}_1.\boldsymbol{r}}}{|\boldsymbol{k}_1|(\boldsymbol{k}_1^2+m^2)}\\ &
%=\frac{1}{16}\rmint_{\boldsymbol{k}_1}\frac{k_1^k \,e^{i\boldsymbol{k}_1.\boldsymbol{r}}}{\Vec{|k_1|}(\boldsymbol{k}_1^2+m^2)}\\ &
%=\Bigg[\frac{i}{16} \pi^{3/2}\, m \,G_{1,3}^{2,1}\left(\frac{m^2 r^2}{4}\Big|
%\begin{array}{c}
% -\frac{1}{2} \\
% -\frac{1}{2},-\frac{1}{2},0 \\
%\end{array}
%\right)-\frac{i\,\pi^{3/2}\, G_{1,3}^{2,1}\left(\frac{m^2 r^2}{4}\Big|
%\begin{array}{c}
 %\frac{1}{2} \\
 %\frac{1}{2},\frac{1}{2},0 \\
%\end{array}
%\right)}{8 \,m\, r^2}\Bigg]n^k
%    \end{split}
%%\end{align}
%}
%\vspace{-0.63cm}
\subsection{The pure scalar sector}\label{Sec4.3}
\textcolor{black}{The contributing terms in this case is simply given by expanding the 3rd and 4th term of the Lagrangian mentioned in (\ref{2.11}).}
The non-relativistic decomposition of the action is given by,
\begin{align}
    \begin{split}
         S[{\varphi}]_{\text{int}}&
        =\frac{1}{m_p}\rmint d^3x dt \,[-2\psi \partial_{0}\varphi\partial_{0}\varphi-2\Hat{\mathcal{A}}^{i}\partial_{i}\varphi\partial_{0}\varphi-\frac{1}{4}\sigma^{lm}\delta_{lm}\partial_{i}\varphi\partial_{i}\phi+\frac{1}{2}\sigma^{ij}\partial_{i}\varphi\partial_{j}\varphi+\psi m^2 \varphi^2]\,.\label{4.24}
    \end{split}
\end{align}
Note that all the interaction terms except the last one will contribute to the effective action at the order of $L\,v^4$, i.e., of 2PN order. The last term in (\ref{4.24}) corresponds to scalar-graviton interaction at 1PN. There are two diagrams contributing to the 1PN effective action that comes from the 3-point vertex as shown in Fig.~(\ref{ch1:fig:scalar-graviton-1pn}). Next, we give the details of the amplitudes corresponding to these diagrams.
\begin{figure}[t!]
    \centering
    \begin{subfigure}[t]{0.22\textwidth}
\thesisdiagrampanel{\scalebox{0.34}{\begin{feynman}
    \electroweak[flip=true, label=$\varphi$, lineWidth=6, color=0693e3]{4.80, 4.00}{5.60, 4.80}
    \fermion[lineWidth=6]{4.00, 5.60}{7.20, 5.60}
    \electroweak[label=$\varphi$, lineWidth=4, color=0693e3]{5.60, 4.80}{6.40, 4.00}
    \electroweak[lineWidth=6, label=$\Tilde{\psi}$]{5.60, 4.80}{5.60, 5.60}
    \fermion[lineWidth=6]{4.00, 4.00}{7.20, 4.00}
\end{feynman}}}
        \caption{}
        \label{fig5aa}
    \end{subfigure}
    \begin{subfigure}[t]{0.22\textwidth}
        \centering
\thesisdiagrampanel{\scalebox{0.34}{\begin{feynman}
    \electroweak[flip=true, label=$\varphi$, lineWidth=6, color=0695e3]{4.80, 4.00}{5.60, 4.80}
    \fermion[lineWidth=6]{4.00, 5.60}{7.20, 5.60}
    \electroweak[lineWidth=6, label=$\varphi$, color=0693e3]{5.60, 4.80}{5.60, 5.60}
    \electroweak[lineWidth=6, label=$\Tilde{\psi}$]{5.60, 4.80}{6.40, 4.00}
    \fermion[lineWidth=6]{4.00, 4.00}{7.20, 4.00}
\end{feynman}
}}
        \caption{}
        \label{fig5bb}
    \end{subfigure}
    \caption{1PN diagrams for the scalar sector with 3 point vertex coming from scalar-graviton interaction.}
    \label{ch1:fig:scalar-graviton-1pn}
\end{figure}
\begin{itemize}
    \item The amplitude corresponds to the diagram in Fig.~(\ref{fig5aa}) with the bulk interaction vertex $\rmint d^4 x \,\Tilde{\psi}\varphi\varphi(x)$ where two scalar fields are contracted with the same worldline:
    \begin{align}
        \begin{split}
\mathcal{S}_{\text{eff}}\Big |_{\text{fig}(\ref{fig5aa})}&=\frac{m^2m_1 m_2^2 s_2^2}{m_{p}^4}\rmint\Bar{\mathcal{D}}\Hat{\xi}\rmint dt_1 \,\Tilde{\psi}(\boldsymbol{x}_1(t_1))\rmint dt_2 \,\varphi(\boldsymbol{x}_2(t_2))\rmint dt_3 \varphi(\boldsymbol{x}_2(t_3))\rmint d^4 x \,\Tilde{\psi}\varphi\varphi(x)\,,\\ &
            =\frac{m^2 m_1 m_2^2 s_2^2}{m_p^4} \rmint dt\prod_{i=1}^3 dt_i \Big\langle\Tilde{\psi}(\boldsymbol{x}_1(t_1))\Tilde{\psi}(x)\Big\rangle \Big\langle\varphi(\boldsymbol{x}_2(t_2))\varphi(x) \Big\rangle\Big\langle\varphi(\boldsymbol{x}_2(t_3))\varphi(x) \Big\rangle\,,\\ &
            =\frac{m^2 m_1 m_2^2 s_2^2}{2m_p^4} \rmint dt \rmint_{\boldsymbol{k},\boldsymbol{k}_1}\,\frac{e^{i\boldsymbol{k}\cdot \boldsymbol{r}}}{\boldsymbol{k}^2(\boldsymbol{k}_1^2+m^2)[(\boldsymbol{k}-\boldsymbol{k}_1)^2+m^2]}\,,\\ &
           % =\frac{m^2 m_1 m_2^2 s_2^2}{2 m_p^4}\rmint dt\,
           % \rmint_{k}\frac{e^{i \boldsymbol{k}\cdot\boldsymbol{r}}}{\boldsymbol{k}^2}
            %\begin{cases}
              %   \frac{1}{8|\boldsymbol{k}|}\Big[1-\frac{2}{\pi}\arctan(2m/|\boldsymbol{k}|)\Big]\, & k^2> 4m^2\\
       % \frac{1}{4\pi|\boldsymbol{k}|}\arctan(|\boldsymbol{k}|/2m)\,, & k^2<4m^2
           % \end{cases}\,\\ &
            =\frac{m^2 m_1 m_2^2 s_2^2}{8\pi m_p^4}\rmint dt \rmint \frac{e^{i\boldsymbol{k}\cdot \boldsymbol{r}}}{\boldsymbol{k}^3}\arctan(|\boldsymbol{k}|/2m)\,,\\&
            =\frac{m_1 m_2^2 s_2^2}{64\pi^2 m_p^4}\rmint dt \,\frac{m-2m^2r\,\text{Ei}(-2mr)-m\,e^{-2mr}}{r}+(1\leftrightarrow 2)\sim\mathcal{O}(Lv^2)\,.
            \label{4.25}
        \end{split}
    \end{align}
    %\vspace{-0.75cm}
    %\vspace{0.25 cm}
    In (\ref{4.25}), $\text{Ei}(z)=-\rmint_{-z}^{\infty}dt\,\frac{e^{-t}}{t}$ is the exponential integral function.
 \item The amplitude corresponds to the diagram in Fig.~(\ref{fig5bb}) with the bulk interaction vertex $\rmint d^4 x \,\Tilde{\psi}\varphi\varphi(x)$ where two scalar fields are contracted with two different worldlines:
    \begin{align}
        \begin{split}
            \mathcal{S}_{\text{eff}}\Big |_{\text{fig}(\ref{fig5bb})}&=\frac{m^2m_1 m_2^2 s_2^2}{m_{p}^4}\rmint\Bar{\mathcal{D}}\Hat{\xi}\rmint dt_1 \,\varphi(\boldsymbol{x}_1(t_1))\rmint dt_2 \,\Tilde{\psi}(\boldsymbol{x}_2(t_2))\rmint dt_3 \varphi(\boldsymbol{x}_2(t_3))\rmint d^4 x \,\Tilde{\psi}\varphi\varphi(x)\,,\\ &=\frac{m^2m_1s_1 m_2^2 s_2}{m_p^4} \rmint d^3x\,dt\prod_{i=1}^3 dt_i \Big\langle\Tilde{\psi}(\boldsymbol{x}_2(t_2))\Tilde{\psi}(x)\Big\rangle \Big\langle\varphi(\boldsymbol{x}_1(t_1))\varphi(x) \Big\rangle\Big\langle\varphi(\boldsymbol{x}_2(t_3))\varphi(x) \Big\rangle\,,\\ &
            =\frac{m^2m_1s_1 m_2^2 s_2}{m_p^4}\rmint_{k}\frac{e^{i\boldsymbol{k}\cdot \boldsymbol{r}}}{\boldsymbol{k}^2+m^2}\rmint_{k_1}\frac{1}{\boldsymbol{k}_1^2[(\boldsymbol{k}_1-\boldsymbol{k})^2+m^2]}\,,\\ &
           % =\frac{m^2m_1s_1 m_2^2 s_2}{m_p^4}\rmint_{k}\frac{e^{i\boldsymbol{k}\cdot\boldsymbol{r}}}{\boldsymbol{k}^2+m^2}
           % \begin{cases}
             %   \frac{1}{8|\boldsymbol{k}|}[1-\frac{2}{\pi}\arctan(m/|\boldsymbol{k}|)]\, & \boldsymbol{k}^2>m^2\\
              %  \frac{1}{4\pi|\boldsymbol{k}|}\arctan(|\boldsymbol{k}|/m)\, & \boldsymbol{k}^2<m^2
           % \end{cases}\,,\\ &
          %  =\frac{m^2m_1s_1 m_2^2 s_2}{4\pi m_p^4}\rmint_{k}\frac{e^{i\boldsymbol{k}\cdot\boldsymbol{r}}}{(\boldsymbol{k}^2+m^2)|\boldsymbol{k}|}\arctan(|\boldsymbol{k}|/m)\,,\\ &
           = \frac{m^2 m_1s_1 m_2^2 s_2}{4\pi m_p^4}\rmint dt \,\beta(m;r)+(1\leftrightarrow 2)\,\sim \mathcal{O}(Lv^2).\label{4.26}
        \end{split}
    \end{align}
    In (\ref{4.26}), $\beta(m;r)$ has no closed-form expression. Its behaviour is shown in Fig.~(\ref{lastfig}) in Appendix~(\ref{ch1:app:B}).
%\begin{figure}
 %   \centering
 %   \includegraphics[scale=0.4]{beta_vs_mr.pdf}
%    \caption{$\beta(m;r)$ vs $mr$ plot}
%    \label{beta}
%\end{figure}
\end{itemize}

%But the 3-point vertices does not contribute because of the following reason: We have the action,
%\begin{align}
 %   \begin{split}
  %      S_{\phi}&\sim\rmint d^4x \underbrace{\Big(\frac{1}{2}\eta^{\mu\nu}h^{\alpha}_{\alpha}-h^{\mu\nu}\Big)}\partial_{\mu}\phi\partial_{\nu}\phi\\ &
  %  \end{split}
%\end{align}

%5\begin{align}
%\begin{split}
      %\mathcal{S}_{\text{eff}}&\sim\rmint dt_1 h_{00}\rmint d^4x {\Big(\frac{1}{2}\eta^{\mu\nu}h^{\alpha}_{\alpha}-h^{\mu\nu}\Big)}\partial_{\mu}\phi\partial_{\nu}\phi\\ &
     % \sim \frac{\eta^{\mu\nu}}{2}\Big\langle h _{00}h_{\alpha\beta}\Big\rangle\eta^{\alpha\beta}-\Big\langle h_{00}h^{\mu\nu}\Big\rangle\sim 0
    %  \end{split}
   % \end{align}
   \begin{figure}[t!]
    \centering
    \begin{subfigure}[t]{0.25\textwidth}
       \centering
\thesisdiagrampanel{\scalebox{0.3}{\begin{feynman}
    \electroweak[label=$\varphi$, lineWidth=6, color=0693e3]{5.60, 4.00}{5.60, 5.00}
    \electroweak[label=$\varphi$, lineWidth=6, color=0693e3]{5.60, 6.40}{5.60, 5.40}
    \fermion[lineWidth=6]{4.00, 6.40}{7.20, 6.40}
    \fermion[lineWidth=6]{4.00, 4.00}{7.20, 4.00}
    \parton{5.60,5.20}{0.20}
\end{feynman}}}
        \caption{}
        \label{fig6ma}
    \end{subfigure}
    \begin{subfigure}[t]{0.25\textwidth}
        \centering
\thesisdiagrampanel{\scalebox{0.3}{\begin{feynman}
    \electroweak[lineWidth=6, color=0693e3, label=$\varphi$]{5.60, 6.00}{6.80, 4.00}
    \fermion[lineWidth=6]{4.00, 4.00}{7.20, 4.00}
    \electroweak[flip=true, lineWidth=6, label=$\Tilde{\psi}$]{4.40, 4.00}{5.60, 6.00}
    \fermion[lineWidth=6]{4.00, 6.00}{7.20, 6.00}
\end{feynman}}}
        \caption{}
        \label{fig6mb}
    \end{subfigure}
    \begin{subfigure}[t]{0.25\textwidth}
        \centering
\thesisdiagrampanel{\scalebox{0.3}{\begin{feynman}
    \electroweak[lineWidth=6, color=0693e3, label=$\varphi$]{5.60, 6.00}{6.80, 4.00}
    \fermion[lineWidth=6]{4.00, 4.00}{7.20, 4.00}
    \electroweak[flip=true, lineWidth=6, label=$\varphi$, color=0693e3]{4.40, 4.00}{5.60, 6.00}
    \fermion[lineWidth=6]{4.00, 6.00}{7.20, 6.00}
\end{feynman}}}
        \caption{}
        \label{fig6mc}
    \end{subfigure}
    \caption{1PN scalar diagrams that contribute to the bound sector.}
    \label{fig6m}
\end{figure}
There are three more diagrams contributing to 1PN coming purely from the worldline coupling operators, one is from the correction to the scalar propagator, and another is from the $\Tilde{\psi}\varphi$ coupling at any of the two worldlines. 
\begin{itemize}
\item The amplitude corresponds to the diagram in Fig.~(\ref{fig6ma}) with worldline vertices $\rmint dt\,\varphi$ and with the corrected scalar propagator:
\begin{align}
\begin{split}
    \mathcal{S}_{\text{eff}}\Big|_{\text{fig}(\ref{fig6ma})}&=\frac{m_1s_1m_2s_2}{m_{p}^2}\rmint\Bar{\mathcal{D}}\Hat{\xi}\rmint dt_1 \,\varphi(\boldsymbol{x}_1(t_1)) \rmint dt_2 \,\varphi(\boldsymbol{x}_2(t_2)),\\ &
    =\frac{m_1s_1m_2s_2}{m_p^2}\rmint dt_1 dt_2\,\Big\langle\varphi(\boldsymbol{x}_1(t_1))\varphi(\boldsymbol{x}_2(t_2))\Big\rangle_{\varoslash}\,,\\ &
    =\frac{m_1s_1m_2s_2}{m_p^2}\rmint dt_1 dt_2\, \partial_{t_1}\partial_{t_2}\delta(t_1-t_2)\rmint_{k}\frac{e^{i\boldsymbol{k}\cdot(\boldsymbol{x}_1-\boldsymbol{x}_{2})}}{(\boldsymbol{k}^2+m^2)^2}\,,\\ &
    %=\frac{m_1s_1m_2s_2}{m_p^2}\rmint dt_1 dt_2 \delta(t_1-t_2)\partial_{t_1}\partial_{t_2}\Big\{\frac{e^{-m|\boldsymbol{x}_1(t_1)-\boldsymbol{x}_2(t_2)|}}{8m\pi}\Big\}\,,\\ &
    =\frac{m_1s_1m_2s_2}{8\pi m_p^2}\rmint dt \,e^{-mr}\Big[\frac{\boldsymbol{v}_1\cdot\boldsymbol{v}_2}{r}-\frac{(\boldsymbol{v}_1\cdot\Hat{n})(\boldsymbol{v}_2\cdot\Hat{n})}{r}(1+mr)\Big]\,\sim \mathcal{O}(Lv^2).
    \end{split}
\end{align}
\item The amplitude corresponds to the diagram as shown in Fig.~(\ref{fig6mb}) with the worldline vertices $\rmint dt \,\Tilde{\psi}\varphi,$  $\rmint dt \,\Tilde{\psi}$ and  $\rmint dt \,\varphi$ : 
\begin{align}
    \begin{split}
        \mathcal{S}_{\text{eff}}\Big |_{\text{fig}(\ref{fig6mb})}&=\frac{m_1s_1m_2^2s_2}{m_{p}^4}\rmint\Bar{\mathcal{D}}\Hat{\xi}\rmint dt_1 \,\varphi(\boldsymbol{x}_1(t_1))\Tilde{\psi}(\boldsymbol{x}_1(t_1)) \rmint dt_2 \,\varphi(\boldsymbol{x}_2(t_2))\rmint dt_3 \Tilde{\psi}(\boldsymbol{x}_2(t_3)) \,,\\ &=\frac{m_1 s_1 m_2^2s_2}{m_p^4}\rmint dt_1 dt_2 dt_3 \Big\langle\varphi(\boldsymbol{x}_1(t_1))\varphi(\boldsymbol{x}_2(t_2))\Big\rangle \Big\langle\Tilde{\psi}(\boldsymbol{x}_1(t_1))\Tilde{\psi}(\boldsymbol{x}_2(t_3))\Big\rangle\,,\\ &
        =\frac{m_1 s_1 m_2^2s_2}{2 m_p^4}\rmint dt \rmint_{k_1}\frac{e^{i\boldsymbol{k}_1\cdot\boldsymbol{r}}}{\boldsymbol{k}_1^2+m^2}\rmint_{k_2}\frac{e^{i\boldsymbol{k}_2\cdot\boldsymbol{r}}}{\boldsymbol{k}_2^2}\,,\\ &
        =\frac{m_1 s_1 m_2^2s_2}{32 \pi^2 m_p^4}\rmint dt \frac{e^{-mr}}{r^2}+(1\leftrightarrow 2)\,\,\sim \mathcal{O}(Lv^2).
    \end{split}
\end{align}
\item The amplitude corresponds to the diagram in Fig.~(\ref{fig6mc}) with the worldline vertices $\rmint dt\, \varphi\varphi$ and $\rmint dt \,\varphi$ : 
\begin{align}
    \begin{split}
\mathcal{S}_{\text{eff}}\Big |_{\text{fig}(\ref{fig6mc})}&=\frac{m_1g_1m_2^2s_2^2}{m_{p}^4}\rmint\Bar{\mathcal{D}}\Hat{\xi}\rmint dt_1 \,\varphi(\boldsymbol{x}_1(t_1))\varphi(\boldsymbol{x}_1(t_1)) \rmint dt_2 \,\varphi(\boldsymbol{x}_2(t_2))\rmint dt_3 \varphi(\boldsymbol{x}_2(t_3)) \,,\\ & =\frac{m_1 g_1 m_2^2 s_2^2}{m_p^4}\rmint dt_1 dt_2 dt_3 \Big\langle\varphi(\boldsymbol{x}_1(t_1))\varphi(\boldsymbol{x}_2(t_2))\Big\rangle \Big\langle\varphi(\boldsymbol{x}_1(t_1))\varphi(\boldsymbol{x}_2(t_3))\Big\rangle\,,\\ &
         =\frac{m_1 g_1 m_2^2 s_2^2}{16\pi^2 m_p^4}\rmint dt \frac{e^{-2mr}}{r^2}+(1\leftrightarrow 2)\,\,\sim \mathcal{O}(Lv^2).
    \end{split}
\end{align}
\end{itemize}
\begin{figure}
\centering
\begin{subfigure}[t]{.22\textwidth}
  \centering
  \thesisdiagrampanel{\scalebox{0.3}{\begin{feynman}
    \fermion[lineWidth=6]{4.00, 6.20}{8.00, 6.20}
    \fermion[lineWidth=6]{4.00, 4.00}{8.00, 4.00}
    \electroweak[lineWidth=6, color=0693e3, label=$\varphi$]{6.00, 5.00}{7.00, 4.00}
    \electroweak[flip=true, lineWidth=6, color=0693e3, label=$\varphi$]{5.00, 4.00}{6.00, 5.00}
    \gluon[lineWidth=6, color=ff6900, label=$\zeta_{ij}$]{6.00, 5.00}{6.00, 6.20}
\end{feynman}}}
  \caption{}
  \label{fig3a}
\end{subfigure}%
\begin{subfigure}[t]{.22\textwidth}
  \centering
   \thesisdiagrampanel{\scalebox{0.3}{\begin{feynman}
    \fermion[lineWidth=6]{4.00, 6.20}{8.00, 6.20}
    \fermion[lineWidth=6]{4.00, 4.00}{8.00, 4.00}
    \electroweak[lineWidth=6, color=0693e3, label=$\varphi$]{6.00, 5.00}{7.00, 4.00}
    \electroweak[flip=true, lineWidth=6, color=0693e3, label=$\varphi$]{5.00, 4.00}{6.00, 5.00}
    \gluon[lineWidth=6, color=ff6900, label=$Tr(\zeta)$]{6.00, 5.00}{6.00, 6.20}
\end{feynman}
}}
  \caption{}
  \label{fig3b}
\end{subfigure}
\begin{subfigure}[t]{.22\textwidth}
\centering
\thesisdiagrampanel{\scalebox{0.3}{\begin{feynman}
    \electroweak[lineWidth=6, color=0693e3, label=$\varphi$]{5.80, 4.80}{6.60, 4.00}
    \fermion[lineWidth=6]{4.00, 6.20}{7.80, 6.20}
    \fermion[lineWidth=6]{4.00, 4.00}{7.80, 4.00}
    \electroweak[flip=true, lineWidth=6, color=0693e3, label=$\varphi$]{5.00, 4.00}{5.80, 4.80}
    \gluon[lineWidth=6, label=$\Hat{\mathcal{A}}_{i}$, color=fcb900]{5.80, 4.80}{5.80, 6.20}
\end{feynman}
}}
    \caption{}
  \label{fig3c}
\end{subfigure}
\begin{subfigure}[t]{0.22\textwidth}
  \centering
  \thesisdiagrampanel{\scalebox{0.3}{\begin{feynman}
    \electroweak[lineWidth=6, color=0693e3, label=$\varphi$]{5.80, 4.80}{6.60, 4.00}
    \fermion[lineWidth=6]{4.00, 6.20}{7.80, 6.20}
    \fermion[lineWidth=6]{4.00, 4.00}{7.80, 4.00}
    \electroweak[flip=true, lineWidth=6, color=0693e3, label=$\varphi$]{5.00, 4.00}{5.80, 4.80}
    \electroweak[lineWidth=6, label=$\Tilde{\psi}$]{5.80, 4.80}{5.80, 6.20}
\end{feynman}}}
\caption{}
\label{fig3d}
\end{subfigure}
\caption{Scalar diagrams coming from the 3-point vertices that contribute at 2 PN order in scalar effective action.}
\label{ch1:fig:scalar-2pn}
\end{figure}
\textcolor{black}{For the scalar sector, we find the effective action up to 2PN order. At 2PN order, corrections appear due to the scalar-graviton interaction vertices shown in Fig.~(\ref{ch1:fig:scalar-2pn}).} Below we list their contributions.
\begin{itemize}
\item The amplitude corresponds to the diagram in Fig.~(\ref{fig3a}) with the following bulk interaction vertex $\rmint d^4 x\frac{1}{2}\boldsymbol{\zeta}^{lm}\partial_{l}\varphi\partial_{m}\varphi(x)$:
\begin{align}
    \begin{split}
        \mathcal{S}_{\text{eff}}\Big|_{\text{fig}(\ref{fig3a})}&=-\frac{m_1}{2m_p^4}\rmint\bar {\mathcal{D}}{\hat\xi}\,\rmint dt_1 \boldsymbol{\zeta}_{ij}(\boldsymbol{x}_1(t_1))v_1^i v_1^j\rmint dt_2 m_2 s_2 \varphi(\boldsymbol{x}_2(t_2))\rmint dt_2 m_2 s_2 \varphi(\boldsymbol{x}_2(t_3))\\&\hspace{5cm}\rmint d^3x \,dt\frac{1}{2}\boldsymbol{\zeta}^{lm}\partial_{l}\varphi\partial_{m}\varphi(x)\,,\\ &
        =-\frac{ m_1 m_2^2 s_2^2}{4m_p^4}\rmint \prod_{i=1}^3 dt_i dt \,v_1^{i}(t_1)v_1^{j}(t_1)\rmint d^3x \,\partial_{l}\Big\langle\varphi(\boldsymbol{x}_2(t_2))\varphi(x)\Big\rangle \partial_{m}\Big\langle\varphi(\boldsymbol{x}_2(t_3))\varphi(x)\Big\rangle \\& \hspace{8.5cm}
        \Big\langle\boldsymbol{\zeta}_{ij}(\boldsymbol{x}_1(t_1))\boldsymbol{\zeta}^{lm}(x)\Big\rangle\,,\\ &
        =-\frac{2 m_1m_2^2 s_2^2}{m_p^4}\rmint dt\, v_1^{i}(t)v_1^{j}(t)\rmint_{\boldsymbol{k},\boldsymbol{k}_1}\frac{e^{i\boldsymbol{k}\cdot \boldsymbol{r}}k_{1l}(k_{1m}-k_m)}{(\boldsymbol{k}_1^2+m^2)[(\boldsymbol{k}-\boldsymbol{k}_1)^2+m^2]\boldsymbol{k}^2}\\&\,\,\,\,\,\,\,\,\,\,\,\,\,\,\,\,\,\,\,\,\,\,\,\,\,\,\,\,\,\,\,\,\,\,\,\,\,\,\,\,\,\,\,\,\,\,\,\,\,\,\,\,\,\,\,\,\,\,\,\,\,\,\,\,\, \,\,\,\,\,\,\,\,\,\,\,\,\,\,\,\,\,\,\,\,\,\,\,\,\,\,\,\,\,\,\,\,\,\,\,\,\,\,\,\,\,\,\,\,\,\,\,\,\,\,\,\,\,\,\,\,\,\,\,\,\,\,\,\,\,\,\,\,\,\,\,\,\,\,\,\,(\delta_{i}^{l}\delta_j^m+\delta_{i}^m\delta_j^l-2\delta_{ij}\delta^{lm})\,,\\ &
       % =\frac{m_1 m_2 s_2^2}{m_{p}^4}\rmint dt v_1^i v_1^j\Big[(\lambda_1-\chi_{1})\delta_{ij}+(\lambda_{2}-\chi_{2})n_i n_j-\delta_{ij}\Big(3(\lambda_{1}-\chi_1)+(\lambda_2-\chi_{2})\Big)\Big]\\ &
        =-\frac{4m_1 m_2^2 s_2^2}{m_p^4}\rmint dt \Big[(-2\lambda_1+2\chi_1-\lambda_2+\chi_2)v_1^2+(\lambda_2-\chi_2)(\boldsymbol{v}_1\cdot\Hat{n})^2\Big]\\&\,\,\,\,\,\,\,\,\,\,\,\,\,\,\,\,\,\,\,\,\,\,\,\,\,\,\,\,\,\,\,\,\,\,\,\,\,\ +(1\leftrightarrow2)\sim \mathcal{O}(Lv^4)\,.
        \label{4.38}
    \end{split}
\end{align}
\vspace{-0.4cm}
The calculation and the functions $\lambda_{1,2}(r)$, $\chi_{1,2}(r)$ are explicitly given in Appendix~ (\ref{ch1:app:B}).
\vspace{0.2cm}
\item The amplitude corresponds to the diagram in Fig.~(\ref{fig3b}) with the following bulk interaction vertex $\rmint d^4x\,\frac{1}{4}\boldsymbol{\zeta}^{lm}\delta_{lm}\partial_{k}\varphi\partial_{k}\varphi(x)$:
\begin{align}
    \begin{split}
          \mathcal{S}_{\text{eff}}\Big |_{\text{fig}(\ref{fig3b})}&=-\frac{m_1}{2m_p^4}\rmint\Bar{\mathcal{D}}\Hat{\xi}\rmint dt_1 \boldsymbol{\zeta}_{ij}(\boldsymbol{x}_1(t_1))v_1^i v_1^j\rmint dt_2 m_2 s_2 \varphi(\boldsymbol{x}_2(t_2))\rmint dt_2 m_2 s_2 \varphi(\boldsymbol{x}_2(t_3))\\&\hspace{5cm}\rmint d^3x \,dt\frac{1}{4}\boldsymbol{\zeta}^{lm}\delta_{lm}\partial_{k}\varphi\partial_{k}\varphi(x)\,,\\ &
        =-\frac{m_1 m_2^2 s_2^2}{8 m_p^4}\rmint dt_1 dt_2 dt_3 dt v_{1}^i(t_1)v_1^j(t_1) \rmint d^3x \,\partial_{k}\Big\langle\varphi(\boldsymbol{x}_2(t_2))\varphi(x)\Big\rangle \partial_{k}\Big\langle\varphi(\boldsymbol{x}_2(t_3))\varphi(x)\Big\rangle
       \\&\hspace{8.5cm} \Big\langle\boldsymbol{\zeta}_{ij}(\boldsymbol{x}_1(t_1))\boldsymbol{\zeta}^{lm}(x)\Big\rangle \delta_{lm}\,,\\ &
        =-\frac{m_1 m_2^2 s_2^2}{ m_p^4}\rmint dt v_1^{i}(t)v_1^j(t)\rmint_{k,k_1}\frac{e^{i\boldsymbol{k}\cdot \boldsymbol{r}}(\boldsymbol{k}_1^2-\boldsymbol{k}\cdot \boldsymbol{k}_1)}{(\boldsymbol{k}_1^2+m^2)[(\boldsymbol{k}_1-\boldsymbol{k})^2+m^2]\boldsymbol{k}^2}\delta_{ij}\,,\\ &
        =-\frac{m_1 m_2^2 s_2^2}{m_p^4}\rmint dt\, v_1^2\Big(3(\lambda_1-\chi_1)+(\lambda_2-\chi_2)\Big)+(1\leftrightarrow 2)\sim \mathcal{O}(Lv^4)\,.
    \end{split}
\end{align}
Again, all the details of the integrals are given in Appendix~(\ref{ch1:app:B}). The functions $\lambda_{1,2},\chi_{1,2}$ are defined in  (\ref{B.5m}) of Appendix~(\ref{ch1:app:B}).
\item The amplitude corresponds to the diagram in Fig.~(\ref{fig3c}) with the bulk interaction vertex $\rmint d^4x\,\Hat{\mathcal{A}}^{k}\partial_{k}\varphi\partial_{0}\varphi(x)$:
\begin{align}
    \begin{split}
          \mathcal{S}_{\text{eff}}\Big|_{\text{fig}(\ref{fig3c})}&=2m_1\rmint\Bar{\mathcal{D}}\Hat{\xi}\rmint dt_1\, \frac{\Hat{\mathcal{A}}_{i}(\boldsymbol{x}_1(t_1))}{m_{p}}v_1^{i}(t_1)\rmint dt_2\,\frac{m_2 s_2}{m_p}\varphi(\boldsymbol{x}_2(t_2))\rmint dt_3\,\frac{m_2 s_2}{m_p}\varphi(\boldsymbol{x}_2(t_3))\\&\hspace{5cm}\rmint d^3x dt \,\frac{\Hat{\mathcal{A}}^{k}}{m_p}\partial_{k}\varphi\partial_{0}\varphi(x)\,,\\ &
        =\frac{2m_1 m_2^2 s_2^2}{m_p^4}\rmint dt dt_1 dt_2 dt_3 v_1^i(t_1)\rmint d^3x \Big\langle\Hat{\mathcal{A}}_{i}(\boldsymbol{x}_1(t_1))\mathcal{\Hat{A}}^{k}(x)\Big\rangle \partial_{k}\Big\langle\varphi(\boldsymbol{x}_2(t_2))\varphi(x)\Big\rangle \\&\hspace{9.5cm}\partial_{0}\Big\langle\varphi(\boldsymbol{x}_2(t_3))\varphi(x  )\Big\rangle\,,\\ &
       % =\frac{2 m_2 m_1^2 s_1^2}{m_p^4}\rmint dt dt_3 \partial_{t}\delta(t-t_3)\rmint_{k_{1,2,3}}\delta(\boldsymbol{k}_1+\boldsymbol{k}_2+\boldsymbol{k}_3)\frac{ik_2^k\,e^{i\boldsymbol{k}_1\cdot \boldsymbol{x}_1(t)}e^{i\boldsymbol{k}_2\cdot \boldsymbol{x}_2(t)}e^{i\boldsymbol{k}_3\cdot \boldsymbol{x}_2(t_3)}}{\boldsymbol{k}_1^2(\boldsymbol{k}_2^2+m^2)(\boldsymbol{k}_3^2+m^2)}\delta_i^k\\ &
        =-\frac{4m_2 m_1^2 s_1^2}{m_p^4}\rmint dt v_1^i(t)v_2^{l}(t)\rmint_{k_2,k_3}\frac{k_2^i\,k_3^l\,e^{-i(\boldsymbol{k}_2+\boldsymbol{k}_3)\cdot \boldsymbol{r}}}{(\boldsymbol{k}_2+\boldsymbol{k}_3)^2(\boldsymbol{k}_2^2+m^2)(\boldsymbol{k}_3^2+m^2)}\,,\\ &
      %  =\frac{2m_2 m_1^2 s_1^2}{m_p^4}\rmint dt\,v_1^{i}(t)v_2^{l}(t)[(\lambda_1+\chi_1)\delta_{il}+(\lambda_2+\chi_2)n_l n_i]\\ &
        =-\frac{4m_2 m_1^2s_1^2}{m_p^4}\rmint dt[(\lambda_1+\chi_1)\boldsymbol{v}_1\cdot\boldsymbol{v}_2+(\lambda_2+\chi_2)(\boldsymbol{v}_1\cdot \Hat{n})(\boldsymbol{v}_2\cdot \Hat{n})]+(1\leftrightarrow 2)\sim \mathcal{O}(L v^4)\,.
    \end{split}
\end{align}
Again all the  details of the computation and the functions $\lambda_{1,2},\chi_{1,2}$ are given in Appendix~(\ref{ch1:app:B}).
%\newpage
\item The amplitude corresponds to the diagram in Fig.~(\ref{fig3d}) with the bulk interaction vertex $\rmint d^4x\, \frac{\Tilde{\psi}}{m_p}\partial_{0}\varphi\partial_{0}\varphi$:
\begin{align}
    \begin{split}
          \mathcal{S}_{\text{eff}}\Big |_{\text{fig}(\ref{fig3d})}&=2m_1\rmint\Bar{\mathcal{D}}\Hat{\xi}\rmint dt_1 \frac{\Tilde{\psi}(\boldsymbol{x}_1(t_1))}{m_p}\rmint dt_2 \frac{m_2 s_2}{m_p}\varphi(\boldsymbol{x}_2(t_2))\rmint dt_3 \frac{m_2 s_2}{m_p}\varphi(\boldsymbol{x}_2(t_3))\\& \,\,\,\,\,\,\,\,\,\,\,\,\,\,\,\,\,\,\,\,\,\,\,\,\,\,\,\,\,\,\,\,\,\,\,\,\,\,\,\,\,\,\,\,\,\,\,\,\,\,\,\,\,\,\,\,\,\,\,\,\,\,\,\,\,\,\,\,\,\,\,\,\,\,\,\,\,\,\,\,\,\,\,\,\,\,\,\,\,\,\,\,\,\,\,\,\,\,\,\,\,\,\,\,\,\,\,\,\,\,\,\,\,\,\,\,\,\,\,\,\,\,\,\, \rmint d^3x\,dt\, \frac{\Tilde{\psi}}{m_p}\partial_{0}\varphi\partial_{0}\varphi\,.\\ &
        =\textstyle{\frac{2m_1 m_2^2 s_2^2}{m_p^4}\rmint dt_1 dt_2 dt_3 dt\rmint d^3x \Big\langle\Tilde{\psi}(\boldsymbol{x}_1(t_1))\Tilde{\psi}(x)\Big\rangle\partial_{0}\Big\langle\varphi(\boldsymbol{x}_2(t_2))\varphi(x)\Big\rangle \partial_{0}\Big\langle\varphi(\boldsymbol{x}_2(t_3))\varphi(x  )\Big\rangle}\,,\\ &
        =\frac{m_1 m_2^2 s_2^2}{m_p^4}\rmint dt \, v_2^a(t)v_2^b(t)\rmint_{k,k_1}\frac{k_{1a}(k_{1b}-k_b)e^{i\boldsymbol{k}\cdot \boldsymbol{r}}}{\boldsymbol{k}^2(\boldsymbol{k}_1^2+m^2)[(\boldsymbol{k}-\boldsymbol{k}_1)^2+m^2]}\,,\\ &
        =\frac{m_1 m_2^2 s_2^2}{m_p^4}\rmint dt \, [v_2^2(\lambda_1-\chi_1)+(\boldsymbol{v}_2\cdot \Hat{n})^2(\lambda_2-\chi_2)]+(1\leftrightarrow 2)\sim \mathcal{O}(Lv^4)\,.
    \end{split}
\end{align}
\end{itemize}
At 2PN there are several more terms that are of the order $\mathcal{O}(Gv^4)$. Below we list out the contribution of those diagrams to the effective action. 
\begin{figure}[t!]
    \centering
    \begin{subfigure}[t]{0.25\textwidth}
        \centering
        \thesisdiagrampanel{\scalebox{0.3}{\begin{feynman}
    \electroweak[lineWidth=4, label=$\varphi$, color=0693e3]{5.80, 4.00}{5.80, 6.20}
    \fermion[lineWidth=4, label=$v^4$]{4.00, 6.20}{7.60, 6.20}
    \fermion[lineWidth=4, label=$v^0$]{4.00, 4.00}{7.60, 4.00}
\end{feynman}}}
\caption{}
\label{fig4a}
    \end{subfigure}
    \begin{subfigure}[t]{0.25\textwidth}
        \centering
        \thesisdiagrampanel{\scalebox{0.3}{\begin{feynman}
    \electroweak[lineWidth=4, label=$\varphi$, color=0693e3]{5.80, 4.00}{5.80, 6.20}
    \fermion[lineWidth=4, label=$v^2$]{4.00, 6.20}{7.60, 6.20}
    \fermion[lineWidth=4, label=$v^2$]{4.00, 4.00}{7.60, 4.00}
\end{feynman}}}
\caption{}
\label{fig4b}
    \end{subfigure}
    \caption{Diagrams contributing to scalar bound sector at 2 PN order coming from worldline vertices. }
    \label{fig4}
\end{figure}
\begin{itemize}
    \item One term that contributes consists of the worldline vertex $\frac{1}{8}m_1 s_1\rmint dt \,v^4\frac{\varphi}{m_p}$ as shown in the Fig.~(\ref{fig4a}). The corresponding amplitude is given by,
    \begin{align}
        \begin{split}
              \mathcal{S}_{\text{eff}}\Big|_{\text{fig}(\ref{fig4a})}&=-\frac{m_1 s_1 m_2 s_2 }{8m_p^2}\rmint\Bar{\mathcal{D}}\Hat{\xi}\rmint dt_1\,v_1^4(t_1) \varphi(\boldsymbol{x}_1(t_1))\rmint dt_2 \varphi(\boldsymbol{x}_2(t_2))\,,\\ &
            =-\frac{m_1 s_1 m_2 s_2 }{8m_p^2}\rmint dt_1 \,dt_2\,v_1^4(t_1)\Big\langle\varphi(\boldsymbol{x}_1(t_1))\varphi(\boldsymbol{x}_2(t_2))\Big\rangle\,,\\ &
            =-\frac{m_1 s_1 m_2 s_2}{32\pi m_p^2}\rmint dt \,v_1^4(t)\,\frac{e^{-mr}}{r}+(1\leftrightarrow2)\sim \mathcal{O}(Lv^4)\,.
        \end{split}
    \end{align}
   % \newpage
    \item Another term that contributes at this order consists of the worldline vertex vertex:$\rmint dt \frac{1}{2m_p}v^2\,m_a s_a \varphi$ as shown in Fig.~(\ref{fig4b}). The corresponding amplitude is given by,
    \begin{align}
        \begin{split}
              \mathcal{S}_{\text{eff}}\Big|_{\text{fig}(\ref{fig4b})}&=\frac{m_1 s_1 m_2 s_2}{4m_p^2}\rmint\Bar{\mathcal{D}}\Hat{\xi}\rmint dt_1 v_1^2(t_1)\varphi(\boldsymbol{x}_1(t_1))\rmint dt_2 v_2^2(t_2)\varphi(\boldsymbol{x}_2(t_2))\,,\\ &
            =\frac{m_1 s_1 m_2 s_2}{4m_p^2}\rmint dt_1 \, dt_2 v_1^2(t_1)v_2^2(t_2)\Big\langle\varphi(\boldsymbol{x}_1(t_1))\varphi(\boldsymbol{x}_2(t_2))\Big\rangle\,,\\ &
            =\frac{m_1 s_1 m_2 s_2}{16 \pi m_p^2}\rmint dt\,v_1^2(t)v_2^2(t)\frac{e^{-mr}}{r}+(1\leftrightarrow 2)\,\sim \mathcal{O}(Lv^4)\,.
        \end{split}
    \end{align}
\end{itemize} 
\subsection{Proca-electromagnetic interaction sector}
Now we focus on the Proca-electromagnetic interaction sector. This sector contains an interaction term of the form $F_{\mu\nu}B^{\mu\nu}$. This term can be expanded \textcolor{black}{up to 1PN order }in the following fashion.
\begin{align}
    \begin{split}
  \frac{\gamma}{2} F_{\mu\nu}B^{\mu\nu}&=\frac{\gamma}{2}(\nabla_\mu A_\nu-\nabla_\nu A_\mu)(\nabla^\mu B^\nu-\nabla^\nu B^\mu)\\ &
%=\frac{\gamma}{2}g^{\mu a }g^{\nu b}(\nabla_\mu A_\nu-\nabla_\nu A_\mu)(\nabla_a B_b-\nabla_b B_a)\\ &
%=\gamma\,g^{\mu a }g^{\nu b}[\nabla_\mu A_\nu\nabla_a B_b-\nabla_\mu A_\nu\nabla_bB_a]\\ &
=\gamma(\partial_0A_i\partial_0B_i-\partial_0A_i\partial_iB_0+\partial_lA_i\partial_lB_i-\partial_lA_i\partial_iB_l-\partial_kA_0\partial_kB_0+\partial_kA_0\partial_0B_k)\,.\label{4.35}
      \end{split}
\end{align}
In the action mentioned in (\ref{4.35}), the two fields are different (although they are quadratic), and they provide us with two-point interaction vertices. 
\begin{figure}[t!]
    \centering
    \begin{subfigure}[t]{0.22\textwidth}
    \centering
    \thesisdiagrampanel{\scalebox{0.3}{ \begin{feynman}
    \electroweak[lineWidth=6, color=9900ef, label=$\mathcal{A}_i$]{6.00, 5.00}{6.00, 6.00}
    \dashed[showArrow=false, color=9900ef, lineWidth=6, label=$\mathcal{B}_0$]{6.00, 4.00}{6.00, 5.00}
    \fermion[lineWidth=6]{4.00, 4.00}{8.00, 4.00}
    \fermion[lineWidth=6]{4.00, 6.00}{8.00, 6.00}
\end{feynman}}}
        \caption{}
        \label{fig9am}
    \end{subfigure}
    \begin{subfigure}[t]{0.22\textwidth}
        \centering
        \thesisdiagrampanel{\scalebox{0.3}{\begin{feynman}
    \electroweak[lineWidth=6, color=9900ef, label=$\mathcal{A}_i$]{6.00, 5.00}{6.00, 6.00}
    \gluon[flip=true, lineWidth=6, color=eb144c, label=$\mathcal{B}_i$]{6.00, 4.00}{6.00, 5.00}
    \fermion[lineWidth=6]{4.00, 4.00}{8.00, 4.00}
    \fermion[lineWidth=6]{4.00, 6.00}{8.00, 6.00}
\end{feynman}
}}
\caption{}
\label{fig9bm}
    \end{subfigure}
    \begin{subfigure}[t]{0.22\textwidth}
        \centering
        \thesisdiagrampanel{\scalebox{0.3}{\begin{feynman}
    \electroweak[lineWidth=6, color=9900ef, label=$\mathcal{A}_i$]{6.00, 5.00}{6.00, 6.00}
    \gluon[flip=true, lineWidth=6, color=eb144c, label=$\mathcal{B}_l$]{6.00, 4.00}{6.00, 5.00}
    \fermion[lineWidth=6]{4.00, 4.00}{8.00, 4.00}
    \fermion[lineWidth=6]{4.00, 6.00}{8.00, 6.00}
\end{feynman}
}}
\caption{}
\label{fig9cm}
    \end{subfigure}
    \begin{subfigure}[t]{0.22\textwidth}
        \centering
        \thesisdiagrampanel{\scalebox{0.3}{\begin{feynman}
    \dashed[showArrow=false, lineWidth=6, color=9900ef, label=$\mathcal{B}_0$]{6.00, 4.00}{6.00, 5.00}
    \fermion[lineWidth=6]{4.00, 4.00}{8.00, 4.00}
    \fermion[lineWidth=6]{4.00, 6.00}{8.00, 6.00}
    \dashed[showArrow=false, lineWidth=6, label=$\mathcal{A}_0$]{6.00, 5.00}{6.00, 6.00}
\end{feynman}}}
\caption{}
\label{fig9dm}
    \end{subfigure}
    \begin{subfigure}[t]{0.22\textwidth}
        \centering
        \thesisdiagrampanel{\scalebox{0.3}{\begin{feynman}
    \gluon[flip=true, color=eb144c, lineWidth=6, label=$\mathcal{B}_k$]{6.00, 4.00}{6.00, 5.00}
    \fermion[lineWidth=6]{4.00, 4.00}{8.00, 4.00}
    \fermion[lineWidth=6]{4.00, 6.00}{8.00, 6.00}
    \dashed[showArrow=false, lineWidth=6, label=$\mathcal{A}_0$]{6.00, 5.00}{6.00, 6.00}
\end{feynman}}}
\caption{}
\label{fig9em}
    \end{subfigure}
    \caption{Diagrams contributing to the bound sector at 1PN order due to the Proca-electromagnetic interaction term.}
    \label{fig:my_label2}
\end{figure}
Now we need to calculate the amplitude of all the diagrams shown in Fig.~(\ref{fig:my_label2}) that contribute to 1PN order in the effective action. Below we list all such contributions. 
%\newpage
\begin{itemize}
\item The amplitude corresponding to the diagram in Fig.~(\ref{fig9am}) with bulk interaction vertex $\rmint d^4 x\,\partial_0\mathcal{A}_i\partial_i\mathcal{B}_0$:
\begin{align}
    \begin{split}
   \mathcal{S}_{\text{eff}}\Big |_{{\text{fig}(\ref{fig9am})}}&=\gamma \,Q_1 Q_2'\rmint\Bar{\mathcal{D}}\Hat{\xi}\rmint dt_1v_1^j(t_1)\mathcal{A}_j(\boldsymbol{x}_1(t_1))\rmint dt_2\, \mathcal{B}_0(\boldsymbol{x}_2(t_2))\rmint d^4x\,\partial_0\mathcal{A}_i(x)\partial_i\mathcal{B}_0(x),\\ &=\gamma \,Q_1 Q_2'\rmint dt_1dt_2v_1^j(t_1)\rmint d^4x\,\partial_0\Big{\langle}\mathcal{A}_j(\boldsymbol{x}_1(t_1))\mathcal{A}_i(x)\Big{\rangle}\partial_i\Big{\langle}\mathcal{B}_0(\boldsymbol{x}_2(t_2))\mathcal{B}_0(x)\Big{\rangle},\\ &
%=\rmint dt_1\, dt_2\,v_1^i(t_1)\rmint d^4x\partial_0\rmint\delta(t-t_1)\frac{e^{-i\boldsymbol{k}\cdot [\boldsymbol{x}(t)-\boldsymbol{x}_1(t_1)]}}{k^2}\delta_{ij}\partial_i\rmint\delta(t-t_2)\frac{e^{-i\boldsymbol{k}_1\cdot[\boldsymbol{x}(t)-\boldsymbol{x}_2(t_2)]}}{k_1^2+\mu^2}\\&
%=-i\rmint dt_1\, dt_2\,dt\,v_1^i(t_1)\partial_0\delta(t-t_1)\rmint_{\boldsymbol{k},\boldsymbol{k}_1}\delta(\boldsymbol{k}+ \boldsymbol{k}_1)\delta(t-t_2)\frac{k_1^ie^{i\boldsymbol{k}\cdot \boldsymbol{x}_1(t_1)}e^{i\boldsymbol{k}_1\cdot \boldsymbol{x}_2(t_2)}}{k^2(k_1^2+\mu_{\gamma}^2)}\\ &
%=-i\rmint dt\, \,\partial_{t_1}\Big[v_1^i(t_1)\rmint \frac{e^{-i\boldsymbol{k}\cdot [\boldsymbol{x}_2(t)-\boldsymbol{x}_1(t_1)]}\,k^i}{\boldsymbol{k}^2(\boldsymbol{k}^2+\mu^2)}\Big]_{t_1=t}\\ &
=i\gamma\,Q_1 Q_2'\rmint dt \Bigg[a^{i}_{1}(t)\rmint_{\boldsymbol{k}}\frac{k^i\,e^{-i\,\boldsymbol{k}\cdot\boldsymbol{r}}}{\boldsymbol{k}^2(\boldsymbol{k}^2+\mu_{\gamma}^2)}+v_1^{i}(t)v_{1}^{a}(t)\frac{i\,k^a\,k^i\,e^{-i\boldsymbol{k}\cdot\boldsymbol{r}}}{\boldsymbol{k}^2(\boldsymbol{k}^2+\mu_{\gamma}^2)}\Bigg]\,,\\ &
%=-i\rmint dt \partial^{t_1}\Big[v_1^i(t_1)\Big(\rmint \frac{e^{-i\boldsymbol{k}\cdot (x_2(t)-x_1(t_1)}k^i}{k^2}-\frac{e^{-i\boldsymbol{k}\cdot \vec r}k^i}{k^2+\mu^2}\Big)\Big]_{t_1=t}\\&
=-\frac{\gamma\,Q_1 Q_2'}{2\pi}\rmint dt \Bigg[(\Vec{a}_1\cdot \Hat{n})\Big\{\frac{  e^{-\mu_{\gamma} r}}{2 \mu_{\gamma} \,r}-\frac{1 -  e^{-\mu_{\gamma} r}}{2\, \mu_{\gamma}^2\, r^2}\Big\}+\boldsymbol{v}_1^2 \,\frac{e^{-r \mu _{\gamma }} \left(-r \mu _{\gamma }+e^{r \mu _{\gamma }}-1\right)}{  r^3 \mu _{\gamma }^2}\\ &
\,\,\,\,\,\,\,\,\,\,\,\,\,\,\,\,\,\,\,\,\,\,\,\,\,\,\,\,\,
+(\boldsymbol{v}_1\cdot \Hat{n})^2\,\frac{e^{-r \mu _{\gamma }} \Big(r \mu _{\gamma } \left(r \mu _{\gamma }+3\right)-3 e^{r \mu _{\gamma }}+3\Big)}{  r^3 \mu _{\gamma }^2}
\Bigg]+(1\leftrightarrow 2)\,\sim \mathcal{O}(L v^2)\,.
\end{split}
\end{align}
%The first part of the integration can be easily done using the formula given in appendix A.
%The result it gives is the following:
%\begin{align}
%\begin{split}
%\partial^{t_1}(\rmint \frac{e^{-i\boldsymbol{k}
%\cdot \boldsymbol{r}}\,k^i}{k^2})&=(2\pi)^3\frac{2i}{(4\pi)^{3/2}}\pi^{1/2}\partial^{t_1}\frac{n^i}{r^2} \\ &%=(2\pi)^3 \frac{2i}{(4\pi)^{3/2}}\pi^{1/2}n^i\frac{(2\boldsymbol{v}_1\cdot \boldsymbol{r})}{r^4}
%\end{split}
%\end{align}
%So, the result now simplifies to:
%\begin{align}
 %  \begin{split}
%  \mathcal{S}_{\text{eff}}(\partial^0A^i\partial^iB^0)&=-i\rmint dt \frac{\boldsymbol{v}_1\cdot\hat n}{\mu^2}\Bigg[\frac{4i\pi^2(\boldsymbol{v}_1\cdot\hat{r})}{r^3}-\partial^{t_1}\rmint \frac{e^{-i\boldsymbol{k}\cdot \vec r}k^i}{(k^2+\mu^2)}\Bigg]_{t_1=t}\\ &\label{3.27}
%\end{split}
%\end{align}
%Now for calculating the second term we do the following trick:
%\begin{align}
 %  \begin{split}

%\text{so},
%\partial^t_1(\frac{1}{r}2\pi^2e^{-\mu r})\\&
%=\Bigg[\frac{-1}{r^2}\boldsymbol{v}_1+\frac{1}{r} \boldsymbol{v}_1(-\mu)\Bigg]2\pi^2e^{(-\mu r)}\\&
%\end{split}
%\end{align}
%Again acting $-\frac{1}{i}\partial^{r_i}$ we get,
%\begin{align}
 %   \begin{split}
%-\frac{1}{i}\partial^{r_i}[\frac{-1}{r^2}\boldsymbol{v}_1+\frac{1}{r} \boldsymbol{v}_1(-\mu)]2\pi^2e^{(-\mu r)}&=-\frac{1}{i}\Bigg[\frac{2}{r^3}(\vec  v_1\cdot\hat r)e^{-\mu r}-\frac{2}{r^2}(\boldsymbol{v}_1\cdot \hat r)e^{-\mu r}(-\mu)+\frac{1}{r}e^{-\mu r}\mu^2( \boldsymbol{v}_1\cdot \hat r)\Bigg]2\pi^2\label{3.29}
%\end{split}
%\end{align} 
%Now we can see that exactly the first term of (\ref{3.27}) cancel against the first term arising from the second integral. again if we expand the first term of (\ref{3.29}) we can easily see that the second term also cancels with the next term of the expansion.So the whole expansion of the numerator begins at $\mathcal{O}(\mu^2)$ and that cancels the $\mu^2$ in the denomitaor.
%therefore we get the resultant expression of (\ref{3.27}) as,
%\begin{align}
 %  \begin{split}
  %    \mathcal{S}_{\text{eff}}(\partial^0A^i\partial^iB^0)&=
%-\rmint dt {2\pi^2}\frac{1}{\mu^2}\Bigg[\frac{2}{r}\mu^2(v_1\cdot \hat r)+\frac{2}{r^3}(v_1\cdot \hat r)\sum_{n=3}^{\infty}\frac{(-\mu r)^n}{n!}+\frac{1}{r}\mu^2(v_1\cdot \hat r)\sum_{k=1}^{\infty}\frac{(-\mu r)^k}{k!}+\\&\frac{2\mu}{r^2}(v_1\cdot \hat r)\sum_{k=1}^{\infty}\frac{(-\mu r)^k}{k!}\Bigg]( v_1\cdot \hat n)\\&+1\longleftrightarrow2
%\end{split}
%\end{align}
\vspace{-0.5cm}
\item The amplitude corresponding to the diagram in Fig.~(\ref{fig9bm}) with bulk interaction vertex $\rmint d^4x\,\partial_l\mathcal{A}_i\partial_l\mathcal{B}_i$:
\begin{align}
    \begin{split}
   \mathcal{S}_{\text{eff}}\Big|_{\text{fig}(\ref{fig9bm})}&=\gamma \,Q_1 Q_2'\rmint\Bar{\mathcal{D}}\Hat{\xi}\rmint dt_1v_1^j(t_1)\mathcal{A}_j(\boldsymbol{x}_1(t_1))\rmint dt_2\, \mathcal{B}_m(\boldsymbol{x}_2(t_2))v_2^m(\boldsymbol{x}_2(t_2))\\ & \hspace{6.5 cm}\rmint d^4x\,\partial_l\mathcal{A}_i(x)\partial_l\mathcal{B}_i(x),\\&=\gamma Q_1 Q_2'\rmint dt_1dt_2v_1^j(t_1)v_2^m(t_2)\rmint d^4x\partial_l\Big{\langle}\mathcal{A}_j(\boldsymbol{x}_1(t_1))\mathcal{A}_i(x)\Big{\rangle}\partial_l\Big{\langle}\mathcal{B}_m(\boldsymbol{x}_2(t_2))\mathcal{B}_i(x)\Big{\rangle}\,,\\ &
%=\textstyle{Q_1Q_2'\rmint dt_1\, dt_2\,v_1^i(t_1)v_2^m(t_2)\rmint d^4x\partial_l\rmint\delta(t-t_1)\frac{e^{-i\boldsymbol{k}\cdot [\boldsymbol{x}(t)-\boldsymbol{x}_1(t_1)]}}{k^2}\delta_{im}\partial_l\small{\rmint\delta(t-t_2)\frac{e^{-i\boldsymbol{k}_1\cdot[\boldsymbol{x}(t)-\boldsymbol{x}_2(t_2)]}}{k_1^2+\mu^2}}}\\&
%=-\scriptsize{Q_1Q_2'\rmint dt_1\, dt_2\,dt\,v_1^i(t_1)v_2^i(t_2)\delta(t-t_1)\rmint_{\boldsymbol{k},\boldsymbol{k}_1}\delta(\boldsymbol{k}+ \boldsymbol{k}_1)\delta(t-t_2)\frac{k^lk_1^le^{i\boldsymbol{k}\cdot \boldsymbol{x}_1(t_1)}e^{i\boldsymbol{k}_1\cdot \boldsymbol{x}_2(t_2)}}{k^2(k_1^2+\mu^2)}}\\ &
%=Q_1Q_2'\rmint dt\, \,\Big[\boldsymbol{v}_1(t)\cdot \boldsymbol{v}_2(t)\rmint \frac{e^{-i\boldsymbol{k}\cdot [\boldsymbol{x}_2(t)-\boldsymbol{x}_1(t)]}}{(\boldsymbol{k}^2+\mu^2)}\Big]\\ &
%=\gamma Q_1Q_2'\rmint dt\, \,\Bigg[\boldsymbol{v}_1(t)\cdot \boldsymbol{v}_2(t)\rmint \frac{e^{-i\boldsymbol{k}\cdot \boldsymbol{r}}}{(\boldsymbol{k}^2+\mu^2)}\Bigg]    \,,             \\ &
=\frac{\gamma Q_1 Q_2'}{4\pi}\rmint dt\, (\boldsymbol{v}_1\cdot \boldsymbol{v}_2) \frac{e^{-\mu_{\gamma}r}}{r}+(1\leftrightarrow 2)\,\sim \mathcal{O}(L v^2)\,.
\end{split}
\end{align}
\item The amplitude corresponding to the diagram in Fig.~(\ref{fig9cm}) with bulk interaction vertex $\rmint d^4x\,\partial_l\mathcal{A}_i\partial_i\mathcal{B}_l$:
\begin{align}
    \begin{split}
   \mathcal{S}_{\text{eff}}\Big|_{\text{fig}(\ref{fig9cm})}&=\gamma \,Q_1 Q_2'\rmint\Bar{\mathcal{D}}\Hat{\xi}\rmint dt_1v_1^j(t_1)\mathcal{A}_j(\boldsymbol{x}_1(t_1))\rmint dt_2\, \mathcal{B}_m(\boldsymbol{x}_2(t_2))v_2^m(\boldsymbol{x}_2(t_2))\\ &\hspace{6.5 cm}\rmint d^4x\,\partial_l\mathcal{A}_i(x)\partial_i\mathcal{B}_l(x),\\&=\gamma Q_1 Q_2'\rmint dt_1dt_2v_1^j(t_1)v_2^m(t_2)\rmint d^4x\partial_l\Big{\langle}\mathcal{A}_j(\boldsymbol{x}_1(t_1))\mathcal{A}_i(x)\Big{\rangle}\partial_i\Big{\langle}\mathcal{B}_m(\boldsymbol{x}_2(t_2))\mathcal{B}_l(x)\Big{\rangle}\,,\\ &
%=Q_1Q_2'\textstyle{\rmint dt_1\, dt_2\,v_1^i(t_1)v_2^m(t_2)\rmint d^4x\partial_l\rmint\delta(t-t_1)\frac{e^{-i\boldsymbol{k}\cdot [\boldsymbol{x}(t)-\boldsymbol{x}_1(t_1)]}}{k^2}\delta_{lm}\partial_i\rmint\delta(t-t_2)\frac{e^{-i\boldsymbol{k}_1\cdot[\boldsymbol{x}(t)-\boldsymbol{x}_2(t_2)]}}{k_1^2+\mu^2}}\\&
%=-Q_1Q_2'\rmint dt_1\, dt_2\,dt\,v_1^i(t_1)v_2^l(t_2)\delta(t-t_1)\rmint_{\boldsymbol{k},\boldsymbol{k}_1}\delta(\boldsymbol{k}+ \boldsymbol{k}_1)\delta(t-t_2)\frac{k^lk_1^ie^{i\boldsymbol{k}\cdot \boldsymbol{x}_1(t_1)}e^{i\boldsymbol{k}_1\cdot \boldsymbol{x}_2(t_2)}}{k^2(k_1^2+\mu^2)}\\ &
%=Q_1Q_2'\rmint dt\, \,\Big[v_1^i(t) v_2^l(t)\rmint \frac{e^{-i\boldsymbol{k}\cdot [\boldsymbol{x}_2(t)-\boldsymbol{x}_1(t)]}}{k^2(\boldsymbol{k}^2+\mu^2)}k^lk^i\Big]\\ &
%=Q_1Q_2'\rmint dt\, \,\Big[v_1^i(t) v_2^l(t)\rmint \frac{e^{-i\boldsymbol{k}\cdot \boldsymbol{r}}}{k^2(\boldsymbol{k}^2+\mu^2)}k^lk^i\Bigg]                 \\ &
=\gamma Q_1Q_2'\rmint dt\, \,\Big[v_1^i(t) v_2^l(t)\
\rmint \frac{k_i k_l\,e^{-i\boldsymbol{k}\cdot \boldsymbol{r}}}{k^2(\boldsymbol{k}^2+\mu^2)}\Bigg]\,,                 \\ &
=\gamma Q_1 Q_2'\rmint dt \, \Big[(\boldsymbol{v}_1\cdot \boldsymbol{v}_2)\frac{e^{-r \mu _{\gamma }} \left(-r \mu _{\gamma }+e^{r \mu _{\gamma }}-1\right)}{2 \pi  r^3 \mu _{\gamma }^2}\\ &
\,\,\,\,\,\,\,\,\,\,\,\,\,\,\,\,\,\,\,\,\,\,\,\,\,\,\,\,\,\,\,\,\,+(\boldsymbol{v}_1\cdot \Hat{n})(\boldsymbol{v}_2\cdot \Hat{n})\frac{e^{-r \mu _{\gamma }} \Big(r \mu _{\gamma } \left(r \mu _{\gamma }+3\right)-3 e^{r \mu _{\gamma }}+3\Big)}{2 \pi  r^3 \mu _{\gamma }^2}\Big]+(1\leftrightarrow 2) \sim \mathcal{O}(L v^2)\,.
\end{split}
\end{align}
\item The amplitude corresponding to the diagram in Fig.~(\ref{fig9dm}) with bulk interaction vertex $\rmint d^4x\,\partial_k \mathcal{A}_0\partial_k \mathcal{B}_0$:
\begin{align}
    \begin{split}
       \mathcal{S}_{\text{eff}}\Big|_{\text{fig}(\ref{fig9dm})}&=\gamma \,Q_1 Q_2'\rmint\Bar{\mathcal{D}}\Hat{\xi}\rmint dt_1 \mathcal{A}_0(\boldsymbol{x}_1(t_1))\rmint dt_2\, \mathcal{B}_0(\boldsymbol{x}_2(t_2))\rmint d^4x\,\partial_k\mathcal{A}_0(x)\partial_k\mathcal{B}_0(x),\\&
       =\gamma Q_1 Q_2'\rmint dt_1dt_2\rmint d^4x\partial_k\Big{\langle}\mathcal{A}_0(\boldsymbol{x}_1(t_1))\mathcal{A}_0(x)\Big{\rangle}\partial_k\Big{\langle}\mathcal{B}_0(\boldsymbol{x}_2(t_2))\mathcal{B}_0(x)\Big{\rangle}\,,\\ &
%=Q_1Q_2'\rmint dt_1\, dt_2\rmint d^4x\partial_k\rmint\delta(t-t_1)\frac{e^{-i\boldsymbol{k}\cdot [\boldsymbol{x}(t)-\boldsymbol{x}_1(t_1)]}}{k^2}\partial_k\rmint\delta(t-t_2)\frac{e^{-i\boldsymbol{k}_1\cdot[\boldsymbol{x}(t)-\boldsymbol{x}_2(t_2)]}}{k_1^2+\mu^2}\\&
%=-Q_1Q_2'\rmint dt_1\, dt_2\,dt\,\delta(t-t_1)\rmint_{\boldsymbol{k},\boldsymbol{k}_1}\delta(\boldsymbol{k}+ \boldsymbol{k}_1)\delta(t-t_2)\frac{k^kk_1^ke^{i\boldsymbol{k}\cdot \boldsymbol{x}_1(t_1)}e^{i\boldsymbol{k}_1\cdot \boldsymbol{x}_2(t_2)}}{k^2(k_1^2+\mu^2)}\\ &
%=Q_1Q_2'\rmint dt\, \,\Big[\rmint \frac{e^{-i\boldsymbol{k}\cdot [\boldsymbol{x}_2(t)-\boldsymbol{x}_1(t)]}}{k^2(\boldsymbol{k}^2+\mu^2)}k^kk^k\Big]\\ &
=\gamma Q_1Q_2'\rmint dt\, \,\Big[\rmint \frac{e^{-i\boldsymbol{k}\cdot \boldsymbol{r}}}{(\boldsymbol{k}^2+\mu^2)}\Bigg]    \,,             \\ &
=\gamma Q_1Q_2'\rmint dt\,\frac{e^{-\mu_{\gamma}r}}{4\pi r}+(1\leftrightarrow 2)\,\sim \mathcal{O}(L v^2)\,.
    \end{split}
\end{align}
\item The amplitude corresponding to the diagram in Fig.~(\ref{fig9em}) with the bulk interaction vertex $\rmint d^4x \,\partial_k\mathcal{A}_0\partial_0 \mathcal{B}_k$:
\begin{align}
    \begin{split}
   \mathcal{S}_{\text{eff}}\Big|_{\text{fig}(\ref{fig9em})}&=\gamma \,Q_1 Q_2'\rmint\Bar{\mathcal{D}}\Hat{\xi}\rmint dt_1\mathcal{A}_0(\boldsymbol{x}_1(t_1))\rmint dt_2\, \mathcal{B}_j(\boldsymbol{x}_2(t_2))v_2^j(\boldsymbol{x}_2(t_2))\rmint d^4x\,\partial_k\mathcal{A}_0(x)\partial_0\mathcal{B}_k(x),\\&=\gamma Q_1 Q_2'\rmint dt_1dt_2v_2^j(t_2)\rmint d^4x\partial_k\Big{\langle}\mathcal{A}_0(\boldsymbol{x}_1(t_1))\mathcal{A}_0(x)\Big{\rangle}\partial_0\Big{\langle}\mathcal{B}_j(\boldsymbol{x}_2(t_2))\mathcal{B}_k(x) \Big{\rangle}\,,\\ &
%=\rmint dt_1\, dt_2\,v_2^k(t_2)\rmint d^4x\partial_k\rmint\delta(t-t_1)\frac{e^{-i\boldsymbol{k}\cdot [\boldsymbol{x}(t)-\boldsymbol{x}_1(t_1)]}}{k^2}\partial_0\rmint\delta(t-t_2)\frac{e^{-i\boldsymbol{k}_1\cdot[\boldsymbol{x}(t)-\boldsymbol{x}_2(t_2)]}}{k_1^2+\mu^2}\\&
%=-i\rmint dt_1\, dt_2\,dt\,v_2^k(t_2)\partial_0\delta(t-t_2)\rmint_{\boldsymbol{k},\boldsymbol{k}_1}\delta(\boldsymbol{k}+ \boldsymbol{k}_1)\delta(t-t_1)\frac{k^ke^{i\boldsymbol{k}\cdot \boldsymbol{x}_1(t_1)}e^{i\boldsymbol{k}_1\cdot \boldsymbol{x}_2(t_2)}}{k^2(k_1^2+\mu^2)}\\ &
%=-i\rmint dt\, \,\partial^{t_2}\Big[v_2^k(t_2)\rmint \frac{e^{-i\boldsymbol{k}\cdot [\boldsymbol{x}_2(t_2)-\boldsymbol{x}_1(t)]}\,k^k}{\boldsymbol{k}^2(\boldsymbol{k}^2+\mu^2)}\Big]_{t_2=t}\\ &
=i\gamma Q_1 Q_2'\rmint dt \Bigg[a^{k}_{2}(t)\rmint_{\boldsymbol{k}}\frac{k^k\,e^{-i\,\boldsymbol{k}\cdot\boldsymbol{r}}}{\boldsymbol{k}^2(\boldsymbol{k}^2+\mu^2)}-v_2^{k}(t)v_{2}^{a}(t)\frac{i\,k^a\,k^k\,e^{-i\boldsymbol{k}\cdot\boldsymbol{r}}}{\boldsymbol{k}^2(\boldsymbol{k}^2+\mu^2)}\Bigg]\,,\\ &
%=-i\rmint dt \partial^{t_1}\Big[v_1^i(t_1)\Big(\rmint \frac{e^{-i\boldsymbol{k}\cdot (x_2(t)-x_1(t_1)}k^i}{k^2}-\frac{e^{-i\boldsymbol{k}\cdot \vec r}k^i}{k^2+\mu^2}\Big)\Big]_{t_1=t}\\&
=-\frac{\gamma Q_1 Q_2'}{2\pi}\rmint dt \Bigg[(\Vec{a}_2\cdot \Hat{n})\Big\{\frac{  e^{-\mu_{\gamma} r}}{2 \mu_{\gamma} \,r}-\frac{1 -  e^{-\mu_{\gamma} r}}{2\, \mu_{\gamma}^2\, r^2}\Big\}+\boldsymbol{v}_2^2 \,\frac{e^{-r \mu _{\gamma }} \left(-r \mu _{\gamma }+e^{r \mu _{\gamma }}-1\right)}{  r^3 \mu _{\gamma }^2}\\ &
\,\,\,\,\,\,\,\,\,\,\,\,\,\,\,\,\,\,\,\,\,\,\,\,\,\,\,\,\,
+(\boldsymbol{v}_2\cdot \Hat{n})^2\,\frac{e^{-r \mu _{\gamma }} \Big(r \mu _{\gamma } \left(r \mu _{\gamma }+3\right)-3 e^{r \mu _{\gamma }}+3\Big)}{  r^3 \mu _{\gamma }^2}
\Bigg]+(1\leftrightarrow 2)\,\sim \mathcal{O}(L v^2)\, \label{4.52}
\end{split}
\end{align}
\end{itemize}

Before we conclude this section, we summarize our results in Table~(\ref{table2mm}). We show the number of diagrams responsible for the corrections to the conservative dynamics due to the various fields and their interactions, apart from the pure gravity contributions. %We also show the order at which $g_{a\gamma\gamma}$ and $\gamma$ enters in the potential correction. \\


\begin{table}[b!]
 \centering
 \scalebox{0.91}{\begin{tabular}{|c|c|c|c|c|}
 \hline
 & $0 PN$ & $1PN$ &  $2PN$ &  $2.5PN$ \\
  \hline
 \textit {scalar} & - & 5 & 6 &-\\
  \hline
 \textit{electromagnetic} & 1 & 3 & - &-\\
 \hline 
 \textit{proca} & 1 & 3 & - &- \\
 \hline 
 \textit{gravitational} & - & 4 & - &-\\
 \hline
 \textit{scalar-em interaction}
 & - & - & - &2($\sim g_{a\gamma\gamma}$)\\
 \hline
 \textit{proca-em interaction}&- & 4 $(\sim \gamma)$ & - &-\\
 \hline
 \end{tabular}}
 \caption{Table showing the number of diagrams contributing to a specific PN order for bound sector. Also, the orders at which terms to due the axion-photon coupling ($g_{a\gamma\gamma}$) and photon-dark photon kinetic mixing term ($\gamma$).}
 \label{table2mm}
 \end{table}
\vspace{0.5cm}
\textcolor{black}{{\bf Corrections to the bound orbit}: \textcolor{black}{Now we can add all the contributions mentioned in 
 (\ref{4.13a})-(\ref{4.52}) and that gives us the real part of the total effective action. To find out the orbit equation, one needs to extremize this.}
\begin{eqnarray}
  &  \frac{\delta}{\delta \boldsymbol{x}_a(t)}\,\text{Re}\,\mathcal{S}_{\text{eff}}^{\text{tot}}[\boldsymbol{x}_a]\Big |_{\mathcal{O}(\hbar^0)}=0\,.\label{4.53m}
\end{eqnarray}}
The leading order contribution to effective action is the Newtonian potential. That gives rise to the Keplerian orbit as the solution of (\ref{4.53m}). This we use to study the radiative dynamics in Sec.~(\ref{Sec5}). Note that we will get corrections to the usual Keplerian orbit solution due to other terms in the effective action at higher-order PN. Although one should consider these corrections and find out the corrected orbit equation and then use it to study the radiative dynamics, in this paper, we will neglect such corrections and work with the leading order solution, i.e. Keplerian orbit solution (in fact, the circular orbit). 
 

\section{The radiative dynamics}\label{Sec5}
Till now, we have studied the conservative dynamics, where no external radiative line is attached to the diagrams. In the radiative sector, we keep one long-wavelength field external and integrate out only the potential modes. For a generic field $\xi=\Hat{\xi}+\Bar{\xi}$, this produces an effective action for the radiation field of the form
\begin{equation}
\mathcal{S}_{\text{eff}}^{\Bar{\xi}}[\boldsymbol{x}_a,\Bar{\xi}]=S_{\text{free}}^{\Bar{\xi}}+\rmint d^4 x\,J(x)\Bar{\xi}(x)\,,
\end{equation}
\noindent where $J(x)$ is the source that encodes the binary dynamics relevant for emission. By comparing with (\ref{2.4m}) one can identify this source with the corresponding Wilson coefficient. Expanding around $\bar\xi=0$,
\begin{align}
    \begin{split}
       & \mathcal{S}_{\text{eff}}^{\Bar{\xi}}[\boldsymbol{x}_a,\Bar{\xi}]= \mathcal{S}_{0}[\boldsymbol{x}_a]+\mathcal{S}_{1}[\boldsymbol{x}_a,\Bar{\xi}]+\mathcal{S}_{2}[\boldsymbol{x}_a,\Bar{\xi}]+\mathcal{O}(\bar\xi^3).\label{5.1m}
    \end{split}
\end{align}
The first term in (\ref{5.1m}) is the conservative action, the second determines the radiation source, and the higher powers describe self-interactions of the radiative field. To compute the emitted power we integrate out the radiation field itself and obtain the effective action for the source,
\begin{align}
\begin{split}
   \mathcal{Z}[J]:= e^{i{\gamma}_{\text{eff}}[J]}&=\rmint \mathcal{D}\Bar{\xi}\rmint \mathcal{D}\Hat{\xi}\,e^{iS[{\xi}]}\,,\\ &
   =\rmint \mathcal{D}\Bar{\xi}\, e^{i\mathcal{S}_{\text{eff}}^{\Bar{\xi}}[\Bar{\xi},J]}\,.
   \label{4.3}
   \end{split}
\end{align}
Again, in the in-out formalism $\mathcal{Z}[J]$ is the overlap between in and out states,
\begin{eqnarray}
    \mathcal{Z}[J]=\Big\langle0_+|0_-\Big\rangle.\label{4.4}
\end{eqnarray}
using the two relations (\ref{4.3}) and (\ref{4.4}) we get \cite{nastase_2019},
\begin{eqnarray}
   | \langle0_+|0_-\rangle|^2=e^{-2\text{Im}[\gamma_{\text{eff}}]}\,.
\end{eqnarray}
Expanding in small $\text{Im}[\gamma_{\text{eff}}]$ we will get,
\begin{align}
    \begin{split}
        2\,\text{Im}[\gamma_{\text{eff}}]=\mathcal{T}\rmint dE\,d\Omega \frac{d^2\Gamma}{dEd\Omega}\label{4.6}
    \end{split}
\end{align}
and the radiated power is given by,
\begin{align}
    \begin{split}
        P=\rmint dE\,d\Omega\,E\,\frac{d^2\Gamma}{dEd\Omega}\label{5.6m}
    \end{split}
\end{align}
where, $\mathcal{T}$ is the interaction time, $d\Gamma$ is the rate of particle emission and $P$ is the radiated power. The practical workflow in every radiative subsection is therefore the following:
\begin{itemize}
    \item identify the diagrams with one external radiation line,
    \item read off the source term or the corresponding multipole moments,
    \item use $\text{Im}\,\gamma_{\text{eff}}$ together with (\ref{5.6m}) to obtain the radiated power.
\end{itemize}
\subsection{Radiative sector PN counting} \label{EFTcounting}
 Now, to compute the radiated power up to a given \textcolor{black}{$N^{(n)}LO$}, we need to systematically assign EFT power-counting rules to the radiative fields. To make things more compact, we display the scaling of the radiation fields $\bar{\xi}\equiv\{\bar\psi,\bar{{\mathcal{A}}}_i,\bar\sigma_{ij},\bar{\phi},\Bar{b}_{\mu}, \Bar{a}_{\mu}\}$ in Table~(\ref{table3mm}). One should note that, for the potential fields, one must perform a partial Fourier transform. Therefore,
$$\hat\xi(x^0,\boldsymbol{x})=\rmint_k e^{i\boldsymbol{k}\cdot \boldsymbol{x}}\,\hat\xi_k(x^0)$$
where $\xi_k\equiv \{\psi_k,\Hat{\mathcal{A}}_{{\boldsymbol{k}}i},\zeta_{{\boldsymbol{k}}ij}, \mathcal{A}_{\boldsymbol{k}i},\mathcal{B}_{\boldsymbol{k}i},\varphi_{\boldsymbol{k}}\}$ are the potential fields. Here in this sector, we define the diagrams as $N^{(n)}LO$, i.e., leading order, next to leading order, and so on. The specification of scaling for the radiative sector is denoted in terms of these notations, as scaling rules are a bit subtle and different in the radiative sector. 
\begin{table}[ht]
\begin{center}
\begin{tabular}{|c|c|}
    \hline
   \multicolumn{2}{|c|}{EFT scaling} \\
\hline
Fields & scaling\\
\hline
t &  $\sim\frac{r}{v}$\\
\hline
  $\delta x$ &$\sim r$\\
  \hline
  $\frac{m}{M_{p}}$ & $\sim L^\frac{1}{2}v^\frac{1}{2}$\\
  \hline
  $Q^2$ & $ \sim L\,v$\\
  \hline
  $\xi_k$ & $\sim r^2v^\frac{1}{2}$\\
  \hline
$\bar{\Xi}$  & $\sim\frac{v}{r}$\\
\hline
$\partial_0\xi_k$ & $\sim\frac{v}{r}\xi_k$\\
\hline
  $\partial_{\mu}\bar{\Xi}$ & $\sim\frac{v}{r}\bar{\Xi}$\\
  \hline 
\end{tabular}
\end{center}
\caption{EFT PN counting for radiative sector}
\label{table3mm}
\end{table} 
\par  \textcolor{black}{Now expanding the metric at the linear level the  relation between NRG fields and the linearized metric components in the radiative sector is given by:$$
\bar g_{\mu\nu}=\eta_{\mu\nu}+\frac{1}{m_p}
\begin{pmatrix}
-2\bar\psi & \bar{\mathcal{A}_j}\\
\bar{\mathcal{A}_i} & \,\,\bar{\sigma_{ij}}-2\delta_{ij}\bar\psi\\
\end{pmatrix}\equiv\eta_{\mu\nu}+\frac{1}{m_p}
\begin{pmatrix}
\bar h_{00} & \bar h_{0j}\\
\bar h_{i0} & \bar h_{ij}\\
\end{pmatrix}
$$
The gravitational radiation formula we used in the later section in (\ref{5.65}) is basically considering the $\bar\psi$ radiation. One subtle point to note here is that when we calculate radiation usually we consider the TT gauge along with the harmonic gauge condition. But from the beginning, we can't set $\bar h_{00}=0$ in the effective Lagrangian for radiation demanded by \textit{Transverse Traceless (TT)} gauge condition.}

\subsection{Scalar interaction and scalar radiation}\label{subsec:radiated_scalars}

We begin with the scalar sector because it already contains all the conceptual ingredients that reappear later: a non-trivial source, a multipole expansion, and the extraction of power through the optical theorem. The strategy is to first derive a compact emission formula in terms of scalar multipoles and then compute the source $J_\phi$ order by order in the PN expansion. In the main text, we therefore keep the multipole logic explicit while compressing the more repetitive tensor algebra. To begin with, we integrate the potential scalar modes and get the effective action for the radiating scalars.  
\begin{equation}
\mathcal{S}_{\text{eff}}^{\Bar{\phi}}=  \rmint d^4x  \left(\underbrace{ -\frac{1}{2} \eta^{\mu \nu}  \partial_\mu \bar{\phi} \partial_\nu \bar{\phi} 
-\frac{1}{2}m^2\bar{\phi}^2 }_{\mathcal{S}_{\rm free}^{\bar\phi}} + \underbrace{\frac{1}{m_p} J   \bar{\phi}}_{\mathcal{S}_{\text{int}}^{\bar\phi}} \right)\;, 
\end{equation}
which leads to the following equation of motion 
\begin{equation}
\textcolor{black}{(\square-m^2 )\bar{\phi} = \frac{J}{m_p} \;, \qquad \square \equiv \eta^{\mu \nu}  \partial_\mu  \partial_\nu \;.}
\end{equation}
We will now calculate the imaginary part of the effective action $\gamma_{\text{eff}}^{\bar\phi}$ for the source obtained by integrating out the radiation scalars as discussed in Eq.~(\ref{4.3}) 
%First, we write down the effective action for the radiating scalar. Then, integrating over radiative modes, we get the dynamics of the source term $J$, a
and then using the optical theorem, we obtain the power radiation for the scalar sector. We will work out the power radiation essentially for circular orbits. Now,
\begin{equation}
\begin{split}
\mathcal{S}_{\rm int}^{\bar\phi} = \rmint d^4 x \frac{J \bar \phi }{m_p}= \rmint dt \, \rmint d^3 x \frac{J (t, \boldsymbol{x})}{m_p}  \bigg( \bar{\phi}(t, \boldsymbol{0}) + x^i \partial_i \bar{\phi}(t,\boldsymbol{0}) + \frac{1}{2} x^i x^j \partial_i \partial_j \bar{\phi}(t, \boldsymbol{0})  \\
+ \frac{1}{3!} x^i x^j x^k \partial_i \partial_j \partial_k \bar{\phi}(t, \boldsymbol{0}) + \ldots \bigg) \;,
\label{5.10}
\end{split}
\end{equation}
From which we can write,
\begin{align}
\begin{split}
S_{\rm int}^{\bar\phi} = \frac{1}{m_p} \rmint dt \left( I_\phi \bar{\phi} + I_\phi^i \partial_i \bar{\phi} + \frac{1}{2} I_\phi^{ij} \partial_i \partial_j \bar{\phi}  + \ldots \right) \;,
\end{split}
\label{eq:multipole_expansion_scalar}
\end{align}
where 
\begin{equation}
\label{5.12mn}
\textcolor{black}{I_\phi  \equiv  \rmint d^3 x \left( J + \frac{1}{6}  (\partial_t^2+m^2) J  x^2 \right) \;, \quad I_\phi^i  \equiv \rmint d^3x \, x^i \left(J + \frac{1}{10} (\partial_t^2+m^2) J x^2 \right) \;, \quad  I_\phi^{ij} \equiv \rmint d^3x J  Q^{ij} \;}
\end{equation}
are respectively the scalar monopole, dipole and quadrupole.
\textcolor{black}{ To get the irreducible forms of the multipole moments, instead of writing $x^ix^j$ one should use,
\begin{align}
    \begin{split}
        {Q}^{ij}=x^ix^j-\frac{1}{3}x^2\delta^{ij}\,.
    \end{split}
\end{align}}
Then we integrate over the scalar radiation field and get the effective action for the source, which is equivalent to evaluating the self-energy diagram shown in Fig.~(\ref{mfig12}).
\begin{figure}
    \centering
\thesisdiagrampanel{\scalebox{0.36}{\begin{feynman}
    \electroweak[flip=true, lineWidth=4]{4.60, 4.10}{5.60, 4.80}
    \fermion[]{4.00, 4.00}{7.20, 4.00}
    \electroweak[lineWidth=4, label=$\bar\phi$]{5.60, 4.80}{6.60, 4.10}
    \fermion[]{4.00, 4.10}{7.20, 4.10}
\end{feynman}
}}
    \caption{Self-energy diagram contributing to the effective action for the source term $J$.}
    \label{mfig12}
\end{figure}
To quadratic order in the source moments, the path integral reduces to
\begin{align}
\begin{split}
i\gamma_{\text{eff}}^{\bar\phi}[J]\sim -\frac{1}{2m_p^2}\rmint dt_1 dt_2 \Big[&
I_{\phi}(t_1)I_{\phi}(t_2)\Big\langle\Bar{\phi}(t_1,\boldsymbol{0})\Bar{\phi}(t_2,\boldsymbol{0})\Big\rangle\\ &
+I_{\phi}^{i}(t_1)I_{\phi}^{j}(t_2)\Big\langle\partial_{i}\Bar{\phi}(t_1,\boldsymbol{0})\partial_{j}\Bar{\phi}(t_2,\boldsymbol{0})\Big\rangle\\ &
+I_{\phi}^{ij}(t_1)I_{\phi}^{kl}(t_2)\Big\langle\partial_{i}\partial_{j}\Bar{\phi}(t_1,\boldsymbol{0})\partial_{k}\partial_{l}\Bar{\phi}(t_2,\boldsymbol{0})\Big\rangle+\ldots\Big]\,.
\label{5.14n}
\end{split}
\end{align}
Odd-point correlators vanish because we expand around a background with $\langle 0 |\bar\phi|0\rangle=0$.\footnote{It would be interesting to revisit this computation when the radiative scalar has a non-zero vacuum expectation value, i.e. $\langle 0 |\bar\phi|0\rangle\ne 0$.}
Using the Feynman propagator
\begin{align}
    \begin{split}
        \Big\langle T \Bar{\phi}(t_1,\boldsymbol{x}_1)\Bar{\phi}(t_2,\boldsymbol{x}_2)\Big\rangle=\rmint \frac{d^4k}{(2\pi)^4}\frac{i}{-k^2-m^2+i\epsilon}e^{-ik\cdot(x_1-x_2)}\,,\label{4.13}
    \end{split}
\end{align}
and rotational invariance, the derivative correlators reduce to
\begin{align}
    \begin{split}
        \Big\langle T \partial_i \Bar{\phi}(t_1,\boldsymbol{0})\partial_j \Bar{\phi}(t_2,\boldsymbol{0})\Big\rangle&=\rmint \frac{d^4k}{(2\pi)^4}\frac{i}{-k^2-m^2+i\epsilon}e^{ik_0(t_1-t_2)}[\partial_{i}^{(x_1)}e^{-i\boldsymbol{k}.\boldsymbol{x}_1}][\partial_{j}^{(x_2)}e^{i\boldsymbol{k}\cdot\boldsymbol{x}_2}]\big|_{\boldsymbol{x}_1,\boldsymbol{x}_2=0}\\ &
        =\rmint \frac{d^4k}{(2\pi)^4}\frac{i}{-k^2-m^2+i\epsilon}e^{ik_0(t_1-t_2)}\langle\langle k_{i}\,k_{j}\rangle\rangle\,,\\ &
        =\frac{1}{3}\rmint \frac{d^4k}{(2\pi)^4}\frac{i}{-k^2-m^2+i\epsilon}e^{ik_0(t_1-t_2)}|\boldsymbol{k}|^2\delta_{ij}\,,\label{4.14}
    \end{split}
\end{align}
\begin{align}
    \begin{split}
          \Big\langle T \,\partial_i\partial_j \Bar{\phi}(t_1,\boldsymbol{0})\,\partial_k\partial_l \Bar{\phi}(t_2,\boldsymbol{0})\Big\rangle&=\rmint \frac{d^4k}{(2\pi)^4}\frac{i}{-k^2-m^2+i\epsilon}e^{ik_0(t_1-t_2)}[\partial_{i}\partial_j^{(x_1)}e^{-i\boldsymbol{k}.\boldsymbol{x}_1}][\partial_{k}\partial_l^{(x_2)}e^{i\boldsymbol{k}\cdot\boldsymbol{x}_2}]\big|_{\boldsymbol{x}_1,\boldsymbol{x}_2=0}\,,\\ &
          =\rmint \frac{d^4k}{(2\pi)^4}\frac{i}{-k^2-m^2+i\epsilon}e^{ik_0(t_1-t_2)}\langle\langle k_ik_jk_kk_l\rangle\rangle\,,\\ &
          =\rmint \frac{d^4k}{(2\pi)^4}\frac{i}{-k^2-m^2+i\epsilon}e^{ik_0(t_1-t_2)}|\boldsymbol{k}|^4 T_{ij,kl}\,,\label{4.15}
    \end{split}
\end{align}
where $T_{ij,kl}$ is constructed from $\langle\langle n_in_jn_kn_l\rangle\rangle$ and $\langle\langle n_in_j \rangle\rangle$ which is symmetric in $ij$ and $kl$ indices \cite{poissonwill}. Then finally we get, 
\begin{equation}
  \gamma_{\text{eff}}^{(\bar\phi)} = \frac{1}{2 m_p^2} \rmint \frac{d^4 k}{(2 \pi)^4} \frac{1}{-k^2-m^2 +i \epsilon} \bigg( | \mathcal{I}_\phi (k_0)|^2 + \frac{1}{3} |\boldsymbol{k}|^2   | \mathcal{I}_\phi^i (k_0)|^2 + \frac{1}{30} |\boldsymbol{k}|^4   | \mathcal{I}_\phi^{ij} (k_0)|^2  \bigg) \;, \label{5.18n}
\end{equation}
where, 
\begin{equation}
\mathcal{I}_\phi(k_0) = \rmint dt \,\mathcal{I}_\phi (t) e^{i k_0 t} \;, \mathcal{I}_{\phi}^{i}(k_0) = \rmint dt \,\mathcal{I}_\phi^i (t) e^{i k_0 t} \;, \mathcal{I}_{\phi}^{ij}(k_0) = \rmint dt \,\mathcal{I}_\phi^{ij} (t) e^{i k_0 t} \;.
\end{equation}
To extract the imaginary part, we need to find  the principal value of the propagator i.e.,
\begin{align}
    \begin{split}
        \frac{1}{-k^2-m^2+i\epsilon}=PV\Big(\frac{1}{-k^2-m^2}\Big)-i\pi\,\delta(-k^2-m^2)\,.
    \end{split}
\end{align}

Therefore, besides the quadrupole,  the monopole and the dipole also contribute to the scalar radiation. 
But for the massive case, we will have some changes in the quadrupole formula compared to \cite{Kuntz:2019zef}.
\begin{align}
    \begin{split}
        P_{\phi}^m=\frac{1}{4\pi^2 m_p^2\mathcal{T}}\sum_{l}\frac{1}{l!(2l+1)!!}\rmint_0^\infty d\omega\,\omega\,(\omega^2-m^2)^{l+1/2}|\mathcal{I}^{m}_{(L)}(\omega)|^2 .
        \label{5.21}
    \end{split}
\end{align}
%\textcolor{black}{After that, we find the power emitted  into scalars for the massless case is given by} \cite{Kuntz:2019zef}\,,
%\begin{equation}
%P_\phi = {2G}\left[  \big\langle \dot{\mathcal{I}}_\phi^2  \big\rangle  + \frac{1}{3} \big\langle  \ddot{\mathcal{I}}_\phi^i \mathcal{\ddot I}_\phi^i  \big\rangle + \frac{1}{30} \big\langle \mathcal{\dddot I}_\phi^{ij}  {\mathcal{\dddot I}}_\phi^{ij}   \big\rangle \right] \;.
%\label{eq:radiated_power_scalar}
%\end{equation}

We will truncate our computation up to $l=1$, i.e., we will compute only the dipole radiation. Here in (\ref{5.21}) $`L$' denotes the \textit{symmetric trace free (STF)} indices $i,j,k,\cdots.$

\subsubsection*{Computing the source term $J_{\phi}$:}
%The effective action can be calculated by integrating out the electromagnetic potential fields while taking the other radiative fields to be zero i.e,
%\begin{align}
%    \begin{split}
 %       \mathcal{Z}[J_\phi]=\rmint \mathcal{D}\tilde{\psi} \mathcal{D}\hat{\mathcal{A}}_{i}\mathcal{D}\varphi \mathcal{D}\mathcal{A}_{\mu}\mathcal{D}\mathcal{B}_{\nu}\mathcal{D}\boldsymbol{\zeta}_{ij}\,e^{iS_{T}+\chi\cdot J}\Big|_{J_{\Bar{a},\Bar{h},\Bar{b}=0}}
%    \end{split}
%\end{align}
We now construct the scalar source only up to the orders needed for the dipole-radiation formula. The diagrams in Fig.~(\ref{Figure10}) can be grouped by PN order, and in the main text we quote the reduced amplitudes after the time and tensor algebra has been simplified. This makes it easier to see which physical effects first enter the source at each order.
\begin{itemize}
\item The $LO$ contribution comes from the worldline vertex $\rmint dt \,m_a s_a \Bar{\phi}.$ The amplitude corresponding to the diagram in Fig.~(\ref{10a}):
\begin{align}
    \begin{split}
\mathcal{S}_{\text{eff}}\Big|_{\text{fig}(\ref{10a})}=-\sum_{a=1}^{2}\rmint dt\,\frac{m_a\,s_a}{m_p}\,\Bar{\phi}(\boldsymbol{x}_a(t))\,\sim\mathcal{O}(L^{1/2}v^{1/2}).
    \end{split}
\end{align}
\begin{figure}
    \centering
  \begin{subfigure}[t]{0.22\textwidth}
\thesisdiagrampanel{\scalebox{0.3}{\begin{feynman}
    \electroweak[lineWidth=6, color=0693e3, label=$\Bar{\phi}$]{6.00, 6.00}{7.00, 7.00}
    \fermion[lineWidth=6]{4.00, 6.00}{8.00, 6.00}
    \fermion[lineWidth=6]{4.00, 4.00}{8.00, 4.00}
\end{feynman}
}}
      \caption{}
      \label{10a}
  \end{subfigure}
   \begin{subfigure}[t]{0.22\textwidth}
 \thesisdiagrampanel{\scalebox{0.3}{\begin{feynman}
    \electroweak[lineWidth=6, color=0693e3, label=$\varphi$]{5.60, 4.00}{5.60, 6.00}
    \electroweak[lineWidth=6, label=$\Bar{\phi}$, color=0693e3]{5.60, 6.00}{6.80, 6.80}
    \fermion[lineWidth=6]{4.00, 6.00}{7.60, 6.00}
    \fermion[lineWidth=6]{4.00, 4.00}{7.60, 4.00}
\end{feynman}}}
      \caption{}
      \label{10b}
  \end{subfigure}
   \begin{subfigure}[t]{0.20\textwidth}
 \thesisdiagrampanel{\scalebox{0.25}{\begin{feynman}
    \fermion[lineWidth=6]{4.00, 4.00}{7.80, 4.00}
    \fermion[lineWidth=6]{4.00, 6.40}{7.80, 6.40}
    \electroweak[lineWidth=6, label=$\tilde{\psi}$]{5.80, 6.40}{5.80, 4.00}
    \electroweak[lineWidth=6, label=$\Bar{\phi}$, color=0693e3]{5.80, 6.40}{7.60, 7.60}
\end{feynman}
}}
      \caption{}
      \label{10c}
  \end{subfigure}
  \begin{subfigure}[t]{0.22\textwidth} \centering
        \thesisdiagrampanel{\scalebox{0.3}{
\begin{feynman}
    \electroweak[label=$\Bar{\phi}$, lineWidth=6, color=0693e3]{5.40, 5.00}{7.60, 5.00}
    \fermion[lineWidth=6]{4.00, 4.00}{6.80, 4.00}
    \electroweak[label=$\varphi$, lineWidth=6, color=0693e3]{5.40, 6.00}{5.40, 5.00}
    \electroweak[label=$\Tilde{\psi}$, lineWidth=6]{5.40, 5.00}{5.40, 4.00}
    \fermion[lineWidth=6]{4.00, 6.00}{6.80, 6.00}
\end{feynman}}}
\caption{}
\label{10h}
\end{subfigure}
  \begin{subfigure}[t]{0.20\textwidth}
 \thesisdiagrampanel{\scalebox{0.25}{\begin{feynman}
    \electroweak[label=$\varphi$, lineWidth=6, color=0693e3]{5.60, 4.00}{5.60, 5.00}
    \electroweak[label=$\bar{\phi}$, color=0693e3, lineWidth=6]{5.60, 6.40}{7.40, 7.80}
    \electroweak[label=$\varphi$, lineWidth=6, color=0693e3]{5.60, 6.40}{5.60, 5.40}
    \fermion[lineWidth=6]{4.00, 6.40}{7.20, 6.40}
    \fermion[lineWidth=6]{4.00, 4.00}{7.20, 4.00}
    \parton{5.60,5.20}{0.20}
\end{feynman}
}}
      \caption{}
      \label{10d}
  \end{subfigure}
%    \caption{}
   % \label{fig:my_label}
     \begin{subfigure}[t]{0.22\textwidth}
        \centering
        \thesisdiagrampanel{\scalebox{0.25}{
\begin{feynman}
    \electroweak[color=0693e3, lineWidth=6, label=$\phi$]{5.60, 5.20}{5.60, 4.00}
    \fermion[lineWidth=6, showArrow=false]{4.00, 6.40}{7.20, 6.40}
    \electroweak[color=0693e3, lineWidth=6, label=$\bar{\phi}$]{5.60, 5.20}{8.20, 5.20}
    \electroweak[lineWidth=6, label=$\psi$]{5.60, 6.40}{5.60, 5.20}
    \fermion[lineWidth=6, showArrow=false]{4.00, 4.00}{7.20, 4.00}
\end{feynman}
}}
\caption{}
\label{10e}
\end{subfigure}
    \begin{subfigure}[t]{0.22\textwidth}
        \centering
        \thesisdiagrampanel{\scalebox{0.25}{\begin{feynman}
    \electroweak[color=0693e3, lineWidth=6, label=$\varphi$]{5.60, 5.20}{5.60, 4.00}
    \fermion[lineWidth=6, showArrow=false]{4.00, 6.40}{7.20, 6.40}
    \gluon[lineWidth=6, label=$\hat{\mathcal{A}_i}$, color=fcb900]{5.60, 6.40}{5.60, 5.20}
    \electroweak[color=0693e3, lineWidth=6, label=$\bar{\phi}$]{5.60, 5.20}{8.20, 5.20}
    \fermion[lineWidth=6, showArrow=false]{4.00, 4.00}{7.20, 4.00}
\end{feynman}
}}
\caption{}
\label{10f}
\end{subfigure}
\begin{subfigure}[t]{0.22\textwidth} \centering
        \thesisdiagrampanel{\scalebox{0.3}{
 \begin{feynman}
    \electroweak[label=$A_m$, color=9900ef, lineWidth=6]{5.40, 5.00}{5.40, 4.00}
    \electroweak[label=$A_i$, color=9900ef, lineWidth=6]{5.40, 6.00}{5.40, 5.00}
    \fermion[lineWidth=6]{4.00, 6.00}{7.00, 6.00}
    \fermion[lineWidth=6]{4.00, 4.00}{7.00, 4.00}
    \electroweak[label=$\bar{\phi}$, lineWidth=6, color=0693e3]{5.40, 5.00}{7.60, 5.00}
\end{feynman}}}
\caption{}
\label{10g}
\end{subfigure}
\caption{Radiative diagrams for the scalar field that contribute up to $N^{(4)}LO$.}
\label{Figure10}
\end{figure}
\item The $N^{(1)}LO$ contribution comes from the worldline vertex $\rmint dt\, \frac{1}{2}\,v_a^2\, m_as_a \frac{\phi}{m_p}$, and its contribution to the effective action has the following form,
\begin{align}
    \begin{split}
        \mathcal{S}_{\text{eff}}=\frac{1}{2}\sum_{a}m_a s_a v_a^2\,\Bar{\phi}(\boldsymbol{x}_a(t_a))\sim \mathcal{O}(L^{1/2}v^{3/2}).
    \end{split}
\end{align}

\item The $N^{(2)}LO$ contribution comes from the vertex: $-\rmint dt\, g_a \varphi\, \Bar{\phi}$ and $-\rmint dt\,s_a \varphi$. The amplitude corresponds to the diagram in Fig.~(\ref{10b}):
\begin{align}
    \begin{split}
          \mathcal{S}_{\text{eff}}\Big|_{\text{fig}(\ref{10b})}&=\frac{m_1 m_2}{m_p^2}\rmint\Bar{\mathcal{D}}\Hat{\xi}\rmint dt_1\,g_1 \varphi(\boldsymbol{x}_1(t_1))\frac{\Bar{\phi}(\boldsymbol{x}_1(t_1))}{m_p}\rmint dt_2\,s_2\,\varphi(\boldsymbol{x}_2(t_2))\,,\\ &
        =\frac{4 \pi  g_1 s_2 m_1 m_2}{m_p^2} \rmint dt\rmint_{\boldsymbol{k}}\frac{e^{i\boldsymbol{k}\cdot \boldsymbol{r}}}{\Vec{k^2}+m^2}\frac{\Bar{\phi}(\boldsymbol{x}_1(t_1))}{m_p}\,,\\ &
        =\frac{g_1 s_2 m_1 m_2}{m_p^2}\rmint dt \frac{e^{-mr}}{r}\frac{\Bar{\phi}(\boldsymbol{x}_1(t))}{m_p}+(1\leftrightarrow 2)\,\sim\mathcal{O}(L^{1/2}v^{5/2}).
    \end{split}
\end{align}

\item \textcolor{black}{Another $N^{(2)}LO$ contribution for the source comes from the worldline vertex:
-$\rmint dts_a m_a\bar{\phi}\tilde{\psi}\,.$}
The amplitude corresponds to the diagram in Fig.~(\ref{10c}):

\begin{align}
    \begin{split}
        \mathcal{S}_{\text{eff}}^{0}\Big|_{\text{fig}(\ref{10c})}&=\frac{s_1m_1m_2}{m_p^3}\rmint\Bar{\mathcal{D}}\Hat{\xi}\rmint dt_1 \,\frac{\Bar{\phi}(\boldsymbol{x}_1(t_1))}{m_p}\tilde{\psi}(\boldsymbol{x}_1(t_1))\rmint dt_2 \tilde{\psi}(\boldsymbol{x}_2(t_2))\,,\\ &
        =\frac{s_1m_1m_2}{m_p^3}\rmint dt_1\,dt_2\,\Bar{\phi}(\boldsymbol{x}_1(t_1))\Big\langle\tilde{\psi}(\boldsymbol{x}_1(t_1))\tilde{\psi}(\boldsymbol{x}_2(t_2))\Big\rangle\,, \\ &
       % =-4\pi G\, s_1\,m_1 (\frac{-m_2}{2})\rmint dt_1\,dt_2\, \delta(t_1-t_2)\,\Bar{\phi}(x_1,t_1)\rmint_{\boldsymbol{k}}\frac{e^{i\boldsymbol{k}\cdot[\boldsymbol{x}_1(t_1)-\boldsymbol{x}_2(t_2)]}}{\boldsymbol{k}^2}\,,\\ &
       % =-2\pi G\, s_1\, m_1(\frac{-m_2}{2})\rmint dt \rmint_{\boldsymbol{k}}\frac{e^{i\boldsymbol{k}\cdot\boldsymbol{r}}}{\boldsymbol{k}^2}\Bar{\phi}(x_1,t)\\ &
       % =-\frac{1}{2} G\, s_1\,m_1(\frac{-m_2}{2})\rmint dt\, \frac{1}{r}\,\Bar{\phi}(x_a,t)\\ &
        =\frac{s_1m_1m_2}{8\pi m_p^2}\rmint dt\, \frac{1}{r} \frac{\Bar{\phi}(\boldsymbol{x}_1(t))}{m_p}+(1\leftrightarrow 2)\,\,\sim\mathcal{O}({L^{1/2}v^{5/2}}).
    \end{split}
\end{align}
\item  The $N^{(3)}LO$ contribution coming from vertex $m^2 \rmint d^4x \Tilde{\psi}\varphi\Bar{\phi}(x)$ corresponding to the diagram in Fig.~(\ref{10h}), contributes to the effective action in the following way:
\begin{align}
    \begin{split}
\mathcal{S}_{\text{eff}}\Big |_{\text{fig}(\ref{10h})}&=\frac{m^2 s_1m_1m_2}{m_p^3}\rmint dt_1dt_2\rmint d^4x \Big\langle \psi(\boldsymbol{x}_2(t_2))\psi(x)\Big\rangle \Big\langle\varphi(\boldsymbol{x}_1(t_1))\varphi(x) \Big\rangle \bar{\phi}(x_0)\,,\\&
=\frac{s_1m_1m_2}{8\pi m_p^2}\rmint dt \frac{(1-e^{-mr})}{r}\frac{\bar{\phi}(\boldsymbol{x}_0(t))}{m_p}+1\leftrightarrow 2\,\sim\mathcal{O}(L^{1/2}v^{7/2}).
    \end{split}
\end{align}
\item \textcolor{black}{One $N^{(4)}LO$ correction comes from the relativistic time correction shown in Fig.~(\ref{10d}). After the time integrations and an integration by parts, the reduced contribution is}
\begin{align}
    \begin{split}
\mathcal{S}_{\text{eff}}\Big|_{\text{fig}(\ref{10d})}
=\frac{m_1 m_2 g_1 s_2}{16 \pi m_p^2}\rmint dt \,
\mathcal{K}_{10d}(t)\,
\frac{\Bar{\phi}(\boldsymbol{x}_1(t))}{m_p}+ T(\Dot{\Bar{\phi}})+ (1 \leftrightarrow 2)\,\sim\mathcal{O}(L^{1/2}v^{9/2})\,.
\label{5.26}
    \end{split}
\end{align}
with
\begin{align}
    \begin{split}
        \mathcal{K}_{10d}(t)=e^{-mr}\Bigg[
        \frac{\boldsymbol{v}_1\cdot\boldsymbol{v}_2}{r}
        -\frac{(\boldsymbol{v}_1\cdot \hat{n})(\boldsymbol{v}_2\cdot \hat{n})}{r}
        +\frac{m}{2}\,(\boldsymbol{v}_2\cdot \Hat{n})(\boldsymbol{v}_1\cdot \hat{n})
        \Bigg].
    \end{split}
\end{align}
Here, 
\begin{align}
    \begin{split}
        T(\Dot{\Bar{\phi}})&=-\frac{m_1 m_2 g_1 s_2}{16\pi m_p^2}\rmint dt\, e^{-mr}(\boldsymbol{v}_2\cdot \Hat{n})\frac{ \Dot{\Bar{\phi}}(\boldsymbol{x}_1(t))}{m_p},\\ &
        =\frac{m_1 m_2 g_1 s_2}{16\pi m_p^2}\rmint dt \frac{e^{-mr}}{r}(\boldsymbol{v}_2\cdot \Dot{\boldsymbol{r}}+\Vec{a}_2\cdot \boldsymbol{r})\,\frac{\Bar{\phi}(\boldsymbol{x}_1(t))}{m_p}\,.
    \end{split}
\end{align}
We have ignored the total derivative term by assuming the time variation of $\Bar{\phi}$ is very small and $|\Dot{\boldsymbol{r}}|=0$ for circular orbits.
\item Another $N^{(4)}LO$ term comes from the bulk interaction $\rmint d^4x \psi\partial_0\phi\partial_0\bar{\phi}$ shown in Fig.~(\ref{10e}). Writing $\Delta\boldsymbol{v}:=\boldsymbol{v}_2-\boldsymbol{v}_1$ and $\Delta\Vec{a}:=\Vec{a}_2-\Vec{a}_1$, its reduced contribution is
\begin{align}
    \begin{split}
\mathcal{S}_{\text{eff}}\Big |_{\text{fig}(\ref{10e})}
=\frac{-s_2m_1m_2}{8\pi m^2 m_p^2}\rmint dt \,
\mathcal{K}_{10e}(t)\,
\frac{\bar\phi(\boldsymbol{x}_0(t))}{m_p}+(1\leftrightarrow 2)\sim\mathcal{O}(L^{1/2}v^{9/2})\,,
\label{5.29}
    \end{split}
\end{align}
where
\begin{align}
    \begin{split}
        \mathcal{K}_{10e}(t)&=\partial_{0}f(r)\,\Delta\boldsymbol{v}\cdot \Hat{n}+f(r)\Bigg\{\Delta\Vec{a}\cdot \Hat{n}
        +\Delta\boldsymbol{v}\cdot
        \Bigg(
        \Hat{n}\frac{(\boldsymbol{v}_1-\boldsymbol{v}_2)\cdot \Hat{n}}{r}
        +\frac{\boldsymbol{v}_2-\boldsymbol{v}_1}{r}
        \Bigg)\Bigg\}.
    \end{split}
\end{align}
where in (\ref{5.29}) the function $f(r)$ is given by,
\begin{align}
    \begin{split}
        f(r)=\frac{e^{-mr}(mr-e^{mr}+1)}{r^2}.
        \label{5.46}
    \end{split}
\end{align}
\vspace{-0.925cm}
\item A third diagram contributes at  $N^{(4)}LO$ through the bulk interaction vertex $\rmint d^4x\,\hat{\mathcal{A}}_i\,\,\partial_i\varphi\,\partial_0\bar{\phi}(x)$ shown in Fig.~(\ref{10f}). In reduced form,
\begin{align}
\begin{split}
S_{\textrm{eff}}\Big |_{\text{fig}(\ref{10f})}
=\frac{-1}{2\pi }\frac{m_1m_2s_2}{m_p^2m^2}\rmint dt\,
\mathcal{K}_{10f}(t)\,
\frac{\Bar{\phi}(\boldsymbol{x}_0(t))}{m_p}  +(1\leftrightarrow 2)\sim\mathcal{O}(L^{1/2}v^{9/2})\,.
\end{split}
\end{align}
with
\begin{align}
\begin{split}
\mathcal{K}_{10f}(t)&=(\boldsymbol{v}_1\cdot \Hat{n}) \partial_0f(r) \\ &
\quad +\Bigg\{
(\Vec{a}_1\cdot \Hat{n})
+\boldsymbol{v}_1\cdot \Bigg[
\frac{(\boldsymbol{v}_1-\boldsymbol{v}_2)\cdot \Hat{n}}{r}\,\Hat{n}
+\frac{\boldsymbol{v}_2-\boldsymbol{v}_1}{r}
\Bigg]
\Bigg\}f(r).
\end{split}
\end{align}
%\item Another interesting bulk 3-point vertex will contribute to the radiation diagrams. The relevant part of the scalar field action is,
%\begin{align}
%\begin{split}
%S_{\Bar{\phi}\tilde{\psi}\Phi}&=-m^2\rmint d^4x\,\varphi\,\tilde{\psi}\,\Bar{\phi} \\ &
%=-m^2\rmint dt\,\Bar{\phi}(\boldsymbol{x}_{cm},t)\rmint_{\boldsymbol{k}}\,\tilde{\psi}_{k}(t)\Phi_{-k}(t)
%\end{split}
%\end{align}
%In order to proceed we need the momentum space propagator of the scalar field,
%\begin{align}
%    \begin{split}
%        \Big\langle\Phi_{k}(t_1)\varphi_{q}(t_2)\Big\rangle=-(2\pi)^3\frac{i}{\boldsymbol{k}^2}\delta(\boldsymbol{k}+\boldsymbol{q})\delta(t_1-t_2)
%    \end{split}
%\end{align}



%Therefore the corresponding amplitude is given by,
%\begin{align}
  %  \begin{split}
  %        \mathcal{S}_{\text{eff}}&=-\frac{m_1s_1m_2}{2(2\pi)^6}m^2\rmint dt_1 \rmint_{\boldsymbol{k}_1}e^{i\boldsymbol{k}_1\cdot\boldsymbol{x}_1}\varphi_{k_1}(t_1)\rmint dt_2\rmint_{\boldsymbol{k}_2}e^{i\boldsymbol{k}_2\cdot\boldsymbol{x}_2}\tilde{\psi}_{k_2}(t_2)\rmint dt \Bar{\phi}(\boldsymbol{x}_{cm},t)\rmint_{\boldsymbol{k}}\,\tilde{\psi}_{k}(t)\Phi_{-k}(t)\\ &
   %     =-\frac{m_1m_2s_1}{2(2\pi)^6}m^2\rmint dt \Bar{\phi}(\boldsymbol{x}_{cm},t)\rmint dt_1 dt_2\rmint_{k,k_1,k_2}e^{i\boldsymbol{k}_1\cdot\boldsymbol{x}_1+i\boldsymbol{k}_2\cdot\boldsymbol{x}_2}\Big\langle\varphi_{k_1}(t_1)\varphi_{-k}(t)\Big\rangle\Big\langle\tilde{\psi}_{k_2}(t_2)\tilde{\psi}_{k}(t)\Big\rangle\\ &
   %     =\textstyle{\frac{-m_1m_2s_1m^2}{2}\rmint dt \Bar{\phi}(\boldsymbol{x}_{cm},t)\rmint dt_1 dt_2\,\delta(t_1-t)\delta(t_2-t)\rmint_{k,k_1,k_2}\delta(\boldsymbol{k}_1-\boldsymbol{k})\delta(\boldsymbol{k}_2+\boldsymbol{k})e^{i\boldsymbol{k}_1\cdot\boldsymbol{x}_1+i\boldsymbol{k}_2\cdot\boldsymbol{x}_2}\frac{1}{\boldsymbol{k}_1^2(\boldsymbol{k}_2^2+m^2)}}\\ &=\textstyle{-\frac{m_1m_2s_1}{2}m^2\rmint dt\Bar{\phi}(\boldsymbol{x}_{cm},t)\rmint_{k}\frac{e^{-i\boldsymbol{k}\cdot \boldsymbol{r}}}{(\boldsymbol{k}^2+m^2)\boldsymbol{k}^2}}\\ &
   %     =\frac{m_1m_2s_1}{8\pi}\rmint dt \, \frac{e^{-mr}-1}{r}\Bar{\phi}(\boldsymbol{x}_{cm},t)
 % &  \end{split}
%\end{align}
Here $f(r)$ is defined in Eq.~(\ref{5.46}).

\item \textcolor{black}{A bulk vertex term causing scalar radiation at $N^{(5)}LO$ is $g_{a\gamma\gamma}\,\epsilon^{0ikm}\rmint d^4x \partial_{0}\mathcal{A}_{i}\partial_{k}\mathcal{A}_{m}\Bar{\phi}$, coming from the \textit{Theta} term. The corresponding reduced amplitude for Fig.~(\ref{10g}) is}
\begin{align}
    \begin{split}
          \mathcal{S}_{\text{eff}}\Big|_{\text{fig}(\ref{10g})}
          =\frac{Q_1 Q_2 g_{a\gamma\gamma}}{8\pi}\rmint dt \,[\Vec{a}_1\cdot (\boldsymbol{v}_2\times \Hat{n})+1\leftrightarrow 2]\,\frac{\Bar{\phi}(\boldsymbol{x}_0(t))}{m_p}\,\sim\mathcal{O}(L^{1/2}v^{11/2})\,.
\label{5.32m}
    \end{split}
\end{align}
%\item \textcolor{black}{Do we need to give this?\textcolor{red}{yes!} The term coming from the tensorial NRG field and scalar coupling is given by,
%\begin{align}
 %   \begin{split}   \mathcal{S}_{\text{eff}}&=\frac{s_2m_1m_2}{8m_p^2}\rmint dt_1dt_2\, v_1^q\,v_1^r\delta^{lm}\rmint d^4x\,\partial_i\Big\langle\varphi(\boldsymbol{x}_2(t_2))\varphi(x,t)\Big\rangle \Big\langle\boldsymbol{\zeta}_{lm}(\boldsymbol{x}_1(t_1))\boldsymbol{\zeta}_{qr}(x)\Big\rangle\partial_i\bar{\phi}+1\leftrightarrow 2 \\ &
   % \end{split}
%\end{align}
%\item  Another term coming from the NRG tensor field and scalar coupling is given by,
%\begin{align}
  %  \begin{split}
  %  S_{eff}&=\frac{s_2m_1m_2}{4m_p^2}\rmint dt_1dt_2 v_1^qv_1^r\rmint d^4x\partial_i\Big\langle\varphi(\boldsymbol{x}_2(t_2))\varphi(x,t)\Big\rangle \Big\langle\zeta_{ij}(\boldsymbol{x}_1(t_1))\zeta_{qr}(x)\Big\rangle\partial_j\bar{\phi}+1\leftrightarrow 2 \\&
  %  \end{split}
%\end{align}
%\item  Another term coming from the NRG tensor field and scalar coupling is given by,
%\begin{align}
  %  \begin{split}
  %  S_{eff}&=\frac{s_2m_1m_2}{4m_p^2}\rmint dt_1dt_2 v_1^qv_1^r\rmint d^4x\partial_j\Big\langle\varphi(\boldsymbol{x}_2(t_2))\varphi(x,t)\Big\rangle \Big\langle\zeta_{ij}(\boldsymbol{x}_1(t_1))\zeta_{qr}(x)\Big\rangle\partial_i\bar{\varphi}+1\leftrightarrow 2 \\&
   % \end{split}
%\end{align}
%\item We don't compute the above three terms as all the above terms have zero contribution as they vanish due to the specific structure of the term at leading order.}
\end{itemize}
Although we get a non-trivial contribution to the radiative effective action from the \textit{theta} term in (\ref{5.32m}), it vanishes for any orbit confined to a \textit{plane}. %We will get non-zero contribution for any orbit lies in 3-\textit{dim}().
\textcolor{black}{Apart from the above-mentioned terms, we can have some more 3-point bulk interaction vertices with a radiating field of the following form,
\begin{equation}
\rmint d^4x \, \boldsymbol{\zeta}_{qr}\partial_{i}\varphi\partial_{j}\Bar{\phi}. \label{5.36n}
\end{equation}
But they would not contribute to the source term for the following reason. To identify the source term, the effective action should have the following form $\rmint d^4x J(x)\Bar{\phi}(\boldsymbol{x}_0,t)$. 
 Now, to write down the interaction action in the specified form, we need to  integrate by parts:
 \begin{equation}
    \rmint dt \rmint d^3x \, \boldsymbol{\zeta}_{qr}\partial_{i}\varphi\partial_{j}\Bar{\phi}=-\rmint dt\Bar{\phi}(\boldsymbol{x}_0(t))\rmint d^3x \partial_{j}\{\boldsymbol{\zeta}_{qr}\partial_{i}\varphi\}\,.
 \end{equation} 
 After that, demanding that the potential fields vanish at spatial infinity, we get the contribution to zero $\implies \text{the source term}\,J=0$ for that specific interaction.}
\\
Now we will calculate the power radiation in the centre of mass frame defined by, $$\boldsymbol{x}_1=\frac{m_2}{m_1+m_2}\boldsymbol{r} \textrm{ and } \boldsymbol{x}_2=\frac{-m_1}{m_1+m_2}\boldsymbol{r}\,.$$
It is evident that the scalar monopole moment does not contribute. We are considering a circular orbit parameterized by, $$x=r \cos(\Omega t),\quad y=r \sin(\Omega t), \quad z=0\,,$$ with the orbital frequency $\Omega$.
The source term can be calculated from the effective action by omitting the radiative fields as defined in (\ref{5.10}). The source term contributing to scalar radiation has the form
\begin{align}
    \begin{split}
        J_{\phi}=& -\sum_{a=1}^{2}m_a\,s_a\delta(\boldsymbol{x}-\boldsymbol{x}_a)+\frac{m_1m_2}{8\pi m_p^2 r}\sum_{a=1}^{2}s_a \delta(\boldsymbol{x}-\boldsymbol{x}_a)+\frac{1}{2}\sum_{a}m_as_av_a^2\delta(\boldsymbol{x}-\boldsymbol{x}_a)\\&
       +\underbrace{ \frac{m_1^2 m_2^2}{16 \pi \,m_p^2(m_1+m_2)^2}}_{\frac{\mu^2}{16\pi}}e^{-mr}r^2\Omega^2\sum_{a\ne b}^{2}g_a s_{b}\,\delta(\boldsymbol{x}-\boldsymbol{x}_a)\\ &+\frac{m_1m_2}{m_p^2} \frac{e^{-mr}}{r}\Big[1+\frac{1}{16}\Big(\frac{m_1}{M}-\nu\Big)\Omega^2\,r^2\Big]\sum_{a\ne b}g_a s_b \delta(\boldsymbol{x}-\boldsymbol{x}_a)\\&+\underbrace{\Big\{[\tilde{c}_1+\tilde{c}_2]r^2\Omega^2+\Big[\frac{\gamma_2m_1-\gamma_1m_2}{M}\Big]r^2\Omega^2+[\tilde{G}_1+\tilde{G}_2]+(\tilde{d}_1-\tilde{d}_2)\Omega^2+\Big[\frac{\Gamma_2m_1-\Gamma_1m_2}{M}\Big]\Omega^2\Big\}}_{\eta(r,\Omega)}\\& \,\,\,\,\,\,\,\,\delta(x-\boldsymbol{x}_0)\,,
    \end{split}
\end{align}
where the functions are defined below,
\begin{align}
\begin{split}
&\tilde{c}_{1,2}=\frac{-s_{2,1}\,\mu\, M f(r)}{8\pi m^2m_p^2r}\,\,\,,
    \tilde{G}_{1,2}=\frac{s_{1,2}\,\mu\, M}{8\pi m_p^2}\frac{(1-e^{-mr})}{r}\,\, ,\,\gamma_{1,2}(r)=-\frac{1}{2\pi}\frac{\mu \,M\, s_{2,1}\,f(r)}{m_p^2 m^2r},\,\\& \tilde{d}_{1,2}=\tilde{c}_{1,2}\, r,\,\,\,\,\,
    \Gamma_{1,2}=\gamma_{1,2}\, r.
    \end{split}
\end{align}
and  $f(r)$ is defined in equation (\ref{5.29}). \par
%\vspace{0.5cm}
 Now the dipole moment is essentially given by,
\begin{align}
    \begin{split}
        I_{\phi}^{i}(t)&=\rmint d^3 x\,x^{i} J_{\phi}+\textcolor{black}{\frac{m^2}{10} \rmint d^3 x\,  x^i x^2\,J_{\phi} }\,,\\ &
        =\sum_{a\ne b}\underbrace{[-m_a s_a(1-\frac{v_a^2}{2})+\frac{ m_1 m_2}{8\pi m_p^2r}s_a +\frac{\mu^2 (g_a s_b)}{16\pi m_p^2}e^{-mr}r^2 \Omega^2+\frac{m_1m_2}{m_p^2}\frac{e^{-mr}}{r}\Big[1+\frac{1}{16}\Big(\frac{m_1}{M}-\nu\Big)\Omega^2\,r^2\Big]g_a s_b]}_{\mathcal{C}_{a}(r)}x_a^{i}\\ &
        +\textcolor{black}{\frac{m^2}{10}\sum_{a\ne b}\mathcal{C}_a x_a^2 x_a^{i}}+\frac{m^2}{10}\eta(r,\Omega){x}_0^i\,\boldsymbol{x}_0^2+\eta(r,\Omega){x}_0^i\,.
    \end{split}
\end{align}
In the frequency domain, the moments are given by,
\begin{align}
    \begin{split}
      &  I_{\phi}^{i}(\omega)=\rmint dt e^{i\omega t}I_{\phi}^{i}(t)\,,\\ &
        \implies I_{\phi}^{x}(\omega)=\Big[ \frac{\mathcal{C}_1 m_2-\mathcal{C}_2 m_1}{m_1+m_2}+\textcolor{black}{\frac{m^2r^2}{10}\frac{\mathcal{C}_1m_2^3-\mathcal{C}_2 m_1^3}{(m_1+m_2)^3} }+\frac{\eta(r,\Omega)}{2}+\frac{m^2}{80}\eta(r,\Omega)r^2\Big]r\,\sqrt{\frac{\pi}{2}}\delta(\omega-\Omega)\,,\\&\quad \text{and}\quad \,I_{\phi}^{y}(\omega)=i I_{\phi}^{x}(\omega),\,.
    \end{split}
\end{align}
Then the radiated power is given by \footnote{The formula for scalar power radiation matches up to some numerical factor with the result of \cite{Huang:2018pbu} due to the different choices of overall normalisation of the gravitational action. Also, the mass-dependent second term, i.e., the $\frac{m^2r^2}{10}(...)$ is not there. This is because of the mass dependence of the multipole moment, coming from the equation of motion, in (\ref{5.12mn}). },\par
\begin{tcolorbox}[thesisresultbox,, title=Scalar Radiation Power]
\restorethesisbodyformat
\begin{equation}
\begin{aligned}
        P_{\phi}&=\frac{1}{12\pi m_p^2 }\Big(\frac{\mathcal{C}_1 m_2-\mathcal{C}_2 m_1}{m_1+m_2}+\textcolor{black}{\frac{m^2 r^2}{10}\frac{\mathcal{C}_1m_2^3-\mathcal{C}_2 m_1^3}{(m_1+m_2)^3} }+\frac{\eta(r,\Omega)}{2}+\frac{m^2}{80}\eta(r,\Omega)r^2\Big)^2\\&\,\,\,\,\,\,\,\,\,\,\,\,\,\,\,\,\,\,\,\,\,\,\,\,\,\,\,\,\,\,\,\,\,\,\,\,\,\,\,\,\,\,\,\,\,\,\,\,\,\,\,\,\,\,\,\,\,\,\,\,\,\,\,\,\,\,\,\,\,\,\,\,\,\,\,\,\,\,\,\,\,\,\,\,\,\,\,\,\,\,\,\,\,\,\,\,\,\,\,\,\,\,\,\,\,\,\,\,\,\,\,\,\,\,\,\,\,\,\,\,\,\,\,\,\,\,\,\,\,\,\,\,\rmint d\omega \,\omega (\omega^2-m^2)^{3/2} r^2\delta(\omega-\Omega)\,,\\ &
        =\frac{1}{12\pi m_p^2 }\Big(\frac{\mathcal{C}_1 m_2-\mathcal{C}_2 m_1}{m_1+m_2}+\textcolor{black}{\frac{m^2r^2}{10}\frac{\mathcal{C}_1m_2^3-\mathcal{C}_2 m_1^3}{(m_1+m_2)^3} }+\frac{\eta(r,\Omega)}{2}+\frac{m^2}{80}\eta(r,\Omega)r^2\Big)^2\,r^2 \, \Omega^4 (1-\frac{m^2}{\Omega^2})^{3/2}\,.
\end{aligned}
\end{equation}
\end{tcolorbox}
%\vspace{0.8 cm}
Before we end this section, some comments are in order. 
\begin{itemize}
   \item  \textit{To the best of our knowledge, this is the first calculation of scalar radiation up to $N^{(4)}LO$. Leading-order computations were carried out in \cite{Huang:2018pbu, Kuntz:2019zef}}.
    \item The contribution coming from the $g_{a\gamma\gamma}$ term drops for the  orbits which are confined to a plane (e.g. for the circular orbit that we have considered in this paper), but it would have contributed otherwise for orbits embedded in 3-dimensions (e.g. some helical orbits embedded in 3d \cite{2021EPJC...81.1048L,Chen:2022qvg,Liu:2020vsy}).
\end{itemize}
\subsection{Electromagnetic interaction and the multipole decomposition}\label{emRad}
Now we consider the radiative photons with the source.  %$S_{int}^{(\bar{a}_{\mu})}$, 
The effective action can be written as
\begin{align}
    \begin{split}
        \mathcal{S}_{\text{eff}}=\rmint d^4 x \Big(-\frac{1}{4}F_{\mu\nu}F^{\mu\nu}+J^{\mu}\Bar{a}_{\mu}\Big)\,.
    \end{split}
\end{align}
We impose the Lorenz gauge: $$\partial_{\mu}\Bar{a}^{\mu}=0.$$ The equation of motion takes the following form,
\begin{eqnarray}
    \Box\Bar{a}_{\mu}=-J_{\mu}\,.
\end{eqnarray}
%We now do a multipole expansion of the vector field $\Bar{a_{\mu}}$ around the centre of mass $\boldsymbol{x}=0$ as $\sum_{M}\frac{x^{M}}{M!}\partial_{M}\Bar{a}_{\mu}(t,\boldsymbol{0})$, the interaction action can be written as,
%\begin{align}
 %   \begin{split}
  %      S_{int}^{\Bar{a_{\mu}}}&=\rmint dt \rmint d^3 x J^{\mu}(t,\boldsymbol{x})\Big(\Bar{a}_{\mu}(t,\boldsymbol{0})+x^i\partial_i\Bar{a}_{\mu}(t,\boldsymbol{0})+\frac{1}{2}x^ix^j\partial_i\partial_{j}\Bar{a}_{\mu}(t,\boldsymbol{0})+\frac{1}{3!}x^ix^jx^k\partial_i\partial_j\partial_{k}\Bar{a}_{\mu}(t,\boldsymbol{0})+...\Big)\\ &
   %    = \rmint dt [\mathcal{I}^{\mu}\Bar{a}_{\mu}+]
    %    \label{4.9}
    %\end{split}
%\end{align}

%As $S_{int}$ is a gauge invariant object one can seperate out the action into electric and magnetic dipole moment[\textcolor{red}{Cite}] as,
%\begin{align}
 %%      S_{int}^{a_{\mu}}=-\rmint dt \rmint d^3x J^{0}\, \Bar{a}^{0}+\rmint dt \sum_{l=1}^{\infty}\frac{1}{l!}\mathcal{I}^{L}\partial_{L-1}E^{k_l}+\rmint dt \sum_{l=1}^{\infty}\frac{l}{(l+1)!}\mathcal{J}^{L}\partial_{L-1}B^{k_l}
   % \end{split}
%\end{align}
%Where the electric and magnetic multipole moments are respectively,
%\begin{align}
 %   \begin{split}
  %      \mathcal{I}^{L}=
   % \end{split}
%\end{align}
To compensate for the extra term we need to add an extra term $\frac{1}{6}\rmint d^3 x\, x^2 J^{\mu}\nabla^2\Bar{a}_{\mu}$ with the monopole term. Now use $\Box \Bar{a}_{\mu}=0$ (as we are measuring outside the source) we can see the term can be written as
\begin{align}
    \begin{split}
        S_{int}'&=\frac{1}{6}\rmint dt \rmint d^3 x\, x^2\,J^\mu (\Box+\partial_{t}^2)\Bar{a}_{\mu}\\ &
        =\frac{1}{6}\rmint dt\rmint d^3x\,\Bar{a}_{\mu}\, \partial_{t}^2 J^{\mu} x^2\,.
    \end{split}
\end{align}
We compute the effective action for the source  following the same approach (in Sec.~(\ref{subsec:radiated_scalars})) as scalar field case. By doing a similar computation like the scalar field case the imaginary part of the effective action for the source in terms of multipole moments can be found as  \cite{Kuntz:2019zef}, 
%\begin{figure}
 %   \centering
  %  \scalebox{0.5}{
%\begin{feynman}
 %   \gluon[color=eb144c]{5.00, 4.00}{6.00, 5.00}
 %   \gluon[color=eb144c]{6.00, 5.00}{7.00, 4.00}
 %   \fermion[]{4.00, 4.00}{8.00, 4.00}
%\end{feynman}
%}
 %   \caption{Caption}
  %  \label{fig:my_label}
%\end{figure}
%The amplitude contributing to the imaginary part of the effective action \cite{Kuntz:2019zef},
%\begin{align}
 %   \begin{split}
  %      iS_{eff}&\sim \rmint dt_1\rmint d t_2 \Big[\mathcal{I}^{(\mu)}_{a}\mathcal{I}^{(\nu)}_{a}\Big\langle T \Bar{a}_{\mu}(t_1,\boldsymbol{0})\Bar{a}_{\nu}(t_2,\boldsymbol{0})\Big\rangle+\mathcal{I}^{(\mu)i}_{a}\mathcal{I}^{(\nu)j}_{a}\Big\langle T \partial_i \Bar{a}_{\mu}(t_1,\boldsymbol{0})\partial_j \Bar{a}_{\nu}(t_2,\boldsymbol{0})\Big\rangle+\\ &
   %     \,\,\,\,\,\,\,\,\,\,\,\,\,\,\,\,\,\,\,\,\,\,\,\,\,\,\,\,\,\,\,\,\,\,\,\,\,\,\,\,\frac{1}{4}\mathcal{I}^{(\mu)ij}_{a}\mathcal{I}^{(\nu)kl}_{a}\Big\langle T \partial_i\partial_j\Bar{a}_{\mu}(t_1,\boldsymbol{0})\partial_k\partial_l\Bar{a}_{\nu}(t_2,\boldsymbol{0})\Big\rangle\Big]\label{4.12}
    %\end{split}
%\end{align}
%Now we have the feynman propagator of the photons,
%\begin{align}
 %   \begin{split}
  %      \Big\langle T \Bar{a}_{\mu}(t_1,\boldsymbol{x}_1)\Bar{a}_{\nu}(t_2,\boldsymbol{x}_2)\Big\rangle=\rmint \frac{d^4k}{(2\pi)^4}\frac{-i\,g_{\mu\nu}}{k^2-i\epsilon}e^{-ik\cdot(x_1-x_2)}\label{4.13}
   % \end{split}
%\end{align}
%Therefore,
%\begin{align}
 %   \begin{split}
  %      \Big\langle T \partial_i \Bar{a}_{\mu}(t_1,\boldsymbol{0})\partial_j \Bar{a}_{\nu}(t_2,\boldsymbol{0})\Big\rangle&=\rmint \frac{d^4k}{(2\pi)^4}\frac{-i\,g_{\mu\nu}}{k^2-i\epsilon}e^{ik_0(t_1-t_2)}[\partial_{i}^{(x_1)}e^{-i\boldsymbol{k}.\boldsymbol{x}_1}][\partial_{j}^{(x_2)}e^{i\boldsymbol{k}\cdot\boldsymbol{x}_2}]\big|_{\boldsymbol{x}_1,\boldsymbol{x}_2=0}\\ &
   %     =\rmint \frac{d^4k}{(2\pi)^4}\frac{-i\,g_{\mu\nu}}{k^2-i\epsilon}e^{ik_0(t_1-t_2)}\langle\langle k_{i}\,k_{j}\rangle\rangle\\ &
    %    =\frac{1}{3}\rmint \frac{d^4k}{(2\pi)^4}\frac{-i\,g_{\mu\nu}}{k^2-%i\epsilon}e^{ik_0(t_1-t_2)}|\boldsymbol{k}|^2\delta_{ij}\label{4.14}
    %\end{split}
%\end{align}
%And the quadrupole contribution has the form,
%\begin{align}
%    \begin{split}
 %         \Big\langle T \,\partial_i\partial_j \Bar{a}_{\mu}(t_1,\boldsymbol{0})\,\partial_k\partial_l \Bar{a}_{\nu}(t_2,\boldsymbol{0})\Big\rangle&=\rmint \frac{d^4k}{(2\pi)^4}\frac{-i\,g_{\mu\nu}}{k^2-i\epsilon}e^{ik_0(t_1-t_2)}[\partial_{i}\partial_j^{(x_1)}e^{-i\boldsymbol{k}.\boldsymbol{x}_1}][\partial_{k}\partial_l^{(x_2)}e^{i\boldsymbol{k}\cdot\boldsymbol{x}_2}]\big|_{\boldsymbol{x}_1,\boldsymbol{x}_2=0}\\ &
   %       =\rmint \frac{d^4k}{(2\pi)^4}\frac{-i\,g_{\mu\nu}}{k^2-i\epsilon}e^{ik_0(t_1-t_2)}\langle\langle k_ik_jk_kk_l\rangle\rangle\\ &
    %      =\rmint \frac{d^4k}{(2\pi)^4}\frac{-i\,g_{\mu\nu}}{k^2-i\epsilon}e^{ik_0(t_1-t_2)}|\boldsymbol{k}|^4 T_{ij,kl}\label{4.15}
    %\end{split}
%\end{align}
%where, $T_{ij,kl}$ is constructed from $\langle\langle n_in_jn_kn_l\rangle\rangle$ and $\langle\langle n_in_j \rangle\rangle$ which is symmetric in $ij$ and $kl$ indices.
%Now putting the results of (\ref{4.13}), (\ref{4.14}) and (\ref{4.15}) into (\ref{4.12}) we will get,
%\begin{align}
 %   \begin{split}
  %      S_{eff}=-\rmint \frac{d^4 k}{(2\pi)^4}\frac{g_{\mu\nu}}{k^2-i\epsilon}[\mathcal{I}^{\mu}(k_0)\mathcal{I}^{*\nu}(k_0)+\frac{1}{3}\mathcal{I}^{(\mu)i}(k_0)\mathcal{I}^{*(\nu)j}(k_0)|\boldsymbol{k}|^2\delta_{ij}+\frac{1}{4}\mathcal{I}^{(\mu)ij}(k_0)\mathcal{I}^{*(\nu)kl}(k_0)|\boldsymbol{k}|^4T_{ij,kl}]
    %\end{split}
%\end{align}
%To extract out the imaginary part we need no bother about the principle value of the propagator i.e,
%\begin{align}
 %   \begin{split}
  %      \frac{1}{k^2}=PV\Big(\frac{1}{k^2}\Big)+i\pi\,\delta(k^2)
   % \end{split}
%\end{align}
%Hence,
\begin{align}
    \begin{split}
       \text{Im} [\gamma_\mathrm{eff}^{(\bar a_\mu)}]&=\frac{1}{8\pi^2}\rmint dk^0\rmint d|\boldsymbol{k}|\,|\boldsymbol{k}|^2 \, \delta(k_0^2-\boldsymbol{k}^2)\Big[|\mathcal{I}^{\mu}(k_0)|^2+\frac{1}{3}|\mathcal{I}^{(\mu)ij}(k_0)|^2|\boldsymbol{k}|^2+\frac{1}{30}|\mathcal{I}^{(\mu)ij}(k_0)|^2|\boldsymbol{k}|^4\Big]\,,\\ &
    =\frac{1}{8\pi^2}\rmint dk^0\rmint d|\boldsymbol{k}|\,|\boldsymbol{k}|^2 \frac{1}{2|k^0|}[\delta(k^0+|\boldsymbol{k}|)+\delta(k^0-|\boldsymbol{k}|)]\Big[|\mathcal{I}^{\mu}(k_0)|^2+\frac{1}{3}|\mathcal{I}^{(\mu)ij}(k_0)|^2|\boldsymbol{k}|^2\\ &\hspace{3cm}
    \,\,\,\,\,\,\,\,\,\,+\frac{1}{30}|\mathcal{I}^{(\mu)ij}(k_0)|^2|\boldsymbol{k}|^4\Big]\,,\\ &
    =\frac{1}{8\pi^2}\rmint_0^{\infty} d\omega\,\Big[|\mathcal{I}^{\mu}(\omega)|^2\omega+\frac{1}{3}|\mathcal{I}^{(\mu)i}(\omega)|^2\omega^3
    +\frac{1}{30}|\mathcal{I}^{(\mu)ij}(\omega)|^2\omega^5\Big]\,.
    \label{5.42}
    \end{split}
\end{align}
%\subsubsection*{Calculating the source term $J_{\mu}$:}
%The effective action can be calculated by integrating out the electromagnetic potential fields while taking the other radiative fields to be zero i.e,
%\begin{align}
 %   \begin{split}
  %      \mathcal{Z}[J_{\Bar{a}}]=\rmint \mathcal{D}\tilde{\psi} \mathcal{D}\hat{\mathcal{A}}_{i}\mathcal{D}\varphi \mathcal{D}\mathcal{A}_{\mu}\mathcal{D}\mathcal{B}_{\nu}\mathcal{D}\boldsymbol{\zeta}_{ij}\,e^{iS_{T}+\chi\cdot J}\Big|_{J_{\Bar{\phi},\Bar{h},\Bar{b}=0}}
   % \end{split}
%\end{align}
 We next compute this effective action up to LO. The diagrams that contribute at this order are shown in Fig.~(\ref{Fig:12}).
\begin{itemize}
\item The $LO$ contribution comes from a diagram with the worldline vertex $\rmint dt Q_a \Bar{a}_{0}$ as shown in Fig.~(\ref{fig12a}) has the following form,
\begin{align}
    \begin{split}
        \mathcal{S}_{\text{eff}}\Big |_{\text{fig}(\ref{fig12a})}=\sum_{a}\rmint dt \,Q_a \,\Bar{a}_{0}\sim \mathcal{O}(L^{1/2}v^{1/2}).
    \end{split}
\end{align}
\begin{figure}
    \centering
    \begin{subfigure}[t]{0.22\textwidth}
      \thesisdiagrampanel{\scalebox{0.3}{\begin{feynman}
    \fermion[lineWidth=6]{4.00, 5.80}{7.40, 5.80}
    \fermion[lineWidth=6]{4.00, 4.00}{7.40, 4.00}
    \dashed[showArrow=false, lineWidth=6, label=$\Bar{a_0}$]{5.60, 5.80}{7.40, 7.20}
\end{feynman}
}}
    \caption{}
    \label{fig12a}
    \end{subfigure}
\begin{subfigure}[t]{0.22\textwidth}
    \centering
\thesisdiagrampanel{\scalebox{0.3}{\begin{feynman}
    \fermion[lineWidth=6]{4.00, 5.80}{7.40, 5.80}
    \fermion[lineWidth=6]{4.00, 4.00}{7.40, 4.00}
    \electroweak[color=9900ef, lineWidth=6, label=$\Bar{a_i}$]{5.60, 5.80}{7.40, 7.20}
\end{feynman}
}}
    \caption{}
    \label{fig12b}
\end{subfigure}
\begin{subfigure}[t]{0.22\textwidth}
    \centering
 \thesisdiagrampanel{\scalebox{0.25}{
\begin{feynman}
    \fermion[lineWidth=6]{4.00, 6.20}{7.40, 6.20}
    \electroweak[lineWidth=6, label=$\Tilde{\psi}$]{5.60, 6.20}{5.60, 4.00}
    \dashed[showArrow=false, lineWidth=6, label=$\Bar{a_0}$]{5.60, 6.20}{7.60, 7.80}
    \fermion[lineWidth=6]{4.00, 4.00}{7.40, 4.00}
\end{feynman}}}
 \caption{}
    \label{fig12d}
\end{subfigure}
\begin{subfigure}[t]{0.22\textwidth}
    \centering
 \thesisdiagrampanel{\scalebox{0.3}{\begin{feynman}
    \electroweak[label=$\tilde{\psi}$, lineWidth=6]{5.60, 5.80}{5.60, 4.00}
    \fermion[lineWidth=6]{4.00, 5.80}{7.40, 5.80}
    \fermion[lineWidth=6]{4.00, 4.00}{7.40, 4.00}
    \electroweak[color=9900ef, label=$\Bar{a}_i$, lineWidth=6]{5.60, 5.80}{7.40, 7.20}
\end{feynman}
}}
    \caption{}
    \label{fig12c}
\end{subfigure}
\caption{Radiative diagrams for electromagnetic field up to $N^{(3)}LO$. }
\label{Fig:12}
\end{figure}


%\item Radiating term from the "Theta" term which gives NGR radiation $\Bar{\psi}$, coming from the worldline vertex $-2\rmint dt\, Q_a\,v_a^{i}A^{i}\,\Bar{\psi}, \rmint dt Q_av_a^i A^{i}, \rmint dt \,m_as_a\varphi$:

%\begin{figure}
%    \centering
 % \thesisdiagrampanel{\scalebox{0.4}{\begin{feynman}
  %  \fermion[lineWidth=4]{4.00, 5.20}{8.00, 5.20}
  %  \electroweak[lineWidth=4, label=$\Bar{\psi}$]{6.00, %4.00}{7.00, 5.20}
  %  \fermion[lineWidth=4]{4.00, 7.80}{8.00, 7.80}
  %  \electroweak[lineWidth=4, color=0693e3, label=$\varphi$]{6.00, 7.80}{6.00, 6.60}
  %  \dashed[lineWidth=4, showArrow=false]{6.00, 6.60}{7.00, 5.20}
  %  \dashed[lineWidth=4, showArrow=false, label=$A_i$]{6.00, 6.60}{6.80, 5.20}
   % \dashed[lineWidth=4, showArrow=false]{6.00, 6.60}{5.00, 5.20}
 %   \dashed[lineWidth=4, showArrow=false]{4.80, 5.20}{6.00, 6.60}
%\end{feynman}}}
 %   \caption{Capt}
 %   \label{fig:my_label}
%\end{figure}
%The corresponding amplitude is given by,
%\begin{align}
  %  \begin{split}
          %\mathcal{S}_{\text{eff}}^{ikm}&=\textstyle{2\rmint dt_2Q_2 v_2^{j}(t_2)A^{j}(x_2,t_2)\Bar{\psi}(x_2,t_2)\rmint dt_3 Q_2\,v_2^{l}(t_3)A^{l}(x_2,t_3)\rmint dt_1\,m_1s_1\varphi(x_1,t_1)\rmint d^4x \varphi\partial^0 A^{i}\partial^{k} A^{m}(x,t)}\\ &    
     %   =2Q_2^2m_1s_1\rmint dt_1dt_2dt_3dt\,v_2^{j}(t_2)v_2^{l}(t_3)\Bar{\psi}(x_2,t_2)\rmint d^3 x\, \partial^{k}\Big\langle A^{j}(x_2,t_2)A^{m}(x,t)\Big\rangle \partial^{0}\Big\langle A^{l}(x_2,t_3)A^{i}(x,t)\Big\rangle\\ &
        %\,\,\,\,\,\,\,\,\,\,\,\,\,\,\,\,\,\,\,\,\,\,\,\,\,\,\,\,\,\,\,\,\,\,\,\,\,\,\Big\langle \varphi(x_1,t_1)\varphi(x_2,t_2)\Big\rangle\\ &
     %   =2Q_2^2m_1s_1\rmint dt\, v_2^{m}(t)a_2^{i}(t)\Bigg[\frac{1}{16} \pi^{3/2}\, m \,G_{1,3}^{2,1}\left(\frac{m^2 r^2}{4}\Big|
%\begin{array}{c}
 %-\frac{1}{2} \\
 %-\frac{1}{2},-\frac{1}{2},0 \\
%\end{array}
%\right)-\frac{\,\pi^{3/2}\, G_{1,3}^{2,1}\left(\frac{m^2 r^2}{4}\Big|
%\begin{array}{c}
 %\frac{1}{2} \\
 %\frac{1}{2},\frac{1}{2},0 \\
%\end{array}
%\right)}{8 \,m\, r^2}\Bigg]n^k\,\\ &
%\Bar{\psi}(x_2,t) +\text{term symmetric in $i,m$}+1\leftrightarrow 2
   % \end{split}
%\end{align}
%The contracted amplitude that contributes to the source term is given by,
%\begin{align}
    %\begin{split}
          %\mathcal{S}_{\text{eff}}&=\frac{\pi^{3/2}}{4}g_{a\gamma\gamma}Q_2^2\,m_1\,s_1\rmint dt\,\Vec{a_2}\cdot(\hat{n}\times\boldsymbol{v}_2)\Bigg[\frac{1}{2}\, m \,G_{1,3}^{2,1}\left(\frac{m^2 r^2}{4}\Big|
%\begin{array}{c}
 %-\frac{1}{2} \\
% -\frac{1}{2},-\frac{1}{2},0 \\
%\end{array}
%\right)-\frac{ %G_{1,3}^{2,1}\left(\frac{m^2 r^2}{4}\Big|
%\begin{array}{c}
 %\frac{1}{2} \\
 %frac{1}{2},\frac{1}{2},0 \\
%\end{array}
%\right)}{ \,m\, r^2}\Bigg]\Bar{\psi}%(x_2,t)\\ &
%+1\leftrightarrow 2
 %   \end{split}
%\end{align}

%\item Another similar type interaction happens from two different electromagnetic worldline coupling: $\rmint dt Q_1v_1^{i}\mathcal{A}^{i}$ and $-2\rmint dt Q_2v_2^{i}\mathcal{A}^{i}\Bar{\psi}$
%\begin{figure}
 %   \centering
 %\thesisdiagrampanel{\scalebox{0.4}{  \begin{feynman}
 %   \gluon[lineWidth=4, color=fcb900, label=$\mathcal{A}_i$]{5.20, 7.80}{6.20, 6.80}
 %   \fermion[lineWidth=4,]{4.00, 5.40}{8.40, 5.40}
    %\electroweak[lineWidth=4,color=0693e3, label=$\phi$]{6.20, 6.80}{7.40, 7.80}
    %\electroweak[lineWidth=4,label=$\Bar{\psi}$]{5.00, 4.20}{6.20, 5.40}
   % \fermion[lineWidth=4,]{4.00, 7.80}{8.40, 7.80}
   % \gluon[lineWidth=4,color=fcb900]%{6.20, 6.80}{6.20, 5.40}
%\end{feynman}
%}
   % \caption{Caption}
   % \label{fig:my_label}
%\end{figure}
%The corresponding amplitude is given by (to calculate this amplitude, we need to invoke the one-loop massive scalar integral \cite{Anastasiou:1999ui,2022arXiv220103593W}),
%\begin{align}
    %\begin{split}
          %\mathcal{S}_{\text{eff}}^{ikm}&=-2\rmint dt_1 Q_1v_1^{j}(t_1)\mathcal{A}^{j}(x_1,t_1)\rmint dt_2 \,m_1s_1 \varphi(x_1,t_2)\rmint dt_3 Q_2 v_2^{l}(t_3)\mathcal{A}^{l}(x_2,t_3)\Bar{\psi}\rmint d^4x \,\varphi\,\partial^0 \mathcal{A}^{i}\partial^{k} \mathcal{A}^{m}(x,t)\\ &
       % =-\textstyle{2Q_1Q_2m_1s_1\rmint dt_1dt_2dt_3dt\,v_1^{j}(t_1)v_2^{l}(t_3)\Bar{\psi}(x_2,t_3)\rmint d^3 x \,\partial^{t}\Big\langle\mathcal{A}^{j}(x_1,t_1)\mathcal{A}^{i}(x,t)\Big\rangle\partial^{k}\Big\langle\mathcal{A}^{l}(x_2,t_3)  \mathcal{A}^{m}(x,t)  \Big\rangle}\\ &
    %\,\,\,\,\,\,\,\,\,\,\,\,\,\,\,\,\,\,\,\,\,\,\,\,\,\,\,\,\,\,\,\,\,\,\,\,\,\,\,\,\,\,\,\,\,\,\,\,\,\,  \Big\langle \varphi(x_1,t_2)\varphi(x,t)\Big\rangle\\ &
   % =-2Q_1Q_2m_1s_1\rmint dt\, dt_1 v_1^{i}(t_1)v_2^{m}(t)\Bar{\psi}(x_2,t)\partial^{t}(\delta(t-t_1))\, \rmint_{\boldsymbol{x}} \rmint_{\boldsymbol{k}_1,\boldsymbol{k}_2,\boldsymbol{k}_3}\frac{-ik_2^{k}}{\boldsymbol{k}_1^2\boldsymbol{k}_2^2(\boldsymbol{k}_3^2+m^2)}e^{i\boldsymbol{k}_1\cdot[\boldsymbol{x}_1(t_1)-\boldsymbol{x}(t)]}\\ &
   %\,\,\,\,\,\,\,\,\,\,\,\,\,\,\,\,\,\,\,\,\,\,\,\,\,\,\,\,\,\,\,\,\,\,\, e^{i\boldsymbol{k}_2\cdot[\boldsymbol{x}_2(t)-\boldsymbol{x}(t)]}
       %e^{i\boldsymbol{k}_3\cdot[\boldsymbol{x}_1(t)-\boldsymbol{x}(t)]}\\ &
     %  =-2Q_1Q_2m_1s_1\rmint dt\,dt_1 v_1^{i}(t_1)v_2^{m}(t)\Bar{\psi}(x_2,t)\partial^{t}(\delta(t-t_1))\rmint_{\boldsymbol{k}_1,\boldsymbol{k}_3}\frac{-i(-k_1^k-k_3^k)}{\boldsymbol{k}_1^2(\boldsymbol{k}_1+\boldsymbol{k}_3)^2(\boldsymbol{k}_3^2+m^2)}\\ &
 % \,\,\,\,\,\,\,\,\,\,\,\,\,\,\,\,\,\,\,\,\,\,\,\,\,\,\,\,\,     e^{i\boldsymbol{k}_1\cdot[\boldsymbol{x}_1(t_1)-\boldsymbol{x}_2(t)]}  e^{i\boldsymbol{k}_3\cdot[\boldsymbol{x}_1(t)-\boldsymbol{x}_2(t)]}\\ &
 %=-2Q_1Q_2m_1s_1\rmint dt\,dt_1 v_2^{m}(t)\delta(t-t_1)\partial^{t_1}\Big[v_1^{i}(t_1)\rmint_{\boldsymbol{k}_1,\boldsymbol{k}_3}\frac{-i(-k_1^k-k_3^k)}{\boldsymbol{k}_1^2(\boldsymbol{k}_1+\boldsymbol{k}_3)^2(\boldsymbol{k}_3^2+m^2)}e^{i\boldsymbol{k}_1\cdot[\boldsymbol{x}_1(t_1)-\boldsymbol{x}_2(t)]} \\ &
%\,\,\,\,\,\,\,\,\,\,\,\,\,\,\,\,\,\,\,\,\,\,\,\,\,\,\,\,\,\,\,\,\,e^{i\boldsymbol{k}_3\cdot[\boldsymbol{x}_1(t)-\boldsymbol{x}_2(t)]}\Big]\\ &
 % =-2Q_1Q_2m_1s_1\rmint dt v_2^{m}(t)a_1^{i}(t)\rmint_{\boldsymbol{k}_1,\boldsymbol{k}_3}\frac{-i(k_1^k+k_3^k)}{\boldsymbol{k}_1^2(\boldsymbol{k}_1+\boldsymbol{k}_3)^2(\boldsymbol{k}_3^2+m^2)}e^{i(\boldsymbol{k}_1+\boldsymbol{k}_3)\cdot\boldsymbol{r}(t)}\\ &
  %=2iQ_1Q_2m_1s_1\rmint dt v_2^{m}(t)a_1^{i}(t)\rmint_{k,k_3}\frac{k^k\,e^{i\boldsymbol{k}\cdot\boldsymbol{r}}}{(\boldsymbol{k}-\V      ec{k_3})^2\boldsymbol{k}^2(\boldsymbol{k}_3^2+m^2)}\\ &
  %=2iQ_1Q_2m_1s_1\rmint dt v_2^{m}(t)a_1^{i}(t)\rmint_{k}\frac{k^k\,e^{i\boldsymbol{k}\cdot\boldsymbol{r}}}{\boldsymbol{k}^2}\textcolor{black}{\rmint_{k_3}\frac{1}{(\boldsymbol{k}-\boldsymbol{k}_3)^2(\boldsymbol{k}_3^2+m^2)}}\bar{\psi}\\ &
 % =\frac{2Q_1 Q_2 m_1 s_1}{4\pi}\rmint dt \, v_2^m (t)a_1^i(t)\,\partial_r\alpha(m;r)n^k\bar{\psi} \label{5.51}
    %\end{split}
%\end{align}
%The details of computations and the the function $\alpha(m;r)$ is given in Appendix~(\ref{app3}).
\item The diagram shown in the Fig.~(\ref{fig12b}) contributes at $N^{(1)}LO.$ It consists of a worldline vertex $\rmint dt\,Q_av_a^{i}\Bar{a}_{i}$ and the corresponding amplitude is given by,
\begin{align}
    \begin{split}
          \mathcal{S}_{\text{eff}}\Big |_{\text{fig}(\ref{fig12b})}=\sum_{a}\rmint dt \, Q_a\,v_a^{i}\, \Bar{a}_{i}(\boldsymbol{x}_a(t))\,\sim\mathcal{O}(L^{1/2}v^{3/2}).
    \end{split}
\end{align}
\item \textcolor{black}{The $N^{(2)}LO$ contribution for the source of $J^{0}$ comes from the diagram shown in Fig.~(\ref{fig12d}) with the worldline vertex:
$-\rmint dt Q_a\,\Bar{a}_{0}\tilde{\psi}$\,.} The amplitude corresponds to the diagram is: 
\begin{align}
    \begin{split} \mathcal{S}_{\text{eff}}\Big |_{\text{fig}(\ref{fig12d})}&=\frac{Q_1m_2}{m_p^2}\rmint\Bar{\mathcal{D}}\Hat{\xi}\rmint dt_1 \Bar{a}_{0}(x_1,t_1)\tilde{\psi}(x_1,t_1)\rmint dt_2 \tilde{\psi}(x_2,t_2)\,,\\ &
        =\frac{Q_1m_2}{m_p^2}\rmint dt_1\,dt_2\,\Bar{a}_{0}(x_1,t_1)\Big\langle\tilde{\psi}(x_1,t_1)\tilde{\psi}(x_2,t_2)\Big\rangle\,, \\ &
      %  =8\pi G Q_a \rmint dt_1\,dt_2\, \delta(t_1-t_2)\,\Bar{a}^{0}(x_1,t_1)\rmint_{\boldsymbol{k}}\frac{e^{i\boldsymbol{k}\cdot[\boldsymbol{x}_1(t_1)-\boldsymbol{x}_2(t_2)]}}{\boldsymbol{k}^2}\,,\\ &
       % =8\pi\,G\,Q_1\rmint dt \rmint_{\boldsymbol{k}}\frac{e^{i\boldsymbol{k}\cdot\boldsymbol{r}}}{\boldsymbol{k}^2}\Bar{a}^{0}(x_1,t)\,,\\ &
        =\frac{1}{8\pi m_p^2}\,\sum_{a\ne b}Q_a m_b\rmint dt\, \frac{1}{r}\,\Bar{a}_{0}(x_a,t)\,\sim\mathcal{O}(L^{1/2}v^{5/2}).
    \end{split}
\end{align}
\item Other radiation diagram consisting of the following world line vertices, $-\rmint dt Q_a v_a^{i}\tilde{\psi}\Bar{a}^{i}$ and $\rmint dt \,\tilde{\psi}$ as shown in Fig.~(\ref{fig12c}) contributes at $N^{(3)}LO\,.$ The corresponding amplitude is given by,
\begin{align}
    \begin{split}
          \mathcal{S}_{\text{eff}}\Big |_{\text{fig}(\ref{fig12c})}&=-Q_1m_2\rmint\Bar{\mathcal{D}}\Hat{\xi}\rmint dt_1\, v_1^{i}(t_1)\tilde{\psi}(\boldsymbol{x}_1(t_1))\Bar{a}_{i}(\boldsymbol{x}_1(t_1))\rmint dt_2 \, \tilde{\psi}(\boldsymbol{x}_2(t_2))\,,\\ &
    =-\frac{Q_1m_2}{2m_p^2}\rmint dt v_1^{i}\rmint_{\boldsymbol{k}}\frac{e^{i\boldsymbol{k}\cdot\boldsymbol{r}}}{\boldsymbol{k}^2}\,\Bar{a}_{i}(\boldsymbol{x}_1(t))\,,\\ &
    =-\frac{Q_1m_2}{2m_p^2}\rmint dt v_1^{i}\, \frac{1}{4\pi r}\,\Bar{a}_{i}(\boldsymbol{x}_1(t))+(1\leftrightarrow 2)\,\sim\mathcal{O}(L^{1/2}v^{7/2}).
    \end{split}
\end{align}
\end{itemize}
By adding all these contributions, we get the total effective action for the radiation sector. Then we can calculate the source term, and it is given by,
\begin{align}
    \begin{split}
      &  J_{a}^{0}=\frac{1}{8\pi m_p^2\,r}\sum_{a\ne b}Q_{a}m_b\delta(\boldsymbol{x}-\boldsymbol{x}_a)+\sum_{a}Q_a\delta(\boldsymbol{x}-\boldsymbol{x}_a)\,,\\ &
J_{a}^{i}=\sum_{a}Q_{a}\,v_a^{i}\delta(\boldsymbol{x}-\boldsymbol{x}_a)-\frac{1}{8\pi m_p^2\,r}\sum_{a\ne b}Q_a\, m_b \,v_a^{i}\delta(\boldsymbol{x}-\boldsymbol{x}_a)\,.
    \end{split}
\end{align}
Finally, we compute the multipole moment using this source term, and then, using (\ref{5.42}) (as well as using the optical theorem), we get the following power radiation formula. Here we have confined ourselves to the dipole term (i.e., $l=1$).\\
\begin{tcolorbox}[thesisresultbox, title=Electromagnetic Dipole Radiation]
\restorethesisbodyformat
\begin{equation}
\begin{aligned}
        P_{em}=&\frac{1}{12\pi}\Bigg\{\Bigg[\frac{(Q_1m_2-Q_2 m_1)}{M}r+\frac{(\mathcal{C}_1m_2^2-\mathcal{C}_2m_1^2)}{M}\Bigg]^2\Omega^4\\ &
        +\Bigg[-\frac{(\mathcal{D}_1m_2^2+\mathcal{D}_2m_1^2)}{M^2}r
+\frac{(\tilde{d_1}m_2^4+\tilde{d_2}m_1^4)}{M^4}r^3\Omega^2
+\frac{\Delta_1m_2^3+\Delta_2m_1^3}{M^2}\\ &
+\frac{(G\tilde{d_1}m_2^5+G\tilde{d_2}m_1^5)}{M^4}r^2\Omega^2
 \Bigg]^2 16r^2\Omega^6\Bigg\}\,.
\end{aligned}
\end{equation}
\end{tcolorbox}
 
 Here, 
 $$ \mathcal{C}_a=\frac{Q_a}{8\pi m_p^2},\mathcal{D}_a=\frac{Q_a}{2},\tilde{d_a}=\frac{Q_a}{20},\Delta_a=\frac{\mathcal{D}_a}{8\pi m_p^2}, $$ with $a=1,2.$\par 
\textit{Again, to the best of our knowledge, the electromagnetic dipole radiation has been computed previously only up to $LO$ \cite{Cardoso:2020iji}. Here we have computed the same up to  $N^{(3)}LO$.}

\subsection{Proca interaction and multipole decomposition}
Now we will focus on the Proca sector. In order to derive the power radiation formula we need to calculate the imaginary part of the effective action of the source. Following the procedure outlined in Sec.~(\ref{subsec:radiated_scalars}) we get, 
\begin{align}
    \begin{split}
        i \gamma_\mathrm{eff}^{(\bar b_\mu)}&\sim \rmint dt_1\rmint d t_2 \Big[\mathcal{I}^{(\mu)}_{b}\mathcal{I}^{(\nu)}_{b}\Big\langle T \Bar{b}_{\mu}(t_1,\boldsymbol{0})\Bar{b}_{\nu}(t_2,\boldsymbol{0})\Big\rangle+\mathcal{I}^{(\mu)i}_{b}\mathcal{I}^{(\nu)j}_{b}\Big\langle T \partial_i \Bar{b}_{\mu}(t_1,\boldsymbol{0})\partial_j \Bar{b}_{\nu}(t_2,\boldsymbol{0})\Big\rangle+\\ &
        \,\,\,\,\,\,\,\,\,\,\,\,\,\,\,\,\,\,\,\,\,\,\,\,\,\,\,\,\,\,\,\,\,\,\,\,\,\,\,\,\frac{1}{4}\mathcal{I}^{(\mu)ij}_{a}\mathcal{I}^{(\nu)kl}_{b}\Big\langle T \partial_i\partial_j\Bar{b}_{\mu}(t_1,\boldsymbol{0})\partial_k\partial_l\Bar{b}_{\nu}(t_2,\boldsymbol{0})\Big\rangle\Big]\,.\label{5.1}
    \end{split}
\end{align}
The main difference between Proca and an electromagnetic propagator is that, in the case of a Proca field, the propagator is a massive propagator. The Feynman propagator for this case is given below,
\begin{align}
    \begin{split}
        \Big\langle T \Bar{b}_{\mu}(t_1,\boldsymbol{x}_1)\Bar{b}_{\nu}(t_2,\boldsymbol{x}_2)\Big\rangle=\rmint \frac{d^4k}{(2\pi)^4}\frac{-i\,\left(g_{\mu\nu}-\frac{k_\mu k_\nu}{\mu_\gamma ^2}\right)}{-k^2-\mu_{\gamma}^2+i\epsilon}e^{-ik\cdot(x_1-x_2)}\,.\label{5.2}
    \end{split}
\end{align}
Therefore, following the same method discussed in Sec.~(\ref{subsec:radiated_scalars}), we get the imaginary part of the effective action.
\begin{align}
    \begin{split}
        \text{Im}[\gamma^{\bar{b}_{\mu}}_{\text{eff}}]\Big|_{(1)}&=\frac{1}{8\pi^2}\rmint dk_{0}\rmint d|\boldsymbol{k}|\,|\boldsymbol{k}|^2\delta(k_0^2-\boldsymbol{k}^2-\mu_{\gamma}^2)\Big[|\mathcal{I}^{\mu}(k_0)|^2+\frac{1}{3}|\mathcal{I}^{(\mu)i}(k_0)|^2|\boldsymbol{k}|^2\\ &
   \hspace{7cm}+\frac{1}{30}|\mathcal{I}^{(\mu)ij}(k_0)|^2|\boldsymbol{k}|^4\Big]
   \,,\\ &
    =\frac{1}{8\pi^2}\rmint dk_0\rmint d|\boldsymbol{k}|\,|\boldsymbol{k}|^2\frac{1}{2\sqrt{k_0^2-\mu_{\gamma}^2}}\Big[\delta(-|\boldsymbol{k}|+\sqrt{k_0^2-\mu_{\gamma}^2})+\delta(|\boldsymbol{k}|+\sqrt{k_0^2-\mu_{\gamma}^2})\Big]\Big[|\mathcal{I}^{\mu}(k_0)|^2\\ &
   \hspace{7cm}+\frac{1}{3}|\mathcal{I}^{(\mu)i}(k_0)|^2|\boldsymbol{k}|^2+\frac{1}{30}|\mathcal{I}^{(\mu)ij}(k_0)|^2|\boldsymbol{k}|^4\Big]\,,\\ &
    =\frac{1}{8\pi^2}\rmint d\omega \, \sum_{L}\frac{1}{l!(2l+1)!!}\,(\omega^2-\mu_{\gamma}^2)^{1/2+l}\,|\mathcal{I}^{(\mu)}_{L}(\omega)|^2\,.
    \label{5.50}
    \end{split}
\end{align}
The other part comes from the second part of the propagator,
\begin{align}
    \begin{split}
\text{Im}[\gamma^{\bar{b}_{\mu}}_{\text{eff}}]\Big|_{(2)}&=-\frac{1}{8\pi^2}\int dt_1 dt_2 I^{\mu}(t_1)I^{\mu}(t_2)\int dk_0\frac{e^{ik_0\,(t_1-t_2)}}{2\sqrt{k_0^2-\mu_{\gamma}^2}}\int d^3k \,\delta\left(-|\boldsymbol{k}|+\sqrt{k_0^2-\mu_{\gamma}^2}\right)\frac{k_\mu k_\nu}{\mu_{\gamma}^2}\\ &
=-\frac{1}{8\pi^2}\int \frac{d\omega}{2\sqrt{\omega^2-\mu_\gamma^2}} I^{\mu}(\omega)I^{*\nu}(\omega)\int d^3k \,\delta\left(-|\boldsymbol{k}|+\sqrt{k_0^2-\mu_{\gamma}^2}\right)\frac{k_\mu k_\nu}{\mu_{\gamma}^2}\\ &
=-\frac{1}{16\pi^2\mu_{\gamma^2}}\int d\omega \,\sqrt{\omega^2-\mu_\gamma^2}\left(\omega^2|I^{0}(\omega)|^2-\frac{1}{3}(\omega^2-\mu_\gamma^2)|I^{i}(\omega)|^2\right)
\end{split}
\end{align}
However, in our following computations, we ignore the second part of the radiative effective action as we are only looking for the leading contribution for this massive mode. Although the correction can be straightforwardly done.
Therefore, the power radiation is given by,
\begin{align}
    \begin{split}
        P_{pr}=\frac{1}{4\pi^2\mathcal{T}}\sum_{l}\frac{1}{l!(2l+1)!!}\rmint_0^\infty d\omega\,\omega\,(\omega^2-\mu_{\gamma}^2)^{l+1/2}|\mathcal{I}^{\mu}_{(L)}(\omega)|^2\,.
    \end{split}
\end{align}
 We need to calculate the two source terms, $J_{b}^{0}$ and $J_{b}^{i}$, and the computation is exactly analogous to the electromagnetic case discussed in Sec.~(\ref{emRad}). Again, we restrict ourselves to \textcolor{black}{$N^{(3)}LO$}. Below, we quote the results for the source terms. 
%\begin{itemize}
 %   \item The lowest order contribution of $J_{b}^{0}$ comes from the worldline vertex: $\rmint dt -2Q_{a}'\,\Bar{b}^{0}\,\tilde{\psi}$. The corresponding amplitude is given by,
%\begin{align}
 %   \begin{split}
  %     &   \mathcal{S}_{\text{eff}(b)}^{0}=2G\sum_{a}Q_{a}'\rmint dt \,\frac{1}{r}\,\Bar{b}^{0}(x_a,t)
   % \end{split}
%\end{align}
%\item The contribution from $J_{b}^{i}$ corresponds to the worldline vertex: $\rmint dt\,Q_a'\,v_a^{i}\Bar{b}^{i}$ with the following amplitude,
%    \begin{split}
 %     &    \mathcal{S}_{\text{eff}(b_i)}=\sum_{a}\rmint dt Q_{a}'\,v_{a}^{i}\Bar{b}^{i}(x_a,t)
  %  \end{split}
%\end{align}
%\item The second order radiation diagram comes from the following world line vertices, $-2\rmint dt Q_a'v_a^{i}\tilde{\psi}\Bar{b}^{i}$ and $\rmint  dt \,\tilde{\psi}$ with the following amplitude:
%\begin{align}
 %   \begin{split}
  %   &     \mathcal{S}_{\text{eff}({b^{i}})}=-2GQ_1'm_2 v_1^{i}\rmint dt \frac{1}{r}\,\Bar{b}^{i}(x_1,t)+(1\leftrightarrow 2)
   % \end{split}
%\end{align}
%\end{itemize}
%Then the leading order source term without any field interaction vertex is given by the sum of the all amplitudes with amputated radiation fields.
\begin{align}
    \begin{split}
      &  J_{b}^{0}=\frac{1}{8\pi m_p^2 r}\sum_{a\neq b}Q_{a}'m_b\delta(\boldsymbol{x}-\boldsymbol{x}_a)+\sum_{a}Q_a'\delta(\boldsymbol{x}-\boldsymbol{x}_a)\,\\ &
      J_{b}^{i}=\sum_{a}Q_{a}'\,v_a^{i}\delta(\boldsymbol{x}-\boldsymbol{x}_a)-\frac{1}{8\pi m_p^2 r}\sum_{a\ne b}Q_a'\, m_b \,v_a^{i}\delta(\boldsymbol{x}-\boldsymbol{x}_a)\,.
    \end{split}
\end{align}
%{For the sake of completeness we show how to calculate the source term for the Proca fields with a single interaction vertex with the electromagnetic field.}\par 
%\begin{align}
 %   \begin{split}
  %      P_{pr}=\textcolor{red}{-}\frac{\textcolor{red}{2}}{4\pi^2\mathcal{T}}\sum_{l}\frac{1}{l!(2l+1)!!}\rmint_0^\infty d\omega\,\omega\,(\omega^2-\mu_{\gamma}^2)^{l+1/2}|\mathcal{I}^{\mu}_{(L)}(\omega)|^2
   % \end{split}
%\end{align}

Now that we have the source terms, we can use (\ref{5.50}) to calculate the multipole moments up to $N^{(3)}LO$ as follows,
 \footnote{ We use $\delta(x-a)^2=\delta(x-a)\delta(0)$.}
\begin{align}
\begin{split}
 |\mathcal{I}^{\mu}_{(L)}(\omega)|^2= &
\Bigg[\frac{({Q}_1'm_2-{Q}_2'm_1)r}{M}+\frac{\mu_{\gamma}^2}{10}\frac{({Q}_1'm_2^3-{Q}_2'm_1^3)}{M^3}r^2+\frac{(\mathcal{C}_1'm_2^2-\mathcal{C}_2'm_1^2)}{M}+\frac{\mu_{\gamma}^2}{10}\frac{(\mathcal{C}_1'm_2^3-\mathcal{C}_2'm_1^3)}{M^3}r^2\Bigg]^2\\&{\frac{\pi\delta(\omega-\Omega)\delta(0)}{2}}+\Bigg[- \frac{(\mathcal{D}_1'm_2^2+\mathcal{D}_2'm_1^2)}{M^2}r
+\frac{(\tilde{d_1}'m_2^4+\tilde{d_2}'m_1^4)}{M^4}r^3\Omega^2\textcolor{black}{(1-\frac{\mu_{\gamma}^2}{\Omega^2}})
\\&+\frac{\Delta_1'm_2^3+\Delta_2'm_1^3}{M^2} +\frac{(\tilde{d_1}'m_2^5+\tilde{d_2}'m_1^5)}{8\pi m_p^2 M^4}r^2\Omega^2\textcolor{black}{(1-\frac{\mu_{\gamma}^2}{\Omega^2}})
 \Bigg]^2\frac{\pi\, r^2\Omega^2}{2}\delta(\omega-2\Omega)\delta(0).
 \end{split}
 \end{align}

\textcolor{black}{Finally, we get the expression for the power radiation  for the Proca sector from the dipole part ($l=1$)}
 \begin{tcolorbox}[thesisresultbox, title=Proca Dipole Radiation]
 \restorethesisbodyformat
 \begin{equation}
 \begin{aligned}
      P_{\text{pr}}=&\frac{1}{12\pi} \Bigg\{\Bigg[\frac{({Q}_1'm_2-{Q}_2'm_1)r}{M}+\frac{\mu_{\gamma}^2}{10}\frac{({Q}_1'm_2^3-{Q}_2'm_1^3)}{M^3}r^2+\frac{(\mathcal{C}_1'm_2^2-\mathcal{C}_2'm_1^2)}{M}+\frac{\mu_{\gamma}^2}{10}\frac{(\mathcal{C}_1'm_2^4-\mathcal{C}_2'm_1^4)}{M^3}r^2\Bigg]^2\\ &
       \Omega^4(1-\frac{\mu_{\gamma}^2}{\Omega^2})^{\frac{3}{2}} +\Bigg[-\frac{(\mathcal{D}_1' m_2^2+\mathcal{D}_2' m_1^2)}{M^2}r
+\frac{(\tilde{d_1}'m_2^4+\tilde{d_2}'m_1^4)}{M^4}r^3\Omega^2\textcolor{black}{(1-\frac{\mu_{\gamma}^2}{\Omega^2}})
+\frac{\Delta'_1m_2^3+\Delta'_2m_1^3}{M^2}\\ &
+\frac{(\tilde{d_1}'m_2^5+\tilde{d_2}' m_1^5)}{8\pi m_p^2 M^4}r^2\Omega^2\textcolor{black}{(1-\frac{\mu_{\gamma}^2}{\Omega^2}})
 \Bigg]^2 16r^2\Omega^6(1-\frac{\mu_{\gamma}^2}{4\Omega^2})^{\frac{3}{2}}\Bigg\}\,,
\end{aligned}
\end{equation}
\end{tcolorbox}
 
 where $$ \mathcal{C}_a'=\frac{Q_a'}{8\pi m_p^2},\mathcal{D}_a'=\frac{Q_a'}{2},\tilde{d_a}'=\frac{Q_a'}{20},\Delta_a'=\frac{\mathcal{D}_a}{8\pi m_p^2}. $$\par 
\textit{Again, to the best of our knowledge, we have computed for the first time the Proca dipole radiation up to  $N^{(3)}LO$.}
\subsection{Gravitational radiation and Multipole decomposition}\label{sec5.5}
Finally, we focus on the gravitational sector. This is already studied in the \cite{Kuntz:2019zef, Levi:2018nxp}. We provide the computation for the gravitational power radiation for completeness, and also compute the corrections coming from different field vertices due to their coupling with the gravitational field at different orders of perturbation, which are explicitly mentioned. 
%To consider gravitational radiation we consider no contraction of $h_{\mu\nu}$, and we are considering only $h_{00}$ in radiation in leading order.Therefore
First, we will discuss the EFT  power counting similar to the one discussed in Sec.~(\ref{EFTcounting}) to understand which diagrams contribute to a \textcolor{black}{specific order}.\par
To see how the multipole moments scale, we need to know the scaling of the pseudo stress-energy tensor $T^{\mu\nu}.$

\begin{equation}\rmint d^3xT^{\mu\nu}\sim \begin{cases}
\begin{array}{ll}
M & (\mu=0,\nu=0),\\
Mv & (\mu=0,\nu=i),\\
Mv^2 & (\mu=i,\nu=j),\\
\end{array}
\end{cases}
\end{equation}
where $M$ denotes the mass-scale. From this we get,

\begin{equation}\rmint d^3x T^{\mu\nu}x^L\sim \begin{cases}
\begin{array}{ll}
L^{\frac{1}{2}}r^lv^{\frac{1}{2}} & (\mu=0,\nu=0),\\
L^{\frac{1}{2}}r^lv^{\frac{3}{2}} & (\mu=0,\nu=i),\\
L^{\frac{1}{2}}r^lv^{\frac{5}{2}} & (\mu=i,\nu=j),\\
\end{array}
\end{cases}
\end{equation}

where $L$ denotes the length-scale.

The formula for power radiation can be obtained in the same way as discussed in (\ref{5.18n}). First, we write down the imaginary part of the effective action.
\begin{align}
    \begin{split}
        i\gamma_{\text{eff}}^{\text{g}}=-\frac{1}{32 m_p^2}\rmint dt_1 dt_2\, I^{ij}_{g}(t_1)I^{kl}_{g}(t_2)\,\Big\langle T\ddot{\Bar{h}}_{ij}^{\text{TT}}(t_1,\boldsymbol{0})\ddot{\Bar{h}}_{kl}^{\text{TT}}(t_2,\boldsymbol{0})\Big\rangle+\cdots,\label{5.56mm}
    \end{split}
\end{align}
where ${\Bar{h}}_{ij}^{\text{TT}}$ is the radiative gravitational field in \textit{transverse-traceless (TT)} gauge and ($\cdots$) denote the higher order moments. Now use the following propagator for ${\Bar{h}}_{ij}^{\text{TT}}$ \cite{Goldberger:2004jt} to evaluate the effective action.
\begin{align}
    \begin{split}
        \Big\langle T\ddot{\Bar{h}}_{ij}^{\text{TT}}(t_1,\boldsymbol{0})\ddot{\Bar{h}}_{kl}^{\text{TT}}(t_2,\boldsymbol{0})\Big\rangle=\frac{8}{5}\Big[\frac{1}{2}(\delta_{ik}\delta_{jl}+\delta_{ik}\delta_{jl}-\frac{1}{3}\delta_{ij}\delta_{kl})\Big]\rmint \frac{d^4k}{(2\pi)^4}\frac{k_0^4}{-k^2-i\epsilon}e^{ik_0(t_1-t_2)},\label{5.57mm}
    \end{split}
\end{align}
Then using (\ref{5.57mm}) into (\ref{5.56mm}) we can write the effective action as \cite{Goldberger:2004jt,Kuntz:2019zef},
\begin{align}
    \begin{split}
        \gamma_{\text{eff}}^{\text{g}}=\frac{1}{20m_p^2}\rmint \frac{d^4 k}{(2\pi)^4}\frac{k_0^4}{-k^2-i\epsilon}|I^{ij}(k_0)|^2+\cdots.
    \end{split}
\end{align}
From this, we can write down the  power radiation in terms of \textit{quadrupole} moment,
\begin{align}
    \begin{split}
        P_{g}=\frac{2G}{5\mathcal{T}}\rmint \frac{d\omega}{2\pi}\omega^6\,|I^{ij}_{g}(\omega)|^2.
    \end{split}
\end{align}
%At this point, however, we have a important issue. For sure we have diagrams with
%the kind of  vertices in which radiated  momentum scales as $k\sim \frac{v}{r}$. Hence if we expand the potential propagator we have no homogeneous correction in v. Its all together a series in different orders of v.
%The analysis of gravitational radiation can be calculated as the way it is done in other field cases. 
In the rest of the section, we mainly use this expression to compute the power radiation. This can be extended systematically by including contributions from higher multipole moments. Interested readers are referred to \cite{Goldberger:2004jt} for more details. For completeness, we also write down the expression for power radiation up to the \textit{octupole} moment in the following way \cite{Goldberger:2004jt},
\begin{align}
    \begin{split}
        \mathcal{P}_{g}=\frac{2G}{5\mathcal{T}}\rmint \frac{d\omega}{2\pi}\, \Bigg\{\omega^6\, |{I}_{h}^{ij}(\omega)|^2+\frac{16\,\omega^6}{45}|J^{ij}(\omega)|^2+\frac{\omega^8}{189}|I^{ijk}(\omega)|^2+.........\Bigg\}\,.
        \label{5.65}
    \end{split}
\end{align}
 \textcolor{black}{ The explicit forms of ${I}^{ij}_{h}, \,J^{ij}$ and $I^{ijk}$ are given by:
\begin{align}
    \begin{split}
      &   {I}^{ij}_{h}(t)=\rmint d^3x \,T^{00}Q^{ij}.\\ &
      J^{ij}=-\frac{1}{2}\rmint d^3x (\epsilon^{ikl}T^{0k}x^ix^j+\epsilon^{jkl}T^{0k}x^ix^l)+.........\\ &
      I^{ijk}=\rmint d^3x (T^{00}+T^{ll})[x^ix^jx^k]^{STF}\,.
    \end{split}
\end{align}
}
%Next we will compute the radiation from the vertex of different field couplings as  shown in Fig.~(\ref{Fig15}). 
%The gravitational radiation sector consists of the following terms:
% The  radiation diagram comes from the following world line vertices, $-2\rmint dt Q_a v_a^{i}\Bar{\psi}{A}^{i}$ and $\rmint  dt \,\tilde{\psi}$.$-2\rmint dtQ_aA^0\Bar{\tilde{\psi}}$,.$-2\rmint dtQ_aA^0\Bar{\tilde{\psi}}^2$, $-2\rmint dt Q_a v_a^{i}\Bar{\Psi}{B}^{i}$ and - $\rmint  dt \,m_a(1+\frac{3}{2}v_a^2)\psi$,$-\frac{1}{2} \rmint m_adt\Bar{\psi}\Psi$,\\$-2\rmint dtQ_aB^0\Bar{\psi}$,.$-2\rmint dtQ_aB^0\Bar{\psi}^2$   with the following diagrams:
\begin{figure}[b!]
\centering
\thesisdiagrampanel[2.3cm]{\scalebox{0.3}{
\begin{feynman}
    \electroweak[lineWidth=4, label=$\Bar{\psi}$]{5.60, 6.00}{7.60, 7.20}
    \fermion[lineWidth=4]{4.00, 4.00}{7.40, 4.00}
    \fermion[lineWidth=4]{4.00, 6.00}{7.40, 6.00}
\end{feynman}}}
\caption{LO radiation coming from the pure gravitational sector.}
\label{figpq}
\end{figure}


\underline{\textit{{$LO$ result for pure gravity}}}:

First, we will briefly discuss the $LO$ power radiation from the pure gravitational sector (i.e., no other field couplings are present).  Only one diagram consists of the vertices $-\sum_a m_a\rmint dt\,\psi$ as shown in Fig.~(\ref{figpq}) contributes to this case. The corresponding amplitude is given by: 
 %The contribution coming from the pure gravitational $LO$ term is given by; 
\begin{align}
    \begin{split}
S_{\text{eff}}\Big|_{{\text{fig}(\ref{figpq})}}=-\sum_{a}m_a\rmint dt\, \frac{\Bar{\psi}(\boldsymbol{x}_a(t))}{m_p}\sim\mathcal{O}(L^{1/2}v^{1/2})\,.
    \end{split}
\end{align}


Now, to compute power radiation, we first define the centre of mass velocities (at leading order) in the following way: $$\boldsymbol{v}_1=\frac{m_2}{m_1+m_2} \boldsymbol{v}\quad \& \quad \boldsymbol{v}_2=-\frac{m_1}{m_1+m_2}  \boldsymbol{v}.$$\\

Then we compute the quadrupole moment,
\begin{align}
\begin{split}
   \mathcal{I}^{ij}_h&=-\rmint d^3 x\rmint dt\underbrace{\sum_a m_a\delta(\boldsymbol{x}-\boldsymbol{x}_a)}_{T^{00}}[x^ix^j]_{STF}\,,\\&
    =-\rmint dt{\sum_a m_a}[x_a^ix_a^j]_{STF}\,.\label{5.59m}
\end{split}
\end{align}
\vspace{-0.1cm}
%\begin{align}
%\begin{split}
%|\mathcal{I}^{ij}_h|^2=(|\mathcal{I}^{xx}|^2+|\mathcal{I}^{yy}|^2+|\mathcal{I}^{xy}|^2+|\mathcal{I}^{yx}|^2)
%\end{split}
%\end{align}

%Implementing this for the above term except the $\frac{3}{2}v_a^2$ 
Using (\ref{5.59m}) one can obtain the formula for $LO$ power radiation as ,
\begin{align}
\begin{split}
{P}_g&=\frac{2G}{5}\rmint \frac{d\omega}{2\pi}\omega^6\underbrace{(\frac{\pi}{2}\mu^2r^4\delta(\omega-2\Omega))}_{|\mathcal{I}^{ij}_h|^2}\,,\\&
=\frac{32G}{5}\mu^2r^4\Omega^6.
\end{split}
\end{align}
which agrees with our known GR radiation formula in leading order \cite{2009GReGr..41.1667H}. As there are couplings of the gravitational field with other fields for our case, we will also have contributions from them to the gravitational power radiation. Next, we calculate those non-trivial corrections of the gravitational power radiation to their respective leading orders.


\vspace{0.2 cm}
\underline{\textit{{Radiation from Pure  gravity at higher order ($N^{(2)}LO$)}}}:
\begin{itemize}
\item \textcolor{black}{The diagram that contributes at $N^{(2)}LO$ is shown Fig.~($\ref{figpq}$). The corresponding amplitude is given by,
\begin{align}
\begin{split}
\mathcal{S}_\text{eff}\Big |_{\text{fig}(\ref{figpq})}=-\sum_a m_a \rmint dt \,\frac{3}{2}v_a^2\,\frac{\Bar{\psi}(\boldsymbol{x}_a(t))}{m_p}\sim\mathcal{O}(L^{1/2}v^{5/2})\,.
\end{split}
\end{align}}
\item The amplitude corresponding to the diagram in Fig.~(\ref{fig13e}) contributing at $N^{(2)}LO$ with worldline coupling $\frac{m_a}{2}\rmint dt\,\tilde{\psi}\Bar{\psi} :$
\begin{align}
    \begin{split}
         \mathcal{S}_{\text{eff}}\Big |_{\text{fig}(\ref{fig13e})}&=\frac{m_1m_2}{2m_p^2}\rmint \Bar{\mathcal{D}}\hat{\xi}\rmint dt_1 \Tilde{\psi}(\boldsymbol{x}_1(t_1))\rmint dt_2\Tilde{\psi}(\boldsymbol{x}_2(t_2))\frac{\Bar{\psi}(\boldsymbol{x}_1(t_1))}{m_p}\,\\&
         =\frac{m_1m_2}{2m_p^2}\rmint dt_1 dt_2 \Big{\langle}\tilde{\psi}(\boldsymbol{x}_1(t_1))\Tilde{\psi}(\boldsymbol{x}_2(t_2))\Big{\rangle}\frac{\Bar{\psi}(\boldsymbol{x}_1(t_1))}{m_p}\,,\\&
      % =\frac{m_1m_2}{m_p^2}\rmint dt \frac{\Bar{\psi}}{m_p}\frac{1}{(4\pi)^{\frac{3}{2}}}\sqrt{\pi}\frac{2}{r}\,,\\ &
      =\frac{m_1m_2}{16m_p^2\pi}\rmint dt\frac{1}{r}\frac{\Bar{\psi}(\boldsymbol{x}_1(t))}{m_p}+(1\leftrightarrow 2)\,,\sim\mathcal{O}(L^{1/2}v^{5/2}).
    \end{split}
\end{align}
\end{itemize}
\underline{\textit{Radiation due to other field couplings with gravity}}\footnote{There are many diagrams contributing to the pure gravitational radiation at different $N^{(n)}LO$. But we don't calculate all of them explicitly as the results are already there in literature \cite{Blanchet:2013haa}.}:
\begin{itemize}
 \item The amplitude corresponding to the diagram in Fig.~(\ref{fig13f}) with worldline interaction vertex $-Q_a'\rmint dt\,\mathcal{B}_0\Bar{\psi}$ at $N^{(2)}LO :$
\begin{align}
    \begin{split}
         \mathcal{S}_{\text{eff}}\Big |_{\text{fig}(\ref{fig13f})}&=-Q_1'Q_2'\rmint\Bar{\mathcal{D}}\Hat{\xi}\rmint dt_1\mathcal{B}_0(\boldsymbol{x}_1(t_1))\rmint dt_2\,\mathcal{B}_0(\boldsymbol{x}_2(t_2))\frac{\Bar{\psi}(\boldsymbol{x}_1(t_1))}{m_p}\\&
       =-Q_1'Q_2'\rmint dt_1 dt_2\Big{\langle}\,\mathcal{B}_0(\boldsymbol{x}_1(t_1))\mathcal{B}_0(\boldsymbol{x}_2(t_2))\Big{\rangle}\frac{\Bar{\psi}(\boldsymbol{x}_1(t_1))}{m_p}\,, \\&
       =Q_1'Q_2'\rmint dt\frac{e^{-\mu_{\gamma} r}}{4\pi r}\frac{\Bar{\psi}(\boldsymbol{x}_1(t))}{m_p}+(1\leftrightarrow 2)\,\,\sim\mathcal{O}(L^{1/2}v^{5/2})\,.
    \end{split}
\end{align}
\end{itemize}
\begin{figure}
    \centering
    \begin{subfigure}[t]{0.2\textwidth}
        \centering
        \thesisdiagrampanel{\scalebox{0.25}{\begin{feynman}
    \fermion[lineWidth=6]{4.00, 6.40}{7.20, 6.40}
    \fermion[lineWidth=6]{4.00, 4.00}{7.20, 4.00}
    \electroweak[lineWidth=6, label=$\tilde{\psi}$]{5.60, 6.40}{5.60, 4.00}
\electroweak[label=$\Bar{\psi}$, lineWidth=6]{5.60, 6.40}{7.20, 8.00}
\end{feynman}
}}
\caption{}
\label{fig13e}
    \end{subfigure}
    \begin{subfigure}[t]{0.2\textwidth}
        \centering
        \thesisdiagrampanel{\scalebox{0.25}{
    \begin{feynman}
    \fermion[lineWidth=6]{4.00, 6.40}{7.20, 6.40}
    \fermion[lineWidth=6]{4.00, 4.00}{7.20, 4.00}
    \dashed[showArrow=false, lineWidth=6, color=9900ef, label=$\mathcal{B}_0$]{5.60, 6.40}{5.60, 4.00}
\electroweak[label=$\Bar{\psi}$, lineWidth=6]{5.60, 6.40}{7.20, 8.00}
\end{feynman}
}}\caption{}
\label{fig13f}
    \end{subfigure}
      \begin{subfigure}[t]{0.2\textwidth}
        \centering
        \thesisdiagrampanel{\scalebox{0.25}{
\begin{feynman}
    \dashed[showArrow=false, lineWidth=6, color=9900ef, label=$\mathcal{B}_0$]{5.60, 5.40}{5.60, 4.00}
    \fermion[lineWidth=6]{4.00, 6.40}{7.20, 6.40}
    \electroweak[lineWidth=6, label=$\Bar{\psi}$]{5.60, 6.40}{7.60, 7.80}
    \fermion[lineWidth=6]{4.00, 4.00}{7.20, 4.00}
    \dashed[lineWidth=6, showArrow=false, label=$\mathcal{A}_0$]{5.60, 6.40}{5.60, 5.40}
\end{feynman}
}}
\caption{}
\label{fig13l}
\end{subfigure}
\begin{subfigure}[t]{0.2\textwidth}
        \centering
        \thesisdiagrampanel{\scalebox{0.25}{
\begin{feynman}
    \fermion[lineWidth=6]{4.00, 6.40}{7.20, 6.40}
    \dashed[showArrow=false, lineWidth=6, label=$\mathcal{A}_0$]{5.60, 6.40}{5.60, 4.00}
    \fermion[lineWidth=6]{4.00, 4.00}{7.20, 4.00}
\electroweak[label=$\Bar{\psi}$, lineWidth=6]{5.60, 6.40}{7.20, 8.00}
\end{feynman}
}}
\caption{}
\label{fig13g}
    \end{subfigure}
    \begin{subfigure}[t]{0.25\textwidth}
    \centering
\thesisdiagrampanel{\scalebox{0.22}{\begin{feynman}
\electroweak[label=$\varphi$, lineWidth=6, color=0693e3]{6.20, 5.10}{6.20, 4.00}
    \electroweak[lineWidth=5, label=$\Bar{\psi}$]{6.50, 5.40}{8.00, 5.40}
    \fermion[lineWidth=6]{4.00, 4.00}{8.40, 4.00}
    \electroweak[lineWidth=6, color=0693e3, label=$\varphi$]{6.20, 5.70}{6.20, 6.80}
    \fermion[lineWidth=6]{4.00, 6.80}{8.40, 6.80}
    \parton[color=eb144c]{6.20,5.40}{0.30}
\end{feynman}}}
    \caption{}
\label{fig13n}
\end{subfigure}
    \begin{subfigure}[t]{0.2\textwidth}
        \centering
        \thesisdiagrampanel{\scalebox{0.25}{
        \begin{feynman}
    \gluon[color=eb144c, label=$\mathcal{B}_i$, lineWidth=6, flip=true]{5.20, 6.40}{5.20, 4.00}
    \fermion[lineWidth=6]{4.00, 4.00}{6.60, 4.00}
    \fermion[lineWidth=6]{4.00, 6.40}{6.60, 6.40}
    \electroweak[lineWidth=6, label=$\bar{\psi}$]{5.20, 6.40}{6.80, 8.00}
\end{feynman}
}}
\caption{}
\label{fig13a}
    \end{subfigure}
    \begin{subfigure}[t]{0.2\textwidth}
        \centering
        \thesisdiagrampanel{\scalebox{0.25}{
        \begin{feynman}
    \fermion[lineWidth=6]{4.00, 4.00}{7.00, 4.00}
    \electroweak[lineWidth=6, color=9900ef, label=$\mathcal{A}_i$]{5.40, 6.40}{5.40, 4.00}
    \electroweak[lineWidth=6, label=$\Bar{\psi}$]{5.40, 6.40}{7.60, 8.00}
    \fermion[lineWidth=6]{4.00, 6.40}{7.00, 6.40}
\end{feynman}
}}\caption{}
\label{fig13b}
    \end{subfigure}
\begin{subfigure}[t]{0.2\textwidth}
        \centering
        \thesisdiagrampanel{\scalebox{0.25}{
\begin{feynman}
    \fermion[lineWidth=6]{4.00, 6.40}{7.20, 6.40}
    \dashed[showArrow=false, label=$\mathcal{B}_0$, color=9900ef, lineWidth=6]{4.40, 4.00}{5.60, 6.40}
    \fermion[lineWidth=6]{4.00, 4.00}{7.20, 4.00}
    \electroweak[lineWidth=6]{5.60, 6.40}{6.80, 4.00}
\electroweak[label=$\Bar{\psi}$, lineWidth=6]{5.60, 6.40}{7.20, 8.00}
\end{feynman}
}}
\caption{}
\label{fig13c}
    \end{subfigure}
\begin{subfigure}[t]{0.2\textwidth}
        \centering
        \thesisdiagrampanel{\scalebox{0.25}{
    \begin{feynman}
    \fermion[lineWidth=6]{4.00, 6.40}{7.20, 6.40}
    \dashed[showArrow=false, lineWidth=6, label=$\mathcal{A}_0$]{4.40, 4.00}{5.60, 6.40}
    \fermion[lineWidth=6]{4.00, 4.00}{7.20, 4.00}
    \electroweak[lineWidth=6]{5.60, 6.40}{6.80, 4.00}
\electroweak[label=$\Bar{\psi}$, lineWidth=6]{5.60, 6.40}{7.20, 8.00}
\end{feynman}
}}\caption{}
\label{fig13d}
    \end{subfigure}
\begin{subfigure}[t]{0.2\textwidth}
        \centering
        \thesisdiagrampanel{\scalebox{0.25}{
    \begin{feynman}
    \fermion[lineWidth=6]{4.00, 6.40}{7.20, 6.40}
    \electroweak[lineWidth=6, color=0693e3, label=$\varphi$]{5.60, 6.40}{5.60, 4.00}
    \fermion[lineWidth=6]{4.00, 4.00}{7.20, 4.00}
\electroweak[label=$\Bar{\psi}$, lineWidth=6]{5.60, 6.40}{7.20, 8.00}
\end{feynman}
}}\caption{}
\label{fig13h}
    \end{subfigure}
\begin{subfigure}[t]{0.2\textwidth}
        \centering
        \thesisdiagrampanel{\scalebox{0.25}{
\begin{feynman}
    \fermion[lineWidth=6]{4.00, 6.40}{7.20, 6.40}
    \electroweak[lineWidth=6, color=0693e3, label=$\varphi$]{5.60, 6.40}{6.80, 4.00}
    \fermion[lineWidth=6]{4.00, 4.00}{7.20, 4.00}
    \electroweak[lineWidth=6, label=$\varphi$, color=0693e3]{4.40, 4.00}{5.60, 6.40}
\electroweak[label=$\Bar{\psi}$, lineWidth=6]{5.60, 6.40}{7.20, 8.00}
\end{feynman}
}}
\caption{}
\label{fig13i}
    \end{subfigure}
\begin{subfigure}[t]{0.2\textwidth}
        \centering
        \thesisdiagrampanel{\scalebox{0.25}{
\begin{feynman}
    \gluon[flip=true, lineWidth=6, label=$\mathcal{B}_i$, color=eb144c]{5.60, 5.20}{5.60, 4.00}
    \fermion[lineWidth=6]{4.00, 6.40}{7.20, 6.40}
    \electroweak[lineWidth=6, label=$\Bar{\psi}$]{5.60, 5.20}{8.20, 5.20}
    \fermion[lineWidth=6]{4.00, 4.00}{7.20, 4.00}
    \electroweak[lineWidth=6, color=9900ef, label=$\mathcal{A}_i$]{5.60, 6.40}{5.60, 5.20}
\end{feynman}
}} 
\caption{}
\label{fig13j}
\end{subfigure}
\begin{subfigure}[t]{0.2\textwidth}
        \centering
        \thesisdiagrampanel{\scalebox{0.25}{\begin{feynman}
    \gluon[lineWidth=6, label=$\mathcal{B}_k$, color=eb144c, flip=true]{5.60, 5.40}{5.60, 4.00}
    \fermion[lineWidth=6]{4.00, 6.40}{7.20, 6.40}
    \dashed[lineWidth=6, label=$\mathcal{A}_0$, showArrow=false]{5.60, 6.40}{5.60, 5.40}
    \electroweak[lineWidth=6, label=$\Bar{\psi}$]{5.60, 6.40}{7.60, 7.80}
    \fermion[lineWidth=6]{4.00, 4.00}{7.20, 4.00}
\end{feynman}
}}
\caption{}
\label{fig13k}
\end{subfigure}
\begin{subfigure}[t]{0.2\textwidth}
        \centering
        \thesisdiagrampanel{\scalebox{0.25}{
\begin{feynman}
    \dashed[showArrow=false, lineWidth=6, color=9900ef, label=$\mathcal{B}_0$]{5.60, 5.20}{5.60, 4.00}
    \electroweak[lineWidth=6, label=$\mathcal{A}_i$, color=9900ef]{5.60, 5.20}{5.60, 6.40}
    \fermion[lineWidth=6]{4.00, 6.40}{7.20, 6.40}
    \electroweak[lineWidth=6, label=$\Bar{\psi}$]{5.60, 6.40}{7.60, 7.80}
    \fermion[lineWidth=6]{4.00, 4.00}{7.20, 4.00}
\end{feynman}
}}
\caption{}
\label{fig13m}
\end{subfigure}
\begin{subfigure}[t]{0.22\textwidth}
    \centering
\thesisdiagrampanel{\scalebox{0.25}{
\begin{feynman}
    \electroweak[lineWidth=6, label=$\bar{\psi}$]{5.60, 6.40}{7.00, 7.60}
    \fermion[lineWidth=6]{4.00, 4.00}{7.20, 4.00}
    \electroweak[color=9900ef, lineWidth=6, label=$\mathcal{A}_l$]{5.60, 5.20}{6.80, 4.00}
    \electroweak[color=0693e3, lineWidth=6, label=$\varphi$]{5.60, 6.40}{5.60, 5.20}
    \electroweak[color=9900ef, lineWidth=6, label=$\mathcal{A}_j$]{4.40, 4.00}{5.60, 5.20}
    \fermion[lineWidth=6]{4.00, 6.40}{7.20, 6.40}
\end{feynman}}}
 \caption{}
\label{fig13o}
\end{subfigure}
\begin{subfigure}[t]{0.22\textwidth}
        \centering
        \thesisdiagrampanel{\scalebox{0.25}{\begin{feynman}
    \electroweak[flip=true, lineWidth=6, label=$\varphi$, color=0693e3]{5.80, 5.80}{6.80, 7.00}
    \electroweak[lineWidth=6, color=9900ef, label=$\mathcal{A}_l$]{5.80, 5.80}{5.80, 4.60}
    \electroweak[color=9900ef, lineWidth=6, label=$\mathcal{A}_j$]{4.80, 7.00}{5.80, 5.80}
    \fermion[lineWidth=6]{4.00, 4.60}{7.40, 4.60}
    \electroweak[lineWidth=6, label=$\bar{\psi}$, flip=true]{5.80, 4.60}{7.20, 4.00}
    \fermion[lineWidth=6]{4.00, 7.00}{7.60, 7.00}
\end{feynman}
}}
 \caption{}
\label{fig13p}
\end{subfigure}
\caption{Radiative diagrams for gravitational fields.}
\label{Fig15}
\end{figure}

\begin{itemize}


\item  The amplitude corresponding to the diagram in Fig.~(\ref{fig13l}) contributing at $N^{(2)}LO$ with the following bulk interaction vertex $\rmint d^4x\,\partial_k \mathcal{A}_0\partial_k \mathcal{B}_0 :$ 
\vspace{-0.5cm}
 \begin{align}
    \begin{split}
\mathcal{S}_{\text{eff}}\Big|_{\text{fig}(\ref{fig13l})}&=-\gamma Q_1 Q_2'\rmint \Bar{\mathcal{D}}\hat{\xi}\rmint dt_1\mathcal{A}_0(\boldsymbol{x}_1(t_1))\rmint dt_2\mathcal{B}_0(\boldsymbol{x}_2(t_2))\rmint d^4x\,\partial_k\mathcal{A}_0(x)\partial_k\mathcal{B}_0(x)\frac{\Bar{\psi}}{m_p}\,,\\&=-\gamma Q_1 Q_2'\rmint dt_1dt_2\rmint d^4x\partial_k\Big{\langle}\mathcal{A}_0(x)\mathcal{A}_0(\boldsymbol{x}_1(t_1))\Big{\rangle}\partial_k\Big{\langle}\mathcal{B}_0(x)\mathcal{B}_0(\boldsymbol{x}_2(t_2))\Big{\rangle}\frac{\Bar{\psi}}{m_p}\,,\\ &
%=Q_1Q_2'\rmint dt_1\, dt_2\rmint d^4x\partial_k\rmint\delta(t-t_1)\frac{e^{-i\boldsymbol{k}\cdot [\boldsymbol{x}(t)-\boldsymbol{x}_1(t_1)]}}{k^2}\partial_k\rmint\delta(t-t_2)\frac{e^{-i\boldsymbol{k}_1\cdot[\boldsymbol{x}(t)-\boldsymbol{x}_2(t_2)]}}{k_1^2+\mu^2}\\&
%=-Q_1Q_2'\rmint dt_1\, dt_2\,dt\,\delta(t-t_1)\rmint_{\boldsymbol{k},\boldsymbol{k}_1}\delta(\boldsymbol{k}+ \boldsymbol{k}_1)\delta(t-t_2)\frac{k^kk_1^ke^{i\boldsymbol{k}\cdot \boldsymbol{x}_1(t_1)}e^{i\boldsymbol{k}_1\cdot \boldsymbol{x}_2(t_2)}}{k^2(k_1^2+\mu^2)}\\ &
%=Q_1Q_2'\rmint dt\, \,\Big[\rmint \frac{e^{-i\boldsymbol{k}\cdot [\boldsymbol{x}_2(t)-\boldsymbol{x}_1(t)]}}{k^2(\boldsymbol{k}^2+\mu^2)}k^kk^k\Big]\\ &
%\nonumber
%\end{split}
%\end{align}
%\begin{align}
%\begin{split}
%=-\gamma Q_1Q_2'\rmint dt\, \,\Big[\rmint \frac{e^{-i\boldsymbol{k}\cdot \boldsymbol{r}}}{(\boldsymbol{k}^2+\mu^2)}\Bigg]                 \frac{\Bar{\psi}}{m_p}\,.\\ &
=-\gamma Q_1Q_2'\rmint dt\,\frac{e^{-\mu_{\gamma}r}}{4\pi r}\Big(\frac{\Bar{\psi}(\boldsymbol{x}_1(t))}{m_p}+\frac{\Bar{\psi}(\boldsymbol{x}_2(t))}{m_p}\Big)+(1\leftrightarrow 2)\,\,\sim \mathcal{O}(L^{1/2}v^{5/2}).
    \end{split}
\end{align}
\item The amplitude corresponding to the diagram in Fig.~(\ref{fig13g})   contributing at $N^{(2)}LO$ with the worldline interaction vertex $-Q_a\rmint dt\,\mathcal{A}_0\Bar{\psi} :$
\begin{align}
    \begin{split}
         \mathcal{S}_{\text{eff}}\Big|_{\text{fig}(\ref{fig13g})}&=-Q_1Q_2\rmint \Bar{\mathcal{D}}\hat{\xi}\rmint dt_1dt_2\mathcal{A}_0(\boldsymbol{x}_1(t_1))\mathcal{A}_0(\boldsymbol{x}_2(t_2))\frac{\Bar{\psi}(\boldsymbol{x}_1(t_1))}{m_p}\,,\\&
          =-Q_1Q_2\rmint dt_1 dt_2\Big{
       \langle}\mathcal{A}_0(\boldsymbol{x}_1(t_1))\mathcal{A}_0(\boldsymbol{x}_2(t_2))\Big{\rangle}\frac{\Bar{\psi}(\boldsymbol{x}_1(t_1))}{m_p} \,,\\&
       =Q_1Q_2\rmint dt\frac{1}{4\pi r}\frac{\Bar{\psi}(\boldsymbol{x}_1(t))}{m_p}\,+(1\leftrightarrow 2)\,\sim\mathcal{O}(L^{1/2}v^{5/2})\,.
    \end{split}
\end{align}
%\vspace{-1cm}
%\begin{align}
%    \begin{split} 
%      \mathcal{S}_{\text{eff}}^{tot}& =
%-\frac{\gamma Q_1Q_2'}{2\pi}\rmint dt \, %\,\Big[\boldsymbol{v}_1(t)\cdot \boldsymbol{v}_2(t) \frac{e^{-\mu r}}{r}\Big](\frac{\Bar{\psi}(x_1,t)}{m_p}+\frac{\Bar{\psi}(x_2,t)+\Bar{\psi}(\boldsymbol{x}_0,t)}{m_p})\\&\textstyle{-2Q_1 %Q_2'\rmint dt \, \Big[\boldsymbol{v}_1\cdot \boldsymbol{v}_2\frac{e^{-r \mu _{\gamma }} \left(-r \mu _{\gamma }+e^{r \mu _{\gamma }}-1\right)}{2 \pi  r^3 \mu _{\gamma }^2}
%+(\boldsymbol{v}_1\cdot \Hat{n})(\boldsymbol{v}_2\cdot \Hat{n})\frac{e^{-r \mu _{\gamma }} \Big(r \mu _{\gamma } \left(r \mu _{\gamma }+3\right)-3 e^{r \mu _{\gamma %}}+3\Big)}{2 \pi  r^3 \mu _{\gamma }^2}\Big]\frac{\Bar{\psi}(\boldsymbol{x}_0,t)+\Bar{\psi}(x_{1},t)+\Bar{\psi}(x_{2},t)}{m_p}})\\&+\frac{m_1m_2^2s_2^2g_1}{16\pi^2m_p^4}\rmint dt(1+\frac{3}{2}v_1^2) \frac{e^{-2mr}}{r^2}\Bar{\psi}(x_1,t)\\&
%+\frac{m_1m_2s_1s_2}{m_p} \rmint dt (\frac{1}{m_p}+\frac{3v_1^2}{2m_p}) \frac{1}{4\pi r}e^{-mr}\Bar{\psi}(x_1,t)-2Q_1Q_2\rmint dt \Bar{\psi}\frac{1}{4\pi r}-2Q_1Q_2\rmint dt \Bar{\psi}\frac{1}{4\pi r}e^{-\mu r}\\&
%+\frac{m_1m_2}{4\pi}\rmint dt \Bar{\psi}(x_1,t)\frac{1}{r}-\frac{m_1m_2Q_1Q_2}{16\pi^2m_p^2}\rmint dt \frac{1}{r^2}\Bar{\psi}(x_1,t)-\frac{m_1m_2Q_1'Q_2'}{16\pi^2m_p^2}\rmint dt \frac{e^{-\mu r}}{r^2}\Bar{\psi}(x_1,t)\\&
%\end{split}
%\end{align}
%\vspace{-0.5cm}
%\begin{align}
%\begin{split}
%&-\frac{2Q_1Q_2}{4\pi }\rmint dt (v_1 \cdot v_2)\frac{1}%{r}\Bar{\psi}(x_1,t)-\frac{2Q_1'Q_2'}{4\pi }\rmint dt %(v_1 \cdot v_2)\frac{e^{-\mu r}}{r}\Bar{\psi}(x_1,t)+\\&
%\frac{2Q_1 Q_2 m_1 s_1\epsilon^{0ikm}}{4\pi}\rmint dt \, v_2^m (t)a_1^i(t)\,\partial_r\alpha(m;r)n^k\bar{\psi}(x_2,t)\\&\textstyle{ +\frac{\pi^{3/2}}{4}g_{a\gamma\gamma}Q_2^2\,m_1\,s_1\rmint dt\,\Vec{a_2}\cdot(\hat{n}\times\boldsymbol{v}_2)\Bigg[\frac{1}{2}\, m \,G_{1,3}^{2,1}\left(\frac{m^2 r^2}{4}\Big|
%\begin{array}{c}
% -\frac{1}{2} \\
% -\frac{1}{2},-\frac{1}{2},0 \\
%\end{array}
%\right)-\frac{ G_{1,3}^{2,1}\left(\frac{m^2 r^2}%{4}\Big|
%\begin{array}{c}
% \frac{1}{2} \\
% \frac{1}{2},\frac{1}{2},0 \\
%\end{array}
%\right)}{ \,m\, r^2}\Bigg]}\\ &\Bar{\psi}(x_2,t)-\frac{1}{\pi}\rmint dt \Bigg[(\Vec{a}_1\cdot \Hat{n})\Big\{\frac{  e^{-\mu_{\gamma} r}}{2 \mu_{\gamma} \,r}-\frac{1 -  e^{-\mu_{\gamma} r}}{2\, \mu_{\gamma}^2\, r^2}\Big\}+v_1^2 \,\frac{e^{-r \mu _{\gamma }} \left(-r \mu _{\gamma }+e^{r \mu _{\gamma }}-1\right)}{  r^3 \mu _{\gamma }^2}\\ &
%\,\,\,\,\,\,\,\,\,\,\,\,\,\,\,\,\,\,\,\,\,\,\,\,\,\,\,\,\,
%\textstyle{(\boldsymbol{v}_1\cdot \Hat{n})^2\,\frac{e^{-r \mu _{\gamma }} \Big(r \mu _{\gamma } \left(r \mu _{\gamma }+3\right)-3 e^{r \mu _{\gamma }}+3\Big)}{  r^3 \mu _{\gamma }^2}
%\Bigg](\frac{\Bar{\psi}(x_1,t)}{m_p}+\frac{\Bar{\psi}(x_2,t)}{m_p})-\textstyle{Q_1Q_2'\rmint dt\,\frac{e^{-\mu_{\gamma}r}}{2\pi r}(\frac{\Bar{\psi}(x_1,t)}{m_p}+\frac{\Bar{\psi}(x_2,t)}{m_p})}}\\&-\frac{1}{\pi}\rmint dt \Bigg[(\Vec{a}_2\cdot \Hat{n})\Big\{\frac{  e^{-\mu_{\gamma} r}}{2 \mu_{\gamma} \,r}-\frac{1 -  e^{-\mu_{\gamma} r}}{2\, \mu_{\gamma}^2\, r^2}\Big\}+v_2^2 \,\frac{e^{-r \mu _{\gamma }} \left(-r \mu _{\gamma }+e^{r \mu _{\gamma }}-1\right)}{  r^3 \mu _{\gamma }^2}\\ &
%\,\,\,\,\,\,\,\,\,\,\,\,\,\,\,\,\,\,\,\,\,\,\,\,\,\,\,\,\,
%(\boldsymbol{v}_2\cdot \Hat{n})^2\,\frac{e^{-r \mu _{\gamma }} \Big(r \mu _{\gamma } \left(r \mu _{\gamma }+3\right)-3 e^{r \mu _{\gamma }}+3\Big)}{  r^3 \mu _{\gamma }^2}
%\Bigg](\frac{\Bar{\psi}(x_1,t)}{m_p}+\frac{\Bar{\psi}(x_2,t)}{m_p})
%+1\leftrightarrow 2\\ &
%  \end{split}
%\end{align}
%\begin{align}
 %   \begin{split} 
 %   \mathcal{T}^{00}& =\textstyle{
%+\frac{Q_1Q_2'}{4\pi} \, \,\Bigg[\frac{-m_1m_2r^2\Omega^2}{M^2} \frac{e^{-\mu r}}{r}\Bigg]\delta(\boldsymbol{x}-\boldsymbol{x}_{\text{cm}})-Q_1Q_2' \Big[\frac{-m_1m_2}{M^2}\left((v\cdot n) (v\cdot n)\right)\left(\frac{\pi \left(1-e^{-\mu r}\right)}{\mu^2 r^2}-\frac{\pi}{2}   e^{-\mu r}-\frac{\pi  e^{-\mu r}}{\mu r}\right) \\&+(\frac{-m_1m_2r^2\Omega^2}{M^2}) \left(\frac{\pi  e^{-\mu r}}{2 \mu r^2}-\frac{\pi \left(1-e^{-\mu r}\right)}{2 \mu^2 r^3}\right) \Big]\delta(\boldsymbol{x}-\boldsymbol{x}_1)+\frac{m_1m_2^2s_2^2g_1}{16\pi^2m_p^4}(1+\frac{3}{2}\frac{m_2^2}{M^2}r^2\Omega^2)\frac{e^{-2mr}}{r^2}\delta(\boldsymbol{x}-\boldsymbol{x}_1)\\&
%+\frac{m_1m_2s_1s_2}{m_p^2}  (\frac{1}{m_p}+\frac{3\frac{m_2^2}{M^2}r^2\Omega^2}{2m_p}) \frac{1}{4\pi r}e^{-mr}\delta(\boldsymbol{x}-\boldsymbol{x}_1)-2Q_1Q_2 \frac{1}{4\pi r}\delta(\boldsymbol{x}-\boldsymbol{x}_1)-2Q_1Q_2 \frac{1}{4\pi r}e^{-\mu r}\delta(\boldsymbol{x}-\boldsymbol{x}_1)\\&
%+\frac{m_1m_2}{4\pi} \frac{1}{r}\delta(\boldsymbol{x}-\boldsymbol{x}_1)-\frac{m_1m_2Q_1Q_2}{16\pi^2} \frac{1}{r^2}\delta(\boldsymbol{x}-\boldsymbol{x}_1)-\frac{m_1m_2Q_1'Q_2'}{16\pi^2} \frac{e^{-\mu r}}{r^2}\delta(\boldsymbol{x}-\boldsymbol{x}_1)\\&
%-\frac{2Q_1Q_2}{4\pi } (\frac{-m_1m_2r^2\Omega^2}{M^2})\frac{1}{r}\delta(\boldsymbol{x}-\boldsymbol{x}_1)-\frac{2Q_1Q_2}{4\pi } (\frac{-m_1m_2r^2\Omega^2}{M^2})\frac{e^{-\mu r}}{r}\delta(\boldsymbol{x}-\boldsymbol{x}_1)+1\leftrightarrow 2}\\ &
 % \end{split}
%\end{align}
%\vspace{-2.5cm}
\item Diagram as shown in Fig.~(\ref{fig13n}) consists of a scalar propagator with relativistic correction, contributing to the gravitational radiation at $N^{(4)}LO.$ The amplitude is given by, 
\begin{align}
    \begin{split}
  \mathcal{S}_{\text{eff}}^{ikm}\Big |_{{\text{fig}(\ref{fig13n})}} &=\frac{m_1m_2 s_1 s_2}{m_p^3}\rmint \Bar{\mathcal{D}}\hat\xi\rmint dt_1 \varphi(\boldsymbol{x}_1(t_1))\bar{\psi}(x_0) \rmint d^4 x\partial_0\varphi(x)\partial_0\varphi(x)\rmint dt_2 \varphi(\boldsymbol{x}_2(t_2))\,, \\&
=\frac{m_1 m_2 s_1 s_2}{16 \pi m_p^2}\rmint dt \,e^{-mr}\Bigg\{\Big(\frac{\boldsymbol{v}_1\cdot\boldsymbol{v}_2}{r}-\frac{(\boldsymbol{v}_1\cdot \hat{n})(\boldsymbol{v}_2\cdot \hat{n})}{r}\Big)+\frac{m}{2}\,(\boldsymbol{v}_2\cdot \Hat{n})(\boldsymbol{v}_1\cdot \hat{n})\Bigg\}\,\frac{\Bar{\psi}(\boldsymbol{x}_0(t)}{m_p}\\&+T(\Dot{\Bar{\psi}})+ (1 \leftrightarrow 2)
\sim \mathcal{O}({L^{1/2}v^{9/2}}),
\end{split}
\end{align}
where $T(\dot{\bar{\psi}})$ is given in  equation (\ref{5.26}) and we just need to replace $\phi$ by $\psi$. 
\begin{align}
    \begin{split}
        T(\Dot{\Bar{\psi}})&=-\frac{m_1 m_2 s_1 s_2}{16\pi m_p^2}\rmint dt\, e^{-mr}(\boldsymbol{v}_2\cdot \Hat{n})\frac{ \Dot{\Bar{\psi}}(\boldsymbol{x}_0(t))}{m_p},\\ &
        =\frac{m_1 m_2 s_1 s_2}{16\pi m_p^2}\rmint dt \frac{e^{-mr}}{r}(\boldsymbol{v}_2\cdot \Dot{\boldsymbol{r}}+\Vec{a}_2\cdot \boldsymbol{r})\,\frac{\Bar{\psi}(\boldsymbol{x}_0(t))}{m_p}\,.
    \end{split}
\end{align}
\item \textcolor{black}{The amplitude corresponding to the diagram with worldline interaction vertex $Q_a'\rmint dt\,v_a\Bar{\psi}\mathcal{B}_i$ in Fig.~(\ref{fig13a}) which contributes at $N^{(4)}LO :$ }
\begin{align}
    \begin{split}
         \mathcal{S}_{\text{eff}}\Big |_{\text{fig}(\ref{fig13a})}&=-Q_1' Q_2'\rmint\Bar{\mathcal{D}}\Hat{\xi}\rmint dt_1\,v_1^i(t_1)\mathcal{B}_i(\boldsymbol{x}_1(t_1))\rmint dt_2 \,v_2^j(t_2)\mathcal{B}_j(\boldsymbol{x}_2(t_2))\frac{\Bar{\psi}(\boldsymbol{x}_1(t_1))}{m_p}\,,\\ &
         =-Q_1'Q_2'\rmint dt_1 dt_2 v_1^i(t_1)v_2^{j}(t_2)\Big\langle\mathcal{B}_i(\boldsymbol{x}_1(t_1))\mathcal{B}_j(\boldsymbol{x}_2(t_2))\Big\rangle\,\frac{\bar\psi(\boldsymbol{x}_1(t_1))}{m_p},\\ &
       %=-2Q_1'Q_2'\rmint dt_1 dt_2v_1^iv_2^j\Big{
       %\langle}B_1^i(\boldsymbol{x}_1(t_1))B_2^j(\boldsymbol{x}_2(t_2))\Big{\rangle}\frac{\Bar{\psi}}{m_p} \\&
       %=-2Q_1'Q_2'\rmint dt_1 v_1^iv_2^j \delta_{ij}\frac{1}{4\pi %r}e^{-\mu r}\frac{\Bar{\psi}}{m_p}\\&
       =\frac{-Q_1'Q_2'}{4\pi }\rmint dt \,(\boldsymbol{v}_1 \cdot \boldsymbol{v}_2)\frac{e^{-\mu r}}{r}\frac{\Bar{\psi}(\boldsymbol{x}_1(t))}{m_p}+(1\leftrightarrow 2)\,\sim \mathcal{O}(L^{1/2}v^{9/2}).
    \end{split}
\end{align}
\item The amplitude corresponding to the diagram with the worldline vertex $Q_a\rmint dt\,v_a\Bar{\psi}\mathcal{A}_i$ in Fig.~(\ref{fig13b}) which contributes at $N^{(4)}LO :$ 
\begin{align}
    \begin{split}
         \mathcal{S}_{\text{eff}}\Big |_{\text{fig}(\ref{fig13b})}&=-Q_1 Q_2\rmint\Bar{\mathcal{D}}\Hat{\xi}\rmint dt_1\,v_1^i(t_1)\mathcal{A}_i(\boldsymbol{x}_1(t_1))\rmint dt_2 \,v_2^j(t_2)\mathcal{A}_j(\boldsymbol{x}_2(t_2))\frac{\Bar{\psi}(\boldsymbol{x}_1(t_1))}{m_p}\,,\\ &
         =-Q_1Q_2\rmint dt_1 dt_2 v_1^i(t_1)v_2^{j}(t_2)\Big\langle\mathcal{A}_i(\boldsymbol{x}_1(t_1))\mathcal{A}_j(\boldsymbol{x}_2(t_2))\Big\rangle\,\frac{\bar\psi(\boldsymbol{x}_1(t_1))}{m_p},\\ &
       %=-2Q_1'Q_2'\rmint dt_1 dt_2v_1^iv_2^j\Big{
       %\langle}B_1^i(\boldsymbol{x}_1(t_1))B_2^j(\boldsymbol{x}_2(t_2))\Big{\rangle}\frac{\Bar{\psi}}{m_p} \\&
       %=-2Q_1'Q_2'\rmint dt_1 v_1^iv_2^j \delta_{ij}\frac{1}{4\pi %r}e^{-\mu r}\frac{\Bar{\psi}}{m_p}\\&
       =\frac{-Q_1Q_2}{4\pi }\rmint dt \,\frac{\boldsymbol{v}_1 \cdot \boldsymbol{v}_2}{r}\frac{\Bar{\psi}(\boldsymbol{x}_1(t))}{m_p}+(1\leftrightarrow 2)\,\sim \mathcal{O}(L^{1/2}v^{9/2}).
    \end{split}
\end{align}
\item The amplitude corresponding to the diagram in Fig.~(\ref{fig13c}) with worldline interaction vertex $Q_a'\rmint dt\,\mathcal{B}_0\Tilde{\psi}\Bar{\psi}$ which contributes at $N^{(4)}LO :$ 
\begin{align}
    \begin{split}
         \mathcal{S}_{\text{eff}}\Big |_{\text{fig}(\ref{fig13c})}&=-\frac{m_2Q_1'Q_2'}{m_p^2}\rmint\Bar{\mathcal{D}}\Hat{\xi}\rmint dt_1\mathcal{B}_0(\boldsymbol{x}_1(t_1))\Tilde{\psi}(\boldsymbol{x}_1(t_1))\rmint dt_2\Tilde{\psi}(\boldsymbol{x}_2(t_2)\rmint dt_3\mathcal{B}_0(\boldsymbol{x}_2(t_3))\frac{\Bar{\psi}(\boldsymbol{x}_1(t_1))}{m_p}\\&
         =-\frac{m_2Q_1'Q_2'}{m_p^2}'\rmint dt_1dt_2dt_3\Big{\langle}\Tilde{\psi}(\boldsymbol{x}_2(t_2)\Tilde{\psi}(\boldsymbol{x}_1(t_1))\Big{\rangle}\Big{\langle}\mathcal{B}_0(\boldsymbol{x}_2(t_3))\mathcal{B}_0(\boldsymbol{x}_1(t_1))\Big{\rangle}\frac{\Bar{\psi}(\boldsymbol{x}_1(t_1))}{m_p}\,,\\&
     %  =-\frac{m_1m_2Q_1'Q_2'}{m_p^2}\rmint dt \frac{1}{(4\pi)^{\frac{3}{2}}}\sqrt{\pi}\frac{2}{r}\frac{1}{4\pi r}e^{-\mu r}\Bar{\psi}(x_1,t_1)\,,\\&
       =-\frac{m_2Q_1'Q_2'}{32m_p^2\pi^2}\rmint dt \frac{e^{-\mu r}}{r^2}\frac{\Bar{\psi}(\boldsymbol{x}_1(t))}{m_p}\,+(1\leftrightarrow 2)\,\sim\mathcal{O}(L^{1/2}v^{9/2}).
    \end{split}
\end{align}
\item The amplitude corresponding to the diagram in Fig.~(\ref{fig13d}) with worldline coupling $Q_a\rmint dt\,\mathcal{A}_0\tilde\psi\Bar{\psi}$ which contributes at $N^{(4)}LO :$ 
\begin{align}
    \begin{split}
         \mathcal{S}_{\text{eff}}\Big |_{\text{fig}(\ref{fig13d})}&=-\frac{m_2Q_1Q_2}{m_p^2}\rmint\Bar{\mathcal{D}}\Hat{\xi}\rmint dt_1\mathcal{A}_0(\boldsymbol{x}_1(t_1))\Tilde{\psi}(\boldsymbol{x}_1(t_1))\rmint dt_2\Tilde{\psi}(\boldsymbol{x}_2(t_2)\rmint dt_3\mathcal{A}_0(\boldsymbol{x}_2(t_3))\frac{\Bar{\psi}(\boldsymbol{x}_1(t_1))}{m_p}\\&=-\frac{m_2Q_1Q_2}{m_p^2}\rmint dt_1dt_2dt_3\Big{\langle}\tilde{\psi}(\boldsymbol{x}_2(t_2)\tilde{\psi}(\boldsymbol{x}_1(t_1))\Big{\rangle}\Big{\langle}\mathcal{A}_0(\boldsymbol{x}_2(t_3)\mathcal{A}_0(\boldsymbol{x}_1(t_1))\Big{\rangle}\frac{\Bar{\psi}(\boldsymbol{x}_1(t_1))}{m_p}\,,\\&
     %  =-\frac{m_1m_2Q_1Q_2}{m_p^2}\rmint dt \frac{1}{(4\pi)^{\frac{3}{2}}}\sqrt{\pi}\frac{2}{r}\frac{1}{4\pi r}\frac{\Bar{\psi}}{m_p}\,,\\&
       =-\frac{m_2Q_1Q_2}{32m_p^2\pi^2}\rmint dt \frac{1}{r^2}\frac{\Bar{\psi}(\boldsymbol{x}_1(t))}{m_p}+(1\leftrightarrow 2)\,\,\sim\mathcal{O}(L^{1/2}v^{9/2})\,.
    \end{split}
\end{align}
 \item The amplitude corresponding to the diagram in Fig.~(\ref{fig13h}) is given by: 
\begin{align}
    \begin{split}
         \mathcal{S}_{\text{eff}}\Big |_{\text{fig}(\ref{fig13h})}&=\frac{m_1m_2s_1s_2}{m_p^2}\rmint \Bar{\mathcal{D}}\hat{\xi}\rmint dt_1  \Big(\frac{1}{m_p}+\frac{3v_1^2}{2m_p}\Big)\varphi(\boldsymbol{x}_1(t_1)){\Bar{\psi}(\boldsymbol{x}_1(t_1))}\rmint dt_2\varphi(\boldsymbol{x}_1(t_2))\,,\\&
       =\frac{m_1m_2s_1s_2}{m_p^2}\rmint dt_1 dt_2 \Big(\frac{1}{m_p}+\frac{3v_1^2}{2m_p}\Big) \Big{\langle}\varphi(\boldsymbol{x}_1(t_1)) \varphi(\boldsymbol{x}_2(t_2))\Big{\rangle}{\Bar{\psi}(\boldsymbol{x}_1(t_1))}\,,\\&
       =\frac{m_1m_2s_1s_2}{m_p^2} \rmint dt_1 dt_2 \Big(\frac{1}{m_p}+\frac{3v_1^2}{2m_p}\Big)\delta(t_1-t_2)\rmint \frac{d^3k}{(2\pi)^3}\frac{e^{i\boldsymbol{k}\cdot\boldsymbol{r}}}{\boldsymbol{k}^2+m^2}{\Bar{\psi}(\boldsymbol{x}_1(t_1))}\,,\\ &
       =\frac{m_1m_2s_1s_2}{m_p^2} \rmint dt \Big(1+\frac{3v_1^2}{2}\Big) \frac{1}{4\pi r}e^{-mr}\frac{\Bar{\psi}(\boldsymbol{x}_1(t))}{m_p}+(1\leftrightarrow 2)\,,\\&\sim \mathcal{O}(L^{1/2}v^{5/2})+\mathcal{O}(L^{1/2}v^{9/2})\,.
    \end{split}
\end{align}
Note that, a part of it contributes at $N^{(2)}LO$ and the other part at $N^{(4)}LO.$ 
 \item Amplitude corresponding to the diagram in Fig.~(\ref{fig13i})  with the radiation coming from a triangle vertex $\phi\phi\Bar{\psi}$  is given by: 
\begin{align}
    \begin{split}
         \mathcal{S}_{\text{eff}}\Big |_{\text{fig}(\ref{fig13i})}&=\frac{m_1m_2^2s_2^2g_1}{m_p^4}\rmint \Bar{\mathcal{D}}\hat{\xi}\rmint dt_1\Big(1+\frac{3}{2}v_1^2\Big)\varphi(\boldsymbol{x}_1(t_1))\varphi(\boldsymbol{x}_1(t_1))\rmint dt_2 \varphi(\boldsymbol{x}_2(t_2)) \,\, \\& \,\,\,\,\,\,\,\,\,\,\,\,\,\,\,\,\,\,\,\,\,\,\,\,\,\,\,\,\,\,\,\,\,\,\,\,\,\,\,\,\,\,\,\,\,\,\,\,\,\,\,\,\,\,\,\,\,\,\,\,\,\,\,\,\,\,\,\,\,\,\,\,\,\,\,\,\,\,\,\,\,\,\,\,\,\,\,\,\,\,\,\,\,\,\,\,\,\,\,\,\,\,\,\,\,\,\,\,\,\,\,\,\,\,\,\,\,\,\,\,\,\,\,\,\,\,\rmint dt_3\varphi(\boldsymbol{x}_2(t_3)) \frac{\Bar{\psi}(\boldsymbol{x}_1(t_1))}{m_p}\,,\\&
         =\frac{m_1m_2^2s_2^2g_1}{m_p^4}\rmint dt_1 dt_2 dt_3\Big(1+\frac{3}{2}v_1^2\Big)\Big{\langle}\varphi(\boldsymbol{x}_1(t_1)) \varphi(\boldsymbol{x}_2(t_2))\Big{\rangle}\Big{\langle}\varphi(\boldsymbol{x}_1(t_1)) \varphi(\boldsymbol{x}_2(t_3))\Big{\rangle}\\& \,\,\,\,\,\,\,\,\,\,\,\,\,\,\,\,\,\,\,\,\,\,\,\,\,\,\,\,\,\,\,\,\,\,\,\,\,\,\,\,\,\,\,\,\,\,\,\,\,\,\,\,\,\,\,\,\,\,\,\,\,\,\,\,\,\,\,\,\,\,\,\,\,\,\,\,\,\,\,\,\,\,\,\,\,\,\,\,\,\,\,\,\,\,\,\,\,\,\,\,\,\,\,\,\,\,\,\,\,\,\,\,\,\,\,\,\,\,\,\,\,\,\,\,\,\, \frac{\Bar{\psi}(\boldsymbol{x}_1(t_1))}{m_p}\,,\\&
    %   =\frac{m_1m_2^2s_2^2g_1}{m_p^4}\rmint dt(1+\frac{3}{2}v_1^2) \frac{e^{-mr}}{4\pi r}\frac{e^{-mr}}{4\pi r}\frac{\Bar{\psi}}{m_p}\,,\\&
       =\frac{m_1m_2^2s_2^2g_1}{16\pi^2m_p^4}\rmint dt \Big(1+\frac{3}{2}v_1^2\Big)\frac{e^{-2mr}}{r^2}\frac{\Bar{\psi}(\boldsymbol{x}_1(t))}{m_p}+(1\leftrightarrow 2)\,,\\&\sim \mathcal{O}(L^{1/2}v^{9/2})+\mathcal{O}(L^{1/2}v^{13/2}).
    \end{split}
\end{align}
Again note that a part of it is contributing at $N^{(4)}LO$ and the other part is contributing at $N^{(6)}LO.$
%\newpage
\item The amplitude corresponds to the diagram in Fig.~(\ref{fig13j}) which contributes at $N^{(4)}LO$ with the bulk vertex $\rmint d^4x\,\psi \partial_l\mathcal{A}_i\partial_l\mathcal{B}_i :$ 
\vspace{-0.3cm}
\begin{align}
    \begin{split}
   \mathcal{S}_{\text{eff}}\Big |_{\text{fig}(\ref{fig13j})}&=-4\gamma Q_1Q_2'\rmint \Bar{\mathcal{D}}\hat{\xi}\rmint dt_1\mathcal{A}_j(\boldsymbol{x}_1(t_1))v_1^j\rmint dt_2\mathcal{B}_m(\boldsymbol{x}_2(t_2))v_2^m(t_2)\,\\&  \,\,\,\,\,\,\,\,\,\,\,\,\,\,\,\,\,\,\,\,\,\,\,\,\,\,\,\,\,\, \,\,\,\,\,\,\,\,\,\,\,\,\,\,\,\,\,\,\,\,\,\,\,\,\,\,\,\,\,\, \,\,\,\,\,\,\,\,\,\,\,\,\,\,\,\,\,\,\,\,\,\,\,\,\,\,\,\,\,\, \,\,\,\,\,\,\,\,\,\,\,\,\,\,\,\,\,\,\,\,\,\,\,\, \rmint d^4x\partial_l\mathcal{A}_i\partial_l\mathcal{B}_i \frac{ \Bar{\psi}(\boldsymbol{x}_0(t))}{m_p}\\&=-4\gamma Q_1Q_2'\rmint \prod_{i=1}^2dt_iv_1^j(t_1)v_2^m(t_2)\rmint d^4x\partial_l\Big{\langle}\mathcal{A}_i(x)\mathcal{A}_j(\boldsymbol{x}_1(t_1))\Big{\rangle}\\&\hspace{6cm}\partial_l\Big{\langle}\mathcal{B}_i(x)\mathcal{B}_m(\boldsymbol{x}_2(t_2))\Big{\rangle}\frac{\Bar{\psi}(\boldsymbol{x}_0(t))}{m_p}\,,\\ &
%=-2Q_1Q_2'\rmint \prod_{i=1}^2dt_iv_1^i(t_1)v_2^m(t_2)\rmint d^4x\partial_l\rmint\delta(t-t_1)\frac{e^{-i\boldsymbol{k}\cdot [\boldsymbol{x}(t)-\boldsymbol{x}_1(t_1)]}}{k^2}\delta_{im}\\& \hspace{6cm}\partial_l
%\rmint\delta(t-t_2)\frac{e^{-i\boldsymbol{k}_1\cdot[\boldsymbol{x}(t)-\boldsymbol{x}_2(t_2)]}}{k_1^2+\mu^2}\frac{\Bar{\psi}}{m_p}\,,\\&
%=+2Q_1Q_2'\rmint dt_1\, %dt_2\,dt\,v_1^i(t_1)v_2^i(t_2)\delta(t-%t_1)\rmint_{\boldsymbol{k},\boldsymbol{k}_1}\delta(\boldsymbol{k}+ \vec %k_1)\delta(t-t_2)\\&\frac{k^lk_1^le^{i\vec %k\cdot \boldsymbol{x}_1(t_1)}e^{i\boldsymbol{k}_1\cdot \vec %x_2(t_2)}}{k^2(k_1^2+\mu^2)}\frac{\Bar{\psi}}%{m_p}\\ &
%=-2Q_1Q_2'\rmint dt\, \,\Big[v_1(t)\cdot v_2(t)\rmint_k \frac{e^{-i\boldsymbol{k}\cdot [\boldsymbol{x}_2(t)-\boldsymbol{x}_1(t)]}}{(\boldsymbol{k}^2+\mu^2)}\Big]\frac{\Bar{\psi}}{m_p}\,,\\ &
=-4\gamma Q_1Q_2'\rmint dt\, \,\Bigg[\boldsymbol{v}_1(t)\cdot \boldsymbol{v}_2(t)\rmint_k \frac{e^{-i\boldsymbol{k}\cdot \boldsymbol{r}}}{(\boldsymbol{k}^2+\mu^2)}\Bigg]                \frac{\Bar{\psi}(\boldsymbol{x}_0(t))}{m_p}\,,\\ &
=-\frac{\gamma Q_1 Q_2'}{\pi}\rmint dt\, \boldsymbol{v}_1\cdot \boldsymbol{v}_2\, \frac{e^{-\mu_{\gamma}r}}{r}\Big(\frac{\Bar{\psi}(\boldsymbol{x}_{0}(t))}{m_p}\Big)+(1\leftrightarrow 2)\,\sim \mathcal{O}(L^{1/2}v^{9/2}).
\end{split}
\end{align}
%\vspace{-0.7 cm}
\item \textcolor{black}{ The amplitude that governs the effective action corresponding to the bulk interaction vertex $\rmint d^4x\,\psi \partial_l\mathcal{A}_i\partial_i\mathcal{B}_l$ at $N^{(4)}LO$,} 
\begin{align}
    \begin{split}
   \mathcal{S}_{\text{eff}}\Big |_{\text{fig}(\ref{fig13j})}&=-4\gamma Q_1Q_2'\rmint \Bar{\mathcal{D}}\hat{\xi}\rmint dt_1\mathcal{A}_j(\boldsymbol{x}_1(t_1))v_1^j\rmint dt_2\mathcal{B}_m(\boldsymbol{x}_2(t_2))v_2^m(t_2)\\& \,\,\,\,\,\,\,\,\,\,\,\,\,\,\,\,\,\,\,\,\,\,\,\,\,\,\,\,\,\, \,\,\,\,\,\,\,\,\,\,\,\,\,\,\,\,\,\,\,\,\,\,\,\,\,\,\,\,\,\, \,\,\,\,\,\,\,\,\,\,\,\,\,\,\,\,\,\,\,\,\,\,\,\,\,\,\,\,\,\,\,\,\,\,\,\,\,\,\,\,\,\,\,\,\,\,\,\,\,\,\,\,\,\,\,\,\,\,\,\,  \rmint d^4x\,\partial_l\mathcal{A}_i\partial_i\mathcal{B}_l \frac{ \Bar{\psi}(\boldsymbol{x}_0(t))}{m_p}\\&
   =-4\gamma\textstyle{Q_1 Q_2'\rmint dt_1dt_2v_1^j(t_1)v_2^m(t_2)\rmint d^4x\partial_l\Big{\langle}\mathcal{A}_i(x)\mathcal{A}_j(\boldsymbol{x}_1(t_1))\Big{\rangle}}\\&\hspace{6cm}\textstyle{\partial_i\Big{\langle}\mathcal{B}_l(x)\mathcal{B}_m(\boldsymbol{x}_2(t_2))\Big{\rangle}\frac{\Bar{\psi}(\boldsymbol{x}_0(t))}{m_p}}\,,\\ &
=-4\gamma Q_1 Q_2'\rmint dt \, \Big[\boldsymbol{v}_1\cdot \boldsymbol{v}_2\frac{e^{-r \mu_{\gamma }} \left(-r \mu _{\gamma }+e^{r \mu _{\gamma }}-1\right)}{2 \pi  r^3 \mu_{\gamma }^2}\\ & \hspace{1cm}
+(\boldsymbol{v}_1\cdot \Hat{n})(\boldsymbol{v}_2\cdot \Hat{n})\frac{e^{-r \mu _{\gamma }} \Big(r \mu _{\gamma } \left(r \mu _{\gamma }+3\right)-3 e^{r \mu _{\gamma }}+3\Big)}{2 \pi  r^3 \mu _{\gamma }^2}\Big]\Big(\frac{\Bar{\psi}(\boldsymbol{x}_0(t))}{m_p}\Big)+(1\leftrightarrow 2)\,\\&\sim\mathcal{O}(L^{1/2}v^{9/2}).
\end{split} 
\end{align}
%\newpage
\vspace{-0.6cm}
\item The amplitude corresponds to the diagram in Fig.~(\ref{fig13k}) that contributes at $N^{(4)}LO$ with the bulk interaction vertex $\rmint d^4x\,\partial_k\mathcal{A}_0\partial_0 \mathcal{B}_k :$
%\vspace{-0.3cm}
\begin{align}
   \mathcal{S}_{\text{eff}}\Big|_{\text{fig}(\ref{fig13k})}
   &=-\gamma Q_1Q_2'\rmint \Bar{\mathcal{D}}\hat{\xi}\rmint dt_1\mathcal{A}_0(\boldsymbol{x}_1(t_1))\rmint dt_2\mathcal{B}_j(\boldsymbol{x}_2(t_2))v_2^j(t_2)\rmint d^4x \partial_k\mathcal{A}_0\partial_0\mathcal{B}_k(x) \frac{ \Bar{\psi}}{m_p}\,,\notag\\
   &=-\gamma Q_1Q_2'\rmint dt_1dt_2v_2^j(t_2)\rmint d^4x\partial_k\Big{\langle}\mathcal{A}_0(x)\mathcal{A}_0(\boldsymbol{x}_1(t_1))\Big{\rangle}\partial_0\Big{\langle}\mathcal{B}_k(x) \mathcal{B}_j(\boldsymbol{x}_2(t_2))\Big{\rangle}\frac{\Bar{\psi}}{m_p}\,,\notag\\
%=\rmint dt_1\, dt_2\,v_2^k(t_2)\rmint d^4x\partial_k\rmint\delta(t-t_1)\frac{e^{-i\boldsymbol{k}\cdot [\boldsymbol{x}(t)-\boldsymbol{x}_1(t_1)]}}{k^2}\partial_0\rmint\delta(t-t_2)\frac{e^{-i\boldsymbol{k}_1\cdot[\boldsymbol{x}(t)-\boldsymbol{x}_2(t_2)]}}{k_1^2+\mu^2}\\&
%=-i\rmint dt_1\, dt_2\,dt\,v_2^k(t_2)\partial_0\delta(t-t_2)\rmint_{\boldsymbol{k},\boldsymbol{k}_1}\delta(\boldsymbol{k}+ \boldsymbol{k}_1)\delta(t-t_1)\frac{k^ke^{i\boldsymbol{k}\cdot \boldsymbol{x}_1(t_1)}e^{i\boldsymbol{k}_1\cdot \boldsymbol{x}_2(t_2)}}{k^2(k_1^2+\mu^2)}\\ &
%=-i\rmint dt\, \,\partial^{t_2}\Big[v_2^k(t_2)\rmint \frac{e^{-i\boldsymbol{k}\cdot [\boldsymbol{x}_2(t_2)-\boldsymbol{x}_1(t)]}\,k^k}{\boldsymbol{k}^2(\boldsymbol{k}^2+\mu^2)}\Big]_{t_2=t}\\ &
%= -i\gamma Q_1Q_2'\rmint dt \Bigg[a^{k}_{2}(t)\rmint_{\boldsymbol{k}}\frac{k^k\,e^{-i\,\boldsymbol{k}\cdot\boldsymbol{r}}}{\boldsymbol{k}^2(\boldsymbol{k}^2+\mu^2)}-v_2^{k}(t)v_{2}^{a}(t)\frac{i\,k^a\,k^k\,e^{-i\boldsymbol{k}\cdot\boldsymbol{r}}}{\boldsymbol{k}^2(\boldsymbol{k}^2+\mu^2)}\Bigg]\frac{\Bar{\psi}}{m_p}\,,\\ &
%=-i\rmint dt \partial^{t_1}\Big[v_1^i(t_1)\Big(\rmint \frac{e^{-i\boldsymbol{k}\cdot (x_2(t)-x_1(t_1)}k^i}{k^2}-\frac{e^{-i\boldsymbol{k}\cdot \vec r}k^i}{k^2+\mu^2}\Big)\Big]_{t_1=t}\\&
   &=\frac{\gamma Q_1Q_2'}{2\pi}\rmint dt \Bigg[(\Vec{a}_2\cdot \Hat{n})\Big\{\frac{  e^{-\mu_{\gamma} r}}{2 \mu_{\gamma} \,r}-\frac{1 -  e^{-\mu_{\gamma} r}}{2\, \mu_{\gamma}^2\, r^2}\Big\}+v_2^2 \,\frac{e^{-r \mu _{\gamma }} \left(-r \mu _{\gamma }+e^{r \mu _{\gamma }}-1\right)}{  r^3 \mu _{\gamma }^2}\notag\\
   &
\,\,\,\,\,\,\,\,\,\,\,\,\,\,\,\,\,\,\,\,\,\,\,\,\,\,\,\,\,
(\boldsymbol{v}_2\cdot \Hat{n})^2\,\frac{e^{-r \mu _{\gamma }} \Big(r \mu _{\gamma } \left(r \mu _{\gamma }+3\right)-3 e^{r \mu _{\gamma }}+3\Big)}{  r^3 \mu _{\gamma }^2}
\Bigg]\Big(\frac{\Bar{\psi}(\boldsymbol{x}_1(t))}{m_p}+\frac{\Bar{\psi}(\boldsymbol{x}_2(t))}{m_p}\Big)\notag\\
   &\quad +(1\leftrightarrow 2)\,\sim\mathcal{O}(L^{1/2}v^{9/2}).
\end{align}
\vspace{-0.7cm}
\item  The amplitude corresponds to the diagram in Fig.~(\ref{fig13m}) that contributes at $N^{(4)}LO$ with bulk interaction vertex $\rmint d^4x\,\partial_0\mathcal{A}_i\partial_i\mathcal{B}_0$ :\footnote{It is important to remember that we have diagrams with asymmetric vertices, i.e., vertices connected with two different fields. Hence, along with shuffling the worldlines, we also have to consider the radiation from the other field vertex (e.g., if we only shuffle the worldlines indices, it will generate a diagram producing radiation from the same field but from different worldlines). But for asymmetric vertices, it demands considering the same diagram with radiation from the other field coupled to one of the worldlines before shuffling the indices).}
\begin{align}
   \mathcal{S}_{\text{eff}}\Big |_{{\text{fig}(\ref{fig13m})}}
   &=-\gamma Q_1Q_2'\rmint \Bar{\mathcal{D}}\hat{\xi}\rmint dt_1v_1^j(t_1)\mathcal{A}_j(\boldsymbol{x}_1(t_1))\rmint dt_2\mathcal{B}_0(\boldsymbol{x}_2(t_2))\rmint d^4x\,\partial_0\mathcal{A}_i(x)\partial_i\mathcal{B}_0(x)\frac{\Bar{\psi}}{m_p}\,,\notag\\
   &=-\gamma Q_1Q_2'\rmint dt_1dt_2v_1^j(t_1)\rmint d^4x\,\partial_0\Big{\langle}\mathcal{A}_i(x)\mathcal{A}_j(\boldsymbol{x}_1(t_1))\Big{\rangle}\partial_i\Big{\langle}\mathcal{B}_0(x)\mathcal{B}_0(\boldsymbol{x}_2(t_2))\Big{\rangle}\frac{\Bar{\psi}}{m_p}\,,\notag\\
%=\rmint dt_1\, dt_2\,v_1^i(t_1)\rmint d^4x\partial_0\rmint\delta(t-t_1)\frac{e^{-i\boldsymbol{k}\cdot [\boldsymbol{x}(t)-\boldsymbol{x}_1(t_1)]}}{k^2}\delta_{ij}\partial_i\rmint\delta(t-t_2)\frac{e^{-i\boldsymbol{k}_1\cdot[\boldsymbol{x}(t)-\boldsymbol{x}_2(t_2)]}}{k_1^2+\mu^2}\\&
%=-i\rmint dt_1\, dt_2\,dt\,v_1^i(t_1)\partial_0\delta(t-t_1)\rmint_{\boldsymbol{k},\boldsymbol{k}_1}\delta(\boldsymbol{k}+ \boldsymbol{k}_1)\delta(t-t_2)\frac{k_1^ie^{i\boldsymbol{k}\cdot \boldsymbol{x}_1(t_1)}e^{i\boldsymbol{k}_1\cdot \boldsymbol{x}_2(t_2)}}{k^2(k_1^2+\mu_{\gamma}^2)}\\ &
%=-i\rmint dt\, \,\partial_{t_1}\Big[v_1^i(t_1)\rmint \frac{e^{-i\boldsymbol{k}\cdot [\boldsymbol{x}_2(t)-\boldsymbol{x}_1(t_1)]}\,k^i}{\boldsymbol{k}^2(\boldsymbol{k}^2+\mu^2)}\Big]_{t_1=t}\\ &
%=i\gamma\,Q_1Q_2'\rmint dt \Bigg[a^{i}_{1}(t)\rmint_{\boldsymbol{k}}\frac{k^i\,e^{-i\,\boldsymbol{k}\cdot\boldsymbol{r}}}{\boldsymbol{k}^2(\boldsymbol{k}^2+\mu_{\gamma}^2)}+v_1^{i}(t)v_{1}^{a}(t)\frac{i\,k^a\,k^i\,e^{-i\boldsymbol{k}\cdot\boldsymbol{r}}}{\boldsymbol{k}^2(\boldsymbol{k}^2+\mu_{\gamma}^2)}\Bigg]\frac{\Bar{\psi}}{m_p}\,,\\ &
%=-i\rmint dt \partial^{t_1}\Big[v_1^i(t_1)\Big(\rmint \frac{e^{-i\boldsymbol{k}\cdot (x_2(t)-x_1(t_1)}k^i}{k^2}-\frac{e^{-i\boldsymbol{k}\cdot \vec r}k^i}{k^2+\mu^2}\Big)\Big]_{t_1=t}\\&
   &=\frac{\gamma Q_1Q_2'}{2\pi}\rmint dt \Bigg[(\Vec{a}_1\cdot \Hat{n})\Big\{\frac{  e^{-\mu_{\gamma} r}}{2 \mu_{\gamma} \,r}-\frac{1 -  e^{-\mu_{\gamma} r}}{2\, \mu_{\gamma}^2\, r^2}\Big\}+v_1^2 \,\frac{e^{-r \mu _{\gamma }} \left(-r \mu _{\gamma }+e^{r \mu _{\gamma }}-1\right)}{  r^3 \mu _{\gamma }^2}\notag\\
   &
\,\,\,\,\,\,\,\,\,\,\,\,\,\,\,\,\,\,\,\,\,\,\,\,\,\,\,\,\,
(\boldsymbol{v}_1\cdot \Hat{n})^2\,\frac{e^{-r \mu _{\gamma }} \Big(r \mu _{\gamma } \left(r \mu _{\gamma }+3\right)-3 e^{r \mu _{\gamma }}+3\Big)}{  r^3 \mu _{\gamma }^2}
\Bigg]\Big(\frac{\Bar{\psi}(\boldsymbol{x}_1(t))}{m_p}+\frac{\Bar{\psi}(\boldsymbol{x}_2(t))}{m_p}\Big)\notag\\
   &\quad +(1\leftrightarrow 2)\,,\sim\mathcal{O}(L^{1/2}v^{9/2})\,.\end{align}
\item The amplitude corresponds to the diagram in Fig.~(\ref{fig13o}) that contributes at $N^{(7)}LO$ with the  3-point bulk interaction vertex $\rmint d^4x\,\varphi\,\partial_{0}\mathcal{A}_{i}\partial_{k}\mathcal{A}_{m}(x) :$ 
\begin{align}
          \mathcal{S}_{\text{eff}}^{ikm}\Big |_{{\text{fig}(\ref{fig13o})}}
          &=-\frac{Q_2^2m_1s_1}{m_p^2}\rmint\Bar{\mathcal{D}}\Hat{\xi}\rmint dt_2\, v_2^{j}(t_2)\mathcal{A}_{j}(\boldsymbol{x}_2(t_2))\rmint dt_3 \,v_2^{l}(t_3)\mathcal{A}_{l}(\boldsymbol{x}_2(t_3))\notag\\
          &\quad \rmint dt_1\,\varphi(\boldsymbol{x}_1(t_1))\frac{\Bar{\psi}}{m_p}(\boldsymbol{x}_1(t_1))\rmint d^4x \,\varphi\,\partial_0 \mathcal{A}_{i}\partial_{k} \mathcal{A}_{m}(x)\,,\notag\\
        &=-\frac{Q_2^2m_1s_1}{m_p^2}\rmint dt_1dt_2dt_3dt\,v_2^{j}(t_2)v_2^{l}(t_3)\frac{\Bar{\psi}(\boldsymbol{x}_1(t_1))}{m_p}\rmint d^3 x\, \partial_{k}\Big\langle\mathcal{A}_{j}(\boldsymbol{x}_2(t_2))\mathcal{A}_{m}(x)\Big\rangle \notag\\
        &\qquad \partial_{0}\Big\langle \mathcal{A}_{l}(\boldsymbol{x}_2(t_3))\mathcal{A}_{i}(x)\Big\rangle \Big\langle \varphi(\boldsymbol{x}_1(t_1))\varphi(\boldsymbol{x}_2(t_2))\Big\rangle\,,\notag\\
    &=-\frac{Q_2^2m_1s_1}{m_p^2}\rmint dt\, v_2^{m}(t)a_2^{i}(t)\Bigg[\frac{1}{16} \pi^{3/2}\, m \,G_{1,3}^{2,1}\left(\frac{m^2 r^2}{4}\Big|
\begin{array}{c}
 -\frac{1}{2} \\
 -\frac{1}{2},-\frac{1}{2},0 \\
\end{array}
\right)\notag\\
&\qquad -\frac{\,\pi^{3/2}\, G_{1,3}^{2,1}\left(\frac{m^2 r^2}{4}\Big|
\begin{array}{c}
 \frac{1}{2} \\
 \frac{1}{2},\frac{1}{2},0 \\
\end{array}
\right)}{8 \,m\, r^2}\Bigg]n^k\,
\frac{\Bar{\psi}(\boldsymbol{x}_1(t))}{m_p}\notag\\
&\quad +\text{term symmetric in $i,m$}+1\leftrightarrow 2\,\sim \mathcal{O}(L^{1/2} v^{15/2}).
\end{align}
 \item \textcolor{black}{A similar type of interaction arises from two different electromagnetic worldline couplings: $\rmint dt Q_1v_1^{i}\mathcal{A}_{i}$ and $-\rmint dt Q_2v_2^{i}\mathcal{A}_{i}\Bar{\psi}$. This is shown in Fig.~(\ref{fig13p}) and it contributes at $N^{(7)}LO.$
Corresponding amplitude is given by (to calculate this amplitude, we need to invoke the one-loop massive scalar integral \cite{Anastasiou:1999ui,2022arXiv220103593W})}: %\vspace{-0.3cm}
\begin{align}
\begin{split}
    \mathcal{S}_{\text{eff}}^{ikm}\Big |_{{\text{fig}(\ref{fig13p})}}&=-\frac{1}{m_p^2}\rmint \Bar{\mathcal{D}}\hat{\xi}\rmint dt_1 Q_1v_1^{j}(t_1)\mathcal{A}_{j}(\boldsymbol{x}_1(t_1))\rmint dt_2 m_1s_1 \varphi(\boldsymbol{x}_1(t_2))\rmint dt_3 Q_2 v_2^{l}(t_3)\mathcal{A}_{l}(\boldsymbol{x}_2(t_3))\\&\,\,\,\,\,\,\,\,\,\,\,\,\,\,\,\,\,\,\,\,\,\,\,\,\,\,\frac{\Bar{\psi}(\boldsymbol{x}_2(t_2))}{m_p}\rmint d^4x \varphi\partial_0 \mathcal{A}_{i}\partial_{k} \mathcal{A}_{m}(x)\,,\\ &
        =-\frac{Q_1Q_2m_1s_1}{m_p^2}\rmint dt_1dt_2dt_3dt\,v_1^{j}(t_1)v_2^{l}(t_3)\frac{\Bar{\psi}(\boldsymbol{x}_2(t_2))}{m_p}\rmint d^3 x \,\partial_{t}\Big\langle\mathcal{A}_{j}(\boldsymbol{x}_1(t_1))\mathcal{A}_{i}(x)\Big\rangle\\ &
\hspace{6cm}\partial_{k}\Big\langle\mathcal{A}_{l}(\boldsymbol{x}_2(t_3))  \mathcal{A}_{m}(x)  \Big\rangle\Big\langle \varphi(x_1(t_2))\varphi(x)\Big\rangle\,,\\ &
    =-\frac{Q_1Q_2m_1s_1}{m_p^2}\rmint dt\, dt_1 v_1^{i}(t_1)v_2^{m}(t)\frac{\Bar{\psi}(\boldsymbol{x}_2(t_2))}{m_p}\partial_{t}\delta(t-t_1)\,\\&\hspace{2cm} \rmint_{\boldsymbol{x}} \rmint_{\boldsymbol{k}_1,\boldsymbol{k}_2,\boldsymbol{k}_3}\frac{-ik_2^{k}}{\boldsymbol{k}_1^2\boldsymbol{k}_2^2(\boldsymbol{k}_3^2+m^2)}e^{i\boldsymbol{k}_1\cdot[\boldsymbol{x}_1(t_1)-\boldsymbol{x}(t)]}e^{i\boldsymbol{k}_2\cdot[\boldsymbol{x}_2(t)-\boldsymbol{x}(t)]} 
       e^{i\boldsymbol{k}_3\cdot[\boldsymbol{x}_1(t)-\boldsymbol{x}(t)]}\,,\\ &
        %=2iQ_1Q_2m_1s_1\rmint dt v_2^{m}(t)a_1^{i}(t)\rmint_{k}\frac{k^k\,e^{i\boldsymbol{k}\cdot\boldsymbol{r}}}{\boldsymbol{k}^2}\textcolor{black}{\rmint_{k_3}\frac{1}{(\boldsymbol{k}-\boldsymbol{k}_3)^2(\boldsymbol{k}_3^2+m^2)}}\frac{\Bar{\psi}}{m_p}\\&
  =\frac{Q_1 Q_2 m_1 s_1}{4\pi\,m_p^2}\rmint dt \, v_2^m (t)a_1^i(t)\,\partial_r\alpha(m;r)n^k \frac{\bar{\psi}}{m_p}(\boldsymbol{x}_2(t))\,\sim \mathcal{O}(L^{1/2}v^{15/2})+(1\leftrightarrow 2)\,.\label{5.85n}
  \end{split}
\end{align}
The details of computations are given in Appendix~(\ref{ch1:app:C}). Also the function  $\alpha(m;r)$ is defined (\ref{C3m}). One can obtain the contribution to the effective action by contracting with $\epsilon^{0ikm}$.
\end{itemize}
%\textcolor{black}{ We can now read off $T^{00}$ as: we collect all the terms with similar delta function structures and combine them into brackets. The $1\longleftrightarrow 2$ means we have to shuffle subscripts to take into consideration couplings to the worldline 2.
%\newpage
%\color{red}
%\begin{align}
   % \begin{split}
 %     \mathcal{T}_{\text{eff}}^{00} =&
%-\frac{\gamma Q_1Q_2'}{2\pi} \,\Big[\boldsymbol{v}_1(t)\cdot \boldsymbol{v}_2(t) %\frac{e^{-\mu r}}{r}\Big](\delta(x-x_1)+\delta(x-x_2)+\delta(\boldsymbol{x}_0-x))\\&\textstyle{-2Q_1 Q_2' %\Big[\boldsymbol{v}_1\cdot \boldsymbol{v}_2\frac{e^{-r \mu _{\gamma }} \left(-r \mu _{\gamma }+e^{r \mu _{\gamma }}-1\right)}{2 \pi  r^3 \mu _{\gamma }^2}
%\Big]\delta(x-x_1)+\delta(x-x_2)+\delta(x-\boldsymbol{x}_0)}\\&+\frac{m_1m_2^2s_2^2g_1}{16\pi^2m_p^4}(1+\frac{3}{2}v_1^2) \frac{e^{-2mr}}{r^2}\delta(x-x_1)\\&
%+\frac{m_1m_2s_1s_2}{m_p} (\frac{1}{m_p}+\frac{3v_1^2}{2m_p}) \frac{1}{4\pi r}e^{-mr}\delta(x-x_1)-2Q_1Q_2 \delta(x-x_1)\frac{1}{4\pi r}-2Q_1Q_2\delta(x-x_1)\frac{1}{4\pi r}e^{-\mu r}\\&
%+\frac{m_1m_2}{4\pi} \delta(x-x_1)\frac{1}{r}-\frac{m_1m_2Q_1Q_2}{16\pi^2m_p^2} \frac{1}{r^2}\delta(x-x_1)-\frac{m_1m_2Q_1'Q_2'}{16\pi^2m_p^2}\frac{e^{-\mu r}}{r^2}\delta(x-x_1)\\&
%-\frac{2Q_1Q_2}{4\pi } (v_1 \cdot v_2)\frac{1}{r}\delta(x-x_1)-\frac{2Q_1'Q_2'}{4\pi } (v_1 \cdot v_2)\frac{e^{-\mu r}}{r}\delta(x-x_1)+\\&
%\textstyle{ -\frac{1}{\pi}\Bigg[v_1^2 %\,\frac{e^{-r \mu _{\gamma }} \left(-r \mu %_{\gamma }+e^{r \mu _{\gamma }}-1\right)}{  %r^3 \mu _{\gamma }^2}}
%\textstyle{
%\Bigg](\delta(x-x_1)+\delta(x-x_2))-%\textstyle{Q_1Q_2'\,\frac{e^{-\mu_{\gamma}r}}%{2\pi r}(\delta(x-x_1)}+\delta(x-x_2)})\\&-%\frac{1}{\pi}\Bigg[v_2^2 \,\frac{e^{-r \mu %_{\gamma }} \left(-r \mu_{\gamma }+e^{r \mu %_{\gamma }}-1\right)}{r^3 \mu _{\gamma }^2}
%\Bigg](\delta(x-x_1)+\delta(x-x_2)
%+1\leftrightarrow 2\\ &
%\end{split}
%\end{align}
%\begin{align}
%\begin{split}
%T^{00}=&\delta(\boldsymbol{x}-\boldsymbol{x}_1)\Bigg\{
%\frac{m_1m_2^2s_2^2g_1}{16\pi^2m_p^4}(1+\frac{3}{2}v_1^2) \frac{e^{-2\textcolor{black}{m}r}}{r^2}
%+\frac{m_1m_2s_1s_2}{m_p} (\frac{1}{m_p}+\frac{3v_1^2}{2m_p}) \frac{1}{4\pi r}e^{-\textcolor{black}{m}r}-\\&Q_1Q_2 \frac{1}{4\pi r}-Q_1'Q_2'\frac{1}{4\pi r}e^{-\mu r}
%+\frac{m_1m_2}{4\pi m_p^2} \frac{1}{r}-\frac{m_1m_2Q_1Q_2}{32\pi^2m_p^2} \frac{1}{r^2}-\frac{m_1m_2Q_1'Q_2'}{32\pi^2m_p^2}\frac{e^{-\mu r}}{r^2}\\&
%-\frac{Q_1Q_2}{4\pi } (v_1 \cdot v_2)\frac{1}{r}-\frac{Q_1'Q_2'}{4\pi } (v_1 \cdot v_2)\frac{e^{-\mu r}}{r}
%\textstyle{ -\frac{\gamma Q_1Q_2'}{2\pi}\Bigg[v_1^2 \,\frac{e^{-r \mu _{\gamma }} \left(-r \mu _{\gamma }+e^{r \mu _{\gamma }}-1\right)}{  r^3 \mu _{\gamma }^2}}
%\Bigg]-\textstyle{\gamma Q_1Q_2'\,\frac{e^{-\mu_{\gamma}r}}{4\pi r}}\\&-\frac{\gamma Q_1Q_2'}{2\pi}\Bigg[v_2^2 \,\frac{e^{-r \mu _{\gamma }} \left(-r \mu _{\gamma }+e^{r \mu _{\gamma }}-1\right)}{  r^3 \mu _{\gamma }^2}
%\Bigg]-\gamma Q_1Q_2'\frac{1}{2\pi} \Bigg[(\Vec{a}_1\cdot \Hat{n})\Big\{\frac{  e^{-\mu_{\gamma} r}}{2 \mu_{\gamma} \,r}-\frac{1 -  e^{-\mu_{\gamma} r}}{2\, \mu_{\gamma}^2\, r^2}\Big\}\Bigg]\\&-\gamma Q_1Q_2'\frac{1}{2\pi} \Bigg[(\Vec{a}_2\cdot \Hat{n})\Big\{\frac{  e^{-\mu_{\gamma} r}}{2 \mu_{\gamma} \,r}-\frac{1 -  e^{-\mu_{\gamma} r}}{2\, \mu_{\gamma}^2\, r^2}\Big\}\Bigg]
%\Bigg\}+\\&\delta(\boldsymbol{x}-\boldsymbol{x}_2)\Bigg\{-\frac{\gamma Q_1Q_2'}{2\pi}\Bigg[v_2^2 \,\frac{e^{-r \mu _{\gamma }} \left(-r \mu _{\gamma }+e^{r \mu _{\gamma }}-1\right)}{  r^3 \mu _{\gamma }^2}
%\Bigg]-\frac{\gamma Q_1Q_2'e^{-\mu_{\gamma}r}}{4\pi r} \\&-\gamma Q_1Q_2'\frac{1}{2\pi} \Bigg[(\Vec{a}_2\cdot \Hat{n})\Big\{\frac{  e^{-\mu_{\gamma} r}}{2 \mu_{\gamma} \,r}-\frac{1 -  e^{-\mu_{\gamma} r}}{2\, \mu_{\gamma}^2\, r^2}\Big\}\Bigg]\\&-\frac{\gamma Q_1Q_2'}{2\pi}\Bigg[v_1^2 \,\frac{e^{-r \mu _{\gamma }} \left(-r \mu _{\gamma }+e^{r \mu _{\gamma }}-1\right)}{  r^3 \mu _{\gamma }^2}\Bigg]-\gamma Q_1Q_2'\frac{1}{2\pi} \Bigg[(\Vec{a}_1\cdot \Hat{n})\Big\{\frac{  e^{-\mu_{\gamma} r}}{2 \mu_{\gamma} \,r}-\frac{1 -  e^{-\mu_{\gamma} r}}{2\, \mu_{\gamma}^2\, r^2}\Big\}\Bigg]+\\&\delta(\boldsymbol{x}-\boldsymbol{x}_{0})\Bigg\{\frac{-\gamma Q_1Q_2'}{\pi} \,\Big[\boldsymbol{v}_1(t)\cdot \boldsymbol{v}_2(t) \frac{e^{-\mu r}}{r}\Big]\textstyle{-4\gamma Q_1 Q_2' \Big[\boldsymbol{v}_1\cdot \boldsymbol{v}_2\frac{e^{-r \mu _{\gamma }} \left(-r \mu _{\gamma }+e^{r \mu _{\gamma }}-1\right)}{2 \pi  r^3 \mu _{\gamma }^2}}
%\Big]+\\&\frac{m_1 m_2 s_1 s_2}{16 \pi m_p^2}\,e^{-mr}\Big(\frac{\boldsymbol{v}_1\cdot\boldsymbol{v}_2}{r}\Big)\Bigg\}\,.
%\end{split}
%\end{align}
\color{black}
Then from this, as discussed in the previous subsections we can read out the source term i.e. $T^{00}$ and from that, we can compute the total power radiation. Before writing that we make the following comment about the contribution coming from axion-photon coupling ($g_
{a\gamma\gamma}$) to the gravitational radiation.  
\vspace{0.2cm}

\underline{{\textit{ The $g_{a\gamma\gamma}$ contribution}}}:\par
The source term coming from the diagrams (\ref{fig13o}) and (\ref{fig13p}), captures the effect of the \textit{theta} term.  %contribution is given by\color{black}
\begin{align}
\begin{split}
T^{00}&\sim
\frac{2Q_1Q_2m_1s_1g_{a\gamma \gamma}}{4\pi\,m_p^2}\, [(\hat n\times \boldsymbol{v}_2)\cdot \Vec{a}_1]
\partial_r\alpha(m;r) \delta(\boldsymbol{x}-\boldsymbol{x}_2)\\&
   + \frac{\pi^{3/2}}{4\,m_p^2}g_{a\gamma\gamma}Q_2^2\,m_1\,s_1\,[\Vec{a_2}\cdot(\hat{n}\times\boldsymbol{v}_2)]\Bigg[\frac{1}{2}\, m \,G_{1,3}^{2,1}\left(\frac{m^2 r^2}{4}\Big|
\begin{array}{c}
 -\frac{1}{2} \\
 -\frac{1}{2},-\frac{1}{2},0 \\
\end{array}
\right)-\frac{ G_{1,3}^{2,1}\left(\frac{m^2 r^2}{4}\Big|
\begin{array}{c}
 \frac{1}{2} \\
 \frac{1}{2},\frac{1}{2},0 \\
\end{array}
\right)}{ \,m\, r^2}\Bigg]\\&\delta(\boldsymbol{x}-\boldsymbol{x}_1)+(1\leftrightarrow 2)\label{5.85}
\end{split}
\end{align}
But the source term in (\ref{5.85}) is zero for any orbit that lies in two dimensions. Hence, the term does not contribute to the gravitational power radiation. This term will only contribute to the power radiation expression for orbits that lie in three dimensions instead.\par

\begin{table}[t!]
 \centering
\scalebox{0.87}{\begin{tabular}{|c|c|c|c|c|c|c|c|c|}
 \hline
  Fields & $LO$ & $N^{(1)}LO$ &  $N^{(2)}LO$ &  $N^{(3)}LO$ &   $N^{(4)}LO$&   $N^{(5)}LO$&   $N^{(6)}LO$ &   $N^{(7)}LO$\\
  \hline
 \textit {scalar} & 1 & 1 & 2 & 1 & 3 & 1 $\sim g_{a\gamma\gamma}$ & - & -\\
  \hline
 \textit{electromagnetic} & 1 & 1 & 1 & 1 & - &-& -& -\\
 \hline 
 \textit{proca} & 1 & 1 & 1 & 1 & - &-& -& -\\
 \hline 
 \textit{gravitational} & 1 & - & 1($\sim \gamma$)+5=6 & - & 4($\sim \gamma$)+7=11 & - & 1 & 2 $\sim g_{a\gamma\gamma}$\\
 \hline
 \end{tabular}}
\caption{Table showing the number of diagrams contributing to each $N^{(n)}LO$ for different radiative fields. We exclude the pure gravitational sector, which is already extensively studied (see \cite{Blanchet:2013haa}), and list only contributions from interactions of other fields with gravity.}
 \label{table2m}
 \end{table}
 
\textit{Before we end this section, we summarize our results in Table~(\ref{table2m}). We show the number of diagrams that contribute to gravitational radiation and arise due to the (minimal) coupling of different fields with the gravitational field. We also highlight the order in which the two couplings constant $g_{a\gamma\gamma}$ and $\gamma$ contribute to the gravitational radiation. Again we emphasize that the pure gravitational sector is already well-studied \cite{Blanchet:2013haa} as mentioned before. Here we mainly focus on the leading order contributions to the gravitational radiation due to the coupling of scalar, em and Proca fields with the gravitational field. To the best of our knowledge, these are some of the new results.} 
Finally, we write down the total gravitational power for \textit{circular orbit} radiation considering the term which arises at $LO$  for pure gravity and terms which arise at their respective leading orders due to the interaction between other fields (scalar, em and Proca) with the gravitational field.
%\begin{align}
   % \begin{split}
   %    \mathcal{ P}_g^{\textrm{LO}}=&\frac{2G}{5\mathcal{T}}\delta(0)\rmint \frac{d\omega}{2\pi}\omega^6\Bigg[\Bigg\{\frac{\mathcal{D}_2m_1^2+\mathcal{D}_1m_2^2}{M^2}\Bigg\}\nu r^4\Omega^2+\Bigg\{\frac{\mathcal{E}_2m_1^2+\mathcal{E}_1m_2^2}{M^2}\Bigg\}(\nu M^2r^2)+\\&\mathcal{L}(r)r^2\frac{(m_2^2+m_1^2)}{M^2}+\nu r^2\Bigg\{G_1(r)m_2^3(1+\frac{3m_2^2r^2\Omega^2}{2M^2})+G_2(r)m_1^3(1+\frac{3m_1^2r^2\Omega^2}{2M^2})\Bigg\}+\\&\nu r^4\Omega^2\Bigg\{A_1(r)(\frac{m_2-m_1}{2M})^2+A_2(r)(\frac{m_2-m_1}{2M})^2\Bigg\}+\Bigg\{(\frac{H_1(r)m_2^2}{ M^2}+\frac{H_2(r)m_1^2}{ M^2})\frac{(m_1^2+m_2^2)}{M^2}(r^4\Omega^2)\Bigg\}+\\&+\Bigg\{(\frac{H_1(r)m_1^2}{ M^2}+\frac{H_2(r)m_2^2}{ M^2})\frac{(m_1^2+m_2^2)}{M^2}(r^4\Omega^2)\Bigg\}+\Bigg\{\frac{\mathcal{P}_2m_1^2+\mathcal{P}_1m_2^2}{M^2}r^2\Bigg\}+\\&\Bigg\{\frac{\mathcal{P}_1m_1^2+\mathcal{P}_2m_2^2}{M^2}r^2\Bigg\}-\Bigg\{\frac{V_1m_2^3}{M^3}+\frac{V_2m_1^3}{M^3}\Bigg\}r^2\Omega^2-\Bigg\{\frac{V_1m_1+V_2m_2}{M}\Bigg\}(\nu r^2 \Omega^2)-\Bigg\{\frac{{V}_1m_1^3}{M^3}+\frac{{V}_2m_2^3}{M^3}\Bigg\}r^2\Omega^2-\\&\Bigg\{\frac{{V}_1m_2+{V}_2m_1}{M}\Bigg\}(\nu r^2 \Omega^2)\Bigg]^2\frac{\pi}{2}\delta(\omega-2\Omega)
   %    \label{5.88}
   % \end{split}
%\end{align}
%\footnote{$\delta(0)$ comes from the following identity of delta function: $|\delta(x-x_0)|^2=\delta(x-x_0)\delta(0)$.}
%\vspace{-0.7cm}
\begin{tcolorbox}[thesisresultbox, title=Gravitational Radiation Power]
\restorethesisbodyformat
\begin{equation}
\begin{aligned}
     P_g=&\frac{32G}{5}\Omega^6\Bigg[\mu r^2+\frac{3}{2}\Bigg\{\frac{m_2^3+m_1^3}{M^2}\Bigg\}\nu r^4\Omega^2+\\&\mathcal{D}\Bigg\{\frac{m_1^2+m_2^2}{M^2}\Bigg\}\nu r^4\Omega^2+\Bigg\{\frac{\mathcal{E}_2m_1^2+\mathcal{E}_1m_2^2}{M^2}\Bigg\}r^2+\mathcal{D}(r)r^2\frac{(m_2^2+m_1^2)}{M^2}+\\&\nu r^2\Bigg\{G_1(r)m_2^3(1+\frac{3m_2^2r^2\Omega^2}{2M^2})+G_2(r)m_1^3(1+\frac{3m_1^2r^2\Omega^2}{2M^2})\Bigg\}+\\&\nu r^4\Omega^2\Bigg\{A_1(r)(\frac{m_2+m_1}{2M})^2+A_2(r)(\frac{m_2+m_1}{2M})^2\Bigg\}+\Bigg\{(\frac{H_1(r)m_2^2}{ M^2}+\frac{H_2(r)m_1^2}{ M^2})\frac{(m_1^2+m_2^2)}{M^2}(r^4\Omega^2)\Bigg\}+\\&\Bigg\{\Big(\frac{H_1(r)m_1^2}{ M^2}+\frac{H_2(r)m_2^2}{ M^2}\Big)\frac{(m_1^2+m_2^2)}{M^2}(r^4\Omega^2)\Bigg\}+\Bigg\{\frac{\mathcal{P}_2m_1^2+\mathcal{P}_1m_2^2}{M^2}r^2\Bigg\}+\\&
       \Bigg\{\frac{\mathcal{P}_1m_1^2+\mathcal{P}_2m_2^2}{M^2}r^2\Bigg\}+\Bigg\{\frac{V_2m_1^3}{M^3}-\frac{V_1m_2^3}{M^3}\Bigg\}r^3\Omega^2-
       \Bigg\{\frac{V_1m_1-V_2m_2}{M}\Bigg\}(\nu r^3 \Omega^2)\\&+\Bigg\{\frac{{V}_1m_1^3}{M^3}-\frac{{V}_2m_2^3}{M^3}\Bigg\}r^3\Omega^2+
       \Bigg\{\frac{{V}_1m_2-{V}_2m_1}{M}\Bigg\}(\nu r^3 \Omega^2)+R(r)\Bigg\{\frac{m_1-m_2}{M}\Bigg\}r^4\Omega^2\Bigg]^2
\end{aligned}
\label{5.81n}
\end{equation}
\end{tcolorbox}

where, 
\begin{align}
\begin{split}
  & \mathcal{D}(r)= \Bigg\{\frac{Q_1Q_2}{4\pi } \frac{1}{r}+\frac{Q_1'Q_2'}{4\pi} \frac{e^{-\mu_{\gamma} r}}{r}\Bigg\}\,,\quad
 %  \mathcal{L}(r)=-2Q_1Q_2 \frac{1}{4\pi r}-2Q_1'Q_2' \frac{1}{4\pi r}e^{-\mu r}\\&
   G_1(r)=\frac{s_2^2g_1}{16\pi^2m_p^4}\frac{e^{-2mr}}{r^2}\,,\\& 
\mathcal{E}_1(r)=\Bigg[\frac{m_1m_2s_1s_2}{m_p}  \Big(\frac{1}{m_p}+\frac{3\frac{m_2^2}{M^2}r^2\Omega^2}{2m_p}\Big) \frac{e^{-mr}}{4\pi r}
+\frac{m_1m_2}{16\pi m_p^2 r}-\frac{m_2 Q_1Q_2}{32\pi^2m_p^2} \frac{1}{r^2}-\frac{m_2 Q_1'Q_2'}{32\pi^2m_p^2} \frac{e^{-\mu_{\gamma} r}}{r^2}\Bigg]\,,\\&
A_1(r)=\Bigg\{\frac{-\gamma Q_1Q_2'}{2\pi}\Big[  \left(\frac{  e^{-\mu_{\gamma} r}}{ r}+\frac{4e^{-r\mu_{\gamma}}\left(-1-r\mu_{\gamma}+e^{\mu_{\gamma} r}\right)}{\mu_{\gamma}^2 r^3}\right) \Big]+\frac{m_1 m_2 s_1 s_2}{16 \pi m_p^2}\,\frac{e^{-m\,r}}{r}\Bigg\}\,,\\&
\mathcal{P}_1(r)=-\frac{\gamma Q_1Q_2'}{4\pi r}e^{-\mu_{\gamma} r}\,,\quad 
H_1(r)=\frac{\gamma Q_1Q_2'}{2\pi}\Bigg[\,\frac{e^{-r \mu _{\gamma }} \left(-r \mu _{\gamma }+e^{r \mu _{\gamma }}-1\right)}{ r^3 \mu _{\gamma }^2}
\Bigg]\,,\\&
V_1(r)=\gamma Q_1Q_2'\frac{1}{2\pi} \Bigg[\Big\{\frac{  e^{-\mu_{\gamma} r}}{2 \mu_{\gamma} \,r}-\frac{1 -  e^{-\mu_{\gamma} r}}{2\, \mu_{\gamma}^2\, r^2}\Big\}\Bigg]\,,\quad 
R(r)=\frac{m_1m_2s_1s_2e^{-mr}}{64\pi m_p^2r}\,.\label{5.81m}
\end{split}
\end{align}
We have used an index subscript for the functions which are not symmetric in two indices (denoting the two worldlines), i.e., one and two, and no indices for the symmetric ones. Also, $\mathcal{E}_2,G_2,A_2,\mathcal{P}_2,H_2,V_2$ can be defined by just replacing the index one by two. There are a few comments that are in order. 

\begin{itemize}
    \item As expected, the total power radiation is a function of the binary separation.
\item In (\ref{5.81m}). Also, we have a term like $\frac{m_1m_2}{16\pi m_p^2 r}$, which doesn't go away even when we set the em charges, dark charges, and screening terms to zero. So it is a pure gravitational radiation correction at $N^{(2)}LO$. We can easily check that it matches the $X^2$ term in  \cite{Huang:2018pbu}. This also serves as a consistency check of our computation. 
\end{itemize}


\section{Discussions and outlooks}\label{disc}
Motivated by recent prospects for searching for dark matter using GW observations \cite{Zhang:2021mks, Coogan:2021uqv, Singh:2022wvw, Becker:2021ivq, Yue:2019ozq, Cardoso:2020iji,Bhattacharya:2023stq}, in this paper we consider a black-hole binary in a dark-matter environment modelled by an \textit{ultra-light particle} (in the form of a massive scalar field with an anomalous coupling to the EM field via a Chern--Simons-type term) and an \textit{ultra-light vector field} (in the form of a massive vector field, or Proca field, that also couples to the EM field through kinetic mixing). Ultimately, we would like to understand whether GW observations can constrain, in addition to the masses of the scalar and Proca fields, the axion-photon coupling parameter $g_{a\gamma\gamma}$ and the Proca-electromagnetic coupling constant $\gamma$, and how those constraints compare with those obtained from other phenomenological considerations mentioned earlier. To that end, we took a first step in this paper by computing the power radiation from a binary system in this dark-matter environment, which is a key ingredient in determining the phase of the gravitational waveform. We use the WEFT formalism to carry out the computation. The study presented in this paper mainly focuses on the regime in which gravity is perturbative. This is approximately true during the adiabatic binary inspiral phase, which is especially important for gravitational-wave detection. A set of length scales is relevant in this regime, ranging from the gravitational radius to the size of the compact objects, their typical orbital separation, and finally the wavelength of the emitted radiation. Keeping these scales in mind and integrating over the different potential fields, we derive the orbital and dissipative dynamics truncated at a given PN order. We then compute the power radiation from a binary moving in a \textit{circular orbit}, which is the first step toward determining the phase of the GW waveform. Below we summarize our results. 
\begin{itemize}
\item We calculated the electromagnetic and Proca bound sectors up to $1PN \sim \mathcal{O}(Lv^2)$. 

\item We also investigated the \textit{em-axion} bound sector, which contributes to the conservative dynamics at $2.5PN\sim \mathcal{O}(Lv^5)$, and the pure scalar bound sector up to $2PN \sim \mathcal{O}(Lv^4)$. \textit{We have provided the details of the 2PN scalar bound sector and of the radiative sector up to $N^{(4)}LO$, which is, to the best of our knowledge, a new result.}

\item \textit{Electromagnetic and Proca radiation has been investigated up to $N^{(3)}LO$. We restricted ourselves to the dipole radiation only. Again this is a new result to the best of our knowledge.} 

\item We also calculated the gravitational radiation due to different field couplings with the gravitational field (results for the pure gravitational sector are well known, and we do not reproduce them here beyond $LO$). \textit{We show that the effect of kinetic term mixing (between photon and dark photon) coupling,  $\gamma$ appears at $N^{(2)}LO (L^{1/2}v^{5/2})$ and $N^{(4)}LO (L^{1/2}v^{9/2}) $ order and
axion-photon, $g_{a\gamma\gamma}$ appears at $\sim N^{(7)}LO(L^{1/2}v^{15/2})$}.  
        
\item \textit{We obtained a novel result showing that the parity-violating axionic term ($g_{a\gamma\gamma}$) leads to detectable power radiation only for genuinely three-dimensional orbits in the case of non-spinning binaries; it vanishes for orbits confined to a two-dimensional plane.} However, it is easy to see that if the two binaries are spinning and their spins are not parallel, this may induce orbital oscillations and hence open up the novel possibility of detecting axions in radiation at infinity.

\end{itemize}

We conclude by mentioning some future directions. First and foremost, we would like to compute the phase of the GW waveform using our results. This will make it possible to estimate various parameters of the theory along the lines of \cite{Zhang:2021mks}, thereby constraining the couplings $g_{a\gamma\gamma}$ and $\gamma$. In turn, this should help constrain this dark-matter model and the parameter space of ultra-light scalar and vector fields. We hope to report on this soon. Furthermore, it will be interesting to generalize our computations to spinning binaries, which would help us build a more complete waveform model. \par 

Another extension of the computation presented in this paper will be to generalize the results for non-Abelian gauge fields and supersymmetric dark matter models. By computing the conservative and radiative dynamics, one will be able to comment on the dependence of our main results on the \textit{color} indices of the gauge fields. One can also introduce finite size effects (tidal deformation etc.)  and investigate the effects of those in our setup. \par 
%As we are interested in different dark matter  candidates, it would be interesting to study different \cite{Jungman:1995df,Ellis:2010kf}. One can give constraint to the different super-symmetric parameters through gravitational wave observations which can be potential verification of `Supersymmetric' theories. 
Last but not least, it will be interesting to investigate the binary dynamics in the Post-Minkowskian (PM) regime. One can use Schwinger-Keldysh `in-in' formalism to study the conservative as well as the dissipative (radiation-reaction) dynamics in a unified framework \cite{Dlapa:2021npj,Kalin:2022hph,Dlapa:2022lmu,Dlapa:2021vgp,Dlapa:2023hsl}. This will enable us to do a full relativistic computation. In recent times, a method known as \textit{Worldline Quantum Field Theory} (WQFT) has been developed to perform such full relativistic computation \cite{Mogull:2020sak,Jakobsen:2022fcj,Jakobsen:2021zvh,Jakobsen:2021lvp,Jakobsen:2021smu,Jakobsen:2022psy} where not only the fields are quantized but also the worldline deflections are quantized. 
%For this purpose we expand the worldline trajectories around the straight lines as, $x_{i}^{\mu}=b_{i}^{\mu}+v_{i}^{\mu}\tau+z_{i}^{\mu}(\tau)$. Inspired by the Eikonal formalism one can write down the WQFT partition function as,
% \begin{align}
  %   \begin{split}
  %       \mathcal{Z}_{\text{WQFT}}=\rmint \mathcal{D}h_{\mu\nu}\,\mathcal{D}\mathcal{A}_{\mu}....\prod_{i=1}^{2}\mathcal{D}z_{i}^{\mu} %\,e^{i[S_{f}+S_{pp}]}
 %    \end{split}
% \end{align}
Using the WQFT method, it would be interesting to compute the full relativistic waveform, i.e., $\langle h_{\mu\nu}\rangle, \langle A_{\mu}\rangle$ etc. and then taking the non-relativistic result one can compare it to our result. It would be a consistency check of the two different formalisms. Another interesting formalism based on scattering amplitude, namely KMOC formalism \cite{Kosower:2018adc}, has been proposed to calculate the gravitational radiation from a binary system. Again in \cite{Mogull:2020sak}, it is shown that the gravitational waveform calculated from KMOC formalism is related to the one-point function in WQFT. Hence it is interesting to find the correspondence between three completely different `new' formalisms using different examples. We hope to report on some of these issues in the near future. 

%\textbf{Future outlooks:} There are many interesting pathway to explore. We  hope to report some of these in near future.
%\begin{itemize}
%\item   This can actually illuminate the aspect of detection of Axion like particles and can shed light on different constraints previously obtained from phenomenological perspectives. Here we don't consider $g_{a\gamma\gamma}$ to be small from the begining and look for possible constraints on them from the pulsar constraints. We are attempting to solve and find out the possible new waveform based on our calculation. It would be interesting to look at the waveforms \cite{Zhang:2021mks} and do fisher metric analysis and calculate several statistical errors and draw inferences via matching those with the existing LIGO data.
 %In WEFT we do not consider any quantization of Wordlines unlike WQFT. 
%\item  Apart from everything our bigger goal is to apply the Effective Field Theory formalism to the alternative and higher derivative theories of gravity which are yet to be explored.
%\end{itemize}


%\newpage
\appendix 
\section{Details of computation used in Section~(\ref{Sec4.2})}\label{ch1:app:A}
In order to evaluate the amplitude in (\ref{Sec4.2}) we need to do two  Feynman loop integrals (\ref{2.15}). Although the integrals are not in the form of Feynman master integrals, one can reduce the integrals by Fourier transforming them into the known master integral form. We have the following two integrals,
\begin{align}
    \begin{split}
       & \mathcal{I}(r)=\rmint_{\boldsymbol{k}_1\Vec{{k_2}}} \frac{k_2^k}{\boldsymbol{k}_1^2\boldsymbol{k}_2^2[(\boldsymbol{k}_1+\boldsymbol{k}_2)^2+\textcolor{black}{m^2}]}e^{i(\boldsymbol{k}_1+\boldsymbol{k}_2).\boldsymbol{r}(t)},\,\Tilde{\mathcal{I}}(r)=\rmint_{\boldsymbol{k}_1\boldsymbol{k}_2}\frac{\,k_2^k\, k_1^a\,e^{i(\boldsymbol{k}_1+\boldsymbol{k}_2).\boldsymbol{r}(t)}}{\boldsymbol{k}_1^2\boldsymbol{k}_2^2[(\boldsymbol{k}_1+\boldsymbol{k}_2)^2+\textcolor{black}{m^2}]}\,.
    \end{split}
\end{align}
We start with the first one,
\begin{align}
    \begin{split}
      &  \mathcal{I}(r)=\rmint_{\boldsymbol{k}_1\Vec{{k_2}}} \frac{k_2^k}{\boldsymbol{k}_1^2\boldsymbol{k}_2^2[(\boldsymbol{k}_1+\boldsymbol{k}_2)^2+\textcolor{black}{m^2}]}e^{i(\boldsymbol{k}_1+\boldsymbol{k}_2).\boldsymbol{r}(t)}\,.
    \end{split}
\end{align}
After doing the Fourier transform we get,
\begin{align}
    \begin{split}
        \hat{I}(\boldsymbol{k})&=\rmint d^3r e^{-i\boldsymbol{k}.\boldsymbol{r}}\,\mathcal{I}(r)\,,\\ &
        =\rmint_{\boldsymbol{k}_1}\frac{k^k-k_1^k}{\boldsymbol{k}_1^2\,(\boldsymbol{k}^2+\textcolor{black}{m^2})\,(\boldsymbol{k}_1-\boldsymbol{k})^2}\,.
    \end{split}
\end{align}
Now do the inverse Fourier transform,
\begin{align}
    \begin{split}
        \mathcal{I}(r)&=\iint_{\boldsymbol{k},\boldsymbol{k}_1}e^{i\boldsymbol{k}.\boldsymbol{r}}\,\Big(\frac{k^k-k_1^k}{\boldsymbol{k}_1^2\,(\boldsymbol{k}^2+\textcolor{black}{m^2})\,(\boldsymbol{k}_1-\boldsymbol{k})^2}\Big)\,,\\ &
    =\iint_{\boldsymbol{k},\boldsymbol{k}_1} e^{i\boldsymbol{k}.\boldsymbol{r}}\frac{k^k}{\boldsymbol{k}_1^2\,(\boldsymbol{k}^2+\textcolor{black}{m^2})\,(\boldsymbol{k}_1-\boldsymbol{k})^2}-\iint_{\boldsymbol{k},\boldsymbol{k}_1} e^{i\boldsymbol{k}.\boldsymbol{r}}\frac{k_1^k}{\boldsymbol{k}_1^2\,(\boldsymbol{k}^2+\textcolor{black}{m^2})\,(\boldsymbol{k}_1-\boldsymbol{k})^2}\,,\\ &
    %=\iint_{\boldsymbol{k},\boldsymbol{k}_1}e^{i\boldsymbol{k}_1.\boldsymbol{r}}\frac{k_1^k}{(\boldsymbol{k}_1^2+\textcolor{black}{m^2})\,\boldsymbol{k}^2\,(\boldsymbol{k}_1-\boldsymbol{k})^2}-\iint_{\boldsymbol{k},\boldsymbol{k}_1}e^{i\boldsymbol{k}_1.\boldsymbol{r}}\frac{k^k}{(\boldsymbol{k}_1^2+\textcolor{black}{m^2})\,\boldsymbol{k}^2\,(\boldsymbol{k}_1-\boldsymbol{k})^2},\,\,(\boldsymbol{k}\longleftrightarrow \boldsymbol{k}_1)\,,\\ &
    %=\rmint_{\Vec{{k_1}}}\frac{k_1^k\,e^{i\boldsymbol{k}_1.\boldsymbol{r}}}{\boldsymbol{k}_1^2+\textcolor{black}{m^2}}\rmint_{\boldsymbol{k}}\frac{1}{\boldsymbol{k}^2(\boldsymbol{k}-\boldsymbol{k}_1)^2}-\rmint_{\Vec{{k_1}}}\frac{e^{i\boldsymbol{k}_1.\boldsymbol{r}}}{\boldsymbol{k}_1^2+\textcolor{black}{m^2}}\rmint_{\boldsymbol{k}}\frac{k^k}{\boldsymbol{k}^2(\boldsymbol{k}-\boldsymbol{k}_1)^2}\,,\\ &
    %=\rmint_{\Vec{{k_1}}}\frac{k_1^k\,e^{i\boldsymbol{k}_1.\boldsymbol{r}}}{\boldsymbol{k}_1^2+\textcolor{black}{m^2}}\Big[\frac{1}{(4\pi)^{3/2}}\Gamma(1/2)^3\frac{1}{|\boldsymbol{k}_1|}\Big]-\rmint_{\Vec{{k_1}}}\frac{e^{i\boldsymbol{k}_1.\boldsymbol{r}}}{\boldsymbol{k}_1^2+\textcolor{black}{m^2}}\Big[\frac{1}{(4\pi)^{3/2
%}}\frac{\Gamma(1/2)\Gamma(3/2)\Gamma(1/2)}{\Gamma(2)}\frac{1}{|\boldsymbol{k}_1|}k_1^k\Big]\,,\\ &
=\frac{1}{8}\rmint_{\boldsymbol{k}_1}\frac{k_1^k \,e^{i\boldsymbol{k}_1.\boldsymbol{r}}}{|\boldsymbol{k}_1|(\boldsymbol{k}_1^2+\textcolor{black}{m^2})}-\frac{1}{16}\rmint_{\boldsymbol{k}_1}\frac{k_1^k \,e^{i\boldsymbol{k}_1.\boldsymbol{r}}}{|\boldsymbol{k}_1|(\boldsymbol{k}_1^2+\textcolor{black}{m^2})}\,,\\ &
=\textstyle{\Bigg[\frac{i}{16} \pi^{3/2}\, m \,G_{1,3}^{2,1}\left(\frac{m^2 r^2}{4}\Big|
\begin{array}{c}
 -\frac{1}{2} \\
 -\frac{1}{2},-\frac{1}{2},0 \\
\end{array}
\right)-\frac{i\,\pi^{3/2}\, G_{1,3}^{2,1}\left(\frac{m^2 r^2}{4}\Big|
\begin{array}{c}
 \frac{1}{2} \\
 \frac{1}{2},\frac{1}{2},0 \\
\end{array}
\right)}{8 \,m\, r^2}\Bigg]n^k}\,.
    \end{split}
\end{align}
If we expand the integral in a small mass limit, we will get,
\begin{align}
    \begin{split}
        -i\mathcal{I}(r)\approx\frac{1}{r}+\frac{1}{3} m^2 r [\log \left(mr\right)+ \gamma_E -\frac{4}{3}]\,.\label{50}
    \end{split}
\end{align}\\
We can see there is a smooth massless limit of the expression: $-i\mathcal{I}(r)\approx\frac{1}{r}$ as expected. Another interesting point to note is that the integral in (\ref{50}) in the small mass limit has the same form as a scalar one-loop integral without the divergent term. That implies that although we are doing some complicated loop integral, the results are convergent. That ensures they are not quantum loop rather we can call them as ``Classical Loops''.\par
%The second intgeral can be evaluated in the same way,
%\begin{align}
    %\begin{split}
        %\mathcal{I}_2(r)&=\iint_{\boldsymbol{k},\boldsymbol{k}_1} e^{i\boldsymbol{k}.\boldsymbol{r}}\frac{k_1^k}{\boldsymbol{k}_1^2\,\boldsymbol{k}^2\,(\boldsymbol{k}_1-\boldsymbol{k})^2}\\ &
    %=\iint_{\boldsymbol{k},\boldsymbol{k}_1}e^{i\boldsymbol{k}_1.\boldsymbol{r}}\frac{k^k}{\boldsymbol{k}_1^2\,\boldsymbol{k}^2\,(\boldsymbol{k}_1-\boldsymbol{k})^2},\,\,(\boldsymbol{k}\longleftrightarrow \boldsymbol{k}_1)\\ &
    %=\rmint_{\Vec{{k_1}}}\frac{e^{i\boldsymbol{k}_1.\boldsymbol{r}}}{\boldsymbol{k}_1^2}\rmint_{\boldsymbol{k}}\frac{k^k}{\boldsymbol{k}^2(\boldsymbol{k}-\boldsymbol{k}_1)^2}-\frac{1}{16}\rmint_{\boldsymbol{k}_1}\frac{k_1^k \,e^{i\boldsymbol{k}_1.\boldsymbol{r}}}{\boldsymbol{k}_1^3}\\ &
    %=\rmint_{\Vec{{k_1}}}\frac{e^{i\boldsymbol{k}_1.\boldsymbol{r}}}{\boldsymbol{k}_1^2}\Big[\frac{1}{(4\pi)^{3/2
%}}\frac{\Gamma(1/2)\Gamma(3/2)\Gamma(1/2)}{\Gamma(2)}\frac{1}{|\boldsymbol{k}_1|}k_1^k\Big]\\ &
%=\frac{1}{16}\rmint_{\boldsymbol{k}_1}\frac{k_1^k \,e^{i\boldsymbol{k}_1.\boldsymbol{r}}}{\boldsymbol{k}_1^3}\\ &
%=\frac{i}{16(4\pi)^{3/2}}\frac{1}{\Gamma(3/2)}\frac{2}{r}\,n^{k}=\frac{i}{32 \pi^2}\frac{1}{r}\,n^k
 %   \end{split}
%\end{align}
%Therefore, 
%\begin{equation}
%\mathcal{I}(r)=\frac{i}{32\pi^2}\frac{1}{r}\,n^k
%\end{equation}%
Now do to the second integral,
\begin{equation}
    \Tilde{\mathcal{I}}(r)=\rmint_{\boldsymbol{k}_1\boldsymbol{k}_2}\frac{\,k_2^k\, k_1^a\,e^{i(\boldsymbol{k}_1+\boldsymbol{k}_2).\boldsymbol{r}(t)}}{\boldsymbol{k}_1^2\boldsymbol{k}_2^2[(\boldsymbol{k}_1+\boldsymbol{k}_2)^2+\textcolor{black}{m^2}]}\,.
\end{equation}
Proceeding as before, we first Fourier transform it.
\begin{align}
    \begin{split}
        \hat{\Tilde{\mathcal{I}}}(\boldsymbol{k})&=\rmint d^3 r\, e^{-i\boldsymbol{k}.\boldsymbol{r}}\, \Tilde{\mathcal{I}}(r)\,,\\ &
        =\rmint_{k_1}\frac{k_1^a(k^k-k_1^k)}{\boldsymbol{k}_1^2\,(\boldsymbol{k}^2+\textcolor{black}{m^2})\,(\boldsymbol{k}_1-\boldsymbol{k})^2}\,.\\ &
    \end{split}
\end{align}
Now do the inverse Fourier transform and get,
\begin{align}
\resizebox{\linewidth}{!}{$
    \begin{aligned}
    \Tilde{\mathcal{I}}(r)&=\iint_{\boldsymbol{k},\boldsymbol{k}_1}\frac{k_1^a\,k^ke^{i\boldsymbol{k}.\boldsymbol{r}}}{\boldsymbol{k}_1^2\,(\boldsymbol{k}^2+\textcolor{black}{m^2})\,(\boldsymbol{k}_1-\boldsymbol{k})^2}-\iint_{\boldsymbol{k},\boldsymbol{k}_1}\frac{k_1^a\,k_1^ke^{i\boldsymbol{k}.\boldsymbol{r}}}{\boldsymbol{k}_1^2\,(\boldsymbol{k}^2+\textcolor{black}{m^2})\,(\boldsymbol{k}_1-\boldsymbol{k})^2}\,,\\
         &=\rmint_{\boldsymbol{k}_1}\frac{k_1^k \,e^{i\boldsymbol{k}_1.\boldsymbol{r}}}{\boldsymbol{k}_1^2+\textcolor{black}{m^2}}\rmint_{\boldsymbol{k}}\frac{k^a}{\boldsymbol{k}^2(\boldsymbol{k}-\boldsymbol{k}_1)^2}-\rmint_{\boldsymbol{k}_1}\frac{e^{i\boldsymbol{k}_1.\boldsymbol{r}}}{\boldsymbol{k}_1^2+\textcolor{black}{m^2}}\rmint_{\boldsymbol{k}}\frac{k^a\,k^k}{\boldsymbol{k}^2(\boldsymbol{k}-\boldsymbol{k}_1)^2}\,,\\
         &=\frac{1}{16}\rmint_{\boldsymbol{k}_1}\frac{k_1^k\,k_1^a\,e^{i\boldsymbol{k}_1.\boldsymbol{r}}}{|\boldsymbol{k}_1|(\boldsymbol{k}_1^2+\textcolor{black}{m^2})}+\frac{1}{16}\rmint_{\boldsymbol{k}_1}\frac{e^{i\boldsymbol{k}_1.\boldsymbol{r}}}{|\boldsymbol{k}_1|(\boldsymbol{k}_1^2+\textcolor{black}{m^2})}\Bigg[\frac{1}{4}\boldsymbol{k}_1^2\delta^{ak}-\frac{3}{4}k_1^a\,k_1^k\Bigg]\,,\\
         &=\frac{1}{64}\rmint_{\boldsymbol{k}_1}\frac{k_1^k\,k_1^a\,e^{i\boldsymbol{k}_1.\boldsymbol{r}}}{|\boldsymbol{k}_1|(\boldsymbol{k}_1^2+\textcolor{black}{m^2})}+\mathcal{F}\,(|\boldsymbol{r}|)\delta^{ak}\,,\\
         &=-\frac{1}{64}\frac{\partial^2}{\partial x^{k}\partial x^{a}}\rmint_{\boldsymbol{k}_1}\frac{e^{i\boldsymbol{k}_1.\boldsymbol{r}}}{|\boldsymbol{k}_1|(\boldsymbol{k}_1^2+\textcolor{black}{m^2})}+\mathcal{F}\,(|\boldsymbol{r}|)\delta^{ak}\,,\\
         &=-\frac{1}{64}\Bigg[ \frac{1}{r}(\delta^{ak}-n^a\,n^k)\Bigg\{\pi ^{3/2} m G_{1,3}^{2,1}\left(\frac{m^2 r^2}{4}|
\begin{array}{c}
 -\frac{1}{2} \\
 -\frac{1}{2},-\frac{1}{2},0 \\
\end{array}
\right)-\frac{2 \pi ^{3/2} G_{1,3}^{2,1}\left(\frac{m^2 r^2}{4}|
\begin{array}{c}
 \frac{1}{2} \\
 \frac{1}{2},\frac{1}{2},0 \\
\end{array}
\right)}{m r^2}\Bigg\}+\\
    &\quad n^a n^k\Bigg\{\frac{\pi ^{3/2} \left(-2 m^2 r^2 G_{1,3}^{2,1}\left(\frac{m^2 r^2}{4}|
\begin{array}{c}
 -\frac{1}{2} \\
 -\frac{1}{2},-\frac{1}{2},0 \\
\end{array}
\right)+8 G_{1,3}^{2,1}\left(\frac{m^2 r^2}{4}|
\begin{array}{c}
 \frac{1}{2} \\
 \frac{1}{2},\frac{1}{2},0 \\
\end{array}
\right)+ m^4 r^4 G_{1,3}^{2,1}\left(\frac{m^2 r^2}{4}|
\begin{array}{c}
 -\frac{3}{2} \\
 -\frac{3}{2},-\frac{3}{2},0 \\
\end{array}
\right)\right)}{2 m r^3}\Bigg\}\\
    &\quad +\mathcal{F}(|\boldsymbol{r}|)\delta^{ak}\Bigg]\,.
    \end{aligned}
$}
\end{align}
%We can write down the amplitude as,
%\begin{align}
 %  \begin{split}
  %        \mathcal{S}_{\text{eff}}^{ikm}&=\textstyle{-i\rmint dt\,v_2^{m}(t)\Big[a_{2}^{i}(t)\rmint_{\boldsymbol{k}_1\Vec{{k_2}}} \frac{k_2^k}{\boldsymbol{k}_1^2\boldsymbol{k}_2^2[(\boldsymbol{k}_1^2+\boldsymbol{k}_2^2)+\textcolor{black}{m^2}]}e^{i(\boldsymbol{k}_1+\boldsymbol{k}_2).\boldsymbol{r}(t)}+v_2^{i}(t) v_2^{a}(t)\rmint_{\boldsymbol{k}_1\boldsymbol{k}_2}\frac{i\,k_2^k\, k_1^a\,e^{i(\boldsymbol{k}_1+\boldsymbol{k}_2).\boldsymbol{r}(t)}}{\boldsymbol{k}_1^2\boldsymbol{k}_2^2[(\boldsymbol{k}_1^2+\boldsymbol{k}_2^2)+\textcolor{black}{m^2}]}\Big]}\\ &
   %     =\textstyle{\frac{\pi^{3/2}}{8}\Bigg[\rmint dt \, v_2^m(t)\,a_2^{i}(t)\Bigg\{\frac{1}{2} \, m \,G_{1,3}^{2,1}\left(\frac{m^2 r^2}{4}\Big|
%\begin{array}{c}
 %-\frac{1}{2} \\
 %-\frac{1}{2},-\frac{1}{2},0 \\
%\end{array}
%\right)-\frac{\, G_{1,3}^{2,1}\left(\frac{m^2 r^2}{4}\Big|
%\begin{array}{c}
% \frac{1}{2} \\
 %\frac{1}{2},\frac{1}{2},0 \\
%\end{array}
%\right)}{m\, r^2}\Bigg\}n^k+}\\ &
%\,\,\,\,\,\,\,\,\,{\frac{1}{8}\,n^a\,n^k\,\rmint dt\,v_2^i(t)v_2^a(t)\Bigg\{-\frac{m G_{1,3}^{2,1}\left(\frac{m^2 r^2}{4}|
%\begin{array}{c}
 %-\frac{1}{2} \\
 %-\frac{1}{2},-\frac{1}{2},0 \\
%\end{array}
%\right)}{r}+\frac{G_{1,3}^{2,1}\left(\frac{m^2 r^2}{4}|
%\begin{array}{c}
 %\frac{1}{2} \\
% \frac{1}{2},\frac{1}{2},0 \\
%\end{array}
%\right)}{m r^3}}\\ &
%\,\,\,\,\,\,\,\,\,{+\frac{1}{2} m^3 r G_{1,3}^{2,1}\left(\frac{m^2 r^2}%{4}|
%\begin{array}{c}
% -\frac{3}{2} \\
 %-\frac{3}{2},-\frac{3}{2},0 \\
%\end{array}
%\right)-\frac{m G_{1,3}^{2,1}\left(\frac{m^2 r^2}{4}|
%\begin{array}{c}
 %-\frac{1}{2} \\
 %-\frac{1}{2},-\frac{1}{2},0 \\
%\end{array}
%\right)}{r}+\frac{2 G_{1,3}^{2,1}\left(\frac{m^2 r^2}{4}|
%\begin{array}{c}
 %\frac{1}{2} \\
 %\frac{1}{2},\frac{1}{2},0 \\
%\end{array}
%\right)}{m r^3}\Bigg\}\Bigg]}
 %  \end{split}
%\end{align}
Finally using these results we can write down the amplitude mentioned in (\ref{4.31}).
\section{ Details of the Feynman integral computation  used in Section~(\ref{Sec4.3})}\label{ch1:app:B}
The integral in (\ref{4.38}) can be evaluated as follows:
\begin{align}
    \begin{split}
         \mathcal{I}_1&= \rmint_{k,k_1}\frac{k_{1i}k_{1j}\,e^{i\boldsymbol{k}\cdot \boldsymbol{r}}}{\boldsymbol{k}^2\,[(\boldsymbol{k}-\boldsymbol{k}_1)^2+m^2](\boldsymbol{k}_{1}^2+m^2)}\,,\\ &
         =\rmint_{k}\frac{e^{i\boldsymbol{k}\cdot \boldsymbol{r}}}{\boldsymbol{k}^2}\underbrace{\rmint_{k_1}\frac{k_{1i}k_{1j}}{(\boldsymbol{k}_1^2+m^2)[(\boldsymbol{k}_1-\boldsymbol{k})^2+m^2]}}_{k_i k_j B_{21}+\delta_{ij}B_{22}}\,.\label{A.1}
    \end{split}
\end{align}
Then the tensor integral in (\ref{A.1}) can be reduced to a scalar integral using Passarino-Veltman reduction \cite{Passarino:1978jh,2022arXiv220103593W} where $B_{21,22}$ is given by,
 \begin{align}
     \begin{split}
          & B_{21}=\frac{1}{2\boldsymbol{k}^2}\Big[\frac{3}{2}\,\boldsymbol{k}^2\,B_{1}+m^2\,B_0-\frac{1}{2}A_0(im)\Big]\,,\\ &
          B_{22}=\frac{1}{4}\Big[-B_1\boldsymbol{k}^2-2m^2B_{0}-A_{0}(im)\Big]\,,
     \end{split}
 \end{align}
 where
 \begin{align}
     \begin{split}
      B_{0}&=\rmint_{k_1}\frac{1}{(\boldsymbol{k}_1^2+m^2)[(\boldsymbol{k}_1-\boldsymbol{k})^2+m^2]}\\ &
     % =\begin{cases}
       %   \frac{1}{8|\boldsymbol{k}|}\Big[1-\frac{2}{\pi}\arctan\Big(\frac{2m}{k}\Big)\Big],& \boldsymbol{k}^2>m^2\\ 
        %  \frac{1}{4\pi|\boldsymbol{k}|}\arctan(|\boldsymbol{k}|/2m), & \boldsymbol{k}^2<m^2
   %   \end{cases}\\ &
      =\frac{1}{4\pi|\boldsymbol{k}|}\arctan(|\boldsymbol{k}|/2m),
    \label{A.3}
     \end{split}
 \end{align}
 and $B_1,A_0(im)$ are given by,
 \begin{align}
     \begin{split}
           &  B_1=\frac{1}{2}B_0=\frac{1}{8\pi|\boldsymbol{k}|}\arctan(|\boldsymbol{k}|/2m)\,,\\ &
       A_0(im)=\frac{m}{4\pi}\,.
     \end{split}
 \end{align}
Therefore the integral in (\ref{A.1}) can be written as,
\begin{align}
    \begin{split}
    \mathcal{I}_{1}&=\rmint_{k}\frac{e^{i\boldsymbol{k}\cdot\boldsymbol{r}}}{\boldsymbol{k}^2}\Big[\frac{3}{4}\frac{1}{8\pi|\boldsymbol{k}|}\arctan(|\boldsymbol{k}|/2m)+\frac{m^2}{8\pi|\boldsymbol{k}|^3}\arctan(|\boldsymbol{k}|/2m)-\frac{m}{16\pi|\boldsymbol{k}|^2}\Big]k_ik_j\\ &
   \,\,\,\, +\rmint_{k}\frac{e^{i\boldsymbol{k}\cdot \boldsymbol{r}}}{\boldsymbol{k}^2}\Big[-\frac{1}{32\pi}|\boldsymbol{k}|\arctan(|\boldsymbol{k}|/2m)-\frac{m^2}{8\pi |\boldsymbol{k}|}\arctan(|\boldsymbol{k}|/2m)-\frac{m}{16\pi}\Big]\delta_{ij}\,,
   \\ & =\frac{3}{32\pi}\Big[\delta_{ij}\frac{1-e^{-2 m r}}{8 \pi  m r^3}+n_in_j\frac{e^{-2m r} \left(2m r-3 e^{2m r}+3\right)}{8 \pi  m r^3}\Big]+\frac{m^2}{8\pi}\Big[\Big\{-\frac{2 m^2 r^2+e^{-2 m r} (2 m r+1)-1}{32 \pi  m^3 r^3}\,n_in_j\Big\}\\ &
  \,\,\,\,\,\, +\Big\{\frac{e^{-2 m r} \left(e^{2 m r} \left(-8 m^3 r^3 \text{Ei}(-2 m r)+6 m^2 r^2-1\right)-4 m^2 r^2+2 m r+1\right)}{96 \pi  m^3 r^3}\delta_{ij}\Big\}\Big]
   \\&-\frac{m}{16\pi}\Big[\frac{1}{8\pi r}\Big(\delta_{ij}-n_i n_j\Big)\Big]
  \,
  +\Big[\frac{e^{-2mr}}{4\pi r^2}-\frac{m}{64 \pi^2 r}\Big(1-e^{-2mr}-2mr \text{Ei}(-2mr)\Big)-\frac{m}{128\pi^2 r}\Big]\delta_{ij}\,,\\ &
   =\lambda_1(r)\delta_{ij}+\lambda_2(r)n_i n_j\,.
    \end{split}\label{B.5m}
\end{align}
The explicit forms of the functions $\lambda_1(r)$ and $\lambda_2(r)$ are given by,
\begin{align}
    \begin{split}
        \lambda_1(r)=&\frac{3-3e^{-2mr}}{256\pi^2m r^3}+\frac{e^{-2 m r} \left(e^{2 m r} \left(-8 m^3 r^3 \text{Ei}(-2 m r)+6 m^2 r^2-1\right)-4 m^2 r^2+2 m r+1\right)}{768\pi^2m r^3}\\ &
        -\frac{m}{128\pi^2 r}+\frac{e^{-2mr}}{4\pi r^2}-\frac{m}{64 \pi^2 r}\Big(1-e^{-2mr}-2mr \text{Ei}(-2mr)\Big)-\frac{m}{128\pi^2 r}\,,
    \end{split}
\end{align}
and,

\begin{align}
    \begin{split}
        \lambda_2(r)=&\frac{e^{-2m r} \left(2m r-3 e^{2m r}+3\right)}{256\pi^2 m r^3}-\frac{2m^2r^2+e^{-2mr}(2mr+1)-1}{256\pi^2m r^3}+\frac{m}{128\pi^2 r}\,.
    \end{split}
\end{align}
The integral in (\ref{B.5m}) can be done as follows,
\begin{align}
    \begin{split}
        \mathcal{I}_{c}(x;r)&=\rmint_{k}\frac{e^{i\boldsymbol{k}\cdot \boldsymbol{r}}}{\boldsymbol{k}^5}\arctan(|\boldsymbol{k}|/x)\,,\\ &
        =\frac{1}{2\pi^2 r x^3}\underbrace{\rmint_{0}^{\infty}dk \, \frac{\sin(\gamma k)}{k^4}\arctan(k)}_{\Hat{\delta}(\gamma)}\,\,,\gamma=x\,r\,.\\ &
    \end{split}\label{B.6m}
\end{align}
The second integral in (\ref{B.5m}) is IR finite. But direct computation is a hard task. So, we first compute the integral in (\ref{B.6m}) and then take darivatives two times w.r.t $x_i.$  But the integral $\Hat{\delta}(\gamma)$ has \textcolor{black}{IR divergence} but one can extract out the finite part by the Feynman trick. For that, we first take double derivative with respect to $\gamma$, and then integrate twice to find out the mother integral with some infinite constant.
\begin{align}
    \begin{split}
         \frac{d^2 \Hat{\delta}(\gamma)}{d\gamma^2}&=-\rmint_{0}^{\infty}dk \, \frac{\sin(\gamma k)}{k^2}\arctan(k)\,,\\ &
         =\frac{1}{2} \pi  \left(\gamma\,  \text{Ei}(-\gamma )+e^{-\gamma }-1\right)\,.
    \end{split}\label{B.7m}
\end{align}
Integrating (\ref{B.7m}) we will have,
\begin{align}
    \begin{split}
\Hat{\delta}(\gamma)&=\rmint^{\gamma}d\gamma'\rmint^{\gamma'}d\gamma''\,\frac{d^2 \Hat{\delta}(\gamma'')}{d\gamma''^2}+C_1 \gamma+C_{2},\,(\text{$C_1, C_2$ are $\gamma$ independent constant})\,,\\ &
        =\frac{\pi}{12}  e^{-\gamma } \Big(\left(1-3 e^{\gamma }\right) \gamma ^2-\gamma +e^{\gamma } \gamma ^3 \text{Ei}(-\gamma )+2\Big)+C_1\,\gamma+C_2.
    \end{split}\label{B.8m}
\end{align}
In order to find the constants, one needs to use the following initial conditions: $\Hat{\delta}(0)=0\,\text{and}\,|\Hat{\delta}'(0)|\rightarrow \infty$. Hence using equation (\ref{B.8m}) one can find that, 
\begin{align}
    \begin{split}
      &  0=\Hat{\delta}(0)=\frac{\pi}{6}+C_2\implies C_2=-\frac{\pi}{6}.\\ &
      \text{and},\,\infty=|\Hat{\delta}'(0)|=|-\frac{\pi}{4}+C_1|\implies |C_1|\rightarrow \infty\,.
    \end{split}
\end{align}
Therefore,
\begin{align}
    \begin{split}
        \mathcal{I}_c=\frac{e^{- x\,r} \Big[(r^2 x^2\, e^{x\,r}\, (x\,r\, \text{Ei}(-x\,r)-3)+r^2 x^2-x\,r+2\Big]}{24 \pi  r x^3}-\frac{1}{12\pi r x^3}+ \text{$`r$' independent infinite constant.}\label{B.9m}
    \end{split}
 \end{align}
Eventually, using the result of (\ref{B.9m}), one can compute the following integral.
\begin{align}
    \begin{split}
     \mathcal{I}_m &= \rmint_{k}\frac{e^{i\boldsymbol{k}\cdot \boldsymbol{r}}}{\boldsymbol{k}^5}\arctan(|\boldsymbol{k}|/x)k_ik_j,\,\,x=2m\\ &
     =-\partial^{x}_{i}\partial^{x}_{j}\,\mathcal{I}_c(2m;r)\,,\\ &
     =\frac{e^{-2 m r} \left(e^{2 m r} \left(-8 m^3 r^3 \text{Ei}(-2 m r)+6 m^2 r^2-1\right)-4 m^2 r^2+2 m r+1\right)}{96 \pi  m^3 r^3}\delta_{ij}\\ &
     \,\,\,\,\,\,\,\,\,\,-\frac{2 m^2 r^2+e^{-2 m r} (2 m r+1)-1}{32 \pi  m^3 r^3}\,n_in_j\,.
    \end{split}
\end{align}
One can obtain the two coefficients by matching the coefficients of $\delta_{ij}$ and $n_i n_j$. Furthermore, we can omit the infinite constant because it does not contribute to the equation of motion obtained from the variation of the effective action.\par 

Another integral relevant for evaluating (\ref{4.38}) is:
\begin{align}
    \begin{split}
         \mathcal{I}_2&= \rmint_{k,k_1}\frac{k_{1i}k_{j}\,e^{i\boldsymbol{k}\cdot \boldsymbol{r}}}{\boldsymbol{k}^2\,[(\boldsymbol{k}-\boldsymbol{k}_1)^2+m^2](\boldsymbol{k}_{1}^2+m^2)}\,,\\ &
         =\rmint_{k}\frac{e^{i\boldsymbol{k}\cdot \boldsymbol{r}}k_j}{\boldsymbol{k}^2}\underbrace{\rmint_{k_1}\frac{k_{1i}}{(\boldsymbol{k}_1^2+m^2)[(\boldsymbol{k}_1-\boldsymbol{k})^2+m^2]}}_{k_{i}B_{1}}\,,\label{A.5}
    \end{split}
\end{align}
where $B_1$ is given in (\ref{A.3}). Therefore the integral in (\ref{A.5}) can be written as,
\begin{align}
    \begin{split} \mathcal{I}_2&=\rmint_{k}\frac{e^{i\boldsymbol{k}\cdot \boldsymbol{r}}k_i k_j}{\boldsymbol{k}^2}\frac{1}{16|\boldsymbol{k}|}\arctan(|\boldsymbol{k}|/2m)\,,\\ &
    =\delta_{ij}\frac{1-e^{-2 m r}}{128 \pi  m r^3}+n_in_j\frac{e^{-2m r} \left(2m r-3 e^{2m r}+3\right)}{128 \pi  m r^3}\,,\\ &
    =\chi_1(r)\delta_{ij}+\chi_2(r)n_i n_j\,.
        \label{ch1:A.6}
    \end{split}
\end{align}
Again the two coefficients $\chi_{1,2}$ can be obtained by matching the coefficients of $\delta_{ij}$ and $n_i n_j$ in (\ref{ch1:A.6}) as,
\begin{align}
    \begin{split}
        &\chi_1(r)=\frac{1-e^{-2mr}}{128\pi m r^3}\,,\\ &
        \text{and,}\\ &
        \chi_2(r)=\frac{e^{-2mr}(2mr-3e^{2mr}+3)}{128\pi m r^3}\,.
    \end{split}
\end{align}
\section{ Details of computation of Feynman integral used in (\ref{Sec4}) and (\ref{Sec5})}\label{ch1:app:C}
In (\ref{4.33mm}) we have the following integral to evaluate:
\begin{align}
    \begin{split}
        I&=\rmint_{k}\frac{k^a e^{i\boldsymbol{k}\cdot \boldsymbol{r}}}{\boldsymbol{k}^2}\underbrace{\rmint_{k_3}\frac{1}{(\boldsymbol{k}-\boldsymbol{k}_3)^2(\boldsymbol{k}_3^2+m^2)}}\\ &
        =\rmint_{k}\frac{k^a e^{i\boldsymbol{k}\cdot \boldsymbol{r}}}{\boldsymbol{k}^2}
        \begin{cases}
        \frac{1}{8|\boldsymbol{k}|}\Big[1-\frac{2}{\pi}\arctan(m/|\boldsymbol{k}|)\Big] & ,k^2> m^2\\
        \frac{1}{4\pi|\boldsymbol{k}|}\arctan(|\boldsymbol{k}|/m) &, k^2<m^2
        \end{cases}\\ &
        =\frac{1}{4\pi}\rmint_{k}\frac{k^a e^{i\boldsymbol{k}\cdot\boldsymbol{r}}}{\boldsymbol{k}^3}\arctan(|\boldsymbol{k}|/m)\,,\\ &
        =-\frac{i}{4\pi}n^{a}\partial_{r}\,\underbrace{\rmint_{k}\frac{e^{i\boldsymbol{k}\cdot\boldsymbol{r}}}{\boldsymbol{k}^3}\arctan(|\boldsymbol{k}|/m)}_{\alpha(m;r)}\,.\label{C.1}
    \end{split}
\end{align}
In (\ref{C.1}) we have used the identity, $\arctan(x)+\arctan(1/x)=\pi/2, \forall x>0$ and $\alpha(m;r)$ has the following form,
\begin{align}
    \begin{split}
        \alpha(m;r)&=\frac{4\pi}{mr}\rmint_{0}^{\infty}dk \frac{\sin(mrk)}{k^2}\arctan(k)\,,\\ &
        =\frac{2 \pi ^2 r \left(-m r \text{Ei}(-m r)-e^{-m r}+1\right)}{m}\,.
    \end{split}
\end{align}
\begin{figure}
    \centering
    \includegraphics[scale=0.5]{int.pdf}
    \caption{$\alpha(m;r)$ vs $mr$ plot}
    \label{fig12}
\end{figure}
 We numerically integrate and show the nature of the integral in Fig~(\ref{fig12}).\par
Another useful relevant integral for (\ref{4.26}) has the following form,
\begin{align}
    \begin{split}
        \beta(m;r)&=\rmint_{k}\frac{e^{i\boldsymbol{k} \cdot \vec r}}{\boldsymbol{k}^2+m^2}\rmint_{k_1}\frac{1}{\boldsymbol{k}_1^2[(\boldsymbol{k}_1-\boldsymbol{k})^2+m^2]}\\ &
        = \rmint_{k}\frac{e^{i\boldsymbol{k} \cdot \vec r}}{\boldsymbol{k}^2+m^2}
        \begin{cases}
                \frac{1}{8|\boldsymbol{k}|}[1-\frac{2}{\pi}\arctan(m/|\boldsymbol{k}|)]\, & \boldsymbol{k}^2>m^2\\
                \frac{1}{4\pi|\boldsymbol{k}|}\arctan(|\boldsymbol{k}|/m)\, & \boldsymbol{k}^2<m^2
            \end{cases}\,,\\ &
 =\rmint_{k}\frac{e^{i\boldsymbol{k}\cdot\boldsymbol{r}}}{(\boldsymbol{k}^2+m^2)|\boldsymbol{k}|}\arctan(|\boldsymbol{k}|/m)\\ &
 =\frac{1}{2\pi^2 mr}\rmint_{0}^{\infty}dk\,\frac{\sin(mrk)}{k^2+1}\arctan(k)\,.\label{C3m}
 \end{split}
\end{align}
The integral in (\ref{C3m}) does not have any closed form expression. One can numerically integrate it and the nature of the function $\beta(m;r)$ is shown in Fig.~(\ref{lastfig}).
\begin{figure}[b!]
    \centering
    \includegraphics[scale=0.4]{beta_vs_mr.pdf}
    \caption{$\beta(m;r)$ vs $mr$ plot}
    \label{lastfig}
\end{figure}


\begin{figure}[t!]
    \centering
    \thesisdiagrampanel{\scalebox{0.36}{\begin{feynman}
    \fermion[label=$$, lineWidth=4]{10.40, 4.80}{11.40, 4.80}
    \fermion[showArrow=false, lineWidth=4]{8.00, 4.40}{7.60, 4.80}
    \fermion[lineWidth=4]{12.60, 4.80}{13.60, 4.80}
    \dashed[showArrow=false, color=eb144c]{4.80, 4.00}{6.00, 4.80}
    \dashed[showArrow=false, color=eb144c]{5.20, 4.00}{6.00, 4.40}
    \dashed[showArrow=false, color=eb144c]{4.80, 4.80}{6.00, 5.60}
    \dashed[showArrow=false, color=eb144c]{4.80, 4.40}{6.00, 5.20}
    \fermion[showArrow=false]{6.00, 4.00}{6.00, 5.60}
    \fermion[showArrow=false, lineWidth=4]{9.20, 5.20}{9.60, 4.80}
    \fermion[lineWidth=4, label=$$]{4.00, 4.00}{6.80, 4.00}
    \fermion[showArrow=false, lineWidth=4]{9.20, 4.40}{9.60, 4.80}
    \fermion[lineWidth=4]{4.00, 5.60}{6.80, 5.60}
    \dashed[showArrow=false, color=eb144c]{4.80, 5.20}{5.60, 5.60}
    \fermion[showArrow=false]{4.80, 4.00}{4.80, 5.60}
    \fermion[showArrow=false, lineWidth=4]{7.60, 4.80}{8.00, 5.20}
    \fermion[lineWidth=5, showArrow=false]{7.60, 4.80}{9.60, 4.80}
    \parton[color=eb144c]{12.00,4.80}{0.60}
\end{feynman}}}
    \caption{Correspondence between PN diagrams (left) and one-loop scalar Feynman diagram (right).}
    \label{fig16}
\end{figure}

%\vspace{-2cm}
\section{Comments on middle vertex contribution to the radiation}\label{ch1:app:D}
In the radiative diagrams, where one radiative field is coming out from the vertex, we implicitly assume that the radiation is coming out from an average position $\boldsymbol{x}_0$. But in general, the radiative field is integrated over all the spacetime i.e, the radiation could be from anywhere in the spacetime. \\

\textbf{Scalar diagrams consisting of radiative bulk vertices:}

1)\,Fig.~(\ref{10h}):

\begin{align}
    \begin{split}
\mathcal{S}_{\text{eff}}\Big |_{\text{fig}(\ref{10h})}&=\frac{m^2 s_1m_1m_2}{m_p^3}\rmint dt_1dt_2\rmint d^4x \Big\langle \psi(\boldsymbol{x}_2(t_2))\psi(x)\Big\rangle \Big\langle\varphi(\boldsymbol{x}_1(t_1))\varphi(x) \Big\rangle \Big\{\bar{\phi}(\vec 0,t)+x^i\partial_{i}\bar\phi(\vec 0,t)+\cdots\Big\}\,\label{E.1}
    \end{split}
\end{align}
The first and second terms in (\ref{E.1}) give the monopole and dipole terms, respectively. First, compute the dipole term. Here, it is difficult to isolate the source term $J$, but one can directly compute the dipole moment by isolating the coefficient of the $\partial_i \bar \phi(\vec 0,t)$. Hence, the effective action has the following, 
\begin{align}
    \begin{split}
        \mathcal{S}_{\text{eff}}\Big |_{\text{fig}(\ref{10h})}&=\frac{m^2 s_1m_1m_2}{2 m_p^3}\rmint dt\,\partial_i\bar\phi(\vec 0,t)\rmint d^3x\,x^i \rmint_{\boldsymbol{k}_1,\boldsymbol{k}_2}\frac{e^{i\boldsymbol{k}_1\cdot (\boldsymbol{x}_2-\boldsymbol{x})}e^{i\boldsymbol{k}_2\cdot (\boldsymbol{x}_1-\boldsymbol{x})}}{\boldsymbol{k}_1^2(\boldsymbol{k}_2^2+m^2)}\,.\label{E.2}
    \end{split}
\end{align}
Now, from \eqref{E.2}, it is clear that the dipole moment has the following form,
\begin{align}
    \begin{split}
        I_{\phi}^{i}\Big|_{\text{fig.} (\ref{10h})}&=\frac{m^2 s_1m_1m_2}{2 m_p^3}\rmint d^3x\,x^i \rmint_{\boldsymbol{k}_1,\boldsymbol{k}_2}\frac{e^{i\boldsymbol{k}_1\cdot (\boldsymbol{x}_2-\boldsymbol{x})}e^{i\boldsymbol{k}_2\cdot (\boldsymbol{x}_1-\boldsymbol{x})}}{\boldsymbol{k}_1^2(\boldsymbol{k}_2^2+m^2)}\,,\\ &
        =\frac{m^2 s_1m_1m_2}{2 m_p^3}\Big[r^i\frac{(m^2r^2-2)+2e^{-mr}(mr)}{4\pi m^4 r^3}+x_2^i\frac{1-e^{-mr}}{4\pi m^2 r}\big]\,.
    \end{split}
\end{align}
2) Fig~(\ref{10e}): Here, the computation is more involved as we have a time derivative over the radiation field.
\begin{align}
    \begin{split}
        \mathcal{S}_{\text{eff}}\Big |_{\text{fig}(\ref{10e})}&=\frac{m_1m_2s_2}{m_p^2}\rmint\Bar{\mathcal{D}}\Hat{\xi}\rmint dt_1dt_2\tilde{\psi}(\boldsymbol{x}_1(t_1))\phi(\boldsymbol{x}_2(t_2))\rmint d^4x \,\frac{\psi}{m_p}\partial_0\phi\partial_0\bar{\phi}\\ &
        =-\frac{m_1m_2s_2}{m_p^2}\rmint \Bar{\mathcal{D}}\Hat{\xi}\rmint dt_1dt_2\tilde{\psi}(\boldsymbol{x}_1(t_1))\phi(\boldsymbol{x}_2(t_2))\rmint d^4x \,\Big[\partial_0\tilde\psi\partial_0 \varphi+\tilde\psi\partial_0^2 \varphi\Big]\Big\{\bar\phi(\vec 0,t)+x^i \partial_i \bar\phi(\vec 0,t)\Big\}
    \end{split}
\end{align}
Again, focus on the dipole term only,
\begin{align}
    \begin{split}
    \mathcal{S}_{\text{eff}}\Big |_{\text{fig}(\ref{10e})}&=\frac{m_1m_2s_2}{m_p^2}\rmint dt_1 dt_2 \rmint d^4x \Big[\partial_0\Big\langle\tilde\psi(\boldsymbol{x}_1(t_1))\tilde\psi(x)\Big\rangle \partial_0\Big\langle\varphi(\boldsymbol{x}_2(t_2))\varphi(x)\Big\rangle\,,\\ &
    \hspace{3 cm}+\Big\langle\tilde\psi(\boldsymbol{x}_1(t_1))\tilde\psi(x)\Big\rangle \partial_0^2\Big\langle\varphi(\boldsymbol{x}_2(t_2))\varphi(x)\Big\rangle\Big]x^{i}\partial_i\bar\phi(\vec 0,t)\,.
    \end{split}
\end{align}
Considering the first term only,
\begin{align}
    \begin{split}
         \mathcal{S}_{\text{eff}}\Big |_{\text{fig}(\ref{10e})}^{(1)}&=\frac{m_1m_2s_2}{m_p^2}\rmint dt_1 dt_2 \rmint d^4x \partial_0\Big\langle\tilde\psi(\boldsymbol{x}_1(t_1))\tilde\psi(x)\Big\rangle \partial_0\Big\langle\varphi(\boldsymbol{x}_2(t_2))\varphi(x)\Big\rangle x^{i}\partial_i \bar\phi(\vec 0,t)\,,\\ &
         =\frac{m_1m_2s_2}{2m_p^2}\rmint dt \,\partial_{i}\bar\phi(\vec 0,t) \rmint d^3x \rmint_{k_1}\frac{(i\boldsymbol{k}_1\cdot \boldsymbol{v}_1) }{\boldsymbol{k}_1^2}\Big[x_1^{i}+i\partial_{k_1^i}\Big]e^{i\boldsymbol{k}_1\cdot (\boldsymbol{x}_1-\boldsymbol{x})}\rmint_{k_2}(i\boldsymbol{k}_2\cdot \boldsymbol{v}_2)\frac{e^{ik_2\cdot (\boldsymbol{x}_2-\boldsymbol{x})}}{\boldsymbol{k}_2^2+m^2}\,.
    \end{split}
\end{align}
Therefore, the dipole moment can be isolated as,
\begin{align}
    \begin{split}
        I_{\phi}^{i}(t)\Big|_{\text{fig.}(\ref{10e})}&=-v_1^a v_2^b\rmint d^3x \rmint_{k_1}\Big[\frac{k_1^a\,x_1^i}{\boldsymbol{k}_1^2} -i\Big(\frac{\delta_{ai}}{\boldsymbol{k}_1^2}-\frac{2k_1^ik_1^a}{\boldsymbol{k}_1^4}\Big)\Big] e^{i\boldsymbol{k}_1\cdot (\boldsymbol{x}_1-\boldsymbol{x})}\rmint_{k_2}k_2^b\,\frac{e^{i\boldsymbol{k}_2\cdot(\boldsymbol{x}_2-\boldsymbol{x})}}{\boldsymbol{k}_2^2+m^2}\\ &
        =v_1^a v_2^b\rmint_{k_1}\Big[\frac{k_1^a\,x_1^i}{\boldsymbol{k}_1^2} -i\Big(\frac{\delta_{ai}}{\boldsymbol{k}_1^2}-\frac{2k_1^ik_1^a}{\boldsymbol{k}_1^4}\Big)\Big]\frac{k_1^b}{\boldsymbol{k}_1^2+m^2}e^{i\boldsymbol{k}_1\cdot \vec r}\\ &
        =v_1^a v_2^b\Bigg[x_1^i\rmint_{\boldsymbol{k}_1}\frac{k_1^a k_1^b}{\boldsymbol{k}_1^2(\boldsymbol{k}_1^2+m^2)}-i\delta_{ai}\rmint_{k_1}\frac{k_1^b}{\boldsymbol{k}_1^2(\boldsymbol{k}_1^2+m^2)}+2i\rmint_{k_1}\frac{k_1^i k_1^a k_1^b}{\boldsymbol{k}_1^4(\boldsymbol{k}_1^2+m^2)}\Bigg]e^{i\boldsymbol{k}\cdot \vec r}
    \end{split}
\end{align}
The integrals can be done by investigating two separate families of Feynman-Fourier integrals that we discuss in the subsequent subsection. However, we are not showing the explicit results of these integrals, as they can be straightforwardly evaluated using the results we provide in the subsequent section.\\\\
\textbf{IBP with exponential and systematic computation of the Fourier integrals.}\\ 
Let's take a moment to systematically compute the Fourier Feynman Integrals.\\
\textbf{\texttt{One-loop Fourier family}:}\,Let us define a family of integrals,
\begin{align}
    \begin{split}
        {J}_{a_1,a_2}= \int d^dk\, \frac{e^{i\boldsymbol{k}\cdot \boldsymbol r}}{(i \boldsymbol{k}\cdot \boldsymbol r)^{a_1}(\boldsymbol{k}^2+m^2)^{a_2}}.
    \end{split}
\end{align}
After solving the IBP relations (with a simple modified version of \textbf{\texttt{LiteRed}}\cite{Lee:2013mka}) we have three master integrals for the first family: ${J}^{}_{0,1},{J}^{}_{0,2}, {J}^{}_{1,1}$. First do the following rescaling, $\boldsymbol{k}\to \boldsymbol{\ell}=\,\boldsymbol{k}/m$ and $\boldsymbol{r}\to \boldsymbol{s}=\boldsymbol{r}m$. We left with the integral,
\begin{align}
    J_{a_1,a_2}=(m^2)^{d/2-a_2}\int d^d\boldsymbol{\ell} \frac{e^{i\boldsymbol{\ell}\cdot\boldsymbol{s}}}{(i\boldsymbol{\ell}\cdot \boldsymbol{s}-i\varepsilon)^{a_1}(\boldsymbol{\ell}^2+1)^{a_2}}
\end{align}
The differential equation satisfied by the master integral is given by (with $x=s^2$),
\begin{align}
    \partial_{x}\boldsymbol{\mathbfcal{J}}=\boldsymbol{\Omega}_{x}\cdot \boldsymbol{\mathbfcal{J}},\qquad   \boldsymbol{\Omega}_{x}=\left(
\begin{array}{ccc}
 -\frac{d-2}{2 x} & -\frac{1}{x} & 0 \\
 -\frac{1}{4} & 0 & 0 \\
 \frac{1}{2 x} & 0 & -\frac{1}{2 x} \\
\end{array}
\right)
\end{align}
Therefore, we have three coupled differential equations to be solved,
\begin{align}
    \begin{split}&\partial_{x}J_{0,1}=\frac{2-d}{2x}J_{0,1}-\frac{1}{x}J_{0,2},\\&
    \partial_{x}J_{0,2}=-\frac{1}{4}J_{0,1},\\ &
    \partial_{x}J_{1,1}=\frac{1}{2x}J_{0,1}-\frac{1}{2x}J_{1,1}.
    \end{split}
\end{align}
The solution is given by,
\begin{align}
\begin{split}
&    J_{0,1}=x^{\frac{2-d}{4}}\left(C_1 I_{\frac{d-2}{2}}(\sqrt{x})+C_2 K_{\frac{d-2}{2}}(\sqrt{x})\right),\\ &
J_{0,2}
=
\frac{1}{2}\,x^{\frac{4-d}{4}}
\left(
-\,C_1 I_{\frac{d-4}{2}}(\sqrt{x})
+
C_2 K_{\frac{d-4}{2}}(\sqrt{x})
\right),\\ &
J_{1,1}(x)
=
\frac{1}{\sqrt{x}}
\left[
C_3
+
\int^{\sqrt{x}}
t^{-\frac{d-2}{2}}
\left(
C_1 I_{\frac{d-2}{2}}(t)
+
C_2 K_{\frac{d-2}{2}}(t)
\right)\,dt
\right].\label{2.201j}
    \end{split}
\end{align}
The constants $C_1, C_2$ can be fixed using the asymptotic expansion around $x\to 0$ of the integral. In this asymptotic, the integral has two distinct momentum scalings,
\begin{align}
   \texttt{hard:} |\boldsymbol{\ell}|\sim\frac{1}{s},\,\,  \texttt{soft:} |\boldsymbol{\ell}|\sim1
\end{align}
There are therefore two regions: from the method of regions \cite{Beneke:1997zp},\\
\textbf{\texttt{hard region}.} In this region the propagator expanded as massless, $\frac{1}{\boldsymbol{\ell}^2+1}\sim \frac{1}{\boldsymbol{\ell}^2}$ and hence,
\begin{align}
    J_{0,1}^{(h)}(x)=2^{d-2}\pi^{d/2} \Gamma\left(\frac{d}{2}-1\right)x^{1-\frac{d}{2}}
\end{align}
\textbf{\texttt{soft region}.} In this region, the exponent has a soft expansion, and considering the leading order, one has,
\begin{align}
    J_{0,1}^{(s)}(x)=\int d^d\boldsymbol{\ell} \frac{1}{\boldsymbol{\ell}^2+1}=\pi^{d/2}\Gamma\left(1-\frac{d}{2}\right)
\end{align}
Collecting all the leading contributions from the asymptotic region,
\begin{align}
    J_{0,1}\overset{x\to 0}{\sim} 2^{d-2}\pi^{d/2}\Gamma\left(\frac{d}{2}-1\right) x^{d/2-1}+\pi^{d/2}\Gamma\left(1-\frac{d}{2}\right)+\cdots\label{2.205a}
\end{align}
Now, expanding $J_{0,1}$ in \eqref{2.201j} we get,
\begin{align}
   J_{0,1}\overset{x\to 0}{\sim}  \frac{2^{-\frac{d}{2}} \left(2 C_1+C_2 \Gamma \left(\frac{d}{2}\right) \Gamma \left(1-\frac{d}{2}\right)\right)}{\Gamma \left(\frac{d}{2}\right)}+C_2 2^{\frac{d}{2}-2} x^{1-\frac{d}{2}} \Gamma \left(\frac{d-2}{2}\right)\label{2.206a}
\end{align}
Now, matching \eqref{2.205a} and \eqref{2.206a} we get,
\begin{align}
    C_1=0,\,C_2=(2\pi)^{d/2}
\end{align}
In a similar way, $C_3$ can be fixed from the asymptotic expansion of  $J_{1,1}$, in soft region,
\begin{align}
\begin{split}
    J_{1,1}^{(s)}&\overset{x\to0}{\sim} \int d^d\ell \frac{1-i\boldsymbol{\ell}\cdot \boldsymbol{s}}{(i\boldsymbol{\ell\cdot s-i\varepsilon})(\boldsymbol{\ell}^2+1)}+\cdots\\ &
    \sim \int d^d\ell \frac{1}{(i\boldsymbol{\ell\cdot s-i\varepsilon})(\boldsymbol{\ell}^2+1)}+\cdots=\frac{\pi^{\frac{d+1}{2}}\Gamma\left(\frac{3-d}{2}\right)}{\sqrt{x}\,}+\cdots
    \end{split}
\end{align}
Comparing we have, $C_3={\pi^{\frac{d+1}{2}}}\Gamma\left(\frac{3-d}{2}\right)$
Therefore, the integrals become,
\begin{align}
\begin{split}
 &   J_{0,1}=(2\pi)^{d/2} \left(\frac{m}{r}\right)^{d/2-1}K_{\frac{d-2}{2}}(mr),\\ &
    J_{0,2}=\frac{1}{2}(m^2)^{d/2}\left(mr\right)^{\frac{4-d}{4}} K_{\frac{d-4}{2}}(mr)\\ &
    J_{1,1}=\frac{1}{\sqrt{m r}}\Bigg[2^{\frac{d}{2}-3} \pi ^d \csc \left(\frac{\pi  d}{2}\right) \sqrt[4]{m r} \Big(4 \pi ^{3/2} \, _1\tilde{F}_2\left(\frac{1}{2};\frac{d}{2},\frac{3}{2};\frac{\sqrt{m r}}{4}\right)\\ &-\pi  2^d \Gamma \left(\frac{3}{2}-\frac{d}{2}\right) (m r)^{\frac{2-d}{4}} \, _1\tilde{F}_2\left(\frac{3}{2}-\frac{d}{2};2-\frac{d}{2},\frac{5}{2}-\frac{d}{2};\frac{\sqrt{m r}}{4}\right)\Big)+\pi ^{\frac{d+1}{2}} \Gamma \left(\frac{3}{2}-\frac{d}{2}\right)\Bigg]
    \end{split}
\end{align}
\paragraph{Two-loop Fourier family.}
We consider the two-loop family
\begin{align}
G_{a_1,a_2,a_3,a_4,a_5}(\boldsymbol{r})
=
\int d^d k\, d^d k_1\,
\frac{e^{D_1}}{D_1^{a_1} D_2^{a_2} D_3^{a_3} D_4^{a_4} D_5^{a_5}},
\end{align}
where
\begin{align}
D_1 = i\,\boldsymbol{k}\!\cdot\!\boldsymbol{r},
\qquad
D_2 = i\,\boldsymbol{k}_1\!\cdot\!\boldsymbol{r},
\qquad
D_3 = \boldsymbol{k}^2,
\qquad
D_4 = \boldsymbol{k}_1^2 + m^2,
\qquad
D_5 = (\boldsymbol{k}_1-\boldsymbol{k})^2 + m^2.
\end{align}
After solving the IBP relations, the family reduces to $17$ master integrals. For our purposes, it is sufficient to consider the subset
\begin{align}
\mathcal G(x)
=
\bigl(
G_1(x),G_2(x),G_3(x),G_4(x),G_5(x),G_6(x)
\bigr)^T,
\end{align}
with
\begin{align}
G_1 &= G_{0,0,0,1,1},
\qquad
G_2 = G_{0,0,0,2,1},
\qquad
G_3 = G_{0,0,0,1,2},
\\
G_4 &= G_{0,0,-1,1,1},
\qquad
G_5 = G_{0,0,1,0,1},
\qquad
G_6 = G_{0,0,1,1,0}.
\end{align}
and we define the dimensionless variable
\begin{align}
x = (m r)^2,
\qquad r = |\boldsymbol{r}|.
\end{align}
The differential system satisfied by these master integrals is
\begin{align}
\partial_x \mathbfcal G(x) = \Delta_x\, \mathbfcal G(x),\quad \Delta_x =
\begin{pmatrix}
\dfrac{2-d}{x} & -\dfrac{1}{x} & -\dfrac{1}{x} & 0 & 0 & 0 \\[6pt]
\dfrac{(d-2)^2}{4x} & 0 & \dfrac{d-2}{2x} & \dfrac{1}{8} & 0 & 0 \\[6pt]
\dfrac{(d-2)^2}{4x} & \dfrac{d-2}{2x} & 0 & \dfrac{1}{8} & 0 & 0 \\[6pt]
\dfrac{2(d-3)}{x} & \dfrac{4}{x} & \dfrac{4}{x} & \dfrac{1-d}{x} & 0 & 0 \\[6pt]
0 & 0 & 0 & 0 & -\dfrac{d-2}{2x} & 0 \\[6pt]
0 & 0 & 0 & 0 & 0 & -\dfrac{d-2}{2x}
\end{pmatrix}.
\end{align}
By symmetry, one has
\begin{align}
G_2(x)=G_3(x),
\qquad
G_5(x)=G_6(x),
\end{align}
and therefore the system reduces to four coupled equations:
\begin{align}
\begin{split}
G_1'(x) &= \frac{2-d}{x}\,G_1(x) - \frac{2}{x}\,G_2(x), \\
G_2'(x) &= \frac{(d-2)^2}{4x}\,G_1(x) + \frac{d-2}{2x}\,G_2(x) + \frac{1}{8}\,G_4(x), \\
G_4'(x) &= \frac{2(d-3)}{x}\,G_1(x) + \frac{8}{x}\,G_2(x) + \frac{1-d}{x}\,G_4(x), \\
G_5'(x) &= \frac{2-d}{2x}\,G_5(x).
\end{split}
\end{align}
Eliminating $G_2$ and $G_4$ from the first three equations, one finds that $G_1$ satisfies the third-order differential equation
\begin{align}
2x^2 G_1'''(x) + 3d\,x\,G_1''(x) + \bigl(d^2-d-2x\bigr) G_1'(x) + (1-d) G_1(x) = 0.
\end{align}
Its general solution can be written as
\begin{align}
G_1(x) = c_1 u_1(x) + c_2 u_2(x) + c_3 u_3(x),
\end{align}
where
\begin{align}
\begin{split}
u_1(x) &= {}_1F_2\!\left(\frac{d-1}{2};\, d-1,\frac{d}{2};\, x\right), \\
u_2(x) &= x^{1-\frac{d}{2}}\,
{}_1F_2\!\left(\frac{1}{2};\, 2-\frac{d}{2},\frac{d}{2};\, x\right), \\
u_3(x) &= x^{2-d}\,
{}_1F_2\!\left(\frac{3-d}{2};\, 3-d,2-\frac{d}{2};\, x\right).
\end{split}
\end{align}
Consequently,
\begin{align}
\begin{split}
G_2(x) &= (2-d)\,G_1(x) - x\,G_1'(x), \\
G_4(x) &= 8\,G_2'(x) - \frac{2(d-2)^2}{x}\,G_1(x) - \frac{4(d-2)}{x}\,G_2(x), \\
G_5(x) &= c_4\,x^{1-\frac{d}{2}}.
\end{split}
\end{align}
Although $G_1$ can be identified with the square of the one-loop master integral discussed previously, let us assume that this relation is not known a priori and determine the constants $c_1,c_2,c_3$ directly from boundary data. To this end, we introduce the dimensionless variables
\begin{align}
\boldsymbol{s} = m \boldsymbol{r},
\qquad
x = \boldsymbol{s}^2,
\end{align}
and rescale the loop momenta according to
\begin{align}
\boldsymbol{q} = \frac{\boldsymbol{k}_1}{m},
\qquad
\boldsymbol{p} = \frac{\boldsymbol{k}_1-\boldsymbol{k}}{m}.
\end{align}
Then the dimensionless master integral can be written as
\begin{align}
G_1(x)
=
m^{4-2d} G_{0,0,0,1,1}(\boldsymbol{r})
=
\int d^d q\, d^d p\,
\frac{e^{i(\boldsymbol{q}-\boldsymbol{p})\cdot \boldsymbol{s}}}{(\boldsymbol{q}^2+1)(\boldsymbol{p}^2+1)}.
\end{align}
To determine the small-$x$ behavior of $G_1$, we use the method of regions in the limit $x \to 0$, or equivalently $s \to 0$. The relevant momentum scalings are
\begin{align}
\texttt{hard:}\quad |\boldsymbol{q}|,|\boldsymbol{p}| \sim \frac{1}{s},
\qquad
\texttt{soft:}\quad |\boldsymbol{q}|,|\boldsymbol{p}| \sim 1.
\end{align}
There are therefore four asymptotic regions.

\paragraph{\texttt{hard-hard} region.}
In this region, both propagators may be expanded as massless,
\begin{align}
\frac{1}{\boldsymbol{q}^2+1} \sim \frac{1}{\boldsymbol{q}^2},
\qquad
\frac{1}{\boldsymbol{p}^2+1} \sim \frac{1}{\boldsymbol{p}^2},
\end{align}
and hence
\begin{align}
G_1^{(hh)}(x)
&=
\int d^d q\,\frac{e^{i\boldsymbol{q}\cdot \boldsymbol{s}}}{\boldsymbol{q}^2}
\int d^d p\,\frac{e^{-i\boldsymbol{p}\cdot \boldsymbol{s}}}{\boldsymbol{p}^2}
\nonumber\\
&=
2^{2d-4}\pi^d
\Gamma\!\left(\frac{d}{2}-1\right)^2
x^{2-d}.
\end{align}

\paragraph{\texttt{hard-soft} and \texttt{soft-hard} regions.}
Consider, for instance, the region $|\boldsymbol{q}| \sim s^{-1}$ and $|\boldsymbol{p}| \sim 1$. The soft exponential is expanded as
\begin{align}
e^{-i\boldsymbol{p}\cdot \boldsymbol{s}}
=
1 - i\,\boldsymbol{p}\cdot \boldsymbol{s} - \frac{1}{2}(\boldsymbol{p}\cdot \boldsymbol{s})^2 + \cdots.
\end{align}
The odd terms vanish upon angular integration, and at leading order, we obtain
\begin{align}
G_1^{(hs)}(x)
&=
\int d^d q\,\frac{e^{i\boldsymbol{q}\cdot \boldsymbol{s}}}{\boldsymbol{q}^2}
\int d^d p\,\frac{1}{\boldsymbol{p}^2+1}.
\end{align}
Including the symmetric \texttt{soft-hard} contribution, one finds
\begin{align}
G_1^{(hs)}(x) + G_1^{(sh)}(x)
=
2^{d-1}\pi^d
\Gamma\!\left(\frac{d}{2}-1\right)
\Gamma\!\left(1-\frac{d}{2}\right)
x^{1-\frac{d}{2}}.
\end{align}

\paragraph{\texttt{soft-soft} region.}
In this case both loop momenta are of order unity, and the exponential can be expanded as
\begin{align}
e^{i(\boldsymbol{q}-\boldsymbol{p})\cdot \boldsymbol{s}} = 1 + \mathcal O(s).
\end{align}
The leading term is therefore
\begin{align}
G_1^{(ss)}(x)
&=
\left(
\int d^d p\,\frac{1}{\boldsymbol{p}^2+1}
\right)^2
=
\pi^d \Gamma\!\left(1-\frac{d}{2}\right)^2.
\end{align}
Collecting the leading contributions from all regions, we obtain the asymptotic expansion
\begin{align}
G_1(x)
\sim
2^{2d-4}\pi^d
\Gamma\!\left(\frac{d}{2}-1\right)^2
x^{2-d}
+
2^{d-1}\pi^d
\Gamma\!\left(\frac{d}{2}-1\right)
\Gamma\!\left(1-\frac{d}{2}\right)
x^{1-\frac{d}{2}}
+
\pi^d
\Gamma\!\left(1-\frac{d}{2}\right)^2
+\cdots.
\end{align}
Matching this expansion against the Frobenius basis $\{u_1,u_2,u_3\}$ immediately yields
\begin{align}
\begin{split}
c_1 &= \pi^d \Gamma\!\left(1-\frac{d}{2}\right)^2, \\
c_2 &= 2^{d-1}\pi^d
\Gamma\!\left(\frac{d}{2}-1\right)
\Gamma\!\left(1-\frac{d}{2}\right), \\
c_3 &= 2^{2d-4}\pi^d
\Gamma\!\left(\frac{d}{2}-1\right)^2.
\end{split}
\end{align}
We consider another the two-loop family that arises in our prebious computations,
\begin{align}
Q_{a_1,a_2,a_3,a_4,a_5}(\boldsymbol{r})
=
\int d^d k\, d^d k_1\,
\frac{e^{D_1}}{D_1^{a_1} D_2^{a_2} D_3^{a_3} D_4^{a_4} D_5^{a_5}},
\end{align}
where
\begin{align}
D_1 = i\,\boldsymbol{k}\!\cdot\!\boldsymbol{r},
\qquad
D_2 = i\,\boldsymbol{k}_1\!\cdot\!\boldsymbol{r},
\qquad
D_3 = \boldsymbol{k}^2+m^2,
\qquad
D_4 = \boldsymbol{k}_1^2 ,
\qquad
D_5 = (\boldsymbol{k}_1-\boldsymbol{k})^2.
\end{align}
Solving IBP identities, we were left with 10 master integrals. However, for our purpose, we need three of them,
\begin{align}
    \mathbfcal{Q}:Q_{0,0,0,1,1},\,Q_{0,0,1,1,1},\, Q_{0,0,2,1,1}.
\end{align}
The differential equation satisfied by them is given by,
\begin{align}
    \partial_{x} \mathbfcal{Q}(x)=\mathbf\Lambda_x \mathbfcal{Q}(x), \qquad \mathbf\Lambda_x =\left(
\begin{array}{ccc}
 \frac{2-d}{x} & 0 & 0 \\
 0 & \frac{3-d}{x} & -\frac{1}{x} \\
 \frac{1}{4} & \frac{2 (d-7) d-x+24}{4 x} & \frac{d-4}{2 x} \\
\end{array}
\right)
\end{align}
Therefore, the system reduces to three coupled differential equations,
\begin{align}
    \begin{split}
        &Q_1'(x)=\frac{2-d}{x}Q_1,\\ &
        Q_2'(x)=\frac{3-d}{x}Q_1-\frac{1}{x}Q_2,\\ &
        Q_3'(x)=\frac{1}{4}Q_1(x)+\frac{2 (d-7) d-x+24}{4 x}Q_2+\frac{d-4}{2x}Q_3.
    \end{split}
\end{align}
As we see, $Q_1$ is decoupled and can be integrated to,
\begin{align}
    Q_1(x)=c_1 x^{2-d}.
\end{align}
After massaging the equations a bit, we find that $Q_2$ satisfies the following differential equation,
\begin{align}
    4 x^2 Q_2'''(x)+(4-6 d)\, x\, Q_2''(x)+(x-2 (d-2) d) Q_2'(x)+(d-2) Q_2(x)=0
\end{align}
 The general solution takes the form,
 \begin{align}
     Q_2(x)=x^{\frac{2-d}{4}}\left(c_1 I_{\frac{d-2}{2}}\left(\sqrt{x}\right)+c_2 K_{\frac{d-2}{2}}(\sqrt{x})\right)+c_3 x^{3-d}\, _1F_2\left(1;4-d,3-\frac{d}{2};\frac{x}{4}\right)
 \end{align}
 Similarly, $Q_3=(3-d)Q_1-x\, Q_2'(x)$. The constants $c_i$ can be fixed by studying the asymptotic expansion of the integrals, as we showed previously.
\section{Useful Feynman integrals} \label{ch1:app:E}
We have seen that there is a direct connection between PN diagrams and the one-loop scalar Feynman diagrams, Fig.~(\ref{fig16}). Below, we list master integrals that will be useful in our context. We take the integrals from \cite{Levi:2011eq}.
%\vspace{-1.3cm}
\begingroup
\small
\begin{eqnarray}
\label{ch1:tensorfourierindentity}
I^i\equiv\rmint\frac{d^d\bf{k}}{(2\pi)^d}\frac{k^ie^{i\bf{k}\cdot\bf{r}}}{({\bf{k}}^2)^\alpha}&=&\frac{i}{(4\pi)^{d/2}}\frac{\Gamma(d/2-\alpha+1)}{\Gamma(\alpha)}\left(\frac{{\bf{r}}^2}{4}\right)^{\alpha-d/2-1/2}n^i,\\
I^{ij}\equiv\rmint\frac{d^d\bf{k}}{(2\pi)^d}\frac{k^ik^je^{i\bf{k}\cdot\bf{r}}}{({\bf{k}}^2)^\alpha}&=&\frac{1}{(4\pi)^{d/2}}\frac{\Gamma(d/2-\alpha+1)}{\Gamma(\alpha)}\left(\frac{{\bf{r}}^2}{4}\right)^{\alpha-d/2-1}\left(\frac{1}{2}\delta^{ij}+(\alpha-1-d/2)n^in^j\right),\\
I^{ijl}\equiv\rmint\frac{d^d\bf{k}}{(2\pi)^d}\frac{k^ik^jk^le^{i\bf{k}\cdot\bf{r}}}{({\bf{k}}^2)^\alpha}&=&\frac{i}{(4\pi)^{d/2}}\frac{\Gamma(d/2-\alpha+2)}{\Gamma(\alpha)}\left(\frac{{\bf{r}}^2}{4}\right)^{\alpha-d/2-3/2}\nonumber\\
&&\quad
\times\left(\frac{1}{2}\left(\delta^{ij}n^l+\delta^{il}n^j+\delta^{jl}n^i\right)+(\alpha-d/2-2)n^in^jn^l\right),\\
I^{ijlm}\textstyle{\equiv\rmint\frac{d^d\bf{k}}{(2\pi)^d}\frac{k^ik^jk^lk^me^{i\bf{k}\cdot\bf{r}}}{({\bf{k}}^2)^\alpha}}&=&\textstyle{\frac{1}{(4\pi)^{d/2}}\frac{\Gamma(d/2-\alpha+2)}{\Gamma(\alpha)}\left(\frac{{\bf{r}}^2}{4}\right)^{\alpha-d/2-2}\left(\frac{1}{4}\left(\delta^{ij}\delta^{lm}+\delta^{il}\delta^{jm}+\delta^{im}\delta^{jl}\right)\right.} \\
&&
\textstyle{+\frac{\alpha-d/2-2}{2}\left(\delta^{ij}n^ln^m+\delta^{il}n^jn^m+\delta^{im}n^jn^l+\delta^{jl}n^in^m+\delta^{jm}n^in^l+\delta^{lm}n^in^j\right) }\nonumber\\
&&\left.\quad
+(\alpha-d/2-2)(\alpha-d/2-3)n^in^jn^ln^m\right).\nonumber
\end{eqnarray}
\endgroup
We use the d-dimensional master formula for one-loop scalar integrals given by ,
\begin{align}
\begin{split}
   & J\equiv\rmint \frac{d^d\bf{k}}{(2\pi)^d}\frac{1}{\left[{\bf{k}}^2\right]^\alpha\left[({\bf{k}-\bf{q}})^2\right]^\beta}=  \frac{1}{(4\pi)^{d/2}}\frac{\Gamma(\alpha+\beta-d/2)}{\Gamma(\alpha)\Gamma(\beta)}\frac{\Gamma(d/2-\alpha)\Gamma(d/2-\beta)}{\Gamma(d-\alpha-\beta)}\left(q^2\right)^{d/2-\alpha-\beta}.\label{ch1:eq:1loop}
    \end{split}
\end{align}
The d-dimensional master formula for one-loop tensor integrals is taken from \cite{Levi:2011eq}. 
Similarly, one can also derive the following d-dimensional formulas for the one-loop tensor integrals: 
\begin{align}
\begin{split}
J^i\equiv\rmint \frac{d^d\bf{k}}{(2\pi)^d}\frac{k^i}{\left[{\bf{k}}^2\right]^\alpha\left[({\bf{k}-\bf{q}})^2\right]^\beta}=\frac{1}{(4\pi)^{d/2}}\frac{\Gamma(\alpha+\beta-d/2)}{\Gamma(\alpha)\Gamma(\beta)}\frac{\Gamma(d/2-\alpha+1)\Gamma(d/2-\beta)}{\Gamma(d-\alpha-\beta+1)}\\ \left(q^2\right)^{d/2-\alpha-\beta}q^i\end{split}
\end{align}
\begin{align}
\begin{split}
J^{ij}\equiv\rmint \frac{d^d\bf{k}}{(2\pi)^d}\frac{k^ik^j}{\left[{\bf{k}}^2\right]^\alpha\left[({\bf{k}-\bf{q}})^2\right]^\beta}=\frac{1}{(4\pi)^{d/2}}\frac{\Gamma(\alpha+\beta-d/2-1)}{\Gamma(\alpha)\Gamma(\beta)}\frac{\Gamma(d/2-\alpha+1)\Gamma(d/2-\beta)}{\Gamma(d-\alpha-\beta+2)}\left(q^2\right)^{d/2-\alpha-\beta}\\
\times\left(\frac{d/2-\beta}{2}q^2\delta^{ij}+(\alpha+\beta-d/2-1)(d/2-\alpha+1)q^iq^j\right),\end{split}
\end{align}
\begin{align}
\begin{split}
J^{ijl}\equiv\rmint \frac{d^d\bf{k}}{(2\pi)^d}\frac{k^ik^jk^l}{\left[{\bf{k}}^2\right]^\alpha\left[({\bf{k}-\bf{q}})^2\right]^\beta}=\frac{1}{(4\pi)^{d/2}}\frac{\Gamma(\alpha+\beta-d/2-1)}{\Gamma(\alpha)\Gamma(\beta)}\frac{\Gamma(d/2-\alpha+2)\Gamma(d/2-\beta)}{\Gamma(d-\alpha-\beta+3)}\left(q^2\right)^{d/2-\alpha-\beta} \\ 
\times\left(\frac{d/2-\beta}{2}q^2\left(\delta^{ij}q^l+\delta^{il}q^j+\delta^{jl}q^i\right)\right.  \left. 
+(\alpha+\beta-d/2-1)(d/2-\alpha+2)q^iq^jq^l\right),
\end{split}
\end{align}


\begin{align}
\begin{split}
&J^{ijlm}\equiv\rmint \frac{d^d\bf{k}}{(2\pi)^d}\frac{k^ik^jk^lk^m}{\left[{\bf{k}}^2\right]^\alpha\left[({\bf{k}-\bf{q}})^2\right]^\beta}=\frac{1}{(4\pi)^{d/2}}\frac{\Gamma(\alpha+\beta-d/2-2)}{\Gamma(\alpha)\Gamma(\beta)}\\& \frac{\Gamma(d/2-\alpha+2)\Gamma(d/2-\beta)}{\Gamma(d-\alpha-\beta+4)}\left(q^2\right)^{d/2-\alpha-\beta}\,\,\,\,
\times\left(\frac{(d/2-\beta)(d/2-\beta+1)}{4}q^4\left(\delta^{ij}\delta^{lm}+\delta^{il}\delta^{jm}+\delta^{jl}\delta^{im}\right)\right.\\
& \textstyle{
+(\alpha+\beta-d/2-2)(d/2-\alpha+2)\frac{d/2-\beta}{2}q^2 
\times\left(\delta^{ij}q^lq^m+\delta^{il}q^jq^m+\delta^{im}q^jq^l+\delta^{jl}q^iq^m+\delta^{jm}q^iq^l+\delta^{lm}q^iq^j\right)}\\
&\left.\,\,\,\,\,\,\,\,\,
+(\alpha+\beta-d/2-2)(\alpha+\beta-d/2-1)(d/2-\alpha+2)(d/2-\alpha+3)\right.
\left.
\times q^iq^jq^lq^m\right).
\end{split}
\end{align}
For more details on the Feynman integral, especially the massive integrals, can be found in \cite{smirnov}.
